% !TeX program = xelatex
% 本文件完整、自包含；可用 XeLaTeX 或 Tectonic 编译。
\documentclass[UTF8,fontset=fandol,12pt,oneside,openany]{ctexrep}
\usepackage[a4paper,left=25mm,right=25mm,top=25mm,bottom=25mm,headheight=18pt,headsep=8mm]{geometry}
\usepackage{amsmath,amssymb,amsthm,mathtools,booktabs,longtable,array}
\usepackage{tikz,enumitem,fancyhdr,needspace,microtype}
\usetikzlibrary{arrows.meta,calc,patterns,positioning}
\usepackage[hidelinks,unicode,bookmarksnumbered=true,hypertexnames=false]{hyperref}
\newcommand{\R}{\mathbb R}
\newcommand{\N}{\mathbb N^{+}}
\newcommand{\dd}{\mathop{}\!\mathrm d}
\newcommand{\ee}{\mathrm e}
\theoremstyle{definition}
\newtheorem{theorem}{定理}[chapter]
\newtheorem{definition}[theorem]{定义}
\newtheorem{proposition}[theorem]{命题}
\newtheorem{lemma}[theorem]{引理}
\newtheorem{corollary}[theorem]{推论}
\newtheorem{example}{例}[chapter]
\newtheorem{exercise}{练习}[chapter]
\newtheorem*{remark}{注}
\renewcommand{\proofname}{证明}
\ctexset{
 chapter={format=\centering\LARGE\bfseries,beforeskip=0pt,afterskip=22pt},
 section={format=\Large\bfseries,beforeskip=20pt,afterskip=10pt},
 subsection={format=\large\bfseries,beforeskip=14pt,afterskip=7pt},
 paragraph={runin=false,format=\normalsize\bfseries,beforeskip=8pt,afterskip=4pt}}
\setlength{\parindent}{2em}
\setlength{\parskip}{.25em}
\linespread{1.18}
\setlist{itemsep=.2em,topsep=.4em,leftmargin=2em}
\setlist[enumerate]{label=\arabic*.,ref=\arabic*}
\setlength{\emergencystretch}{3em}
\setlength{\LTpre}{.5em}
\setlength{\LTpost}{.5em}
\renewcommand{\arraystretch}{1.3}
\widowpenalty=10000
\clubpenalty=10000
\predisplaypenalty=10000
\raggedbottom
\setcounter{tocdepth}{1}
\setcounter{secnumdepth}{2}
\pagestyle{fancy}
\fancyhf{}
\fancyhead[L]{\small 微积分（甲）I}
\fancyhead[R]{\small\nouppercase{\leftmark}}
\fancyfoot[C]{\thepage}
\renewcommand{\headrulewidth}{.35pt}
\renewcommand{\chaptermark}[1]{\markboth{第\thechapter 章\quad #1}{}}
\fancypagestyle{plain}{\fancyhf{}\fancyhead[L]{\small 微积分（甲）I}\fancyhead[R]{\small\nouppercase{\leftmark}}\fancyfoot[C]{\thepage}\renewcommand{\headrulewidth}{.35pt}}

\hypersetup{pdftitle={微积分（甲）I：知识讲义与基础、经典例题},pdfauthor={CangLing-Cerulean}}
\begin{document}
\begin{titlepage}
\centering
\vspace*{28mm}
{\Large 大学知识库\par}
\vspace{15mm}
{\Huge\bfseries 微积分（甲）I\par}
\vspace{12mm}
{\LARGE\bfseries 知识讲义与基础、经典例题\par}
\vspace{12mm}
\rule{.65\textwidth}{.5pt}
\vfill
{\large CangLing-Cerulean\par}
\vspace{5mm}
{2026 年 10 月\par}
\vspace*{18mm}
\end{titlepage}
\pagenumbering{roman}
\chapter*{阅读说明}
\markboth{阅读说明}{}
\addcontentsline{toc}{chapter}{阅读说明}
本册面向第一次学习微积分的读者，从概念动机和严格定义出发，逐步建立证明与计算方法。例题展示方法选择及关键推理，章末练习用于检查理解。拓展内容标明前置知识，首次学习可先完成基础路线，再阅读综合方法。

两册采用同一章节顺序：函数、数列极限、函数极限、连续性与间断点、导数与微分、
微分中值定理、洛必达法则、泰勒公式、导数的应用、不定积分、定积分、定积分的应用、反常积分。
记号 $\N$ 表示正整数集，$\R$ 表示实数集。凹凸性采用“凸（下凸）”表示图像在弦下方的约定，
具体以定义不等式为准。

\tableofcontents
\clearpage
\pagenumbering{arabic}
\chapter{函数}
\label{lec:01:chapter}

微积分研究量怎样变化，以及变化怎样积累。开始计算之前，先要确定变化的对象：哪些输入可以使用，一个输入对应哪个输出，输出能够落在哪个范围。本章只用集合、代数、不等式及高中三角知识，建立后续极限、导数和积分共用的语言。

\section{函数、映射与定义域}
\label{lec:01:function}
\subsection{输入、规则与输出}
圆的半径为 $r$ 时，面积为 $\pi r^2$。这里 $r$ 的取值受几何意义限制，通常取 $r>0$。同一个式子 $\pi x^2$ 若允许 $x\in\R$，描述的就是另一个定义域的函数。函数的意义不能只看一个表达式。

\begin{definition}
设 $A,B$ 为非空集合。若规则 $f$ 给 $A$ 中每个元素 $x$ 指定 $B$ 中唯一的元素 $f(x)$，称 $f:A\to B$ 为映射。$A$ 为定义域，$B$ 为陪域，实际输出组成的集合
\[
 f(A)=\{f(x):x\in A\}
\]
称为值域。定义域为实数的子集、值为实数的映射称为实函数。
\end{definition}
“每个”要求定义域内没有遗漏的输入；“唯一”排除一个输入对应两个输出。例如 $y^2=x$ 在 $x>0$ 时有两个解，关系本身不能指定一个实函数；选定 $y=\sqrt{x}$ 后才确定了一个函数。不同输入可以得到同一输出，$f(x)=x^2$ 中 $f(1)=f(-1)$ 并不违反函数定义。

\begin{definition}
映射 $f:A\to B$ 若满足 $f(x_1)=f(x_2)\Rightarrow x_1=x_2$，称为单射；若 $f(A)=B$，称为满射；同时为单射和满射时称为双射。
\end{definition}
单射问“一个输出能否来自两个输入”，满射问“陪域中的输出是否全部覆盖”。例如 $x\mapsto x^2$ 从 $\R$ 到 $[0,+\infty)$ 是满射但不是单射；限制到 $[0,+\infty)$ 后是双射。

\subsection{求定义域时取条件的交集}
未指定定义域的实函数表达式，通常采用使表达式有实数意义的最大集合。分母不为零，偶次根号的被开方数非负，对数真数严格为正。反三角函数的输入限制见本章后文。多个条件同时出现时必须取交集；化简约去因子不会自动补回原式排除的点。

\begin{example}[逐层检查定义域]
求 $F(x)=\dfrac{\ln\sqrt{1-x^2}}{x-\tfrac12}$ 的定义域，并比较 $\dfrac{x^2-1}{x-1}$ 与 $x+1$。
\end{example}
\paragraph{思路分析} 先检查最外层运算。对数要求整个根式为正，比根式本身要求非负更严格；再排除分母为零的点。
\paragraph{解答} 条件为 $1-x^2>0$ 以及 $x\ne\tfrac12$，故
\[
 D_F=(-1,\tfrac12)\cup(\tfrac12,1).
\]
虽然在 $x=\pm1$ 时根式存在，$\ln0$ 仍无意义。另有
\[
 \frac{x^2-1}{x-1}=x+1\quad(x\ne1).
\]
左式的自然定义域为 $\R\setminus\{1\}$，右式为 $\R$，所以两者在共同定义域上取值相同，但作为具有各自自然定义域的函数并不相同。

\subsection{求值域要说明哪些值能够取到}
求出一个上界或下界，只证明输出落在某个集合中；求值域还需证明该集合中的每个值都有对应输入。配方、换元以及反解方程常能在不使用导数的情况下完成这一步。

\begin{example}[配方与端点]
求 $f(x)=x^2-2x+3$ 在 $[-1,2)$ 上的值域。
\end{example}
\paragraph{思路分析} 配方把问题化为输入到顶点 $x=1$ 的距离。右端点不取到，不意味着最大值一定不取到，还要检查另一端。
\paragraph{解答} $f(x)=(x-1)^2+2$。当 $x\in[-1,2)$ 时，$x-1\in[-2,1)$，故 $(x-1)^2\in[0,4]$。其中 $0$ 在 $x=1$ 取到，$4$ 在 $x=-1$ 取到。对任意 $y\in[2,6]$，取
\[
 x=1-\sqrt{y-2}\in[-1,1],
\]
就有 $f(x)=y$。因此值域恰为 $[2,6]$。

\section{复合函数与反函数}
\label{lec:01:composition}
\subsection{复合的顺序与中间变量}
若先把摄氏温度换成华氏温度，再用仪器的刻度规则显示读数，就经历两次对应。数学中用 $f\circ g$ 表示“先 $g$，后 $f$”。

\begin{definition}
设实函数 $g$、$f$ 的定义域分别为 $D_g,D_f$。在非空集合
\[
 D_{f\circ g}=\{x\in D_g:g(x)\in D_f\}
\]
上，定义 $(f\circ g)(x)=f(g(x))$。$g$ 称为内函数，$f$ 称为外函数，$u=g(x)$ 为中间变量。
\end{definition}
若希望复合在整个 $D_g$ 上定义，就必须有 $g(D_g)\subseteq D_f$；仅有两集合相交，只能在符合条件的部分输入上复合。

\begin{example}[复合顺序决定定义域]
设 $f(u)=\sqrt{u}$，$g(x)=x-1$，分别求 $f\circ g$ 和 $g\circ f$。
\end{example}
\paragraph{思路分析} 写出内函数的输出，再检查它能否进入外函数。顺序不同，检查的对象也不同。
\paragraph{解答} 对 $f\circ g$，先得到 $u=x-1$，须满足 $u\ge0$，故
\[
 (f\circ g)(x)=\sqrt{x-1},\qquad x\in[1,+\infty).
\]
对 $g\circ f$，内函数须满足 $x\ge0$，外函数 $u-1$ 对所有实数有意义，故
\[
 (g\circ f)(x)=\sqrt{x}-1,\qquad x\in[0,+\infty).
\]
这说明复合运算一般不满足交换律。拆解复杂表达式时，也按从外到内的顺序认出每一层；例如 $\ee^{\sqrt{\sin(\ee^x+1)}}$ 可拆成 $w=\ee^x+1$、$v=\sin w$、$u=\sqrt v$、$y=\ee^u$，并要求 $v\ge0$。

\subsection{反函数的存在与求法}
\label{lec:01:inverse}
反函数把输出还原为输入。若一个输出可能来自两个输入，还原就没有唯一结果。

\begin{theorem}
函数 $f:D\to f(D)$ 有反函数，当且仅当 $f$ 是单射。其反函数记为 $f^{-1}:f(D)\to D$，满足
\[
 f^{-1}(f(x))=x\quad(x\in D),\qquad f(f^{-1}(y))=y\quad(y\in f(D)).
\]
\end{theorem}
\begin{proof}
若反函数存在而 $f(x_1)=f(x_2)$，对这个相同输出应用反函数，得到 $x_1=x_2$，故 $f$ 为单射。反过来，对每个 $y\in f(D)$，值域的定义保证至少有一个 $x\in D$ 使 $f(x)=y$；单射保证至多有一个。把该唯一输入指定给 $y$，便得到反函数及上述恒等式。
\end{proof}
$f^{-1}(x)$ 表示反函数值，不能与倒数 $1/f(x)$ 混淆。求反函数时先确认单射，求出值域，再从 $y=f(x)$ 解出 $x$，最后交换变量名称并写清新定义域。图像中点 $(x,y)$ 与反函数的点 $(y,x)$ 对应，所以两图像关于 $y=x$ 对称。

\begin{example}[限制定义域后求反函数]
求 $f(x)=(x-1)^2$ 在 $[1,+\infty)$ 上的反函数；说明在 $\R$ 上为何没有反函数。
\end{example}
\paragraph{思路分析} 解平方方程会出现两个根，给定定义域决定应取哪一根。
\paragraph{解答} 若 $1\le x_1<x_2$，则 $0\le x_1-1<x_2-1$，平方后仍严格小于，故 $f$ 为单射。其值域为 $[0,+\infty)$。由 $y=(x-1)^2$ 和 $x-1\ge0$ 得 $x=1+\sqrt y$，故
\[
 f^{-1}(x)=1+\sqrt x,\qquad x\in[0,+\infty).
\]
在 $\R$ 上，$f(0)=f(2)=1$，同一输出无法唯一还原，所以不存在反函数。若把原定义域改为 $(-\infty,1]$，反函数则为 $1-\sqrt x$。反函数表达式依赖所选分支。

\section{单调性、奇偶性与周期性}
\label{lec:01:properties}
\subsection{单调性是任意两点之间的比较}
\begin{definition}
设 $E\subseteq\R$。若任意 $x_1,x_2\in E$ 满足 $x_1<x_2$ 时都有 $f(x_1)\le f(x_2)$，称 $f$ 在 $E$ 上递增；改成严格不等号时称为严格递增。将输出不等号反向，分别定义递减与严格递减。
\end{definition}
“递增”在本讲义允许相等，称严格递增时才排除相等。单调性必须写明所在集合。证明时任选两点，计算差值、比较比值或使用已知基本函数的保序性；不能只比较几个特殊点。

\begin{example}[差值法与跨区间比较]
证明 $h(x)=\dfrac{x}{1+x}$ 在 $(-1,+\infty)$ 上严格递增，并说明 $q(x)=1/x$ 在其整个定义域上不单调。
\end{example}
\paragraph{思路分析} 有理式作差后可能约成含 $x_2-x_1$ 的乘积。分母的符号由所在区间控制。
\paragraph{解答} 任取 $-1<x_1<x_2$，有
\[
 h(x_2)-h(x_1)=\frac{x_2(1+x_1)-x_1(1+x_2)}{(1+x_1)(1+x_2)}
 =\frac{x_2-x_1}{(1+x_1)(1+x_2)}>0.
\]
故 $h$ 严格递增。对于 $q$，在 $(-\infty,0)$、$(0,+\infty)$ 各自严格递减，但 $-1<1$ 而 $q(-1)<q(1)$，违反整个定义域上的递减条件；$1<2$ 而 $q(1)>q(2)$，又违反递增条件。于是 $q$ 在 $\R\setminus\{0\}$ 上不单调，却是单射且反函数为自身。这也说明存在反函数不必要求整个定义域上严格单调。

\begin{theorem}
严格单调函数存在反函数，反函数与原函数具有相同的严格单调方向。
\end{theorem}
\begin{proof}
严格单调保证不同输入有不同输出，所以存在反函数。以严格递增为例，取值域内 $y_1<y_2$，记 $x_i=f^{-1}(y_i)$。若 $x_1=x_2$ 则 $y_1=y_2$；若 $x_1>x_2$ 则递增性给出 $y_1>y_2$。两种情形均矛盾，故 $x_1<x_2$。递减情形同理，$y_1<y_2$ 时必有 $x_1>x_2$。
\end{proof}
复合函数的单调性也可由两次比较决定：两次保序或两次反序最终保序，一次保序一次反序最终反序。使用这一规律前要确认内函数输出始终处在外函数所讨论的单调区间。例如 $x^2$ 在 $\R$ 上不能当成递增函数使用。

\subsection{奇偶性与定义域的对称}
\begin{definition}
若定义域 $D$ 满足 $x\in D\Rightarrow -x\in D$，且对所有 $x\in D$ 有 $f(-x)=f(x)$，称 $f$ 为偶函数；若 $f(-x)=-f(x)$，称 $f$ 为奇函数。
\end{definition}
偶函数的点 $(x,f(x))$ 与 $(-x,f(x))$ 成对，图像关于 $y$ 轴对称；奇函数的点与 $(-x,-f(x))$ 成对，图像关于原点对称。奇函数若在零点有定义，必有 $f(0)=-f(0)$，故 $f(0)=0$；逆命题不成立，如 $x^2$ 在零点为零却不是奇函数。

\begin{example}[从定义域到代入检验]
判断 $f(x)=\dfrac{x}{1+x^2}$、$g(x)=\ln\dfrac{1+x}{1-x}$ 的奇偶性，并将任意定义在对称集合 $D$ 上的函数 $F$ 拆成偶函数与奇函数之和。
\end{example}
\paragraph{思路分析} 先核对输入 $-x$ 是否始终可用，再比较代入后的式子。要拆出对称部分，可以同时使用 $F(x)$ 与 $F(-x)$。
\paragraph{解答} $f$ 的定义域为 $\R$，且 $f(-x)=-x/(1+x^2)=-f(x)$，所以为奇函数。$g$ 的真数为正等价于 $-1<x<1$，定义域对称。在该区间内，两个比值均为正，故
\[
 g(-x)=\ln\frac{1-x}{1+x}=-\ln\frac{1+x}{1-x}=-g(x).
\]
因此 $g$ 也是奇函数。定义
\[
 F_{\mathrm e}(x)=\frac{F(x)+F(-x)}2,\qquad
 F_{\mathrm o}(x)=\frac{F(x)-F(-x)}2.
\]
直接代入可得 $F_{\mathrm e}(-x)=F_{\mathrm e}(x)$、$F_{\mathrm o}(-x)=-F_{\mathrm o}(x)$，而两式相加为 $F$。例如 $F(x)=x^2+x+1$ 的偶部分为 $x^2+1$、奇部分为 $x$。
\paragraph{易错说明} $x^2$ 在 $[0,+\infty)$ 上的限制不是偶函数，因为定义域不对称。恒零函数在对称定义域上既奇又偶；若函数同时奇偶，则 $f(x)=-f(x)$，所以只能恒为零。

\subsection{周期是全局的平移重复}
\begin{definition}
若 $T>0$ 满足 $D+T=D$，且对所有 $x\in D$ 有 $f(x+T)=f(x)$，称 $T$ 为 $f$ 的一个正周期，$f$ 为周期函数。其中 $D+T=\{x+T:x\in D\}$。若正周期中存在最小者，称为最小正周期。
\end{definition}
一个周期的正整数倍也是周期。周期性只保证重复，单独不能保证有界。例如 $\tan x$ 以 $\pi$ 为周期，却在其定义域上无界。常函数的每个正数都是周期，所以没有最小正周期。

\begin{example}[求周期并证明最小性]
求 $f(x)=\sin(2x+\pi/3)$ 的最小正周期，并说明狄利克雷函数的周期特点。
\end{example}
\paragraph{思路分析} 把输入平移转为相位平移；写出一个可用周期后，还要排除更小者。
\paragraph{解答} 平移 $\pi$ 使相位增加 $2\pi$，故 $\pi$ 为周期。若 $T>0$ 为周期，令 $u=2x+\pi/3$，由于 $u$ 遍历 $\R$，有 $\sin(u+2T)=\sin u$ 对每个 $u$ 成立。取 $u=0$ 得 $\sin2T=0$；取 $u=\pi/2$ 得 $\cos2T=1$。两式共同要求 $2T=2k\pi$，其中 $k$ 为正整数，故最小正周期为 $\pi$。
定义狄利克雷函数 $D(x)$ 在有理数上取 $1$、无理数上取 $0$。任意正有理数 $r$ 平移后不改变输入的有理或无理性质，所以 $r$ 为周期；正无理数 $s$ 不是周期，因为 $D(0)=1$ 而 $D(s)=0$。正有理数中没有最小者，故此函数虽非常值，也无最小正周期。

\section{有界性、最值与确界}
\label{lec:01:bounds}
\subsection{统一的界与无界的量词}
\begin{definition}
若存在实数 $U$ 使每个 $x\in D$ 都满足 $f(x)\le U$，称 $f$ 在 $D$ 上有上界；若存在实数 $L$ 使每个 $x\in D$ 都满足 $f(x)\ge L$，称有下界。同时有上界和下界，称有界。
\end{definition}
有界等价于存在 $M>0$ 使所有输入满足 $|f(x)|\le M$。若已知 $L\le f(x)\le U$，可取 $M=1+\max\{|L|,|U|\}$；反向则取上下界 $\pm M$。这里 $M$ 必须对全部输入共用，不能随 $x$ 改变。每个点的函数值有限，并不代表函数有界。

无界的准确表述为
\[
 \text{对每个 }M>0,\ \text{存在 }x\in D\text{，使 }|f(x)|>M.
\]
只要缺少上界或缺少下界，就无界。例如 $\ee^x$ 有下界 $0$，却没有上界。

\begin{example}[放缩与构造输入]
证明 $f(x)=\dfrac{x\sin x}{1+x^2}$ 在 $\R$ 上有界；证明 $g(x)=1/\sqrt{x}$ 在 $(0,1]$ 上无界。
\end{example}
\paragraph{思路分析} 第一题中正弦已受统一常数控制，只需控制有理因子。第二题要对任意给定的界构造击破它的输入。
\paragraph{解答} 因 $(|x|-1)^2\ge0$，有 $2|x|\le1+x^2$；分母始终为正，故
\[
 |f(x)|\le\frac{|x|}{1+x^2}\le\frac12\qquad(x\in\R).
\]
因此 $M=1/2$ 是可用的界。这个证明没有断言 $1/2$ 是最大值：要在最后一个不等号取等需 $|x|=1$，要在前一个取等又需 $|\sin x|=1$，不能仅凭界直接认定其被取到。
对任意 $M>0$，取 $x_M=1/(M+1)^2$，则 $0<x_M<1$ 且
\[
 g(x_M)=M+1>M.
\]
所以 $g$ 无界。其值域为 $[1,+\infty)$，有下界与最小值 $1$，这与无界没有矛盾。

\subsection{最值必须取到，确界只要求最优的界}
\begin{definition}
若有 $x_0\in D$ 使所有 $x\in D$ 满足 $f(x)\le f(x_0)$，则 $f(x_0)$ 为最大值；最小值类似。值域的所有上界中最小者称为上确界，记 $\sup_{x\in D}f(x)$；所有下界中最大者称为下确界，记 $\inf_{x\in D}f(x)$。
\end{definition}
实数的完备性保证非空、上有界的实数集合有上确界，非空、下有界的实数集合有下确界。这是实数系统的基本性质，本章将其作为依据；本章不把无上界函数的上确界写成一个有限实数。

\begin{proposition}
$S$ 是非空实数集合。实数 $U$ 为 $S$ 的上确界，当且仅当：一是对每个 $s\in S$ 有 $s\le U$；二是对每个 $\varepsilon>0$，存在 $s\in S$ 使 $U-\varepsilon<s$。
\end{proposition}
\begin{proof}
若 $U$ 为最小上界而第二条不成立，就存在 $\varepsilon>0$ 使所有 $s\le U-\varepsilon$，从而 $U-\varepsilon$ 成为更小的上界，矛盾。反之，第一条说明 $U$ 是上界；对任意 $V<U$，在第二条取 $\varepsilon=U-V$，找到 $s>V$，所以 $V$ 不是上界。故 $U$ 是最小上界。
\end{proof}
这一定理将“最小上界”落实为可检查的两步：先证明全部输出不超过它，再证明输出可以从下方任意靠近它。若上确界被取到，它就是最大值；若没有取到，则不存在最大值。

\begin{example}[有确界而没有最大值]
求 $h(x)=\dfrac{x}{1+x}$ 在 $[0,+\infty)$ 上的值域、上下确界及最值。
\end{example}
\paragraph{思路分析} 写成 $1-1/(1+x)$ 得到上界；反解输出能同时验证值域和端点是否取得。
\paragraph{解答} $h(x)=1-1/(1+x)$，故 $0\le h(x)<1$。对任意 $y\in[0,1)$，取 $x=y/(1-y)\ge0$，有 $h(x)=y$，所以值域恰为 $[0,1)$。因此下确界与最小值均为 $0$，在 $x=0$ 取到。对任意 $\varepsilon>0$，取 $x=1/\varepsilon$，则
\[
 1-h(x)=\frac{\varepsilon}{1+\varepsilon}<\varepsilon,
\]
故上确界为 $1$。然而每个有限输入都有 $h(x)<1$，所以没有最大值。

\section{基本初等函数与图像变换}
\label{lec:01:elementary}
\subsection{幂、指数与对数}
常数函数 $y=C$ 的值域是单点集合。幂函数 $y=x^a$ 的定义域随指数改变：非负整数幂通常在 $\R$ 上定义，负整数幂排除零；$x^{1/2}$ 要求 $x\ge0$，$x^{1/3}$ 可在 $\R$ 上定义。任意实指数在 $x>0$ 上统一定义为
\[
 x^a=\ee^{a\ln x}.
\]
讨论零点或负输入时须另外检查指数种类，不能把 $x>0$ 的统一定义直接延伸。

指数函数 $a^x$ 要求 $a>0$ 且 $a\ne1$，定义域为 $\R$，值域为 $(0,+\infty)$；$a>1$ 时严格递增，$0<a<1$ 时严格递减。对数函数 $\log_a x$ 是它的反函数，定义域为 $(0,+\infty)$，值域为 $\R$，具有相同单调方向。$\ln x$ 表示底数为 $\ee$ 的对数。对正数 $u,v$，有
\[
 \ln(uv)=\ln u+\ln v,\qquad \ln(u/v)=\ln u-\ln v,
 \qquad \ln(u^a)=a\ln u.
\]
这些式子来自指数运算法则及反函数关系。若因子为负，不能拆开取实对数；例如 $\ln(x^2)=2\ln|x|$ 对 $x\ne0$ 成立，而 $2\ln x$ 只在 $x>0$ 有意义。

\subsection{三角函数及其取值限制}
三角函数的自变量一律采用弧度。下表给出自然定义域、值域及最小正周期，其中 $k\in\mathbb Z$。
\begin{center}
\begin{tabular}{llll}
\toprule
函数 & 定义域条件 & 值域 & 最小正周期\\
\midrule
$\sin x$ & $x\in\R$ & $[-1,1]$ & $2\pi$\\
$\cos x$ & $x\in\R$ & $[-1,1]$ & $2\pi$\\
$\tan x$ & $x\ne\pi/2+k\pi$ & $\R$ & $\pi$\\
$\cot x$ & $x\ne k\pi$ & $\R$ & $\pi$\\
$\sec x$ & $x\ne\pi/2+k\pi$ & $(-\infty,-1]\cup[1,+\infty)$ & $2\pi$\\
$\csc x$ & $x\ne k\pi$ & $(-\infty,-1]\cup[1,+\infty)$ & $2\pi$\\
\bottomrule
\end{tabular}
\end{center}
其中 $\tan x=\sin x/\cos x$、$\cot x=\cos x/\sin x$、$\sec x=1/\cos x$、$\csc x=1/\sin x$。由 $\sin^2x+\cos^2x=1$，在相应分母非零时分别除以 $\cos^2x$、$\sin^2x$，可得
\[
 1+\tan^2x=\sec^2x,\qquad 1+\cot^2x=\csc^2x.
\]
这些等式也解释 $\sec x$、$\csc x$ 的绝对值至少为 $1$。常用单调区间为：$\sin x$ 在 $[-\pi/2,\pi/2]$ 严格递增，$\cos x$ 在 $[0,\pi]$ 严格递减，$\tan x$ 在 $(-\pi/2,\pi/2)$ 严格递增，$\cot x$ 在 $(0,\pi)$ 严格递减。全实轴上不能把正弦、余弦当作单调函数。

\subsection{反三角函数的主值}
限制到上述单调区间，三角函数便有反函数。选择哪个区间是定义的一部分，本讲义固定如下约定。
\begin{center}
\begin{tabular}{llll}
\toprule
函数 & 定义域 & 值域（主值区间） & 单调方向\\
\midrule
$\arcsin x$ & $[-1,1]$ & $[-\pi/2,\pi/2]$ & 严格递增\\
$\arccos x$ & $[-1,1]$ & $[0,\pi]$ & 严格递减\\
$\arctan x$ & $\R$ & $(-\pi/2,\pi/2)$ & 严格递增\\
$\operatorname{arccot}x$ & $\R$ & $(0,\pi)$ & 严格递减\\
\bottomrule
\end{tabular}
\end{center}
因此 $\sin(\arcsin x)=x$ 对 $|x|\le1$ 成立，而 $\arcsin(\sin t)=t$ 只在 $t\in[-\pi/2,\pi/2]$ 成立。其余反三角函数也必须区分这两个复合方向。

\begin{example}[主值和符号共同决定答案]
计算 $\arcsin(\sin(5\pi/6))$；化简 $\cos(\arcsin x)$ 与 $\tan(\arccos x)$，并写出适用范围。
\end{example}
\paragraph{思路分析} 反三角函数返回主值区间中的角；利用三角恒等式求另一三角值时，开平方的符号由这个角所在区间决定。
\paragraph{解答} $\sin(5\pi/6)=1/2$，主值区间内正弦为 $1/2$ 的角是 $\pi/6$，所以第一问为 $\pi/6$。设 $\theta=\arcsin x\in[-\pi/2,\pi/2]$，则 $\cos\theta\ge0$，故
\[
 \cos(\arcsin x)=\sqrt{1-x^2},\qquad -1\le x\le1.
\]
设 $\varphi=\arccos x\in[0,\pi]$，则 $\sin\varphi\ge0$，因此
\[
 \tan(\arccos x)=\frac{\sqrt{1-x^2}}{x},\qquad x\in[-1,1]\setminus\{0\}.
\]
在 $x=0$ 时角为 $\pi/2$，正切无意义；在 $x<0$ 时结果为负，不能用所有边长均为正的直角三角形图直接忽略符号。

\subsection{初等函数与典型分段函数}
由常数、幂、指数、对数、三角和反三角函数经过有限次四则运算及复合得到的函数，称为初等函数，各运算仍须满足定义域条件。例如 $x^x$ 在 $x>0$ 上等于 $\ee^{x\ln x}$，所以为初等函数。仅凭是否写成分段形式，无法判断是否为初等函数：$|x|=\sqrt{x^2}$ 就是初等函数。

后续反例中还会出现以下三个函数：
\[
 \operatorname{sgn}x=\begin{cases}-1,&x<0,\\0,&x=0,\\1,&x>0,\end{cases}
 \qquad \lfloor x\rfloor=n\quad\Longleftrightarrow\quad n\le x<n+1,
 \quad n\in\mathbb Z,
\]
以及前述狄利克雷函数。$\lfloor x\rfloor$ 是向下取整，例如 $\lfloor-1.2\rfloor=-2$，不能把负数的小数部分直接删掉当作向下取整。分段函数每一段的条件必须覆盖定义域，交界点的取值必须唯一。

\subsection{由点的对应推导图像变换}
\label{lec:01:graph}
若原图像上有点 $(u,f(u))$，对变换后的表达式求出新输入与 $u$ 的关系，便能避免平移方向记忆错误。对
\[
 y=A f(B(x-h))+k,\qquad A\ne0,\quad B\ne0,
\]
令 $u=B(x-h)$，对应点为
\[
 (u,f(u))\longmapsto\left(h+\frac{u}{B},\ k+A f(u)\right).
\]
所以 $h$ 控制水平平移，$k$ 控制竖直平移；$|B|$ 越大，水平方向越压缩，$B<0$ 还包含左右反射；$|A|$ 控制竖直伸缩，$A<0$ 包含上下反射。变换后的定义域为 $h+D_f/B$。$|f(x)|$ 把原图像在 $x$ 轴下方的部分翻到上方；$f(|x|)$ 则保留右半图像并向左镜像，二者作用不同。

\begin{example}[绝对值图像与值域]
从 $y=x^2$ 出发说明 $y=2(x-1)^2-1$ 的图像，并画出 $y=|x^2-1|$ 的示意图、求其值域。
\end{example}
\paragraph{思路分析} 前者追踪顶点及伸缩；后者先找 $x^2-1$ 的正负分界，再翻转负值部分。
\paragraph{解答} 第一条抛物线由 $y=x^2$ 向右平移 $1$，竖直拉伸为原来的 $2$ 倍，再向下平移 $1$，顶点为 $(1,-1)$，对称轴为 $x=1$，值域为 $[-1,+\infty)$。第二个函数为
\[
 |x^2-1|=\begin{cases}1-x^2,&|x|\le1,\\x^2-1,&|x|>1.\end{cases}
\]
在 $[-1,1]$ 上把原来的负值翻到上方。函数总非负，在 $x=\pm1$ 取 $0$；对任意 $y\ge0$，取 $x=\sqrt{y+1}$ 可取到 $y$，故值域为 $[0,+\infty)$。
\begin{figure}[htbp]
\centering
\begin{tikzpicture}[x=1.3cm,y=1cm]
\draw[->] (-2.1,0)--(2.2,0) node[right] {$x$};
\draw[->] (0,-1.35)--(0,3.25) node[above] {$y$};
\draw[dashed,domain=-2:2,samples=65] plot (\x,{\x*\x-1});
\draw[thick,domain=-2:-1,samples=30] plot (\x,{\x*\x-1});
\draw[thick,domain=-1:1,samples=45] plot (\x,{1-\x*\x});
\draw[thick,domain=1:2,samples=30] plot (\x,{\x*\x-1});
\node[below] at (-1,0) {$-1$}; \node[below] at (1,0) {$1$};
\node[above left] at (0,1) {$1$}; \node[below left] at (0,-1) {$-1$};
\node[anchor=west] at (2.15,2.1) {$y=|x^2-1|$};
\end{tikzpicture}
\caption{虚线为原图像的负值部分，实线为取绝对值后的图像。}
\label{lec:01:fig:absolute}
\end{figure}

\section{分层练习与解析}
\label{lec:01:exercises}
\subsection{基础巩固}
\begin{exercise}
求 $F(x)=\dfrac{\sqrt{x+2}}{\ln(3-x)}$ 的定义域。
\end{exercise}
\paragraph{解析} 根式要求 $x\ge-2$，对数要求 $x<3$，分母要求 $\ln(3-x)\ne0$，即 $x\ne2$。交集为 $[-2,2)\cup(2,3)$。不能仅把分母条件写成 $3-x\ne0$。

\begin{exercise}
求 $f(x)=3-2x$ 在 $(-1,2]$ 上的反函数、值域与最值。
\end{exercise}
\paragraph{解析} 任取 $x_1<x_2$，有 $f(x_1)-f(x_2)=2(x_2-x_1)>0$，故严格递减。由 $-1<x\le2$ 得 $-1\le3-2x<5$。反解 $x=(3-y)/2$ 可验证每个 $y\in[-1,5)$ 均取得。故值域为 $[-1,5)$，反函数为 $f^{-1}(x)=(3-x)/2$，定义域 $[-1,5)$；最小值 $-1$ 在原输入 $2$ 取到，上确界为 $5$，无最大值。

\begin{exercise}
计算 $\arccos(\cos(4\pi/3))$、$\sin(\arctan x)$。
\end{exercise}
\paragraph{解析} $\cos(4\pi/3)=-1/2$，而 $[0,\pi]$ 内余弦为 $-1/2$ 的角是 $2\pi/3$，所以第一问为 $2\pi/3$。设 $\theta=\arctan x$，由 $\cos\theta>0$ 和 $1+\tan^2\theta=1/\cos^2\theta$ 得 $\cos\theta=1/\sqrt{1+x^2}$，故 $\sin\theta=\tan\theta\cos\theta=x/\sqrt{1+x^2}$，适用于全部 $x\in\R$。

\subsection{方法迁移}
\begin{exercise}
求 $f(x)=\sqrt{4-x^2}$、$g(x)=x+1$ 的两个复合函数及其定义域。
\end{exercise}
\paragraph{解析} $f$ 的定义域为 $[-2,2]$。$(f\circ g)(x)=\sqrt{4-(x+1)^2}$ 要求 $-2\le x+1\le2$，故定义域 $[-3,1]$。$(g\circ f)(x)=\sqrt{4-x^2}+1$ 的外函数无额外限制，故定义域为 $[-2,2]$。

\begin{exercise}
证明 $f(x)=\sqrt{1+x^2}-|x|$ 在 $\R$ 上有界，求其值域、确界及最值。
\end{exercise}
\paragraph{解析} 乘以共轭式，得到
\[
 f(x)=\frac1{\sqrt{1+x^2}+|x|},\qquad 0<f(x)\le1.
\]
最大值为 $1$，在 $x=0$ 取到。给定 $y\in(0,1]$，取 $t=(1-y^2)/(2y)\ge0$，可验证 $\sqrt{1+t^2}=t+y$，从而 $f(t)=y$。故值域为 $(0,1]$，上确界为 $1$，下确界为 $0$，无最小值。下确界也可直接检验：任意 $\varepsilon>0$，取 $x=1/\varepsilon$，有 $f(x)\le1/|x|=\varepsilon$；若需严格小于，可取 $x=2/\varepsilon$。

\subsection{综合提升}
\begin{exercise}
设 $f:\R\to\R$ 为奇函数，且 $f(x+1)=-f(x)$ 对所有实数成立。证明 $2$ 为周期，并说明能否据此断言最小正周期为 $2$。
\end{exercise}
\paragraph{解析} 对 $x+1$ 再应用条件，有 $f(x+2)=-f(x+1)=f(x)$，所以 $2$ 为周期。条件没有排除更小的周期。恒零函数满足全部条件且无最小正周期；非常值例子 $f(x)=\sin(3\pi x)$ 也满足奇性与 $f(x+1)=-f(x)$，但最小正周期为 $2/3$。因此不能把“找到周期”写成“找到最小正周期”。

\section{本章方法自查}
遇到一个函数，依次检查定义域、值域、运算条件和所讨论的集合。求反函数时检查单射与输出分支；判断单调性时比较任意两点，区分各分支与整个定义域；判断奇偶性时先查定义域对称；求周期时检查全局平移，另证最小性；求确界时先证界，再证可任意逼近，最后检查能否取到。完成这些检查后，后续微积分计算的对象和条件才是确定的。


\chapter{数列极限}
\label{lec:02:chapter}

数列极限研究的是：当项数越来越大时，数列的项能否稳定地接近同一个数。本章从极限的精确定义出发，建立运算与判定工具，再处理根式、乘积、平均值与递推数列。


\textbf{进入本章之前}\quad 需要会绝对值、三角不等式、分式与根式运算、数学归纳法。上确界是一个实数集合所有上界中最小的一个；本章会用它解释单调数列为什么收敛。

\textbf{学完之后}\quad 应当能够读懂并写出 $\varepsilon$--$N$ 证明，区分有界与收敛，选择合适的极限工具，完整证明递推数列的收敛性。

基础路线按本章小节顺序展开；Stolz 定理是补充工具，不需要导数或泰勒公式。


\section{从数列到极限}

数列 $\{a_n\}_{n\ge1}$ 可以看作定义在正整数集上的函数，$n$ 是项数，$a_n$ 是第 $n$ 项。观察下面四个数列：

{\small\begin{longtable}{>{\raggedright\arraybackslash}p{0.23\linewidth}>{\raggedright\arraybackslash}p{0.34\linewidth}>{\raggedright\arraybackslash}p{0.35\linewidth}}
\toprule
\textbf{数列} & \textbf{项的变化} & \textbf{收敛性} \\
\midrule\endfirsthead
\toprule
\textbf{数列} & \textbf{项的变化} & \textbf{收敛性} \\
\midrule\endhead
$1+(-1)^n/n$ & 在 $1$ 两侧交替，偏离幅度越来越小 & 收敛到 $1$ \\
$(-1)^n$ & 在 $-1,1$ 之间持续跳动 & 发散 \\
$n^2$ & 随项数增大而无限增大 & 发散，可记趋于 $+\infty$ \\
$0,1,0,1/2,0,1/3,\ldots$ & 偶数项趋于零，奇数项恒为零 & 收敛到 $0$ \\
\bottomrule\end{longtable}}

“有无振荡”不能决定收敛性。关键是振荡的幅度能否最终小于任意指定的误差。有界也不能单独保证收敛，$(-1)^n$ 就是反例。

\subsection{极限的精确定义}

先认识三个符号：$\forall$ 表示“对任意”，$\exists$ 表示“存在”，$\varepsilon$ 表示预先指定的正误差。它没有固定数值，既可以是 $0.1$，也可以是 $10^{-6}$，还可以更小。

\begin{definition}[数列极限]
对给定数列 $\{a_n\}$ 和实数 $A$，若

\begin{equation}\forall\varepsilon>0,\quad\exists N\in\mathbb N^+,\quad\forall n>N,\quad |a_n-A|<\varepsilon,\label{lec:02:definition}\end{equation}

则称 $a_n$ 收敛到 $A$，记作 $a_n\to A$ 或 $\lim_{n\to\infty}a_n=A$。
\end{definition}

这个定义包含三个依次发生的要求。

\begin{enumerate}
\item \textbf{先给误差。} $\varepsilon$ 可以任意小，不能只检查某个固定误差。
\item \textbf{再选起点。} $N$ 可以随 $\varepsilon$ 改变，误差越小，通常需要越大的 $N$。
\item \textbf{统一控制尾部。} 同一个 $N$ 必须保证之后每一项都满足误差要求。
\end{enumerate}

在几何上，$|a_n-A|<\varepsilon$ 等价于
\[
A-\varepsilon<a_n<A+\varepsilon.
\]

因此，收敛意味着对每一条以 $A$ 为中心的水平误差带，都可以找到一个位置，使其后的所有点都进入带内。仅有无穷多项进入带内仍不足够：还可能有无穷多项不断跑到带外。


\begin{figure}[htbp]
\centering
\begin{tikzpicture}[x=0.36cm,y=3.8cm,>=stealth,font=\small]
\fill[gray!12] (1,.9) rectangle (30,1.1);
\draw[->,gray!75] (0,0)--(31.3,0) node[right] {$n$};
\draw[->,gray!75] (0,-.04)--(0,1.8) node[above] {$a_n$};
\draw[black,dashed] (0,.9)--(30.6,.9) node[right] {$A-\varepsilon$};
\draw[black,dashed] (0,1.1)--(30.6,1.1) node[right] {$A+\varepsilon$};
\draw[black] (0,1)--(30.6,1) node[right] {$A=1$};
\draw[gray,dashed] (10,0)--(10,1.7);
\node[below] at (10,0) {$N=10$};
\foreach \n in {1,...,10} {\fill[black] (\n,{1+pow(-1,\n)/\n}) circle (1.5pt);}
\foreach \n in {11,...,30} {\fill[black] (\n,{1+pow(-1,\n)/\n}) circle (1.5pt);}
\node[align=center,font=\small] at (5,1.74) {有限项可以在带外};
\node[align=center,font=\small,text=black] at (21,1.53) {$n>N$ 的所有项都在带内\\$|a_n-A|<\varepsilon$};
\end{tikzpicture}
\caption{$a_n=1+(-1)^n/n$ 的误差带。这里 $\varepsilon=0.1$；收敛要求对任意正误差都能作同样的尾部控制。}
\label{lec:02:band}
\end{figure}


\textbf{量词易错点。} 如果把定义改成“对每个 $n$，存在一个 $\varepsilon>0$ 使 $|a_n-A|<\varepsilon$”，任意数列、任意 $A$ 都满足，已经失去极限的含义。也不能让 $N$ 随当前要检查的 $n$ 改变，否则尾部无法被统一控制。

\subsection*{极限不存在时应当怎样表述}

“不收敛到给定的 $A$”是上述定义的否定：
\[
\exists\varepsilon_0>0,\quad\forall N\in\mathbb N^+,\quad\exists n>N,\quad |a_n-A|\ge\varepsilon_0.
\]

要证明整个数列发散，需要对每个候选 $A\in\mathbb R$ 排除收敛，即
\[
\forall A\in\R,\quad \exists\varepsilon_A>0,\quad
\forall N\in\N,\quad \exists n>N:\ |a_n-A|\ge\varepsilon_A.
\]
误差可以依赖候选值 $A$。例如常数列 $a_n=1$ 不收敛于 $0$，却收敛于 $1$，所以“不收敛于给定值”和“发散”不能混用。后面的子列判别法通常更方便。这里“无限远处仍能找到坏项”比“前面有坏项”强得多。

\subsection*{无穷大与发散的区别}

记 $a_n\to+\infty$，是指任给 $M>0$，都存在 $N$，使所有 $n>N$ 有 $a_n>M$；$a_n\to-\infty$ 则要求最终 $a_n<-M$。无穷大不是实数极限值，按照本章的有限极限定义，两者都属于发散。

无界数列一定发散，但不一定趋于正无穷或负无穷。例如 $a_n=(-1)^n n$ 一条子列趋于正无穷，另一条趋于负无穷。因此不能用“无界”直接替代任何一个带符号的无穷趋势。

\section{怎样按定义证明极限}

定义证明通常分为两段：先在草稿中估计误差、确定可用的 $N$；再按定义的量词顺序写正式证明。估计时允许把误差适当放大，只要放大后的表达式容易控制。

首先建立 $n^{-p}\to0$（$p>0$）这一基础事实。任给 $\varepsilon>0$，取正整数 $N>\varepsilon^{-1/p}$，则 $n>N$ 时 $|n^{-p}-0|=n^{-p}<\varepsilon$。这里 $p>0$ 保证幂函数在正数上严格递增，也保证倒数随 $n$ 变小；若 $p=0$，数列恒等于 $1$，不能沿用零极限结论。

\begin{example}[直接构造整数起点]
用定义证明
\[
\lim_{n\to\infty}\frac{2n+1}{n+2}=2.
\]
\end{example}

\paragraph{思路分析} 从结论出发计算误差，而不是只说分子、分母同除以 $n$：
\[
\left|\frac{2n+1}{n+2}-2\right|=\frac3{n+2}<\frac3n.
\]

只要 $n>3/\varepsilon$，误差就小于 $\varepsilon$。这给出了选择 $N$ 的依据。

\paragraph{解答} 任给 $\varepsilon>0$，取
\[
N=\max\left\{1,\left\lceil\frac3\varepsilon\right\rceil\right\}.
\]

则当 $n>N$ 时，
\[
\left|\frac{2n+1}{n+2}-2\right|<\frac3n<\varepsilon.
\]

由定义得所求极限为 $2$。这里不需要找最小的 $N$，找到一个有效的整数即可。

记号 $\lceil x\rceil$ 表示大于或等于 $x$ 的最小整数，例如 $\lceil2.3\rceil=3$；它用于把误差估计给出的实数起点改成整数起点。

\begin{example}[几何数列趋于零]
对 $|q|<1$，证明 $q^n\to0$。
\end{example}

\paragraph{解答} 当 $q=0$ 时结论直接成立。下面设 $0<|q|<1$。给定 $\varepsilon>0$，先令 $\eta=\min\{\varepsilon,1/2\}$，于是 $0<\eta<1$。为了使 $|q|^n<\eta$，取对数可得
\[
n\ln|q|<\ln\eta
\quad\Longleftrightarrow\quad
n>\frac{\ln\eta}{\ln|q|}.
\]

因为 $\ln|q|<0$，除法后不等号要反向。取正整数
\[
N>\frac{\ln\eta}{\ln|q|},
\]

则对所有 $n>N$，有 $|q^n|<\eta\le\varepsilon$，从而 $q^n\to0$。

\textbf{方法说明。} $q$ 为负时数列可以交替变号，但其绝对值仍是一个衰减的几何数列；证明中应控制 $|q|^n$。

\begin{example}[用放缩控制阶乘型数列]
对任意固定实数 $c$，证明 $c^n/n!\to0$。
\end{example}

\paragraph{思路分析} 分母是阶乘，逐项把 $|c|$ 与分母因子比较，比直接解 $|c|^n/n!<\varepsilon$ 容易。后续的比值可以统一压在某个小于 $1$ 的数下面。

\paragraph{解答} $c=0$ 时数列各项均为零。设 $c\ne0$，选正整数 $m\ge2|c|$。当 $n>m$ 时，
\[
\frac{|c|^n}{n!}
=\frac{|c|^m}{m!}\prod_{k=m+1}^n\frac{|c|}{k}
\le\frac{|c|^m}{m!}\left(\frac12\right)^{n-m}.
\]

右侧是常数乘以趋于零的几何数列，所以由夹逼思想可得 $|c|^n/n!\to0$，继而 $c^n/n!\to0$。

若题目要求把 $N$ 明确写出，记 $C=|c|^m/m!$，再选择 $N>m$ 使 $C2^{m-N}<\varepsilon$，则对 $n>N$ 有 $|c|^n/n!<\varepsilon$。

\section{收敛数列的基本性质与证明}

\begin{theorem}[极限唯一]
若 $a_n\to A$ 且 $a_n\to B$，则 $A=B$。
\end{theorem}

\paragraph{解答} 假设 $A\ne B$。取 $\varepsilon=|A-B|/3$。由两次收敛定义，分别存在 $N_1,N_2$，使充分大的 $n$ 满足
\[
|a_n-A|<\varepsilon,\qquad |a_n-B|<\varepsilon.
\]

取 $n>\max\{N_1,N_2\}$，由三角不等式得
\[
|A-B|\le|A-a_n|+|a_n-B|<2\varepsilon=\frac23|A-B|,
\]

矛盾。故 $A=B$。

这里的关键是把两套误差要求同时实现，所以需要把两个起点取最大值。

\begin{theorem}[收敛数列必有界]
若 $a_n\to A$，则存在 $M>0$，使所有 $n$ 都满足 $|a_n|\le M$。
\end{theorem}

\paragraph{解答} 在极限定义中取 $\varepsilon=1$。存在 $N$，使 $n>N$ 时
\[
|a_n|\le|a_n-A|+|A|<1+|A|.
\]

尾部已经被统一控制；前 $N$ 项只有有限个，因此可以取
\[
M=\max\{|a_1|,\ldots,|a_N|,1+|A|\}.
\]

这样全部项都满足 $|a_n|\le M$。

\textbf{逆命题不成立。} $a_n=(-1)^n$ 始终有 $|a_n|=1$，但不收敛。有界性负责排除无限增大，不能排除持续振荡。

\begin{theorem}[保号性与保序性]
若 $a_n\to A>0$，则从某项起 $a_n>0$。更具体地，取 $\varepsilon=A/2$，便有
\[
|a_n-A|<A/2\quad\Longrightarrow\quad a_n>A/2>0.
\]
\end{theorem}

同理，若 $A<0$，则数列最终为负。若 $a_n\le b_n$ 最终成立，且两者分别收敛到 $A,B$，则 $A\le B$。

\paragraph{保序性的证明} 若反而有 $A>B$，取 $\varepsilon=(A-B)/3$。充分大的 $n$ 同时满足
\[
a_n>A-\varepsilon>B+\varepsilon>b_n,
\]
与最终的 $a_n\le b_n$ 矛盾。同样可知，如果 $A<B$，则最终 $a_n<b_n$：分别以 $(B-A)/3$ 为误差，两条误差带互不重叠。

\textbf{严格不等号可能消失。} 例如 $0<1/n$ 对每个 $n$ 成立，但两侧极限都为零。因此由 $a_n<b_n$ 一般只能推出 $A\le B$。

\subsection*{有限项与子列}

改变、删去或添加有限项，不会改变数列的收敛性；收敛时也不会改变极限。这是因为极限定义只控制某个位置之后的尾部。

若 $n_1<n_2<\cdots$ 是正整数序列，则 $\{a_{n_k}\}$ 称为子列。由于 $n_k\ge k$，$n_k\to\infty$。当 $a_n\to A$ 时，任何子列也收敛到 $A$。

这一步也可以直接写成定义证明：若 $n>N$ 保证 $|a_n-A|<\varepsilon$，则 $k>N$ 时 $n_k\ge k>N$，故 $|a_{n_k}-A|<\varepsilon$。因原数列本身也是一条子列，得到刻画：原数列收敛当且仅当所有子列均收敛且极限相同。实际判发散时，只需找出两条极限不同的子列。

\begin{example}[两个子列判发散]
证明 $a_n=(-1)^n$ 发散。
\end{example}
\paragraph{思路分析} 若整体收敛，奇偶两条子列必须有同一个极限。
\paragraph{解答} 对 $a_n=(-1)^n$，偶数项子列恒等于 $1$，奇数项子列恒等于 $-1$。如果原数列收敛，两条子列应具有同一个极限，与 $1\ne-1$ 矛盾。所以原数列发散。

子列判别还告诉我们：一条子列收敛，只能说明某一部分项的行为，不能据此判断整个数列收敛。

\begin{proposition}[有限条子列拼合]
若奇数项子列和偶数项子列都收敛到同一个 $A$，则原数列收敛到 $A$。
\end{proposition}
\begin{proof}
任给 $\varepsilon>0$，分别选 $K_1,K_2$，使 $k>K_1$ 时 $|a_{2k-1}-A|<\varepsilon$，$k>K_2$ 时 $|a_{2k}-A|<\varepsilon$。取 $N=2\max\{K_1,K_2\}+1$，则每个 $n>N$ 不论奇偶，都对应一个大于上述起点的 $k$，故 $|a_n-A|<\varepsilon$。
\end{proof}
同样可处理按余数分成的有限条子列。这里“有限条”保证各起点能取一个共同最大值。

\section{极限的四则运算}

设 $a_n\to A$、$b_n\to B$，则
\[
a_n\pm b_n\to A\pm B,\qquad a_nb_n\to AB.
\]

如果 $B\ne0$，则 $b_n$ 从某项起不为零，并且
\[
\frac{a_n}{b_n}\to\frac AB.
\]

\paragraph{加减法的定义证明} 任给 $\varepsilon>0$，选共同起点使两个误差均小于 $\varepsilon/2$，则
\[
|(a_n\pm b_n)-(A\pm B)|
\le|a_n-A|+|b_n-B|<\varepsilon.
\]
固定常数倍也同理：常数 $c\ne0$ 时，将输入误差选为 $\varepsilon/|c|$；$c=0$ 时每项均为零。

\textbf{为什么乘法成立。} 由恒等变形，
\[
a_nb_n-AB=(a_n-A)b_n+A(b_n-B).
\]

由于收敛数列 $\{b_n\}$ 有界，可以用常数控制 $|b_n|$。于是两个误差项都趋于零。这也说明“无穷小乘有界量仍是无穷小”。

\textbf{为什么除法要有 $B\ne0$。} 由保号性，充分大的 $n$ 满足 $|b_n|>|B|/2$，因此
\[
\left|\frac1{b_n}-\frac1B\right|
=\frac{|b_n-B|}{|b_nB|}
\le\frac{2|b_n-B|}{|B|^2}\to0.
\]

再应用乘法即可。若 $B=0$，这个估计无法成立，不能直接把分子、分母的极限相除。

\textbf{加法只能直接用于固定的有限项数。} 对每个固定 $k$，$1/n\to0$，但
\[
\underbrace{\frac1n+\cdots+\frac1n}_{n\text{项}}=1.
\]

这里项数随 $n$ 改变，不能逐项取极限后把结果都算成零。类似的和式需要整体估计、夹逼或平均值方法。

\begin{example}[抓住最高阶项]
求 $\lim_{n\to\infty}(3n^2+n-1)/(2n^2-5n+4)$。
\end{example}

\paragraph{解答} 分子、分母同除以 $n^2$：
\[
\frac{3n^2+n-1}{2n^2-5n+4}
=\frac{3+1/n-1/n^2}{2-5/n+4/n^2}\to\frac32.
\]

除法法则的适用依据是分母的新极限为 $2\ne0$。同类题中，应先比较最高次幂：分子次数较低时极限为零；相同时看首项系数之比；分子次数较高时要另外分析符号与无穷趋势。

\begin{example}[根式之差先有理化]
求 $\lim_{n\to\infty}(\sqrt{n^2+n}-n)$。
\end{example}

\paragraph{思路分析} 两项都趋于无穷，无法用“无穷减无穷”得到结果。相减会抵消主要部分，应先乘共轭式。

\paragraph{解答}
\[
\sqrt{n^2+n}-n
=\frac{n}{\sqrt{n^2+n}+n}
=\frac1{\sqrt{1+1/n}+1}\to\frac12.
\]

有理化把难以控制的差转成了一个可用四则运算求极限的商。

\subsection*{开方与极限交换的依据}

当 $a_n\ge0$ 且 $a_n\to A\ge0$ 时，确实有 $\sqrt{a_n}\to\sqrt A$，但最好知道为什么。若 $A>0$，有理化后得到
\[
|\sqrt{a_n}-\sqrt A|
=\frac{|a_n-A|}{\sqrt{a_n}+\sqrt A}
\le\frac{|a_n-A|}{\sqrt A}\to0.
\]

若 $A=0$，给定 $\varepsilon>0$，使 $a_n<\varepsilon^2$ 即可得到 $\sqrt{a_n}<\varepsilon$。因此后面解开方递推时，取极限有明确的依据。

对任意实数 $a_n\to A$，立方根也满足 $\sqrt[3]{a_n}\to\sqrt[3]A$。当 $A\ne0$，令 $u=\sqrt[3]{a_n},v=\sqrt[3]A$，利用
\[
|u-v|=\frac{|a_n-A|}{u^2+uv+v^2},\qquad
u^2+uv+v^2=\left(u+\frac v2\right)^2+\frac34v^2\ge\frac34v^2>0
\]

即可控制误差；$A=0$ 时使 $|a_n|<\varepsilon^3$ 即可。偶次根要求输入非负，奇次根可以处理任意实数，这个定义域区别不能忽略。

\section{夹逼定理}

\begin{theorem}[夹逼定理]
若从某项起有 $b_n\le a_n\le c_n$，且 $b_n\to A,c_n\to A$，则 $a_n\to A$。
\end{theorem}

\paragraph{解答} 任给 $\varepsilon>0$，两端收敛保证存在 $N_1,N_2$，使充分大的 $n$ 有
\[
A-\varepsilon<b_n,\qquad c_n<A+\varepsilon.
\]

同时选取起点，使原有夹逼不等式也成立，便得 $A-\varepsilon<a_n<A+\varepsilon$。由定义，$a_n\to A$。

常用零极限形式是：若 $|a_n|\le b_n$ 最终成立，且 $b_n\to0$，则 $a_n\to0$。

\begin{example}[变项数和的整体夹逼]
求
\[
c_n=\sum_{k=1}^n\frac{\sqrt{k(k+1)}}{n^2+k}
\]
的极限。
\end{example}
\paragraph{思路分析} 项数随 $n$ 增加，不能逐项取极限。每项根式夹在 $k$ 与 $k+1$ 之间；分母夹在 $n^2$ 与 $n^2+n$ 之间。把每项换为可求和的上下界，并检查分母放缩方向。
\paragraph{解答} 因为各量均为正，对 $1\le k\le n$，
\[
\frac{k}{n^2+n}
<\frac{\sqrt{k(k+1)}}{n^2+k}
<\frac{k+1}{n^2}.
\]
从 $k=1$ 到 $n$ 求和，有
\[
\frac12=\frac{n(n+1)}{2(n^2+n)}
<c_n<\frac{n(n+3)}{2n^2}=\frac12+\frac{3}{2n}.
\]
两端趋于 $1/2$，故 $c_n\to1/2$。
\paragraph{易错说明} 给正分式找上界时，分子可以放大，分母应缩小；若二者同时放大，不能据此判断商的方向。

\begin{example}[振荡因子乘衰减因子]
求 $\lim_{n\to\infty}\sin n/n$。
\end{example}

\paragraph{解答} 因 $|\sin n|\le1$，
\[
\left|\frac{\sin n}{n}\right|\le\frac1n\to0.
\]

故原极限为零。这里无需判断 $\sin n$ 是否有极限；有界性已经足以控制它。

\begin{example}[乘积取 $n$ 次方根]
求 $\lim_{n\to\infty}(1+2^n)^{1/n}$。
\end{example}

\paragraph{解答} 对 $n\ge1$，有 $2^n\le1+2^n\le2\cdot2^n$。取 $n$ 次方根：
\[
2\le(1+2^n)^{1/n}\le2\cdot2^{1/n}.
\]

由于 $2^{1/n}\to1$，两端极限均为 $2$，所以所求极限为 $2$。一个固定的正数因子经 $n$ 次方根后趋于 $1$，经常可以忽略其对最终极限的影响，但夹逼步骤应写明。

\textbf{补足一个基础事实。} 为什么 $2^{1/n}\to1$？给定 $\delta>0$，由二项式展开有 $(1+\delta)^n\ge1+n\delta$。当 $n>1/\delta$ 时，$(1+\delta)^n>2$，因而 $1<2^{1/n}<1+\delta$。这就是趋于 $1$ 的定义证明，不需要提前使用函数极限。

同一证明适用于任意固定 $c>1$：选 $n>(c-1)/\delta$，即可得到 $1<c^{1/n}<1+\delta$。$c=1$ 时各项为 $1$；$0<c<1$ 时对 $1/c>1$ 使用结论再取倒数，所以所有固定 $c>0$ 都有 $c^{1/n}\to1$。

\section{单调有界定理与递推数列}

\begin{theorem}[单调有界定理]
单调递增且有上界的数列收敛；单调递减且有下界的数列收敛。
\end{theorem}

先看上确界与最大值的区别。集合 $S=\{1-1/n:n\ge1\}$ 中每个数都小于 $1$，所以没有最大值；但它的上确界为 $1$。要验证 $A=\sup S$，既要证明 $A$ 是上界，也要证明任意比 $A$ 小的数都不再是上界。后一条等价于：对任意 $\varepsilon>0$，存在 $s\in S$ 满足 $s>A-\varepsilon$。实数的确界原理保证非空、有上界的实数集合具有上确界。

\textbf{证明递增情形。} 设 $A=\sup\{a_n:n\ge1\}$。任给 $\varepsilon>0$，$A-\varepsilon$ 不是该集合的上界，因此存在 $N$ 使 $a_N>A-\varepsilon$。由单调性，$n>N$ 时
\[
A-\varepsilon<a_N\le a_n\le A.
\]

故 $|a_n-A|<\varepsilon$，即 $a_n\to A$。递减情形用下确界同理。

这个证明依赖实数的完备性。它特别适合不知道极限具体值、却能判断项的大小关系的情形。

\begin{corollary}[单调数列的一条收敛子列]
若 $\{a_n\}$ 单调且有一条子列收敛，则原数列收敛，且与该子列有相同极限。
\end{corollary}
\begin{proof}
设数列递增，子列 $\{a_{n_k}\}$ 收敛，因而有界：$a_{n_k}\le M$。任给 $n$，可选 $k$ 使 $n_k\ge n$，于是 $a_n\le a_{n_k}\le M$。原数列递增有上界，故收敛。由收敛数列的子列性质，其极限等于已知子列极限。递减情形用下界同理。
\end{proof}

\subsection*{递推题的四个步骤}

{\small\begin{longtable}{>{\raggedright\arraybackslash}p{0.23\linewidth}>{\raggedright\arraybackslash}p{0.34\linewidth}>{\raggedright\arraybackslash}p{0.35\linewidth}}
\toprule
\textbf{步骤} & \textbf{要解决的问题} & \textbf{常见做法} \\
\midrule\endfirsthead
\toprule
\textbf{步骤} & \textbf{要解决的问题} & \textbf{常见做法} \\
\midrule\endhead
找界 & 后续项能否留在某个区间内 & 数学归纳法、不动点候选值 \\
证单调 & 相邻项的大小关系是否稳定 & 作差、作比、利用递推函数 \\
证收敛 & 能否合法引入极限 $L$ & 单调有界定理或收缩估计 \\
求极限 & $L$ 满足哪个方程、选哪一根 & 连续性、符号和已有区间 \\
\bottomrule\end{longtable}}

可以先解形式上的不动点方程来寻找界，但正式证明不能把“候选极限”当作已经存在的极限。

\begin{example}[开方递推的完整证明]
设 $a_1=1$，$a_{n+1}=\sqrt{2+a_n}$，求 $\lim a_n$。
\end{example}

\textbf{第一步：找不变区间。} 我们猜测极限为 $2$，于是尝试证明 $1\le a_n<2$。初值满足条件。若 $1\le a_n<2$，则
\[
1<\sqrt3\le a_{n+1}=\sqrt{2+a_n}<2.
\]

由归纳法，该区间限制对全部 $n$ 成立。

\textbf{第二步：证递增。} 因两项都为正，可以比较平方：
\[
a_{n+1}^2-a_n^2=2+a_n-a_n^2=(2-a_n)(a_n+1)>0.
\]

所以 $a_{n+1}>a_n$。

\textbf{第三步：证收敛。} 数列递增且有上界 $2$，由单调有界定理，存在实数极限 $L$，且 $1\le L\le2$。

\textbf{第四步：求极限。} 利用前面已证明的平方根极限性质，
\[
L=\sqrt{2+L},\qquad L^2-L-2=0.
\]

候选根为 $2,-1$，结合 $L\ge1$，得 $L=2$。

\textbf{进一步看收敛速度。} 有理化给出
\[
2-a_{n+1}=\frac{2-a_n}{2+\sqrt{2+a_n}}
\le\frac{2-a_n}{3}.
\]

由此 $0<2-a_n\le3^{-(n-1)}$，不仅说明收敛，还给出了误差的几何衰减估计。

\begin{example}[不单调时用收缩估计]
设 $a_1=1$，$a_{n+1}=1/(1+a_n)$，证明收敛并求极限。
\end{example}

\paragraph{思路分析} 项的大小交替变化，不能直接使用单调有界定理。先找一个能同时约束所有项的区间，再考察与不动点的距离。

\paragraph{解答} $a_2=1/2$。对 $1/2\le a_n\le1$，有 $1/2\le a_{n+1}\le2/3\le1$，故该区间对全部项成立。设
\[
\alpha=\frac{\sqrt5-1}{2},\qquad\alpha=\frac1{1+\alpha}.
\]

对 $n\ge1$，
\[
|a_{n+1}-\alpha|
=\frac{|a_n-\alpha|}{(1+a_n)(1+\alpha)}
\le\frac49|a_n-\alpha|,
\]

因为 $a_n,\alpha\ge1/2$。反复应用得到
\[
|a_n-\alpha|\le\left(\frac49\right)^{n-1}|a_1-\alpha|\to0.
\]

故 $a_n\to\alpha$。这个例子展示了有界振荡也可能收敛；需要控制的是振荡幅度。

\section{常见量级及重要数列}

\begin{example}[有理递推与因式分解]
设 $x_0=2$，$x_{n+1}=(2x_n^3+1)/(3x_n^2)$，证明收敛并求极限。
\end{example}
\paragraph{思路分析} 分母要求先证 $x_n>0$。形式不动点方程提示目标为 $1$；将 $x_{n+1}-1$ 因式分解，得到不变下界，再用相邻项差证明递减。
\paragraph{解答} 若 $x_n\ge1$，则
\[
x_{n+1}-1=\frac{2x_n^3-3x_n^2+1}{3x_n^2}
=\frac{(x_n-1)^2(2x_n+1)}{3x_n^2}\ge0.
\]
从 $x_0=2$ 出发，归纳得 $x_n\ge1$。再有
\[
x_{n+1}-x_n=\frac{1-x_n^3}{3x_n^2}\le0.
\]
数列递减且有下界 $1$，故收敛到某个 $L\ge1$。分母极限为 $3L^2>0$，可对递推关系取极限，得 $3L^3=2L^3+1$，即 $L^3=1$，所以 $L=1$。

以下极限常用于比较大小，范围条件不可省略：
\[
\frac1{n^p}\to0\quad(p>0),\qquad q^n\to0\quad(|q|<1),\qquad\frac{c^n}{n!}\to0\quad(c\in\mathbb R).
\]

\subsection*{多项式与指数的比较}

当 $a>1,p>0$ 时，$n^p/a^n\to0$。取正整数 $k>p$，将 $a$ 写成 $1+h$，其中 $h>0$。由二项式展开，$n\ge2k$ 时
\[
a^n=(1+h)^n\ge\binom nk h^k
\ge\frac{(n/2)^k}{k!}h^k.
\]

因此
\[
0\le\frac{n^p}{a^n}\le\frac{2^kk!}{h^k}n^{p-k}\to0.
\]

固定底数大于 $1$ 的指数增长最终超过任意固定正次幂，但这句话必须配合条件使用：底数为 $1$ 或小于 $1$ 时行为完全不同。

\subsection{高次方根消除多项式增长}

还有一个常用事实是 $n^{1/n}\to1$。因为 $n\ge1$，有 $n^{1/n}\ge1$。给定 $\delta>0$，对 $n\ge2$ 用二项式展开的二次项估计，得到
\[
(1+\delta)^n\ge\frac{n(n-1)}2\delta^2>n
\]

只需令 $n>1+2/\delta^2$。于是 $1\le n^{1/n}<1+\delta$，证明了极限为 $1$。该论证不需要提前使用对数的函数极限。

\subsection{自然常数的数列定义}
\[
u_n=\left(1+\frac1n\right)^n
\]

递增且有上界，因此收敛，通常以其极限定义 $e$。从二项式展开可以直接看到单调性：第 $k$ 项写成
\[
\binom nk\frac1{n^k}
=\frac1{k!}\prod_{j=0}^{k-1}\left(1-\frac jn\right),\qquad 0\le k\le n.
\]

对固定 $k$，这些因子随 $n$ 增大而不减，而下一次展开还多出一个正项，故 $u_{n+1}>u_n$。另一方面，
\[
u_n\le\sum_{k=0}^n\frac1{k!}<3,
\]

其中对 $k\ge2$ 用 $k!\ge2^{k-1}$ 控制尾部即可。于是数列收敛。

这里建立的是 $e$ 的数列来源；函数形式的 $(1+x)^{1/x}\to e$ 及其变形，在第3章讨论。

现在还可由前面的幂与指数比较推出 $(\ln n)^p/n\to0$（$p>0$），无需导数：令 $m=\lfloor\ln n\rfloor$，则 $\ee^m\le n<\ee^{m+1}$ 且 $m\to+\infty$，故充分大 $n$ 时
\[
0\le\frac{(\ln n)^p}{n}\le\frac{(m+1)^p}{\ee^m}\to0.
\]
末尾仍是固定正次幂与固定底数大于 $1$ 的指数之比。这里使用第1章的对数、指数单调性和反函数关系。

\begin{example}[根式平均的共同界]
求 $\displaystyle\lim_{n\to\infty}\frac{1^{1/n}+2^{1/n}+\cdots+n^{1/n}}n$。
\end{example}
\paragraph{思路分析} 求和中的第 $k$ 项也依赖 $n$，不能直接套用算术平均定理。用每项的共同下界 $1$ 和共同上界 $n^{1/n}$ 控制整体。
\paragraph{解答} 对 $1\le k\le n$，有 $1\le k^{1/n}\le n^{1/n}$，所以
\[
1\le\frac1n\sum_{k=1}^n k^{1/n}\le n^{1/n}\to1.
\]
由夹逼定理，所求极限为 $1$。

\section{平均值极限：固定前段与可控尾部}

\begin{theorem}[算术平均保持极限]
若 $a_n\to A$，则
\[
\frac{a_1+\cdots+a_n}{n}\to A.
\]
\end{theorem}

\paragraph{解答} 任给 $\varepsilon>0$，选 $N$，使 $k>N$ 时 $|a_k-A|<\varepsilon/2$。记有限前段误差之和为
\[
C=\sum_{k=1}^N|a_k-A|.
\]

当 $n>N$ 时，
\[
\left|\frac1n\sum_{k=1}^na_k-A\right|
\le\frac Cn+\frac1n\sum_{k=N+1}^n|a_k-A|
<\frac Cn+\frac\varepsilon2.
\]

再取 $n>2C/\varepsilon$，右侧小于 $\varepsilon$。故结论成立。

\textbf{思路总结。} 固定前段可能很大，但只含有限项，除以 $n$ 后贡献趋零；后段项数虽多，每一项却都已经接近 $A$。这是一种后续积分平均值与部分和估计中反复出现的拆分方法。

\begin{remark}[Cesàro 定理的逆命题]
算术平均保持极限也称 Cesàro 定理。逆命题不成立：$a_n=(-1)^n$ 发散，但其前 $n$ 项之和的绝对值不超过 $1$，所以算术平均的绝对值不超过 $1/n$，仍趋于零。平均会削弱某些振荡，不能由平均收敛反推原数列收敛。
\end{remark}

\begin{example}[调和平均对应的简单应用]
已知 $a_n\to A>0$ 且所有 $a_n>0$，求
\[
\lim_{n\to\infty}\frac n{1/a_1+\cdots+1/a_n}.
\]
\end{example}

\paragraph{解答} 因 $A\ne0$，有 $1/a_n\to1/A$。对倒数数列应用算术平均定理，得到
\[
\frac1n\sum_{k=1}^n\frac1{a_k}\to\frac1A.
\]

再次取倒数，所求极限为 $A$。这里先证明分母经归一化后的极限非零，才可以使用除法法则。

\section{柯西收敛准则：不先猜极限}

极限定义要求数列尾部接近某个固定实数 $A$。有时我们还不知道 $A$ 是多少，却能比较任意两项的距离，这时可以使用柯西收敛准则。

\begin{theorem}[柯西收敛准则]
实数列 $\{a_n\}$ 收敛，当且仅当
\[
\forall\varepsilon>0,\quad\exists N,\quad\forall m,n>N,\quad |a_m-a_n|<\varepsilon.
\]
\end{theorem}

\textbf{必要性。} 若 $a_n\to A$，使 $m,n>N$ 时都有 $|a_m-A|,|a_n-A|<\varepsilon/2$，则三角不等式给出 $|a_m-a_n|<\varepsilon$。

\textbf{充分性。} 先取误差 $1$。固定一个足够靠后的项，其他尾部项与它的距离小于 $1$，前段又只有有限项，所以整个数列有界。对每个 $n$，定义尾部项的下确界与上确界：
\[
l_n=\inf\{a_k:k\ge n\},\qquad u_n=\sup\{a_k:k\ge n\}.
\]

随着 $n$ 增大，尾部集合缩小，因此 $l_n$ 递增、$u_n$ 递减，两者都受原数列的界约束。由单调有界定理，它们分别收敛到实数 $l,u$。

任给 $\varepsilon>0$，柯西条件保证充分大的 $n$ 之后，任何两项的距离小于 $\varepsilon$。于是尾部的上下确界也满足 $0\le u_n-l_n\le\varepsilon$。具体地，从所有 $a_j\le a_k+\varepsilon$ 先对 $k$ 取下确界，再对 $j$ 取上确界即可得到这一估计。取极限有 $0\le u-l\le\varepsilon$，由 $\varepsilon$ 任意，得 $u=l$。最后利用 $l_n\le a_n\le u_n$ 和夹逼定理，得到 $a_n\to l$。

\textbf{必须控制任意两项。} $a_{n+1}-a_n\to0$ 仅控制相邻项，比柯西条件弱，不能保证收敛。

柯西条件的否定为
\[
\exists\varepsilon_0>0,\quad \forall N\in\N,\quad
\exists m,n>N:\ |a_m-a_n|\ge\varepsilon_0.
\]
实数列满足这个否定就发散，不需要逐个排除候选极限。

\begin{example}[正项和的柯西估计]
证明 $S_n=\sum_{k=1}^n1/k^2$ 收敛。
\end{example}
\paragraph{思路分析} 不需要计算和的极限值，只需把任意长的尾部控制为一个随起点趋零的上界。将 $1/k^2$ 放大为可消去相邻项的 $1/[k(k-1)]$。
\paragraph{解答} 若 $m>n\ge1$，则
\[
0<S_m-S_n=\sum_{k=n+1}^m\frac1{k^2}
\le\sum_{k=n+1}^m\left(\frac1{k-1}-\frac1k\right)
=\frac1n-\frac1m<\frac1n.
\]
任给 $\varepsilon>0$，取 $N>1/\varepsilon$。任意 $m,n>N$，若相等则差为零；若不等则交换下标使大者在前，由上述估计得 $|S_m-S_n|<\varepsilon$。由柯西准则，$S_n$ 收敛。这个证明同时表明，尾部界应与终点 $m$ 无关。

\begin{example}[调和数列为什么发散]
证明 $H_n=1+1/2+\cdots+1/n$ 发散。
\end{example}

\paragraph{思路分析} 相邻项差为 $1/(n+1)\to0$，但这不足以判断收敛。我们寻找相隔较远的两项，使其距离始终不会变得很小。

\paragraph{解答} 对任意 $n\ge1$，
\[
H_{2n}-H_n=\sum_{k=n+1}^{2n}\frac1k\ge n\cdot\frac1{2n}=\frac12.
\]

取固定误差 $\varepsilon_0=1/2$。不论起点 $N$ 多大，都能选 $n>N$，此时 $H_n,H_{2n}$ 两个尾部项的距离至少为 $1/2$，违反柯西条件。因此 $H_n$ 发散。由于 $H_n$ 递增，若它有上界就会收敛，所以实际上它无上界并趋于正无穷。

注意 $H_n$ 发散与平均值小节所用的 $H_n/n\to0$ 并不矛盾；两个表达式的增长行为不同。

\subsection{补充：收缩递推的一般保证}

前面的倒数递推已经给出直接控制误差的方法。若不动点还不知道，可以由柯西准则先建立它的存在。

\begin{theorem}[闭区间上的收缩递推]
设 $F$ 把闭区间 $[u,v]$ 映入自身，且存在 $0<q<1$，使对所有 $x,y\in[u,v]$ 有
\[
|F(x)-F(y)|\le q|x-y|.
\]
对任意初值 $x_0\in[u,v]$，递推 $x_{n+1}=F(x_n)$ 收敛到区间内唯一不动点 $\alpha$，并有
\[
|x_n-\alpha|\le\frac{q^n}{1-q}|x_1-x_0|.
\]
\end{theorem}
\begin{proof}
区间封闭性保证每一步都有定义且仍在 $[u,v]$ 内。递推估计给出
$|x_{n+1}-x_n|\le q^n|x_1-x_0|$。当 $m>n$ 时，用三角不等式求和：
\[
|x_m-x_n|\le |x_1-x_0|\sum_{k=n}^{m-1}q^k
\le\frac{q^n}{1-q}|x_1-x_0|.
\]
右侧与 $m$ 无关且趋零，所以由柯西准则数列收敛。保序性保证其极限 $\alpha\in[u,v]$。又
$|F(x_n)-F(\alpha)|\le q|x_n-\alpha|\to0$，故 $x_{n+1}\to F(\alpha)$；同时删去首项不改变极限，于是 $F(\alpha)=\alpha$。若还有不动点 $\beta$，则
$|\alpha-\beta|\le q|\alpha-\beta|$，因 $q<1$，只能有 $\alpha=\beta$。最后在上面的尾部估计中固定 $n$ 令 $m\to\infty$，得到误差界。
\end{proof}
两项条件缺一不可。$F(x)=x+1$ 的迭代没有区间封闭性；$F(x)=1-x$ 虽映 $[0,1]$ 入自身，却只能取 Lipschitz 常数 $1$，从 $0$ 出发持续振荡。后续第6章的中值定理可把一致导数界 $|F'|\le q<1$ 转为本定理的收缩不等式，本章直接使用代数估计。

\section{补充工具：Stolz 定理}

基础题通常先考虑极限定义、四则运算、夹逼和单调有界。部分和与增长量之比还可以使用下面的差分工具。

\begin{theorem}[Stolz 定理的无穷分母形式]
设 $b_n$ 严格递增，且 $b_n\to+\infty$。如果
\[
\frac{a_{n+1}-a_n}{b_{n+1}-b_n}\to L\in\mathbb R,
\]
则 $a_n/b_n\to L$。
\end{theorem}

\paragraph{解答} 任给 $\eta>0$，选 $N$ 使 $k\ge N$ 时
\[
L-\eta<\frac{a_{k+1}-a_k}{b_{k+1}-b_k}<L+\eta.
\]

由于每个分母增量为正，乘回并从 $k=N$ 累加到 $n-1$，得
\[
(L-\eta)(b_n-b_N)<a_n-a_N<(L+\eta)(b_n-b_N).
\]

充分大的 $n$ 满足 $b_n>0$。除以 $b_n$ 后，固定项 $a_N/b_n,b_N/b_n$ 都趋零，因此 $a_n/b_n$ 最终可以被控制在 $L-2\eta$ 与 $L+2\eta$ 之间。由于 $\eta$ 任意，结论成立。

这个证明也解释了为什么需要分母严格递增，以及为什么不能在分母趋于一个有限数时套用该形式。

\begin{example}[平方和除以立方]
求 $\lim_{n\to\infty}(1^2+\cdots+n^2)/n^3$。
\end{example}

\paragraph{解答} 令 $a_n=\sum_{k=1}^nk^2$、$b_n=n^3$。分母严格递增并趋于无穷，且
\[
\frac{a_{n+1}-a_n}{b_{n+1}-b_n}
=\frac{(n+1)^2}{3n^2+3n+1}\to\frac13.
\]

由 Stolz 定理，原极限为 $1/3$。也可以使用平方和公式直接算出；第11章还会把它解释为 $\int_0^1x^2\,dx$ 对应的黎曼和。

\begin{theorem}[Stolz 定理的零分母形式]
若 $a_n\to0$，$b_n>0$ 严格递减且 $b_n\to0$，以及
\[
\frac{a_{n+1}-a_n}{b_{n+1}-b_n}\to L\in\mathbb R,
\]
则同样有 $a_n/b_n\to L$。
\end{theorem}

\paragraph{证明} 由于 $b_k-b_{k+1}>0$，差分比的估计可等价写为对 $a_k-a_{k+1}$ 的估计。从充分大的 $n$ 开始累加差分不等式到 $m$，再令 $m\to\infty$：
\[
(L-\varepsilon)(b_n-b_{m+1})\le a_n-a_{m+1}\le(L+\varepsilon)(b_n-b_{m+1}).
\]

两数列都趋于零，故得到 $(L-\varepsilon)b_n\le a_n\le(L+\varepsilon)b_n$，最后除以正数 $b_n$ 即可。此形式的条件 $a_n\to0$ 不可删除：例如 $a_n=1+2^{-n},b_n=2^{-n}$，差分之比恒等于 $1$，但原来的商为 $2^n+1$，趋于正无穷。

差分之比没有极限时，只能说明无法通过上述定理得到结论，不能据此断定原商发散。

\paragraph{无穷极限的扩展} 两种 Stolz 形式中，差分之比也可以趋于 $+\infty$ 或 $-\infty$。以无穷分母形式的 $+\infty$ 为例，任给 $M>0$，充分大 $k$ 有差分比 $>2M$。累加得
\[
\frac{a_n}{b_n}>
2M\left(1-\frac{b_N}{b_n}\right)+\frac{a_N}{b_n}.
\]
右侧趋于 $2M$，因而充分大 $n$ 时原商 $>M$。负无穷情形改用上界；零分母形式则按前述尾部累加证明。每次使用仍须检查严格单调、分母趋势以及零分母形式的分子趋零条件。

\begin{example}[零分母形式]
求 $\displaystyle\lim_{n\to\infty}\frac{\sum_{k=n}^{\infty}2^{-k}}{2^{-n}}$，用 Stolz 零分母形式验证。
\end{example}
\paragraph{思路分析} 这里无限和仅为已知几何尾和，定义为有限尾和的极限；由几何求和公式等于 $2^{1-n}$。故分子确实趋于零，可以使用零分母形式。
\paragraph{解答} 记 $a_n=\sum_{k=n}^{\infty}2^{-k}$、$b_n=2^{-n}$。$b_n>0$ 严格递减趋零，$a_n\to0$，而
\[
\frac{a_{n+1}-a_n}{b_{n+1}-b_n}
=\frac{-2^{-n}}{-2^{-n-1}}=2.
\]
由零分母形式，原商趋于 $2$，与几何尾和直接求商的结果一致。

\section{方法选择与自检}

{\small\begin{longtable}{>{\raggedright\arraybackslash}p{0.23\linewidth}>{\raggedright\arraybackslash}p{0.34\linewidth}>{\raggedright\arraybackslash}p{0.35\linewidth}}
\toprule
\textbf{题面特征} & \textbf{先考虑的方法} & \textbf{解答中必须说明的环节} \\
\midrule\endfirsthead
\toprule
\textbf{题面特征} & \textbf{先考虑的方法} & \textbf{解答中必须说明的环节} \\
\midrule\endhead
明确要求“用定义证明” & 控制 $|a_n-A|$，构造 $N$ & 任意 $\varepsilon$、统一起点、全部尾部项 \\
多项式之比 & 除以最高次幂 & 分母新极限非零 \\
两个大根式相减 & 有理化 & 差中主项抵消后的余量 \\
有界振荡量乘衰减量 & 夹逼 & 对振荡量给统一界 \\
给定递推关系 & 不变区间、单调或收缩 & 先证收敛，再解不动点方程 \\
部分和除以项数 & 前段、尾段拆分 & 前段固定，尾段统一小 \\
部分和除以增长量 & Stolz 或后续的黎曼和 & 检查差分定理的分母条件 \\
\bottomrule\end{longtable}}

本章最重要的逻辑区分是：\textbf{猜测极限、证明极限存在、计算极限值}分别解决不同问题。一个候选值满足递推方程，并不自动说明原数列收敛到它。


\section{分层练习}
\subsection*{A组：基础理解与计算}
\begin{enumerate}[label=A\arabic*.]
\item 用 $\varepsilon$--$N$ 定义证明 $1/n^2\to0$，并写出一个可用的 $N(\varepsilon)$。
\item 判断 $(-1)^n/(n+1)$ 是否收敛，说明所用定理。
\item 求 $\lim_{n\to\infty}(\sqrt{n^2+4n}-n)$，解释为什么应先有理化。
\item 求 $\lim_{n\to\infty}(5n^3-n)/(2n^3+7)$，指出除法法则的适用条件。
\item “正数项数列有界就一定收敛”是否正确？若不正确，给出一个每项都严格为正的反例。
\end{enumerate}
\subsection*{B组：方法选择与逻辑辨析}
\begin{enumerate}[label=B\arabic*.]
\item 设 $a_1=2$，$a_{n+1}=\sqrt{6+a_n}$。按“找界、证单调、证收敛、求极限”的顺序解答。
\item 求 $\lim_{n\to\infty}\frac1n\sum_{k=1}^n1/k$，分别说明固定前段与尾段的作用。
\item 若 $a_{n+1}-a_n\to0$，能否推出 $a_n$ 收敛？提示：考虑 $a_n=\sqrt n$。
\item 设 $a_1=0$，$a_{n+1}=1-a_n$。形式上解 $L=1-L$ 会得到 $L=1/2$；为什么这不是原数列的极限？
\end{enumerate}
\subsection*{C组：补充工具}
\begin{enumerate}[label=C\arabic*.]
\item 用 Stolz 定理求 $\lim_{n\to\infty}(1+2+\cdots+n)/n^2$，逐项检查分母条件。
\end{enumerate}

\section{练习解析}
\paragraph{A1.} 任给 $\varepsilon>0$，取 $N=\max\{1,\lceil\varepsilon^{-1/2}\rceil\}$。当 $n>N$ 时，$|1/n^2-0|=1/n^2<\varepsilon$。起点不必最优，但必须统一控制全部后续项。
\paragraph{A2.} 因 $|(-1)^n/(n+1)|=1/(n+1)\to0$，由零极限形式的夹逼定理可得极限为零。交替变号不妨碍收敛。
\paragraph{A3.} 两项分别趋于无穷，不能使用差的运算法则。有理化得
\[
\sqrt{n^2+4n}-n=\frac{4n}{\sqrt{n^2+4n}+n}=\frac4{\sqrt{1+4/n}+1}\to2.
\]
\paragraph{A4.} 同除以 $n^3$ 后，分式变为 $(5-1/n^2)/(2+7/n^3)$，分母极限为非零数 $2$，故所求极限为 $5/2$。
\paragraph{A5.} 不正确。取 $a_n=1+(-1)^n/2$，全部项均在 $[1/2,3/2]$，但奇、偶数项子列分别收敛到 $1/2,3/2$，所以原数列发散。
\paragraph{B1.} 初值 $a_1=2$ 满足 $2\le a_1<3$；若 $2\le a_n<3$，则 $\sqrt8\le a_{n+1}<3$，故 $2\le a_{n+1}<3$。由归纳法，这一不变区间对全部项成立。因为两项均为正，且
\[
a_{n+1}^2-a_n^2=6+a_n-a_n^2=(3-a_n)(a_n+2)>0,
\]
数列递增。由单调有界定理收敛，设极限为 $L\in[2,3]$。平方根极限性质给出 $L=\sqrt{6+L}$，即 $L^2-L-6=0$。候选根为 $3,-2$，由 $L\ge2$ 选得 $L=3$。
\paragraph{B2.} 将 $a_k=1/k\to0$ 代入算术平均定理，得到极限为零。按定义写证明时，给定 $\varepsilon>0$，选 $N$ 使 $k>N$ 时 $1/k<\varepsilon/2$；记 $C=\sum_{k=1}^N1/k$，则
\[
0\le\frac1n\sum_{k=1}^n\frac1k\le\frac Cn+\frac\varepsilon2.
\]
再使 $n>2C/\varepsilon$ 即可。
\paragraph{B3.} 不能。$a_n=\sqrt n$ 无界，故不收敛；但
\[
a_{n+1}-a_n=\frac1{\sqrt{n+1}+\sqrt n}\to0.
\]
相邻项越来越接近，并不意味着所有足够靠后的项都接近同一个固定数。
\paragraph{B4.} 实际数列为 $0,1,0,1,\ldots$，两条子列极限不同，所以发散。$L=1-L$ 只是在假设已经收敛时才能得到的必要条件，不能代替收敛性的证明。
\paragraph{C1.} 令 $a_n=\sum_{k=1}^nk$、$b_n=n^2$。$b_n$ 严格递增且趋于正无穷，差分之比为 $(n+1)/(2n+1)\to1/2$，故所求极限为 $1/2$。


\textbf{离开本章之前}\quad 请尝试不看答案，写出开方递推例题的四个证明步骤，并解释 $(-1)^n$ 为什么有界却不收敛。如果只能记住“先设极限为 $L$”，还需要回到单调有界小节重新检查证明顺序。


\chapter{函数极限}
\label{lec:03:chapter}

数列极限考察整数指标越来越大时的变化；函数极限允许自变量沿任意实数点趋近。前置知识是实数的绝对值、函数的定义域、数列极限及其运算法则。本章要解决三个问题：怎样准确表达“无限接近”，怎样证明极限存在或不存在，以及怎样在不使用导数和泰勒公式的条件下计算基本极限。

\section{从趋近现象到严格定义}
\label{lec:03:definition}
例如当 $x\ne1$ 时，$(x^2-1)/(x-1)=x+1$。虽然原式在 $1$ 无定义，靠近 $1$ 的函数值却靠近 $2$。极限关注邻近点的函数值，而非中心点的取值。把中心值改为任意数，都不会改变这个极限。

\begin{definition}[有限点处的函数极限]
设 $f$ 在 $x_0$ 的某个去心邻域 $0<|x-x_0|<\delta_0$ 内有定义。若对任意 $\varepsilon>0$，存在 $\delta>0$，使得
\[
0<|x-x_0|<\delta\quad\Longrightarrow\quad |f(x)-A|<\varepsilon,
\]
则称 $\lim_{x\to x_0}f(x)=A$。可以总取 $\delta\le\delta_0$。
\end{definition}
这里 $\varepsilon$ 是别人提出的函数值误差要求，$\delta$ 是我们提供的自变量控制范围。先给 $\varepsilon$，再选 $\delta$，最后对范围内的\emph{所有} $x$ 成立；$\delta$ 可依赖 $\varepsilon$，不能随被检验的 $x$ 临时改变。对于一般定义域 $D$，要求 $x_0$ 是 $D$ 的聚点，并把量词限制为 $x\in D$；本章双侧极限的基本讨论采用完整去心邻域。

\begin{figure}[htbp]
\centering
\begin{tikzpicture}[scale=0.85]
\draw[->] (-0.3,0)--(6,0) node[right] {$x$};
\draw[->] (0,-0.3)--(0,4.3) node[above] {$y$};
\draw[dashed] (0,1.5) node[left] {$A-\varepsilon$}--(5.7,1.5);
\draw[dashed] (0,2.5) node[left] {$A+\varepsilon$}--(5.7,2.5);
\draw[dotted] (0,2) node[left] {$A$}--(5.7,2);
\draw[dashed] (2,0) node[below left] {$x_0-\delta$}--(2,3.7);
\draw[dashed] (4,0) node[below right] {$x_0+\delta$}--(4,3.7);
\draw[dotted] (3,0) node[below] {$x_0$}--(3,3.7);
\draw[thick,domain=1.3:4.7,samples=40] plot (\x,{2+0.32*(\x-3)});
\draw[fill=white] (3,2) circle (0.065);
\node[anchor=west] at (4.2,3.4) {$y=f(x)$};
\end{tikzpicture}
\caption{竖直带内的去心部分必须映入给定的水平误差带。中心点的值不参与检验。}
\end{figure}

\paragraph{定义的否定}
“$f$ 不趋于指定的 $A$”准确写成
\[
\exists\varepsilon_0>0,\ \forall\delta>0,\ \exists x:
\quad0<|x-x_0|<\delta,\quad |f(x)-A|\ge\varepsilon_0.
\]
要说没有任何有限极限，还须对每个 $A\in\R$ 否定上述极限定义。证明 $f$ 不趋于 $0$ 并不排除它趋于 $1$。

\begin{example}[先固定范围，再控制误差]
按定义证明 $\lim_{x\to2}x^2=4$。
\end{example}
\paragraph{思路分析}误差是 $|x-2||x+2|$。第一因子由 $\delta$ 控制，第二因子也变化，先把它限制在一个固定上界内。
\paragraph{证明}任给 $\varepsilon>0$，取 $\delta=\min\{1,\varepsilon/5\}$。若 $0<|x-2|<\delta$，则
\[
|x+2|\le |x-2|+4<5,\qquad
|x^2-4|<5|x-2|<5\delta\le\varepsilon.
\]
因此定义成立。这里 $1$ 负责保证估计可用，$\varepsilon/5$ 负责达到最终精度，两项要求必须同时满足。

\begin{example}[分母远离零]
按定义证明 $\lim_{x\to1}1/(x+1)=1/2$。
\end{example}
\paragraph{思路分析}通分后出现 $|x+1|$，需要给分母一个\emph{正的下界}，而非上界。
\paragraph{证明}任给 $\varepsilon>0$，取 $\delta=\min\{1,2\varepsilon\}$。当 $0<|x-1|<\delta$ 时，$0<x<2$，所以 $x+1>1$，从而
\[
\left|\frac1{x+1}-\frac12\right|
=\frac{|x-1|}{2|x+1|}<\frac{|x-1|}{2}<\varepsilon.
\]
若误把 $x+1<3$ 用来放大这个分式，就会把不等号方向弄反。

\subsection{左右极限和无穷远处的极限}
右极限 $\lim_{x\to x_0^+}f(x)=A$ 把条件改为 $0<x-x_0<\delta$；左极限把条件改为 $0<x_0-x<\delta$。在相应一侧有定义是讨论的前提。

\begin{theorem}[双侧与单侧的关系]
函数在两侧去心邻域有定义时，$\lim_{x\to x_0}f(x)=A$ 当且仅当两个单侧极限都为 $A$。
\end{theorem}
\begin{proof}
双侧定义直接限制到任一侧，即得必要性。反过来，给定 $\varepsilon>0$，分别由左右定义取得 $\delta_-,\delta_+$，取 $\delta=\min\{\delta_-,\delta_+\}$。满足 $0<|x-x_0|<\delta$ 的点必在某一侧，该侧的误差估计适用。
\end{proof}

\begin{example}[先拆左右两侧]
判断 $\lim_{x\to0}|x|/x$ 是否存在。
\end{example}
\paragraph{解答}当 $x>0$ 时值为 $1$，当 $x<0$ 时值为 $-1$。右极限是 $1$，左极限是 $-1$，所以双侧极限不存在。仅从正数代入得到 $1$，只检验了右侧。

\begin{definition}[无穷远处的有限极限]
设 $f$ 在充分大的正数处有定义。$\lim_{x\to+\infty}f(x)=A$ 指
\[
\forall\varepsilon>0,\ \exists M>0,\ \forall x>M:
\quad |f(x)-A|<\varepsilon.
\]
$x\to-\infty$ 时把 $x>M$ 改为 $x<-M$；$|x|\to\infty$ 时改为 $|x|>M$，并要求两端均有定义。
\end{definition}
两端趋于同一个 $A$，才能写 $|x|\to\infty$ 的有限极限。例如 $x/\sqrt{x^2+1}$ 在正无穷处趋于 $1$，在负无穷处趋于 $-1$，没有两端共同的有限极限。原因是 $\sqrt{x^2}=|x|$，不能在负数处写成 $x$。

\paragraph{有理函数在无穷远处}
设 $P,Q$ 分别为 $k,m$ 次多项式，首项系数分别为非零的 $a,b$。提出最高次幂得
\[
\frac{P(x)}{Q(x)}
=x^{k-m}\frac{a+O(1/x)}{b+O(1/x)}.
\]
这里也可暂把 $O(1/x)$ 理解为有限个趋于零的低次项之和，正式记号见\ref{lec:03:orders}节。于是 $k<m$ 时趋于 $0$，$k=m$ 时趋于 $a/b$。$k>m$ 时绝对值趋于无穷；确定符号还要检查 $a/b$ 和 $x^{k-m}$。正无穷处符号由 $a/b$ 决定，负无穷处还受差次数的奇偶性影响。

\begin{example}[最高次项和两端符号]
分别考察 $\displaystyle\frac{2x^3-x}{x^2+1}$ 在 $x\to+\infty$ 与 $x\to-\infty$ 时的变化。
\end{example}
\paragraph{解答}表达式为 $x(2-1/x^2)/(1+1/x^2)$，后一个因子趋于 $2$，充分远处可限制在 $(1,3)$ 内。正数端表达式大于 $x$，趋于 $+\infty$；负数端表达式小于 $x$，趋于 $-\infty$。两端共同成立的是绝对值趋于无穷，不能把共同结论写成趋于 $+\infty$。

\begin{example}[在无穷远处选择阈值]
按定义证明 $\lim_{x\to+\infty}(2x+1)/(x+3)=2$。
\end{example}
\paragraph{解答}对 $x>0$，误差为 $5/(x+3)<5/x$。给定 $\varepsilon>0$，取 $M=5/\varepsilon$，则 $x>M$ 蕴含误差小于 $\varepsilon$。有限点处寻找小半径，正无穷处寻找大阈值，两者都是把误差放大为可控制的简单量。

\section{极限性质、归结原则与存在判据}
\label{lec:03:criteria}
\begin{proposition}[唯一性、局部有界性与保号性]
有限极限唯一。若 $f(x)\to A$，则 $f$ 在充分小的去心邻域有界；若 $A>0$，则该邻域内 $f(x)>A/2>0$，若 $A<0$，则该邻域内 $f(x)<A/2<0$。
\end{proposition}
\begin{proof}
若同时趋于不同的 $A,B$，取 $\varepsilon=|A-B|/3$，在两个定义所给邻域的交集中有
$|A-B|\le|A-f(x)|+|f(x)-B|<2|A-B|/3$，矛盾。取 $\varepsilon=1$，得 $|f(x)|<|A|+1$，证明有界性。取 $\varepsilon=|A|/2$ 得保号性。
\end{proof}
这是局部结论，不约束远离 $x_0$ 的函数值。邻域内 $f(x)>0$ 只能推出极限 $A\ge0$，不能推出 $A>0$；例如 $x^2>0$（$x\ne0$），但极限为 $0$。

\begin{theorem}[Heine 定理，归结原则]
在去心邻域内，$\lim_{x\to x_0}f(x)=A$ 当且仅当对每个满足 $x_n\ne x_0$、$x_n\to x_0$ 的数列（各项在定义域内），都有 $f(x_n)\to A$。
\end{theorem}
\begin{proof}
若函数极限成立，对任给 $\varepsilon>0$ 选相应 $\delta$。由 $x_n\to x_0$，充分大的 $n$ 满足 $0<|x_n-x_0|<\delta$，故 $|f(x_n)-A|<\varepsilon$。

反过来，若函数不趋于 $A$，由定义的否定有固定 $\varepsilon_0>0$，使任意小去心邻域内都存在误差至少为 $\varepsilon_0$ 的点。在半径 $\min\{\delta_0/2,1/n\}$ 的去心邻域里选一个这样的点 $x_n$。于是 $x_n\to x_0$，而 $|f(x_n)-A|\ge\varepsilon_0$ 对每项成立，违背假设。
\end{proof}
单侧和无穷远处也有对应版本，例如 $x\to+\infty$ 时要求任意 $x_n\to+\infty$。验证某一个数列的函数值收敛不足以证明函数极限；但找到两个极限不同的函数值数列，就足以否定函数极限。

\begin{example}[把振荡相位固定]
证明 $\sin(1/x)$ 在 $x\to0$ 时无极限。
\end{example}
\paragraph{思路分析}直接观察“振荡”还不是证明。让 $1/x$ 分别落在正弦取得 $1$ 与 $-1$ 的位置，就能获得两个稳定而不同的函数值数列。
\paragraph{证明}取
\[
x_n=\frac1{2n\pi+\pi/2},\qquad
y_n=\frac1{2n\pi+3\pi/2}.
\]
两列均趋于 $0$ 且避开 $0$，但 $\sin(1/x_n)=1$，$\sin(1/y_n)=-1$。由 Heine 定理，甚至右极限也不存在。

\begin{theorem}[四则运算与夹逼]
若 $f\to A$、$g\to B$，则 $f\pm g\to A\pm B$、$fg\to AB$；若 $B\ne0$，则 $f/g\to A/B$。若邻域内 $p\le h\le q$，且 $p,q\to A$，则 $h\to A$。
\end{theorem}
\begin{proof}
任取 $x_n\to x_0$、$x_n\ne x_0$，由 Heine 定理得到 $f(x_n)\to A$、$g(x_n)\to B$。数列的四则法则再给出对应结果；$B\ne0$ 时保号性保证分母最终非零。由 Heine 定理返回函数极限。夹逼定理也可以这样归结；直接证明时，对给定 $\varepsilon$ 同时保证 $A-\varepsilon<p$ 与 $q<A+\varepsilon$，便有 $|h-A|<\varepsilon$。
\end{proof}
分母趋于零时，除法法则失效，$0/0$ 只是提醒我们继续变形的未定式。$x/x$、$x^2/x$、$x/x^2$ 在 $0$ 附近分别表现为 $1$、$0$ 和无有限极限，不能给 $0/0$ 指定统一结果。

\begin{example}[因式分解与有理化]
求 $\lim_{x\to1}(x^2-1)/(x^3-1)$ 及 $\lim_{x\to1}(\sqrt{x}-1)/(\sqrt[3]{x}-1)$。
\end{example}
\paragraph{思路分析}分子分母同时为零，先找相同的消失因子。根式差可通过平方差、立方差制造 $x-1$。
\paragraph{解答}在 $x\ne1$ 的邻域内，第一式为
\[
\frac{x+1}{x^2+x+1}\longrightarrow\frac23.
\]
令 $a=\sqrt{x}$、$b=\sqrt[3]{x}$，由 $a-1=(x-1)/(a+1)$、$b-1=(x-1)/(b^2+b+1)$，第二式为
\[
\frac{b^2+b+1}{a+1}\longrightarrow\frac32.
\]
这里只在去心邻域约掉非零的 $x-1$，无须补定义中心点。根式的趋近可用恒等式证明，例如 $|\sqrt{x}-1|=|x-1|/(\sqrt{x}+1)\le|x-1|$。

\begin{example}[夹逼含取整的表达式]
求 $\lim_{x\to0^+}x\lfloor1/x\rfloor$。
\end{example}
\paragraph{解答}由 $1/x-1<\lfloor1/x\rfloor\le1/x$，因为 $x>0$，相乘不改变不等号方向，得到
\[
1-x<x\lfloor1/x\rfloor\le1.
\]
两端趋于 $1$，故极限为 $1$。不要先要求取整函数本身有极限；这里起作用的是整体误差小于 $x$。

\paragraph{单调有界与柯西判据}
函数若在 $(a,b)$ 单调递增且有上界，则 $x\to b^-$ 时存在有限极限。证明可取 $L=\sup\{f(x):a<x<b\}$：给定 $\varepsilon>0$，有 $c<b$ 使 $L-\varepsilon<f(c)$，当 $c<x<b$ 时 $L-\varepsilon<f(x)\le L$。其他单侧情形用下确界或改变方向处理。此法适合证明存在性，通常不能直接给出数值。

函数极限的柯西判据为：有限极限存在当且仅当
\[
\forall\varepsilon>0,\ \exists\delta>0:\quad
0<|x-x_0|,|y-x_0|<\delta\ \Longrightarrow\ |f(x)-f(y)|<\varepsilon.
\]
必要性由分别把 $f(x),f(y)$ 控制在极限的 $\varepsilon/2$ 范围内得到。充分性：先取任意一列 $z_n\to x_0$，$z_n\ne x_0$，由判据得到 $f(z_n)$ 是柯西数列，由实数完备性它趋于某 $A$；再给定 $\varepsilon$，对充分靠近的 $x$ 和充分大的 $n$，有 $|f(x)-f(z_n)|<\varepsilon/2$、$|f(z_n)-A|<\varepsilon/2$，即得定义。此判据无需预先猜出 $A$。

\paragraph{复合极限的条件}
若 $g(x)\to b$，且 $f(u)\to A$（$u\to b$），并且 $g(x)$ 在充分靠近时始终不等于 $b$，则 $f(g(x))\to A$。证明按外层误差先选 $u$ 的半径，再由内层极限保证 $g(x)$ 落入该去心范围。若 $g$ 可能等于 $b$，需额外知道 $f(b)=A$。反例：$f(u)=1$（$u\ne0$）、$f(0)=0$，$g(x)\equiv0$；外层去心极限为 $1$，复合函数却恒为 $0$。第4章连续性的定义正是消除这一障碍的条件。

\section{无穷小、无穷大与阶的比较}
\label{lec:03:orders}
\begin{definition}[无穷小与无穷大]
在指定的趋近过程中，$\alpha(x)\to0$ 称为无穷小。$f(x)\to+\infty$ 表示对任意 $K>0$，充分靠近时都有 $f(x)>K$；$f(x)\to-\infty$ 改为 $f(x)<-K$。$|f(x)|\to+\infty$ 只要求 $|f(x)|>K$，不保证符号固定。
\end{definition}
无穷小是一个变化的量，不能当成“很小的固定常数”。同一个 $x$ 在 $x\to0$ 时是无穷小，在 $x\to+\infty$ 时是无穷大。有限极限 $f\to A$ 等价于 $f=A+\alpha$、$\alpha\to0$。

\begin{proposition}[有界量乘无穷小]
若 $|h(x)|\le C$ 在充分靠近时成立，且 $\alpha(x)\to0$，则 $h(x)\alpha(x)\to0$。
\end{proposition}
\begin{proof}
取 $C_1=\max\{C,1\}$。给定 $\varepsilon>0$，在足够小的邻域内同时满足有界条件与 $|\alpha|<\varepsilon/C_1$，于是 $|h\alpha|\le C_1|\alpha|<\varepsilon$。
\end{proof}

\begin{example}[振荡项的整体误差]
求 $\lim_{x\to0}\bigl(3+x\sin(1/x)+x^2\cos(1/x)\bigr)$。
\end{example}
\paragraph{解答}两种振荡函数均无极限，但
$|x\sin(1/x)|\le|x|$、$|x^2\cos(1/x)|\le x^2$。故整体与 $3$ 的差的绝对值不超过 $|x|+x^2\to0$，极限为 $3$。不能对无极限的振荡因子直接套乘积法则。

若 $\alpha$ 最终非零，则 $\alpha\to0$ 当且仅当 $|1/\alpha|\to+\infty$，证明只需把误差阈值 $\varepsilon$ 与大小阈值 $1/\varepsilon$ 对应。由 $1/x$ 的左右符号可见，这不保证 $1/\alpha$ 趋于某个固定符号的无穷大。

有限个绝对值趋于无穷的量之积仍有绝对值趋于无穷；若要确定乘积是正无穷还是负无穷，还须检查符号。两个无穷大的和可能相消，例如 $x+(-x)=0$。无穷大与有界量的乘积不能统一判断：当 $x\to+\infty$ 时，$x\cdot1\to+\infty$，$x\cdot(1/x)=1$，$x\cdot(1/x^2)\to0$。这正是 $0\cdot\infty$ 被称为未定式的原因。相反，绝对值趋于无穷的量加上有界量仍绝对值趋于无穷，因为 $|f+h|\ge|f|-|h|$。

“局部无界”要求每个去心邻域内都\emph{能找到}任意大的绝对值；无穷大要求充分靠近后\emph{全部}绝对值任意大。$\sin(1/x)/x$ 在 $0$ 的每个邻域内无界，但沿 $x_n=1/(n\pi)$ 恒为零，因此绝对值也不趋于无穷。$\sin(1/x)$ 则有界而无极限。三种概念应分开。

\begin{definition}[阶、等价与大、小 $O$]
设 $\alpha,\beta$ 为无穷小，且 $\beta$ 最终非零。若 $\alpha/\beta\to0$，记 $\alpha=o(\beta)$；若比值趋于非零常数 $c$，称二者同阶；若 $c=1$，记 $\alpha\sim\beta$，称等价无穷小。若 $|\alpha/\beta|$ 最终有界，记 $\alpha=O(\beta)$，此定义也可用于一般函数。
\end{definition}
例如 $x^3=o(x^2)$，$2x^2$ 与 $x^2$ 同阶但不等价；$x\sin(1/x)=O(x)$，却不是 $o(x)$。记号 $o(\beta)$ 代表满足条件的一类余项；同一式里出现的两个 $o(\beta)$ 不一定是同一个函数。由比值运算可证明：有限个 $o(\beta)$ 的和仍为 $o(\beta)$，$o(x^p)o(x^q)=o(x^{p+q})$。

\begin{theorem}[乘除因子的等价替换]
若 $\alpha\sim\widetilde\alpha$、$\beta\sim\widetilde\beta$，相关分母最终非零，则
\[
\frac{\alpha/\beta}{\widetilde\alpha/\widetilde\beta}
=\frac\alpha{\widetilde\alpha}\frac{\widetilde\beta}\beta\longrightarrow1.
\]
因此当替换后的商有有限或确定的无穷极限时，原商具有同一极限。多个乘除因子的情形相同。
\end{theorem}
差式的有效规则是\emph{误差相对于所保留尺度可以忽略}。若要计算 $(a-b)/d$，替换为 $(\widetilde a-\widetilde b)/d$，充分条件为
$a-\widetilde a=o(d)$、$b-\widetilde b=o(d)$。只知 $a\sim\widetilde a$、$b\sim\widetilde b$，没有这一保证。

\section{两个重要极限及其推论}
\label{lec:03:important}
\begin{theorem}[第一重要极限]
采用弧度制时，$\displaystyle\lim_{x\to0}\frac{\sin x}{x}=1$。
\end{theorem}
\begin{proof}
先设 $0<x<\pi/2$。单位圆上 $A=(1,0)$、$B=(\cos x,\sin x)$，过 $A$ 的切线与射线 $OB$ 相交于 $C=(1,\tan x)$。三角形 $OAB$、扇形 $OAB$、三角形 $OAC$ 的面积依次增大，面积分别为 $\sin x/2$、$x/2$、$\tan x/2$，所以
\[
\sin x<x<\tan x,\qquad \cos x<\frac{\sin x}{x}<1.
\]
又 $0\le1-\cos x=2\sin^2(x/2)\le x^2/2$，故 $\cos x\to1$，夹逼得到右极限为 $1$。对 $x<0$ 令 $t=-x>0$，$\sin x/x=\sin t/t$，左极限也为 $1$。
\end{proof}
上述内接三角形必须是 $OAB$；若使用垂足 $D$ 构成的 $OBD$，其面积是 $\sin x\cos x/2$，不能写成 $\sin x/2$。面积与角度的公式也是使用弧度制的原因。

若 $u(x)\to0$ 且 $u$ 在去心邻域最终非零，复合极限给出 $\sin u(x)/u(x)\to1$。若 $u$ 可能为零，则原商在这些点无定义；可以把函数 $\sin t/t$ 在 $t=0$ 补定义为 $1$，再讨论这个补定义函数的复合极限。第二重要极限的复合形式也须检查相应的非零条件。

\begin{figure}[htbp]
\centering
\begin{tikzpicture}[scale=2.2]
\draw[->] (0,0)--(1.5,0);
\draw[->] (0,0)--(0,1.15);
\draw (1,0) arc (0:60:1);
\draw[thick] (0,0)--(1,0)--(1,0.839)--cycle;
\draw (0,0)--(0.766,0.643)--(1,0);
\draw (0.22,0) arc (0:40:0.22);
\node at (0.3,0.1) {$x$};
\node[below left] at (0,0) {$O$};
\node[below] at (1,0) {$A$};
\node[left] at (0.63,0.67) {$B$};
\node[right] at (1,0.839) {$C$};
\end{tikzpicture}
\caption{内接三角形、单位圆扇形和外切三角形给出 $\sin x<x<\tan x$。}
\end{figure}

\begin{theorem}[第二重要极限]
以第2章的数列极限 $\ee=\lim_{n\to\infty}(1+1/n)^n$ 为基础，有
\[
\lim_{s\to+\infty}\left(1+\frac1s\right)^s=\ee,
\qquad \lim_{t\to0}(1+t)^{1/t}=\ee.
\]
第二式在 $t\ne0$、$1+t>0$ 的邻域内讨论。
\end{theorem}
\begin{proof}
对 $s\ge1$ 令 $n=\lfloor s\rfloor$。利用正底数幂的单调性，
\[
\left(1+\frac1{n+1}\right)^n
\le\left(1+\frac1s\right)^s
\le\left(1+\frac1n\right)^{n+1}.
\]
左端为 $(1+1/(n+1))^{n+1}/(1+1/(n+1))$，右端为 $(1+1/n)^n(1+1/n)$，两端均趋于 $\ee$，故中间式趋于 $\ee$。取 $s=1/t$ 得 $t\to0^+$ 的结论。

若 $t\to0^-$，令 $s=-1/t\to+\infty$。对 $s>1$，
\[
(1+t)^{1/t}=\left(1-\frac1s\right)^{-s}
=\left(1+\frac1{s-1}\right)^{s-1}
 \left(1+\frac1{s-1}\right)\longrightarrow\ee.
\]
左右极限相同，双侧结论成立。
\end{proof}

\subsection{常用等价量怎样得到}
在对极限取对数或作反三角换元之前，先独立证明所需的趋近性质。这里只使用第1章的单调性、指数运算法则与第2章的伯努利不等式，不借用导数或尚未建立的连续性结论。

\begin{lemma}[指数、对数与反三角换元的趋近性质]
\label{lec:03:elementary-approach}
当 $h\to0$ 时，$\ee^h\to1$；因此对任意实数 $c$，当 $x\to c$ 时有 $\ee^x\to\ee^c$。
对任意 $y_0>0$，当 $y\to y_0$ 时有 $\ln y\to\ln y_0$。
当 $x\to0$ 时，$\arcsin x\to0$、$\arctan x\to0$。
\end{lemma}
\begin{proof}
令 $r_n=\ee^{1/n}>1$。由伯努利不等式，
\[
\ee=r_n^n=[1+(r_n-1)]^n\ge1+n(r_n-1),
\qquad 1\le r_n\le1+\frac{\ee-1}{n}.
\]
故 $r_n\to1$，其倒数也趋于 $1$。给定任意误差，先选充分大的 $n$，再限制 $|h|<1/n$；指数单调性给出
\[
\ee^{-1/n}\le\ee^h\le\ee^{1/n}.
\]
两端可以同时任意接近 $1$，这就证明了 $\ee^h\to1$。由 $\ee^x=\ee^c\ee^{x-c}$ 得到任意点的结论。

对数用单调反解处理。给定 $\varepsilon>0$，令
\[
y_-=\ee^{\ln y_0-\varepsilon}<y_0,
\qquad y_+=\ee^{\ln y_0+\varepsilon}>y_0,
\qquad \delta=\min\{y_0-y_-,y_+-y_0\}>0.
\]
若 $|y-y_0|<\delta$，则 $y_-<y<y_+$。取单调递增的对数，得
$\ln y_0-\varepsilon<\ln y<\ln y_0+\varepsilon$，证明完毕。

最后，先取 $0<\varepsilon<\pi/2$。当 $|x|<\sin\varepsilon$ 时，由正弦在主值区间上的严格递增性，有 $|\arcsin x|<\varepsilon$；当 $|x|<\tan\varepsilon$ 时，同理有 $|\arctan x|<\varepsilon$。任意较大的误差也由这样的较小误差控制，因此两个原点极限成立。
\end{proof}

现在可以合法地使用这些趋近性质。第4章会把“函数值随输入趋近”的性质统一称为连续性。由上面的三角极限及恒等式得
\[
\sin x\sim x,\qquad \tan x=\frac{\sin x}{\cos x}\sim x,
\qquad 1-\cos x=2\sin^2(x/2)\sim\frac{x^2}{2}.
\]
设 $t=\arcsin x\to0$，则 $x=\sin t$，$\arcsin x/x=t/\sin t\to1$；设 $t=\arctan x\to0$，同理 $t/\tan t\to1$。这两种换元使用反函数在原点的趋近性质。

对第二重要极限取自然对数，得到
\[
\frac{\ln(1+x)}x=\ln\bigl((1+x)^{1/x}\bigr)\longrightarrow1.
\]
令 $u=\ee^x-1\to0$，当 $x\ne0$ 时 $u\ne0$，由 $x=\ln(1+u)$ 得
\[
\frac{\ee^x-1}x=\frac u{\ln(1+u)}\longrightarrow1.
\]
因此 $\ln(1+x)\sim x$、$\ee^x-1\sim x$。若 $a>0$、$a\ne1$，则 $a^x-1=\ee^{x\ln a}-1\sim x\ln a$。对固定 $\alpha\ne0$，设 $v=\alpha\ln(1+x)\to0$，于是
\[
\frac{(1+x)^\alpha-1}{\alpha x}
=\frac{\ee^v-1}v\frac{\ln(1+x)}x\longrightarrow1.
\]
$\alpha=0$ 时差恒为零，不能称其与 $\alpha x$ 等价。所有式子均要求内部变量趋于 $0$，对数及实幂的底数在讨论邻域内为正。

\begin{example}[整体换元与乘除替换]
求 $\displaystyle\lim_{x\to0}\frac{\ln(1+2x)\tan(3x)}{1-\cos(2x)}$。
\end{example}
\paragraph{思路分析}三个内部变量 $2x,3x,2x$ 均趋于零；分子是乘积，分母是完整的 $1-\cos(2x)$，可分别替换，须保留系数和平方。
\paragraph{解答}在充分小的邻域内 $1+2x>0$，且分母对 $x\ne0$ 非零。由于
\[
\ln(1+2x)\sim2x,\quad \tan(3x)\sim3x,
\quad 1-\cos(2x)\sim\frac{(2x)^2}{2}=2x^2,
\]
原极限为 $\lim 6x^2/(2x^2)=3$。

\begin{example}[相消后先改成乘积]
求 $\displaystyle\lim_{x\to0}\frac{\tan x-\sin x}{x^3}$。
\end{example}
\paragraph{思路分析}$\tan x\sim x$、$\sin x\sim x$ 只知道一阶近似，两项相减后这一阶完全消失。利用恒等式把差化成无需更高阶展开的乘积。
\paragraph{解答}在原点附近 $\cos x\ne0$，
\[
\frac{\tan x-\sin x}{x^3}
=\frac{\sin x}{x}\frac{1-\cos x}{x^2}\frac1{\cos x}
\longrightarrow1\cdot\frac12\cdot1=\frac12.
\]
直接用 $(x-x)/x^3$ 得零丢失了真正的三阶量。加减问题有时靠因式分解或有理化就能解决；需要更精细展开的情形留到第8章。

\section{幂指型极限与方法选择}
\label{lec:03:power}
对正底数 $u(x)>0$，实幂定义为
\[
u(x)^{v(x)}=\exp\bigl(v(x)\ln u(x)\bigr).
\]
先独立证明 $v\ln u$ 的极限，再由指数函数连续性得到原极限。不要先设原极限 $L$ 存在、写 $\ln L$，然后把存在性当作已知。$1^\infty$、$0^0$、$\infty^0$ 是未定式标签，不是运算结果。

\begin{example}[把幂指型化为乘积]
求 $\lim_{x\to0}(1+2x)^{3/x}$ 及 $\lim_{x\to0}(\cos x)^{1/(1-\cos x)}$。
\end{example}
\paragraph{解答}第一式在小邻域内底数为正，
\[
\frac3x\ln(1+2x)=6\frac{\ln(1+2x)}{2x}\longrightarrow6,
\]
故极限为 $\ee^6$。第二式在小去心邻域内 $0<\cos x<1$，令 $t=\cos x-1\to0$，$t\ne0$，则对数为
\[
\frac{\ln(\cos x)}{1-\cos x}
=-\frac{\ln(1+t)}t\longrightarrow-1,
\]
所以极限为 $\ee^{-1}$。

\begin{example}[无穷远处的参数幂指型]
对任意实常数 $a$，求 $\displaystyle\lim_{x\to+\infty}\left(\frac{x-a}{x+a}\right)^x$。
\end{example}
\paragraph{思路分析}充分大时分子分母均为正，底数趋于 $1$。先写成 $1$ 加一个无穷小，随后检查指数与小量的乘积。
\paragraph{解答}若 $a=0$，表达式恒为 $1=\ee^{-2a}$。若 $a\ne0$，令 $t=-2a/(x+a)$，则 $t\to0$、$t\ne0$，对数为
\[
x\ln(1+t)=xt\frac{\ln(1+t)}t
=\frac{-2ax}{x+a}\frac{\ln(1+t)}t\longrightarrow-2a.
\]
故极限为 $\ee^{-2a}$。单独处理 $a=0$ 避免在中途除以零。

\paragraph{计算流程}
先检查定义域与趋近方向，再试直接代入。若分母极限非零，四则运算足够；若出现 $0/0$，检查因式分解、有理化和基本极限；若有有界振荡因子，估计其整体误差；若为正底数幂指型，取对数。涉及相消时先找恒等变形，不能用一阶等价量覆盖所需的高阶信息。判断不存在时，左右极限和 Heine 定理通常比反复计算更直接。

\begin{center}
\begin{tabular}{p{0.25\textwidth}p{0.31\textwidth}p{0.31\textwidth}}
\toprule
方法 & 识别特征 & 必查条件\\
\midrule
四则法则 & 各部分已有有限极限 & 除法的分母极限非零\\
因式分解、有理化 & 多项式或根式同时为零 & 约去的因子在去心邻域非零\\
等价替换 & 基本小量的乘除组合 & 内部变量趋零，误差精度足够\\
夹逼估计 & 振荡、取整、难直接算的量 & 两个界趋于同一个值\\
对数法 & 正底数的变量指数 & 底数为正，先证明对数极限\\
Heine 定理 & 需否定极限或转成数列 & 所选点属于定义域且避开中心\\
\bottomrule
\end{tabular}
\end{center}

\section{分层练习与完整解析}
下列均为自编练习，不作历年真题标注。
\begin{exercise}[基础：定义与方向]
按定义证明 $\lim_{x\to3}(2x-5)=1$；判断 $\lim_{x\to0}1/x$ 的有限极限及单侧无穷极限。
\end{exercise}
\paragraph{解析}给定 $\varepsilon>0$ 取 $\delta=\varepsilon/2$，则 $|(2x-5)-1|=2|x-3|<\varepsilon$。对任意 $K>0$ 取 $\delta=1/K$，当 $0<x<\delta$ 时 $1/x>K$，当 $-\delta<x<0$ 时 $1/x<-K$，故右侧为 $+\infty$、左侧为 $-\infty$；双侧没有有限极限，且 $|1/x|\to+\infty$。

\begin{exercise}[基础：局部估计]
用定义证明 $\lim_{x\to0}x\sin(1/x)=0$，并说明 $\sin(1/x)$ 无极限与本结论为何相容。
\end{exercise}
\paragraph{解析}任给 $\varepsilon>0$ 取 $\delta=\varepsilon$。若 $0<|x|<\delta$，则 $|x\sin(1/x)|\le|x|<\varepsilon$。振荡因子的幅度经乘以 $x$ 后缩到零；乘积收敛无需每个因子都收敛。

\begin{exercise}[迁移：根式与基本极限]
求 $\displaystyle\lim_{x\to0}\frac{\sqrt{1+3x}-1}{\sin(2x)}$。
\end{exercise}
\paragraph{解析}在 $1+3x>0$ 的邻域有
\[
\frac{\sqrt{1+3x}-1}{\sin(2x)}
=\frac{3x}{\sin(2x)(\sqrt{1+3x}+1)}
=\frac32\frac{2x}{\sin(2x)}\frac1{\sqrt{1+3x}+1}
\longrightarrow\frac34.
\]
分母的两因子极限分别为基本比例中的非零值和 $2$，各步合法。

\begin{exercise}[迁移：同阶与等价]
比较 $x\to0$ 时 $1-\cos x$ 与 $x^2$，以及 $\ln(1+x^2)$ 与 $\sin^2x$ 的阶；求后两者之商的极限。
\end{exercise}
\paragraph{解析}$(1-\cos x)/x^2\to1/2$，故同阶而不等价。$\ln(1+x^2)/x^2\to1$，$\sin^2x/x^2\to1$，因此二者等价，商趋于 $1$。内部变量为 $x^2$ 时，对数的标准等价量是 $x^2$。

\begin{exercise}[综合：相消与幂指]
求 $\displaystyle\lim_{x\to0}\frac{\sqrt{1+x^2}-\cos x}{x^2}$ 及 $\displaystyle\lim_{x\to0}(1+x^2)^{\sin x/x^3}$。
\end{exercise}
\paragraph{解析}第一式拆为
\[
\frac{\sqrt{1+x^2}-1}{x^2}+\frac{1-\cos x}{x^2}
=\frac1{\sqrt{1+x^2}+1}+\frac{1-\cos x}{x^2}
\longrightarrow\frac12+\frac12=1.
\]
这次拆分后两项各自存在极限，和法则适用。第二式底数恒正，对数为
\[
\frac{\sin x}{x^3}\ln(1+x^2)
=\frac{\sin x}{x}\frac{\ln(1+x^2)}{x^2}\longrightarrow1,
\]
故原极限为 $\ee$。

\begin{exercise}[综合：用两列否定极限]
设 $f(x)=1$（$x$ 为有理数），$f(x)=0$（$x$ 为无理数）。证明它在 $x\to0$ 时无极限，而 $xf(x)\to0$。
\end{exercise}
\paragraph{解析}取 $x_n=1/n$ 和 $y_n=\sqrt2/n$，两列均趋于零，前列有理、后列无理，所以函数值分别恒为 $1,0$，由 Heine 定理极限不存在。另一方面 $|xf(x)|\le|x|$，故乘积趋于零。

\paragraph{自查}能否不看答案写出定义的量词顺序和否定？能否指出每次约分、替换、取对数的条件？遇到相消时，能否说明所保留的尺度是什么？能否用明确数列把“振荡”改写成极限不存在的证明？这些检查比只记住结果更能保证方法可迁移。


\chapter{连续性与间断点}
\label{lec:04:chapter}

本章的前置知识是函数极限、单侧极限、极限的运算法则以及数列的收敛子列。极限描述点附近的趋势，连续性进一步要求这个趋势与该点的函数值衔接。学习本章后，应能判断分段函数是否连续、辨别间断的原因，并从连续性证明方程有解、最值存在和函数值的统一控制。全章先讨论一点附近的局部性质，再讨论整个闭区间上的性质；后者依靠区间的闭性与有界性，不能只凭每一点连续就直接断言。

\section{连续：极限与函数值衔接}
\label{lec:04:continuity}

\subsection{定义、增量与单侧连续}
对函数 $f(x)=x^2$，当 $x$ 接近 $2$ 时，$f(x)$ 接近 $4=f(2)$。如果只将 $f(2)$ 改成 $7$，附近的极限仍为 $4$，但图像在该点无法衔接。因此连续性同时涉及附近的值和中心点的值。

\begin{definition}
设 $x_0$ 属于函数 $f$ 的定义域 $D$。若对任意 $\varepsilon>0$，存在 $\delta>0$，使得对所有 $x\in D$，
\[
 |x-x_0|<\delta\quad\Longrightarrow\quad
 |f(x)-f(x_0)|<\varepsilon,
\]
则称 $f$ 在 $x_0$ 连续。在通常的内点情形，这等价于
\[
 \lim_{x\to x_0}f(x)=f(x_0).
\]
\end{definition}

这里 $\delta$ 可以依赖 $\varepsilon$ 和所考察的点 $x_0$，但不能在选定后再随 $x$ 改变。连续性的定义允许 $x=x_0$；此时函数值之差为零。令 $\Delta x=x-x_0$，$\Delta y=f(x_0+\Delta x)-f(x_0)$，便有
\[
 f\text{ 在 }x_0\text{ 连续}\quad\Longleftrightarrow\quad
 \Delta y\to0\quad(\Delta x\to0).
\]
连续只要求函数值增量趋零，没有要求增量比 $\Delta y/\Delta x$ 有极限，后者是下一章的导数问题。

\begin{definition}
若 $\lim_{x\to x_0^-}f(x)=f(x_0)$，称 $f$ 在 $x_0$ 左连续；若 $\lim_{x\to x_0^+}f(x)=f(x_0)$，称右连续。函数在 $[a,b]$ 连续，指它在 $(a,b)$ 每一点连续，在 $a$ 右连续，在 $b$ 左连续，记作 $f\in C[a,b]$。
\end{definition}

由双侧极限与单侧极限的关系，在内点 $x_0$ 连续当且仅当左右连续同时成立。求分段函数的连续参数时，必须把左极限、右极限和点值放在一起比较。

\begin{example}[直接使用定义]
用 $\varepsilon$--$\delta$ 定义证明 $f(x)=x^2$ 在任意 $x_0\in\R$ 连续。
\end{example}
\paragraph{思路分析}
将函数值之差分解为 $(x-x_0)(x+x_0)$。前一因子可以由邻域大小控制，后一因子则先用一个固定的小邻域建立上界。
\paragraph{证明}
当 $|x-x_0|<1$ 时，$|x|\le |x_0|+1$，于是
\[
 |x^2-x_0^2|=|x-x_0|\,|x+x_0|
 \le(2|x_0|+1)|x-x_0|.
\]
给定 $\varepsilon>0$，取
\[
 \delta=\min\left\{1,\frac{\varepsilon}{2|x_0|+1}\right\}.
\]
则 $|x-x_0|<\delta$ 推出 $|x^2-x_0^2|<\varepsilon$，证明完成。
\paragraph{方法回顾}
先限制在 $|x-x_0|<1$ 内，再确定精度所需的邻域，是定义证明中常用的两步控制。此处的 $\delta$ 随 $x_0$ 改变，不能据此认为全实轴上存在统一的同一 $\delta$。

\begin{example}[分段函数拼接]
求 $a,b$，使
\[
 f(x)=\begin{cases}
 \dfrac{\sin x}{x},&x<0,\\
 a,&x=0,\\
 \ee^x+b,&x>0
 \end{cases}
\]
在 $0$ 连续。
\end{example}
\paragraph{思路分析}
左、右两段各自连续，唯一需要拼接的点为 $0$。连续要求三项相等，不只是两侧极限相等。
\paragraph{解答}
由重要极限和指数函数在零点的极限，
\[
 f(0^-)=\lim_{x\to0^-}\frac{\sin x}{x}=1,
 \qquad f(0^+)=\lim_{x\to0^+}(\ee^x+b)=1+b.
\]
因此连续的充要条件是 $1=a=1+b$，即 $a=1,b=0$。若仅取 $b=0$ 而 $a\ne1$，两侧极限存在且相等，函数仍不连续。

\subsection{连续性的数列判别}
\begin{theorem}[数列判别]
函数 $f$ 在 $x_0\in D$ 连续，当且仅当对任意 $x_n\in D$ 且 $x_n\to x_0$ 的数列，都有 $f(x_n)\to f(x_0)$。
\end{theorem}
\begin{proof}
若 $f$ 连续，给定 $\varepsilon>0$，先从定义取 $\delta>0$，再由 $x_n\to x_0$ 取 $N$，使 $n>N$ 时 $|x_n-x_0|<\delta$，便有 $|f(x_n)-f(x_0)|<\varepsilon$。

反过来，若 $f$ 在 $x_0$ 不连续，定义的否定为：存在 $\varepsilon_0>0$，对每个 $\delta>0$，都能找到 $x\in D$，使 $|x-x_0|<\delta$ 而 $|f(x)-f(x_0)|\ge\varepsilon_0$。对每个 $n$ 取 $\delta=1/n$，得到 $x_n\to x_0$，但函数值始终与 $f(x_0)$ 相差至少 $\varepsilon_0$，与数列条件矛盾。
\end{proof}
这个判别既可用来否定连续性，也会在闭区间最值定理的证明中把收敛子列传递到函数值。

\section{连续函数的运算与局部性质}
\label{lec:04:operations}

\subsection{四则运算和复合函数}
\begin{theorem}
若 $f,g$ 在 $x_0$ 连续，则 $f+g,f-g,fg$ 在 $x_0$ 连续。若另有 $g(x_0)\ne0$，则 $f/g$ 在 $x_0$ 连续。若 $\varphi$ 在 $x_0$ 连续，$F$ 在 $u_0=\varphi(x_0)$ 连续，且附近的 $\varphi(x)$ 属于 $F$ 的定义域，则 $F\circ\varphi$ 在 $x_0$ 连续。
\end{theorem}
\begin{proof}
四则结论由极限运算法则得到，例如
\[
 \lim_{x\to x_0}f(x)g(x)=f(x_0)g(x_0).
\]
商法则还需要分母在附近非零：由 $g(x)\to g(x_0)\ne0$，可使 $|g(x)-g(x_0)|<|g(x_0)|/2$，于是 $|g(x)|>|g(x_0)|/2$。

对复合函数，给定 $\varepsilon>0$，由外函数连续性取 $\eta>0$，使 $|u-u_0|<\eta$ 时 $|F(u)-F(u_0)|<\varepsilon$。再由内函数连续性取 $\delta>0$，使 $|x-x_0|<\delta$ 时 $|\varphi(x)-u_0|<\eta$。代入 $u=\varphi(x)$ 就得到
\[
 |F(\varphi(x))-F(\varphi(x_0))|<\varepsilon.
\]
\end{proof}

若仅知 $\varphi(x)\to u_0$，即使 $\varphi(x_0)$ 未定义，也仍可由外函数连续性得到
\[
 \lim_{x\to x_0}F(\varphi(x))=F(u_0).
\]
这里必须先证明内层有极限且落在外函数的连续点。不能把极限符号穿过间断的外函数。

基本初等函数在其定义域内连续。因而多项式处处连续，有理函数在分母非零处连续，通常的初等函数在表达式有意义的定义区间上连续。三角函数连续性也可直接估计，例如
\[
 |\sin x-\sin y|
 =2\left|\cos\frac{x+y}{2}\sin\frac{x-y}{2}\right|
 \le |x-y|.
\]
指数函数在任意点的连续性由第3章引理\ref{lec:03:elementary-approach}独立证明的 $\ee^h\to1$ 得到：$\ee^{x_0+h}-\ee^{x_0}=\ee^{x_0}(\ee^h-1)\to0$。该引理也已经证明对数函数在每个正数处的连续性；反三角函数的全定义域连续性可由后面的反函数定理建立。实指数幂在 $x>0$ 上由 $x^\alpha=\ee^{\alpha\ln x}$ 连续；负底数是否有实值必须另行检查。

\begin{example}[定义域先于代入]
求 $f(x)=\ln(2-x^2)/\sqrt{1-x}$ 的定义域，并说明它在何处连续。
\end{example}
\paragraph{解答}
对数真数需满足 $2-x^2>0$，故 $-\sqrt2<x<\sqrt2$。平方根位于分母，需 $1-x>0$，故 $x<1$。两条件相交给出定义域 $(-\sqrt2,1)$。在该区间内，内层多项式、外层对数与平方根都连续，分母非零，所以 $f$ 在此区间每一点连续。端点 $-\sqrt2,1$ 不在定义域内，不能把端点值直接代入原式。

\begin{example}[先求内层极限]
求 $\displaystyle\lim_{x\to0}\left(2+\frac{\sin x}{x}\right)^{\cos x}$。
\end{example}
\paragraph{思路分析}
底数趋于 $3>0$，可在充分小的去心邻域中取对数；先求指数中的有限极限，再用指数函数连续性。
\paragraph{解答}
记 $u(x)=2+\sin x/x$，则 $u(x)\to3$，且附近 $u(x)>0$。于是
\[
 \cos x\ln u(x)\longrightarrow1\cdot\ln3=\ln3.
\]
由恒等式 $u(x)^{\cos x}=\exp(\cos x\ln u(x))$ 和指数函数连续性，所求极限为 $\ee^{\ln3}=3$。先求对数的极限避免了先假设原幂式极限存在的循环论证。

\subsection{局部有界和局部保号}
\begin{proposition}
若 $f$ 在 $x_0$ 连续，则它在 $x_0$ 的某个邻域内有界。若 $f(x_0)>0$，则在某个邻域内 $f(x)>0$；若 $f(x_0)<0$，则在某个邻域内 $f(x)<0$。
\end{proposition}
\begin{proof}
取 $\varepsilon=1$，连续性给出附近 $|f(x)-f(x_0)|<1$，从而 $|f(x)|<|f(x_0)|+1$，这证明局部有界。

当 $f(x_0)>0$ 时取 $\varepsilon=f(x_0)/2$，得到 $f(x)>f(x_0)/2>0$。负值情形应用于 $-f$ 即得。
\end{proof}
保号性的条件是中心点值严格非零。函数 $f(x)=x$ 在零点连续，但零点任意双侧邻域都同时有正值和负值。局部有界也不等于全区间有界，例如 $f(x)=x$ 在全实轴上每点附近有界，但在全实轴上无界。

\subsection{反函数为什么连续}
\label{lec:04:inverse}
\begin{theorem}[反函数连续性]
设 $f$ 在区间 $I$ 上连续且严格单调。则值域 $J=f(I)$ 是区间，反函数 $f^{-1}:J\to I$ 存在、严格单调且连续。
\end{theorem}
\begin{proof}
若 $I$ 只有一个点，结论直接成立。其余情形中，严格单调保证每个函数值只对应一个原像，所以反函数存在且有相应的单调性。任取 $y_1<y_2$ 属于 $J$，它们的原像之间是一个位于 $I$ 的闭区间；介值定理说明所有介于 $y_1,y_2$ 的值都能取到，故 $J$ 是区间。

证明连续性时先设 $f$ 严格递增，固定 $y_0=f(x_0)$。若 $x_0$ 是 $I$ 的内点，给定 $\varepsilon>0$，可选 $x_-,x_+\in I$，使
\[
 x_0-\varepsilon<x_-<x_0<x_+<x_0+\varepsilon.
\]
取 $\delta=\min\{y_0-f(x_-),f(x_+)-y_0\}>0$。对 $y\in J$，若 $|y-y_0|<\delta$，则
\[
 f(x_-)<y<f(x_+)\quad\Longrightarrow\quad
 x_-<f^{-1}(y)<x_+.
\]
因此 $|f^{-1}(y)-x_0|<\varepsilon$。若 $x_0$ 为区间中包含的端点，只选区间内一侧的辅助点，另一侧用反函数值必在 $I$ 内限制，得到同样的相对定义域连续性。严格递减情形对 $-f$ 作相同论证。
\end{proof}
此证明中的介值定理在本章后面建立。连续和严格单调共同保证反函数的定义域是区间，不能仅凭一一对应就声称反函数连续。

\begin{example}[反函数无需显式公式]
设 $f(x)=x^3+x$，证明其反函数在 $\R$ 上存在且连续。
\end{example}
\paragraph{思路分析}
无须解三次方程写出反函数。先检查连续、严格单调，再由两端无穷极限证明值域为全实轴。
\paragraph{解答}
多项式 $f$ 连续。若 $x_2>x_1$，则
\[
 f(x_2)-f(x_1)
 =(x_2-x_1)(x_2^2+x_1x_2+x_1^2+1)>0,
\]
因为 $x_2^2+x_1x_2+x_1^2=(x_2+x_1/2)^2+3x_1^2/4\ge0$。故 $f$ 严格递增。又 $f(x)\to\pm\infty$ 当 $x\to\pm\infty$。给定 $y\in\R$，可取 $A<B$ 使 $f(A)<y<f(B)$，由介值定理存在唯一 $x$ 使 $f(x)=y$。故反函数定义在整个 $\R$ 上，并由反函数连续性定理连续。

\section{间断点：先算两侧，再看点值}
\label{lec:04:discontinuity}

讨论间断点时，通常要求函数在该点的某个去心邻域有定义，该点本身可以无定义。定义域之外且没有这种邻近结构的点，不应一概称为间断点。端点则按对应的单侧连续性处理。

\begin{definition}
设 $x_0$ 是可考察双侧极限的点。若 $f$ 在 $x_0$ 不连续，而左右极限都存在为有限实数，则称第一类间断点。其中：左右极限相等的为可去间断点；左右极限不相等的为跳跃间断点，其跳跃度为 $|f(x_0^+)-f(x_0^-)|$。至少一个单侧有限极限不存在的为第二类间断点，常见的情形有无穷型和振荡型。
\end{definition}

分类操作次序是：先列出分界点、分母零点等可疑点；在其他位置利用初等函数连续性；对每个可疑点分别求左右极限；若两侧有限且相同，再与点值比较。无穷极限 $+\infty,-\infty$ 不属于第一类所需的有限极限。

\begin{example}[约分与补定义]
讨论 $f(x)=(x^2-1)/(x-1)$ 的间断点，并使它成为全实轴上的连续函数。
\end{example}
\paragraph{解答}
当 $x\ne1$ 时，$f(x)=x+1$，所以该函数除 $1$ 外连续。约分并没有给原函数增加点 $1$ 的定义；但它给出
\[
 \lim_{x\to1}f(x)=\lim_{x\to1}(x+1)=2.
\]
因此 $1$ 是可去间断点。定义 $f(1)=2$ 后得到连续延拓；这个值是唯一可使该点连续的补充值。

\begin{example}[绝对值和指数引起的跳跃]
分别讨论 $g(x)=(x^2-1)/|x-1|$ 在 $1$ 以及 $h(x)=1/(1+\ee^{1/x})$ 在 $0$ 的间断类型。
\end{example}
\paragraph{思路分析}
绝对值要按左右两侧去掉；指数项需要根据 $1/x$ 的单侧符号判定趋势。
\paragraph{解答}
当 $x>1$ 时 $g(x)=x+1$，当 $x<1$ 时 $g(x)=-(x+1)$。故
\[
 g(1^- )=-2,\qquad g(1^+)=2,
\]
为跳跃间断点，跳跃度为 $4$。
对 $h$，当 $x\to0^+$ 时 $\ee^{1/x}\to+\infty$，所以 $h(x)\to0$；当 $x\to0^-$ 时 $\ee^{1/x}\to0$，所以 $h(x)\to1$。故零点为跳跃间断点，跳跃度为 $1$。两个函数的两侧极限不相同，任意补定义都无法消除跳跃。

\begin{example}[振荡、无穷与振幅衰减]
比较 $\sin(1/x)$、$1/x$ 和 $x\sin(1/x)$ 在 $0$ 附近的情形；对第三个函数补定义 $f(0)=0$。
\end{example}
\paragraph{解答}
令
\[
 x_n=\frac1{\pi/2+2n\pi},\qquad
 y_n=\frac1{3\pi/2+2n\pi}.
\]
两列都从右侧趋零，而 $\sin(1/x_n)=1$、$\sin(1/y_n)=-1$。因此右极限不存在，$\sin(1/x)$ 在零点有第二类振荡间断，虽有界仍没有极限。

对 $1/x$，右极限为 $+\infty$，左极限为 $-\infty$，零点为第二类无穷间断。对 $x\sin(1/x)$，
\[
 |x\sin(1/x)|\le |x|\longrightarrow0,
\]
由夹逼定理得极限零，恰与补定义相等，故在零点连续。振荡因子出现不等于函数一定间断，需要检查它的振幅是否趋零。

\begin{remark}
有界性与极限存在性不能相互替代。$\sin(1/x)$ 有界但无极限；$\sin(1/x)/x$ 无界振荡，但既不趋于 $+\infty$ 也不趋于 $-\infty$。因此“第二类间断”等同于“趋于无穷”的说法错误。
\end{remark}

\section{闭区间上的全局性质}
\label{lec:04:compact}

\subsection{有界与最值：两个不同的结论}
局部连续性在每一点给出不同的邻域。要控制整个区间，必须把这些局部信息连接起来。本节证明使用已学的实数完备性和有界数列必有收敛子列：在闭区间中，收敛子列的极限仍落在区间内。

\begin{theorem}[有界与最值定理]
若 $f\in C[a,b]$，其中 $a\le b$ 为有限实数，则 $f$ 有界，且存在 $x_m,x_M\in[a,b]$，使
\[
 f(x_m)=m=\min_{x\in[a,b]}f(x),\qquad
 f(x_M)=M=\max_{x\in[a,b]}f(x).
\]
\end{theorem}
\begin{proof}
若 $a=b$，结论直接成立。下面设 $a<b$。

先证明有界。若 $f$ 无界，可对每个正整数 $n$ 取 $x_n\in[a,b]$，使 $|f(x_n)|>n$。有界数列 $x_n$ 有子列 $x_{n_k}\to x_*\in[a,b]$。连续性给出 $f(x_{n_k})\to f(x_*)$，因此该函数值子列有界；但 $|f(x_{n_k})|>n_k\to\infty$，矛盾。

现在值域非空且有界，令 $M=\sup\{f(x):x\in[a,b]\}$。上确界定义保证每个 $n$ 都存在 $z_n\in[a,b]$，使
\[
 M-\frac1n<f(z_n)\le M.
\]
取收敛子列 $z_{n_k}\to x_M\in[a,b]$。由夹逼定理 $f(z_{n_k})\to M$，由连续性又有 $f(z_{n_k})\to f(x_M)$，所以 $f(x_M)=M$。对 $-f$ 应用已证的最大值结论，就得到最小值。
\end{proof}

上确界只说明存在最紧的上界，最大值还要求这个上界被取到。证明最值时不能写“连续所以令最大值为 $M$”后就结束；需要闭区间条件或其他保证取到的论证。

\begin{example}[闭区间条件的边界]
比较 $x$ 在 $(0,1)$、$1/x$ 在 $(0,1)$ 以及在 $[0,1]$ 上定义的函数
\[
 q(x)=\begin{cases}x,&0\le x<1,\\0,&x=1\end{cases}
\]
是否有界、是否取到最大值和最小值。
\end{example}
\paragraph{解答}
$x$ 在 $(0,1)$ 上连续且有界，其上下确界为 $1,0$，但二者都取不到。$1/x$ 在 $(0,1)$ 上连续，却向零点趋近时无上界；其下确界为 $1$，也取不到。$q$ 的定义域为闭区间且有界，最小值 $0$ 可取到，但上确界 $1$ 取不到，因为它在右端点不连续。这三个反例分别说明：开区间可能取不到最值、开区间可能无界、闭区间若缺连续性仍可能取不到最值。

\subsection{零点定理与介值定理}
\label{lec:04:ivt}
\begin{theorem}[零点定理]
若 $f\in C[a,b]$，$a<b$，且 $f(a)f(b)<0$，则存在 $\xi\in(a,b)$ 使 $f(\xi)=0$。
\end{theorem}
\paragraph{证明分析}
端点异号说明根被夹在两端之间。每次二分保留异号的半段，区间长度趋零，两个端点逼近同一点；连续性使这一点的函数值同时非正、非负。
\begin{proof}
不妨设 $f(a)<0<f(b)$。从 $[a_0,b_0]=[a,b]$ 开始，取中点 $c_n=(a_n+b_n)/2$。若 $f(c_n)=0$，已得到根；否则保留端点函数值异号的一半，记为 $[a_{n+1},b_{n+1}]$。若过程一直继续，则
\[
 f(a_n)<0<f(b_n),\qquad b_n-a_n=\frac{b-a}{2^n}.
\]
左端点递增有上界、右端点递减有下界，因此各自收敛。长度趋零保证二者极限相同，记为 $\xi\in[a,b]$。连续性给出 $f(a_n)\to f(\xi)$、$f(b_n)\to f(\xi)$，由极限保序性得 $f(\xi)\le0$ 且 $f(\xi)\ge0$，于是 $f(\xi)=0$。端点值不为零，所以 $\xi\ne a,b$。
\end{proof}

\begin{corollary}[介值定理]
若 $f\in C[a,b]$，则介于 $f(a)$ 与 $f(b)$ 之间的任意实数 $C$ 都能被 $f$ 在 $[a,b]$ 内取到。若 $C$ 严格介于两端点值之间，则存在 $\xi\in(a,b)$ 使 $f(\xi)=C$。
\end{corollary}
\begin{proof}
若 $C$ 等于某个端点值，取相应端点即可。否则构造 $F(x)=f(x)-C$，它连续且端点异号，应用零点定理得到 $F(\xi)=0$。
\end{proof}

\begin{figure}[htbp]
\centering
\begin{tikzpicture}[x=1cm,y=1cm]
\draw[->] (0,0)--(5.3,0) node[right] {$x$};
\draw[->] (0,-1)--(0,2.4) node[above] {$y$};
\draw[dashed] (0,1)--(5,1) node[right] {$y=C$};
\draw[thick] (0.6,-0.5) .. controls (1.4,-0.3) and (1.8,1.6) .. (2.7,1.3) .. controls (3.6,0.9) and (4.1,1.5) .. (4.7,2);
\fill (0.6,-0.5) circle (1.5pt) node[below] {$(a,f(a))$};
\fill (4.7,2) circle (1.5pt) node[above left] {$(b,f(b))$};
\draw[dashed] (0.6,-0.5)--(0.6,0) node[above] {$a$};
\draw[dashed] (4.7,2)--(4.7,0) node[below] {$b$};
\end{tikzpicture}
\caption{连续曲线从低于 $C$ 走向高于 $C$，必须与水平线 $y=C$ 相交。}
\label{lec:04:fig-ivt}
\end{figure}

将最值定理与介值定理结合，$f\in C[a,b]$ 的值域恰为 $[m,M]$：最值给出两端可取，介值给出中间没有缺口。零点定理仅保证至少一个根；端点同号不能否定根的存在，例如 $x^2$ 在 $[-1,1]$ 的端点均为正，但零点存在。

\subsection{存在性题目中的辅助函数}
\begin{example}[存在与唯一分别证明]
证明方程 $x^3+x-1=0$ 在 $(0,1)$ 有唯一实根。
\end{example}
\paragraph{思路分析}
把方程左边作为连续函数，用端点异号证明存在；再用严格单调排除两个根。两项工作不能互相替代。
\paragraph{解答}
令 $F(x)=x^3+x-1$，它在 $[0,1]$ 连续，且 $F(0)=-1,F(1)=1$。零点定理给出 $\xi\in(0,1)$。对 $x_2>x_1$，
\[
 F(x_2)-F(x_1)=(x_2-x_1)(x_2^2+x_1x_2+x_1^2+1)>0,
\]
所以 $F$ 严格递增，不能在两个不同点都取值零，根唯一。若按二分法估计根，每次保留异号区间，$n$ 次后二分区间长度为 $2^{-n}$；取中点近似时误差不超过 $2^{-n-1}$。

\begin{example}[开区间中的零点]
设 $f$ 在 $(a,b)$ 连续，且 $\lim_{x\to a^+}f(x)=A$、$\lim_{x\to b^-}f(x)=B$ 为有限值，$AB<0$。证明 $f$ 在 $(a,b)$ 内有零点。
\end{example}
\paragraph{思路分析}
原函数在端点未定义，不能直接声称它属于 $C[a,b]$。有两种合法途径：按极限补定义，或在内部选取异号端点。这里采用后一种。
\paragraph{证明}
不妨设 $A<0<B$。由单侧极限的保号性，在靠近 $a$ 的右侧有 $f(x)<0$，在靠近 $b$ 的左侧有 $f(x)>0$。各选 $a_1,b_1\in(a,b)$，并使 $a_1<b_1$，满足 $f(a_1)<0<f(b_1)$。函数在 $[a_1,b_1]$ 连续，故存在 $\xi\in(a_1,b_1)\subset(a,b)$ 使 $f(\xi)=0$。
\paragraph{条件辨析}
若只给出两个端点极限存在而不给 $AB<0$，结论不成立，例如常函数 $f(x)=1$ 没有零点。若 $AB=0$，也未必有内部零点，例如 $f(x)=x$ 在 $(0,1)$ 上。

\begin{example}[固定点]
设 $f\in C[a,b]$ 且对所有 $x\in[a,b]$ 有 $a\le f(x)\le b$。证明存在 $\xi\in[a,b]$ 使 $f(\xi)=\xi$。
\end{example}
\paragraph{思路分析}
目标是两函数相等，构造它们的差 $F(x)=f(x)-x$。值域限制刚好给出端点的相反方向不等式。
\paragraph{证明}
$F$ 在 $[a,b]$ 连续，且 $F(a)=f(a)-a\ge0$、$F(b)=f(b)-b\le0$。若任一等号成立，对应端点就是固定点；否则端点严格异号，由零点定理得内部零点。因此必有 $\xi\in[a,b]$ 满足 $F(\xi)=0$。结论包括端点，不能无条件改成 $\xi\in(a,b)$，因为 $f(x)\equiv a$ 只有端点固定点。

\begin{example}[正值的统一下界]
设 $f\in C[a,b]$ 且处处 $f(x)>0$。证明存在 $m>0$，使 $f(x)\ge m$，并说明 $1/f$ 为什么有界。
\end{example}
\paragraph{证明}
最值定理给出 $x_m\in[a,b]$ 使 $m=f(x_m)=\min f$。因为该点也满足严格正值条件，所以 $m>0$。于是
\[
 0<\frac1{f(x)}\le\frac1m\qquad(x\in[a,b]).
\]
这一步的关键是最小值在区间内取到；在 $(0,1)$ 上 $f(x)=x>0$ 却没有正的统一下界。

\begin{proposition}[奇次多项式必有实根]
次数为奇数的实系数多项式 $P$ 至少有一个实零点。
\end{proposition}
\begin{proof}
设 $P(x)=a_0x^{2k+1}+a_1x^{2k}+\cdots+a_{2k+1}$，$a_0\ne0$。乘以 $-1$ 不改变零点，故可设 $a_0>0$。提出最高次幂，括号内为 $a_0+a_1/x+\cdots+a_{2k+1}/x^{2k+1}\to a_0>0$。因此 $P(x)\to+\infty$ 当 $x\to+\infty$，$P(x)\to-\infty$ 当 $x\to-\infty$。选有限实数 $A<B$ 使 $P(A)<0<P(B)$，在 $[A,B]$ 应用零点定理即得实根。
\end{proof}

\section{补充：一致连续与斜率控制}
\label{lec:04:uniform}

本节为补充内容。普通连续性在每一点各选一个邻域；一致连续要求对给定误差，在整个集合中选同一个控制距离。

\begin{definition}
函数 $f:D\to\R$ 在 $D$ 上一致连续，是指
\[
 \forall\varepsilon>0\ \exists\delta>0\ \forall x,y\in D:
 \quad |x-y|<\delta\Longrightarrow |f(x)-f(y)|<\varepsilon.
\]
若存在常数 $L\ge0$，使所有 $x,y\in D$ 都满足 $|f(x)-f(y)|\le L|x-y|$，则称 $f$ 满足 Lipschitz 条件。
\end{definition}

普通连续性允许 $\delta$ 依赖固定的 $x_0$；一致连续中的 $\delta$ 只能依赖 $\varepsilon$，对任意两点都有效。Lipschitz 条件直接给出一致连续：取 $\delta=\varepsilon/(L+1)$ 即可。一致连续的否定则是
\[
 \exists\varepsilon_0>0\ \forall\delta>0\ \exists x,y\in D:
 \quad |x-y|<\delta,\quad |f(x)-f(y)|\ge\varepsilon_0.
\]
因此，构造两点距离趋零、函数值距离却始终有正下界的两列点，就能证明不一致连续。

\begin{theorem}[闭区间上的一致连续性]
若 $f\in C[a,b]$，则 $f$ 在 $[a,b]$ 上一致连续。
\end{theorem}
\begin{proof}
若不一致连续，存在 $\varepsilon_0>0$，每个正整数 $n$ 都可取 $x_n,y_n\in[a,b]$，满足
\[
 |x_n-y_n|<1/n,\qquad |f(x_n)-f(y_n)|\ge\varepsilon_0.
\]
取 $x_n$ 的收敛子列 $x_{n_k}\to x_*\in[a,b]$。由
\[
 |y_{n_k}-x_*|\le |y_{n_k}-x_{n_k}|+|x_{n_k}-x_*|\to0
\]
知 $y_{n_k}\to x_*$。连续性给出两列函数值都趋于 $f(x_*)$，因而其差趋零，与固定正下界矛盾。
\end{proof}

\begin{example}[全实轴上的统一距离]
证明 $\sin x$ 在 $\R$ 上一致连续，而 $x^2$ 在 $\R$ 上不一致连续。
\end{example}
\paragraph{解答}
前面已经用三角恒等式证明 $|\sin x-\sin y|\le |x-y|$，取 $\delta=\varepsilon$，就对所有实数点对同时有效。

对 $x^2$，取 $x_n=n$、$y_n=n+1/n$。它们距离为 $1/n\to0$，而
\[
 y_n^2-x_n^2=2+\frac1{n^2}\ge2.
\]
若存在对 $\varepsilon=1$ 有效的统一 $\delta$，选 $n>1/\delta$ 就会违背定义，故不一致连续。$x^2$ 仍在每一点连续，其失败来自远处增长过快。

\begin{example}[端点附近的两点构造]
证明 $1/x$ 在 $(0,1)$ 上不一致连续；证明 $\sqrt{x}$ 在 $[0,+\infty)$ 上一致连续。
\end{example}
\paragraph{解答}
对 $1/x$，取 $x_n=1/(n+1)$、$y_n=1/(n+2)$，则
\[
 |x_n-y_n|=\frac1{(n+1)(n+2)}\to0,
 \qquad \left|\frac1{x_n}-\frac1{y_n}\right|=1.
\]
因此 $\varepsilon_0=1$ 给出不一致连续的证据。

对平方根，若 $x\ge y\ge0$，则
\[
 (\sqrt{x}-\sqrt{y})^2=x+y-2\sqrt{xy}\le x-y,
\]
因为 $y\le\sqrt{xy}$。交换两点也得 $|\sqrt{x}-\sqrt{y}|\le\sqrt{|x-y|}$。给定 $\varepsilon>0$，取 $\delta=\varepsilon^2$ 即得一致连续。这个例子说明定义域无界并不妨碍一致连续，Lipschitz 条件也非必要条件：在零点附近 $\sqrt{x}/x=1/\sqrt{x}$ 无界，不能有固定的 Lipschitz 常数。

\subsection{割线斜率控制连续性}
\label{lec:04:secants}
若题目给出三个点的割线斜率次序，目标却是连续或 Lipschitz 估计，应固定区间外侧的点，把内部的可变斜率夹在两个固定斜率之间。令
\[
 S(u,v)=\frac{f(v)-f(u)}{v-u}\qquad(u<v).
\]
三个点 $u<v<w$ 满足恒等式
\[
 S(u,w)=\frac{v-u}{w-u}S(u,v)+\frac{w-v}{w-u}S(v,w).
\]
右边是权重为正、和为一的平均值。所以若 $S(u,v)\le S(v,w)$，就有 $S(u,v)\le S(u,w)\le S(v,w)$。这些代数事实不用导数，能直接处理本章两道真题。

\begin{example}[从割线得到右连续]
设实值函数 $f$ 在区间 $I$ 上对任意 $u<v<w$ 满足 $S(u,v)\le S(v,w)$。证明 $f$ 在 $I$ 的每个内点右连续。
\end{example}
\paragraph{思路分析}
不能先假设 $f$ 可导。固定 $a<x_0<b$，用涉及 $a,b$ 的斜率分别控制短区间 $[x_0,x]$ 的斜率，然后乘以趋零的长度。
\paragraph{证明}
取 $a,x_0,b\in I$ 且 $a<x_0<b$。当 $x_0<x<b$ 时，由三点条件与上述加权平均关系，
\[
 S(a,x_0)\le S(x_0,x)\le S(x_0,b).
\]
令 $C=\max\{|S(a,x_0)|,|S(x_0,b)|\}$，得到
\[
 |f(x)-f(x_0)|\le C(x-x_0).
\]
给定 $\varepsilon>0$，取 $\delta=\min\{b-x_0,\varepsilon/(C+1)\}$，则 $0<x-x_0<\delta$ 时函数值差小于 $\varepsilon$，右连续得证。左侧可以用相应的固定割线证明，因此这种三点条件实际还能推出内点连续。

\section{分层练习与完整解析}
\label{lec:04:exercises}
以下均为自编练习，按定义、方法迁移和综合证明递进。

\begin{exercise}[基础：补定义与分类]
设 $f(x)=(1-\cos x)/x^2$（$x\ne0$）。求连续补定义的值。再判断 $g(x)=\lfloor x\rfloor$ 在整数 $k$ 的连续性与跳跃度。
\end{exercise}
\paragraph{解析}
由 $1-\cos x=2\sin^2(x/2)$，
\[
 \frac{1-\cos x}{x^2}=\frac12\left(\frac{\sin(x/2)}{x/2}\right)^2\to\frac12.
\]
所以唯一补充值是 $f(0)=1/2$。对取整函数，当 $x\to k^-$ 时 $\lfloor x\rfloor=k-1$，当 $x\to k^+$ 时 $\lfloor x\rfloor=k$，且 $g(k)=k$。因此 $k$ 是跳跃度为 $1$ 的跳跃间断点，右连续但不左连续。非整数点附近函数为常数，故连续。

\begin{exercise}[基础：振荡因子的边界]
对 $\alpha\in\R$，定义 $f_\alpha(0)=0$，$f_\alpha(x)=|x|^\alpha\sin(1/x)$（$x\ne0$）。判定零点连续的充要条件。
\end{exercise}
\paragraph{解析}
若 $\alpha>0$，则 $|f_\alpha(x)|\le |x|^\alpha\to0$，故连续。若 $\alpha=0$，用本章构造的两列点使正弦取 $1$ 和 $-1$，极限不存在。若 $\alpha<0$，在 $x_n=1/(\pi/2+2n\pi)$ 处函数值为 $x_n^\alpha\to+\infty$，在 $y_n=1/(3\pi/2+2n\pi)$ 处函数值为 $-y_n^\alpha\to-\infty$，所以不可能趋零，且发生无界振荡。因此充要条件为 $\alpha>0$。

\begin{exercise}[迁移：对称点的等值]
设 $f\in C[0,1]$。判定以下结论是否成立：存在 $\xi\in[0,1/2]$，使 $f(\xi)=f(\xi+1/2)$。若不成立，给出反例，并讨论增加条件 $f(0)=f(1)$ 后的情形。
\end{exercise}
\paragraph{解析}
原条件不足，不能证明。取 $f(x)=x$，则对所有 $\xi\in[0,1/2]$ 有 $f(\xi+1/2)-f(\xi)=1/2$。若增加 $f(0)=f(1)$，令 $F(x)=f(x)-f(x+1/2)$，则 $F$ 在 $[0,1/2]$ 连续，且
\[
 F(0)=f(0)-f(1/2),\qquad F(1/2)=f(1/2)-f(0)=-F(0).
\]
若 $F(0)=0$，可取 $\xi=0$；否则端点异号，由零点定理得到 $\xi\in(0,1/2)$。本题的训练目标是先核验条件，再构造辅助函数。

\begin{exercise}[迁移：无界区间上的最小值]
设 $f$ 在 $\R$ 上连续，且 $\lim_{|x|\to\infty}f(x)=+\infty$。证明 $f$ 在 $\R$ 上取到最小值。
\end{exercise}
\paragraph{解析}
先把可能的最小值限制到一个有限闭区间。由无穷远极限，存在 $R>0$，使 $|x|>R$ 时 $f(x)>f(0)+1$。在 $[-R,R]$ 上，最值定理给出 $x_m$ 使 $f(x_m)=m=\min_{[-R,R]}f$。由于 $0$ 在这个区间内，$m\le f(0)$；区间外却有 $f(x)>f(0)+1>m$，故 $m$ 也是全实轴最小值。这里没有直接在无界区间上套用闭区间定理，而是先建立外侧估计。

\begin{exercise}[综合：反函数的定量控制]
设 $f(x)=x^3+x$，反函数为 $g$。证明对所有 $u,v\in\R$，$|g(u)-g(v)|\le |u-v|$，从而 $g$ 一致连续。
\end{exercise}
\paragraph{解析}
已知 $f$ 是严格递增的连续双射。令 $x=g(u),y=g(v)$，则
\[
 |u-v|=|f(x)-f(y)|
 =|x-y|(x^2+xy+y^2+1)\ge |x-y|,
\]
因为 $x^2+xy+y^2\ge0$。代回即得 $|g(u)-g(v)|\le|u-v|$，取 $\delta=\varepsilon$ 证明一致连续。

\begin{exercise}[综合：有界不保证一致连续]
证明 $\sin(1/x)$ 在 $(0,1)$ 上有界且连续，但不一致连续。
\end{exercise}
\paragraph{解析}
函数是连续函数的复合且分母非零，所以在 $(0,1)$ 每一点连续；又其绝对值不超过 $1$。取
\[
 x_n=\frac1{2n\pi+\pi/2},\qquad
 y_n=\frac1{2n\pi+3\pi/2}.
\]
它们均属 $(0,1)$，且
\[
 |x_n-y_n|=\frac{\pi}{(2n\pi+\pi/2)(2n\pi+3\pi/2)}\to0,
\]
但 $|\sin(1/x_n)-\sin(1/y_n)|=2$。固定 $\varepsilon_0=1$，这两列点否定一致连续。

\paragraph{章末自查}
遇到连续性问题，先分清考察点、定义域和单侧方向；遇到间断分类，先比较两个有限单侧极限，再看点值；遇到存在性证明，写出辅助函数并逐项核验连续和端点符号；遇到最值或统一下界，检查最值是否取到；遇到一致连续，检查控制距离能否对整个集合中的所有点对同时使用。每条结论都要连同它的条件使用。


\chapter{导数与微分}
\label{lec:05:chapter}

本章的前置知识是函数极限、连续性和基本初等函数。极限描述变量怎样接近一个数，导数进一步研究：自变量发生很小的变化时，函数值以怎样的速度变化。我们先由差商建立定义，再从定义证明求导法则，最后把导数组织成可执行的计算方法。中值定理、泰勒公式和积分尚未建立，本章的基础推导不使用这些工具。

\section{从平均变化率到导数}
\label{lec:05:definition}

设物体的位置为 $s(t)=t^2$。从时刻 $t_0$ 到 $t_0+h$ 的平均速度为
\[
 \frac{s(t_0+h)-s(t_0)}h=2t_0+h.
\]
当时间间隔趋于零时，平均速度趋于 $2t_0$。这使我们能用一个极限定义某一时刻的速度，而无需把“瞬间”当作一个非零时间段。

几何上，$[f(x_0+h)-f(x_0)]/h$ 是曲线上两点之间的割线斜率；若其具有有限极限，就用该极限作为切线斜率。

\begin{definition}[导数]
设 $f$ 在 $x_0$ 的某个开邻域内有定义。若有限极限
\[
 f'(x_0)=\lim_{h\to0}\frac{f(x_0+h)-f(x_0)}h
       =\lim_{x\to x_0}\frac{f(x)-f(x_0)}{x-x_0}
\]
存在，则称 $f$ 在 $x_0$ 可导。这个数称为导数。若在区间内每一点均可导，把各点的导数组成的函数称为导函数，记作 $f'$。
\end{definition}

导数 $f'(x_0)$ 是一个数，导函数 $f'(x)$ 是一个函数。定义中的 $h$ 可正可负，分子必须减去实际点值 $f(x_0)$。差商趋于无穷时，本书称该点不可导；曲线仍可能有垂直切线，例如 $y=x^{1/3}$ 在原点。

\begin{definition}[单侧导数]
分别以 $h\to0^-$ 和 $h\to0^+$ 定义 $f'_-(x_0)$、$f'_+(x_0)$。两者均为有限数且相等，是内点可导的充要条件。在闭区间端点只讨论相应的单侧导数。
\end{definition}

\begin{example}[定义法与几何意义]
用定义求 $f(x)=x^2$ 在 $x_0=1$ 的导数，并写出切线与法线。
\end{example}
\paragraph{思路分析}先把差商中的零除零消去，再取极限；直线还需要经过曲线上的实际点。
\paragraph{解答}
\[
 f'(1)=\lim_{h\to0}\frac{(1+h)^2-1}h
       =\lim_{h\to0}(2+h)=2.
\]
切线为 $y-1=2(x-1)$，法线为 $y-1=-\tfrac12(x-1)$。一般切线为 $y-f(x_0)=f'(x_0)(x-x_0)$；当 $f'(x_0)=0$ 时法线是竖直线 $x=x_0$，不能写 $-1/f'(x_0)$。

\begin{figure}[htbp]
\centering
\begin{tikzpicture}[scale=1.65]
 \draw[->] (-.15,0)--(2.25,0) node[right] {$x$};
 \draw[->] (0,-.2)--(0,3.1) node[above] {$y$};
 \draw[domain=0:1.72,samples=70] plot (\x,{\x*\x}) node[right] {$y=x^2$};
 \draw[dashed,domain=.4:1.9] plot (\x,{2*\x-1});
 \draw[domain=.55:1.75] plot (\x,{2.5*\x-1.5});
 \fill (1,1) circle (.035) node[above left] {$P(1,1)$};
 \fill (1.5,2.25) circle (.035) node[above left] {$Q$};
 \draw[dotted] (1.5,0)--(1.5,2.25);
 \node[below] at (1.5,0) {$1+h$};
 \node at (1.75,.9) {切线};
 \node[anchor=west] at (2.05,2.25) {割线};
\end{tikzpicture}
\caption{固定 $P$，令 $Q$ 趋向 $P$；割线斜率 $2+h$ 趋于切线斜率 $2$。}
\end{figure}

\section{可导、连续与分段点}
\label{lec:05:continuity}

\begin{theorem}
若 $f$ 在 $x_0$ 可导，则在 $x_0$ 连续。
\end{theorem}
\begin{proof}
对 $h\ne0$，有
\[
 f(x_0+h)-f(x_0)=\frac{f(x_0+h)-f(x_0)}h\,h.
\]
右边第一个因子趋于有限数 $f'(x_0)$，第二个趋于零，故函数增量趋于零，这正是连续性。
\end{proof}

逆命题不成立。对 $f(x)=|x|$，零点处左右差商分别为 $-1$ 与 $1$；它连续但不可导。对 $f(x)=x^{1/3}$，差商 $h^{-2/3}\to+\infty$；连续性仍不能保证有限导数。

\paragraph{分段点的操作顺序}
先查点值与左右极限是否一致；若不连续，立即判不可导。若连续，再计算左右差商；相等时才得到导数。分别求分段表达式的导数，通常只能得到分界点两侧的导函数，不能把其值直接当作分界点导数。

\begin{example}[连续条件与可导条件]
设
\[
 f(x)=\begin{cases}x^2+x+b,&x<0,\\ \arctan(ax),&x\ge0.\end{cases}
\]
求使 $f$ 在 $0$ 可导的 $a,b$。
\end{example}
\paragraph{思路分析}连续性决定函数值能否接上，可导性决定两侧斜率能否接上；两个条件不能合为一个。
\paragraph{解答}因 $f(0)=0$，左极限为 $b$，故必须 $b=0$。此时
\[
 f'_-(0)=\lim_{x\to0^-}(x+1)=1,
 \qquad f'_+(0)=\lim_{x\to0^+}\frac{\arctan(ax)}x=a.
\]
当 $a=0$ 时右差商为零；当 $a\ne0$ 时用 $\arctan u/u\to1$ 得最后等式。因此 $a=1,b=0$，并且这组数满足两侧导数相等，充分性也得到验证。

\begin{example}[振荡函数在特殊点求导]
设 $f(0)=0$，$f(x)=x^2\sin(1/x)$（$x\ne0$）。求 $f'(0)$，并判断 $f'$ 在 $0$ 是否连续。
\end{example}
\paragraph{思路分析}函数表达式在零点无意义，必须用差商；有界振荡乘以趋零因子可夹逼。判断导函数连续则需要另外计算附近的导数。
\paragraph{解答}由
\[
 \left|\frac{f(h)-f(0)}h\right|=|h\sin(1/h)|\le |h|
\]
得 $f'(0)=0$。在 $x\ne0$ 处求导得
\[
 f'(x)=2x\sin(1/x)-\cos(1/x).
\]
取 $x_n=1/(2n\pi)$，有 $f'(x_n)=-1$；取 $z_n=1/((2n+1)\pi)$，有 $f'(z_n)=1$。所以 $f'$ 在零点不连续。可导并不保证导函数连续。

\begin{example}[识别导数定义]
已知 $f$ 在 $a$ 可导，求
\[
 \lim_{h\to0}\frac{f(a+2h)-f(a-3h)}h.
\]
\end{example}
\paragraph{思路分析}原分子跨越基点 $a$，插入 $f(a)$ 才能组成两个定义差商。
\paragraph{解答}
\[
 \frac{f(a+2h)-f(a-3h)}h
 =2\frac{f(a+2h)-f(a)}{2h}
  +3\frac{f(a-3h)-f(a)}{-3h}\longrightarrow5f'(a).
\]
一般系数为 $\alpha+\beta$。这里已知可导是前提；仅知对称差商有极限不够。例如 $|x|$ 的 $[f(h)-f(-h)]/(2h)$ 恒为零，但 $f'(0)$ 不存在。

\section{由定义建立基本公式和运算法则}
\label{lec:05:rules}

\subsection{最初的一组公式}
由极限章节已建立的 $\sin u/u\to1$、$\ln(1+u)/u\to1$、$(\ee^u-1)/u\to1$，可直接证明下列公式；它们不依赖泰勒展开。

对常数 $c$，差商恒为零，所以 $c'=0$。对 $\ee^x$，
\[
 \frac{\ee^{x+h}-\ee^x}h=\ee^x\frac{\ee^h-1}h\longrightarrow\ee^x.
\]
对 $a>0$，将 $a^x=\ee^{x\ln a}$ 的差商写成
\[
 a^x\ln a\,\frac{\ee^{h\ln a}-1}{h\ln a}
\]
得到 $(a^x)'=a^x\ln a$；$a=1$ 时直接按常数处理。
对 $x>0$，
\[
 \frac{\ln(x+h)-\ln x}h
 =\frac1x\frac{\ln(1+h/x)}{h/x}\longrightarrow\frac1x.
\]
故 $(\log_a x)'=1/(x\ln a)$（$a>0,a\ne1$）。

利用三角恒等式，
\[
 \frac{\sin(x+h)-\sin x}h
 =\cos(x+h/2)\frac{\sin(h/2)}{h/2}\longrightarrow\cos x,
\]
而余弦差商为 $-\sin(x+h/2)\sin(h/2)/(h/2)$，故 $(\cos x)'=-\sin x$。角度必须使用弧度制。

\subsection{四则运算及其证明}
\begin{theorem}[四则求导法则]
若 $u,v$ 在考察点可导，则
\[
 (u\pm v)'=u'\pm v',\quad (cu)'=cu',\quad
 (uv)'=u'v+uv'.
\]
若另有 $v\ne0$，则
\[
 \left(\frac uv\right)'=\frac{u'v-uv'}{v^2}.
\]
\end{theorem}
\begin{proof}
和与常数倍直接由极限的线性性质得到。为证明乘积，在基点记函数值为 $u,v$，增量为 $\Delta u,\Delta v$。展开乘积差得
\[
 \frac{(u+\Delta u)(v+\Delta v)-uv}h
 =v\frac{\Delta u}h+u\frac{\Delta v}h
  +\frac{\Delta u}h\Delta v.
\]
可导保证 $\Delta v\to0$，三个项依次趋于 $vu',uv',0$。对于商，连续性保证邻近点分母仍不为零，且
\[
 \frac1h\left(\frac{u+\Delta u}{v+\Delta v}-\frac uv\right)
 =\frac{v(\Delta u/h)-u(\Delta v/h)}{v(v+\Delta v)}
 \longrightarrow\frac{vu'-uv'}{v^2}.
\]
\end{proof}

由商法则导出
\[
 (\tan x)'=\sec^2x,\quad(\cot x)'=-\csc^2x,
 \quad(\sec x)'=\sec x\tan x,\quad
 (\csc x)'=-\csc x\cot x,
\]
各式均要求原三角函数有定义。例如
$ (\tan x)'=(\cos^2x+\sin^2x)/\cos^2x$，这同时解释分母为何不能为零。

\section{复合、反函数和对数求导}
\label{lec:05:chain}

\begin{theorem}[链式法则]
若 $g$ 在 $x_0$ 可导，$f$ 在 $u_0=g(x_0)$ 可导，则 $f\circ g$ 在 $x_0$ 可导，且
\[
 (f\circ g)'(x_0)=f'(g(x_0))g'(x_0).
\]
\end{theorem}
\begin{proof}
由 $f$ 可导，可以写成
\[
 f(u_0+k)-f(u_0)=[f'(u_0)+\varepsilon(k)]k,
 \qquad\varepsilon(k)\to0.
\]
补定义 $\varepsilon(0)=0$，该等式对 $k=0$ 也成立。取
$k=g(x_0+h)-g(x_0)$，由 $g$ 连续知 $k\to0$，所以
\[
 \frac{f(g(x_0+h))-f(g(x_0))}h
 =[f'(u_0)+\varepsilon(k)]\frac{g(x_0+h)-g(x_0)}h
 \longrightarrow f'(u_0)g'(x_0).
\]
这种写法也覆盖 $k=0$ 的情形，避免直接除以函数增量时遗漏零增量。
\end{proof}

应用链式法则，$x^\alpha=\ee^{\alpha\ln x}$（$x>0$）给出
$ (x^\alpha)'=\alpha x^{\alpha-1}$。正整数幂在全实轴成立；任意实指数的统一公式使用 $x>0$。零点和负半轴若属于某一特定幂函数的定义域，须另按其实际定义检查。
对 $x<0$，$\ln|x|=\ln(-x)$，所以同样有 $(\ln|x|)'=1/x$；不能为方便把负半轴上的对数改写成 $\ln x$。

\begin{example}[逐层求导]
求 $y=\ln(1+\sin^2(3x))$ 的导数。
\end{example}
\paragraph{思路分析}先列层次 $x\mapsto3x\mapsto\sin(3x)\mapsto\sin^2(3x)\mapsto\ln(1+\sin^2(3x))$，每层只求自己对应的导数。
\paragraph{解答}对数的真数至少为 $1$，全实轴有定义。链式法则给出
\[
 y'=\frac1{1+\sin^2(3x)}\cdot2\sin(3x)\cdot\cos(3x)\cdot3
 =\frac{6\sin(3x)\cos(3x)}{1+\sin^2(3x)}.
\]
每个因子对应一层；漏掉最后的 $3$ 是把关于中间变量的导数误当成关于 $x$ 的导数。

\begin{example}[内层不可导仍须检查复合函数]
求 $f(x)=\cos(x^{2/3})$ 在 $0$ 的导数。
\end{example}
\paragraph{思路分析}内函数在零点不可导，链式法则的条件不满足，但不能由此断言复合函数不可导；外函数的变化可能更小。
\paragraph{解答}利用 $1-\cos u=2\sin^2(u/2)$ 及 $|\sin v|\le|v|$，
\[
 \left|\frac{\cos(h^{2/3})-1}h\right|
 \le\frac{|h|^{4/3}}{2|h|}=\frac12|h|^{1/3}\to0.
\]
故 $f'(0)=0$。附近导数公式中出现 $x^{-1/3}$，只表明该公式不能直接代入零点。

\subsection{反函数求导与反三角函数}
\label{lec:05:inverse}
\begin{theorem}[反函数求导]
设 $f$ 在区间 $I$ 上连续且严格单调，$x_0$ 是内点，$f$ 在 $x_0$ 可导且 $f'(x_0)\ne0$。则反函数 $\varphi=f^{-1}$ 在 $y_0=f(x_0)$ 可导，并且
\[
 \varphi'(y_0)=\frac1{f'(x_0)}.
\]
\end{theorem}
\begin{proof}
连续严格单调函数的反函数连续。令 $x=\varphi(y)$，则 $y\to y_0$ 时 $x\to x_0$；一一对应性保证 $y\ne y_0$ 时 $x\ne x_0$。于是
\[
 \frac{\varphi(y)-\varphi(y_0)}{y-y_0}
 =\frac1{[f(x)-f(x_0)]/(x-x_0)}\longrightarrow\frac1{f'(x_0)}.
\]
\end{proof}

若 $f'(x_0)=0$，上述倒数极限法失效。例如 $f(x)=x^3$ 的反函数为 $x^{1/3}$，其在零点没有有限导数。

在主值区间 $y=\arcsin x\in[-\pi/2,\pi/2]$ 上，$x=\sin y$。当 $|x|<1$，$\cos y=\sqrt{1-x^2}>0$，故
\[
 (\arcsin x)'=\frac1{\sqrt{1-x^2}}.
\]
类似地，$y=\arccos x\in[0,\pi]$，其内部 $\sin y>0$，得
\[
 (\arccos x)'=-\frac1{\sqrt{1-x^2}},\quad |x|<1.
\]
由 $x=\tan y$ 和 $x=\cot y$ 分别得到
\[
 (\arctan x)'=\frac1{1+x^2},\qquad
 (\operatorname{arccot}x)'=-\frac1{1+x^2}.
\]
后式采用主值 $\operatorname{arccot}x\in(0,\pi)$。端点 $\pm1$ 的反正弦、反余弦差商无有限极限，不能沿用内部公式称其端点可导。

\begin{example}[在对应点求反函数导数]
设 $f(x)=x^3+x$，$\varphi$ 是其反函数。求 $\varphi'(2)$。
\end{example}
\paragraph{思路分析}先解 $f(x_0)=2$；倒数公式中的导数应在 $x_0$，而非直接在 $2$ 求。
\paragraph{解答}若 $b>a$，则
$f(b)-f(a)=(b-a)(b^2+ab+a^2+1)>0$，故 $f$ 严格递增，且连续、值域为 $\R$。因为 $f(1)=2$，$f'(1)=4\ne0$，所以 $\varphi'(2)=1/4$。

\subsection{幂指函数和乘除结构}
当 $u(x)>0$ 时，$u(x)^{v(x)}=\ee^{v(x)\ln u(x)}$，从而
\[
 \frac{\dd}{\dd x}u^v=u^v\left(v'\ln u+v\frac{u'}u\right).
\]
只有底数与指数均为变量时需要这两个项；固定实指数可直接使用幂法则。

\begin{example}[对数求导与定义域]
求 $y=x^x$（$x>0$）以及 $z=(x-1)^2\sqrt{x+2}/(x+3)$（$x>-2$）的导数。
\end{example}
\paragraph{思路分析}对数把乘除转成加减，把幂指数降到前面，但要求取对数的量非零；原函数的零点要另行检查。
\paragraph{解答}由 $\ln y=x\ln x$ 得
\[
 y'/y=\ln x+1,\qquad y'=x^x(\ln x+1).
\]
对 $z$，在 $x\ne1$ 时
\[
 \ln z=2\ln|x-1|+\tfrac12\ln(x+2)-\ln(x+3),
\]
所以
\[
 z'=z\left(\frac2{x-1}+\frac1{2(x+2)}-\frac1{x+3}\right).
\]
在 $x=1$，不能取 $\ln z$。直接用 $z=(x-1)^2 q(x)$ 且 $q$ 在 $1$ 连续，得
$[z(1+h)-z(1)]/h=hq(1+h)\to0$，故 $z'(1)=0$。对数不能把和拆开：例如 $\ln(x^x+2^x)$ 不等于 $x\ln x+x\ln2$。

\section{隐函数：先检查点，再沿关系求导}
\label{lec:05:implicit}

方程 $F(x,y)=0$ 描述一条关系；它可能有多条分支，甚至在指定横坐标处无实数点。求导时须先确定考察的 $(x_0,y_0)$ 确实满足方程，再检查附近是否能写为 $y=y(x)$。

\paragraph{存在性条件与使用范围}
一个常用的充分条件是：$F$ 在 $(x_0,y_0)$ 邻域内有连续的一阶偏导数，$F(x_0,y_0)=0$，且 $F_y(x_0,y_0)\ne0$。隐函数定理保证该点附近唯一的可微分支 $y(x)$。此定理属于多元微积分，本章引用其存在性结论；这里 $F_y$ 表示固定 $x$ 对 $y$ 求导。对于二阶导数，通常要求 $F$ 在邻域内有连续的二阶偏导数。若题目直接给出可导分支，则可在该分支上使用链式法则。

对 $F(x,y(x))=0$ 求导得到 $F_x+F_y y'=0$，从而 $y'=-F_x/F_y$；再求导得到
\[
 F_{xx}+2F_{xy}y'+F_{yy}(y')^2+F_y y''=0.
\]
刚学习时逐项求导通常更清楚。特别是 $\dd(y^2)/\dd x=2yy'$，而对 $yy'$ 求导要同时得到 $(y')^2$ 和 $yy''$。

\begin{example}[隐函数二阶导数]
在单位圆的上半支 $x^2+y^2=1,y>0$，求 $y',y''$。
\end{example}
\paragraph{思路分析}上半支保证 $F_y=2y\ne0$，也确定根号取正。二阶导数宜对一阶关系继续求导，避免把商展开得过长。
\paragraph{解答}依次得到
\[
 2x+2yy'=0,\qquad y'=-\frac xy,
\]
\[
 2+2(y')^2+2yy''=0.
\]
代入 $y'$ 并用 $x^2+y^2=1$，得
\[
 y''=-\frac{1+x^2/y^2}y=-\frac1{y^3}.
\]
在 $y=0$ 的圆上两点，不能除以 $y$；它们是以 $x$ 为自变量的端点，切线竖直。

\begin{example}[求点值的顺序]
由 $y-\ee^{xy}=0$ 确定原点横坐标附近的分支，求 $y(0),y'(0),y''(0)$。
\end{example}
\paragraph{思路分析}指数的内层 $xy$ 自身是乘积。先求 $y(0)$，再求 $y'(0)$，最后代入二阶关系，每次代值都保留足够信息。
\paragraph{解答}当 $x=0$，有 $y=1$，且 $F_y=1-x\ee^{xy}$ 在 $(0,1)$ 为 $1$，故分支存在并光滑。求导得
\[
 y'=\ee^{xy}(y+xy'),
\]
在 $0$ 得 $y'(0)=1$。再求导得
\[
 y''=\ee^{xy}(y+xy')^2+\ee^{xy}(2y'+xy''),
\]
因此 $y''(0)=1+2=3$。不能在求导前把 $x=0$ 代入整个关系，因为代值后已失去附近变化的信息。

\section{参数方程及二阶导数}
\label{lec:05:parametric}

设 $x=x(t),y=y(t)$ 在 $t_0$ 附近连续可导，$x(t)$ 在该邻域严格单调，且 $x'(t_0)\ne0$。严格单调性与连续性保证局部连续反函数 $t=t(x)$ 存在；再由本章的反函数求导和链式法则，得到
\[
 \frac{\dd y}{\dd x}=\frac{y'(t)}{x'(t)}.
\]
若 $x,y$ 二阶可导，则继续运用链式法则：
\[
 \frac{\dd^2y}{\dd x^2}
 =\frac1{x'(t)}\frac{\dd}{\dd t}\left(\frac{y'(t)}{x'(t)}\right)
 =\frac{y''(t)x'(t)-y'(t)x''(t)}{[x'(t)]^3}.
\]
二阶导数是“一阶斜率对 $x$ 的变化率”，每次由对 $t$ 求导转成对 $x$ 求导，都需要除以 $x'(t)$。它不等于 $y''(t)/x''(t)$。这里由 $x'$ 连续且 $x'(t_0)\ne0$，可缩小邻域使分母始终非零。学完第6章后，还可由导数保号推出严格单调，因此在常用的连续可导条件下，单独检查 $x'(t_0)\ne0$ 就足以保证局部反解；这一步的依据是中值定理。

\begin{example}[正规参数点与失效点]
设 $x=t^2,y=t^3$，求 $t>0$ 时的一、二阶导数，并讨论 $t=0$。
\end{example}
\paragraph{思路分析}先查 $x'=2t$，因此公式只可直接用于 $t\ne0$；零点须用消参后的函数单独判断。
\paragraph{解答}在 $t>0$ 时
\[
 \frac{\dd y}{\dd x}=\frac{3t^2}{2t}=\frac32t,\qquad
 \frac{\dd^2y}{\dd x^2}=\frac1{2t}\frac{\dd}{\dd t}(3t/2)=\frac3{4t}.
\]
上支消参为 $y=x^{3/2},x\ge0$。在端点 $x=0$，右导数是
$\lim_{x\to0^+}x^{3/2}/x=0$，但其二阶右导数差商为 $(3/2)\sqrt{x}/x\to+\infty$，没有有限二阶导数。若把 $t<0$ 分支也纳入，同一 $x>0$ 对应正负两个 $y$，整条曲线不能写成单值 $y(x)$。

\begin{example}[变量代换中的二阶导数]
在 $x>0$ 时令 $x=\ee^t$，记 $Y(t)=y(\ee^t)$。把表达式 $x^2y''+a_1xy'+a_2y$ 改写为关于 $t$ 的导数。
\end{example}
\paragraph{思路分析}此例展示参数公式的用途：以指数代换消去导数前的 $x$ 次幂。$Y'$ 是对 $t$ 的导数，$y'$ 是对 $x$ 的导数，须分开记号。
\paragraph{解答}因为 $\dd x/\dd t=\ee^t=x>0$，
\[
 y'=\frac{Y'}x,\qquad
 y''=\frac1x\frac{\dd}{\dd t}(Y'\ee^{-t})
 =\frac{Y''-Y'}{x^2}.
\]
所以
\[
 x^2y''+a_1xy'+a_2y=Y''+(a_1-1)Y'+a_2Y.
\]
这称为欧拉型表达式的变量代换，求解相应微分方程属于后续课程；这里仅验证导数变换。漏掉 $-Y'$ 会把二阶求导误当作单纯除以 $x^2$。

\section{高阶导数与结构化计算}
\label{lec:05:higher}

若 $f'$ 在点处可导，则其导数称为 $f''$；递次定义 $f^{(n)}=(f^{(n-1)})'$。$f^{(n)}$ 表示求导 $n$ 次，而 $(f')^n$ 是一阶导数的 $n$ 次方。一阶可导不保证二阶可导，例如 $f(x)=x|x|$ 在 $0$ 的一阶导数为零，附近导函数为 $2|x|$，其左右导数为 $-2$ 和 $2$。

重复基本公式或归纳可得
\begin{align*}
 (\ee^{ax})^{(n)}&=a^n\ee^{ax},\\
 (x^\alpha)^{(n)}&=\alpha(\alpha-1)\cdots(\alpha-n+1)x^{\alpha-n},\quad x>0,\\
 (\ln x)^{(n)}&=(-1)^{n-1}(n-1)!x^{-n},\quad n\ge1,\\
 [\sin(ax+b)]^{(n)}&=a^n\sin(ax+b+n\pi/2),\\
 [\cos(ax+b)]^{(n)}&=a^n\cos(ax+b+n\pi/2).
\end{align*}
例如对数公式在 $n=1$ 为 $1/x$；若第 $n$ 阶式成立，再求导会乘 $-n$ 并使幂次降低一，正好得到第 $n+1$ 阶式。三角函数每次求导都增加相位 $\pi/2$。

\begin{theorem}[Leibniz 公式]
若 $u,v$ 在区间上 $n$ 阶可导，约定 $u^{(0)}=u$，则
\[
 (uv)^{(n)}=\sum_{k=0}^n\binom nk u^{(k)}v^{(n-k)}.
\]
\end{theorem}
\begin{proof}
$n=1$ 是乘积法则。假设对 $n$ 成立，再对每项求导，分别产生
$\binom nk u^{(k+1)}v^{(n-k)}$ 和 $\binom nk u^{(k)}v^{(n-k+1)}$。把第一类下标改为 $j=k+1$，第二类改为 $j=k$，相同的 $u^{(j)}v^{(n+1-j)}$ 在内部的系数为
$\binom n{j-1}+\binom nj=\binom{n+1}j$，两个端项系数为 $1$。得到 $n+1$ 阶公式，归纳完成。
\end{proof}

\paragraph{方法选择}多项式乘指数或三角函数时，用 Leibniz 公式，因为多项式高阶导数会消失。可裂项的分式先裂项；三角幂先降幂；只求某一点的高阶导数且函数满足简单微分关系时，构造递推。第八章将讲有限阶泰勒公式，使用系数求导数时必须先确认足够阶的可导性，不能把任意形式的渐近式倒推为高阶可导。

\begin{example}[多项式乘指数]
求 $f(x)=x^2\ee^x$ 的 $n$ 阶导数（$n\ge2$）与 $f^{(10)}(0)$。
\end{example}
\paragraph{思路分析}$x^2$ 只有零、一、二阶导数非零，因此求和只保留三项。
\paragraph{解答}
\[
 f^{(n)}(x)=x^2\ee^x+\binom n1(2x)\ee^x+\binom n2(2)\ee^x
 =\ee^x[x^2+2nx+n(n-1)].
\]
在 $n=10,x=0$ 得 $f^{(10)}(0)=90$。因组合数计数的是导数分配方式，不能漏掉二阶项前的 $\binom n2$。

\begin{example}[先裂项再求高阶导数]
求 $f(x)=1/[(x+1)(x+2)]$ 的 $n$ 阶导数。
\end{example}
\paragraph{思路分析}反复使用商法则会不断增加项数，裂项后则只有两个相同类型的幂函数。
\paragraph{解答}在 $x\ne-1,-2$，
\[
 f(x)=\frac1{x+1}-\frac1{x+2},\qquad
 f^{(n)}(x)=(-1)^n n!\left[\frac1{(x+1)^{n+1}}-\frac1{(x+2)^{n+1}}\right].
\]
公式对 $n=0$ 也成立；每次求导将分母指数加一，并乘以其负指数。

\begin{example}[由微分关系建立递推]
设 $f(x)=\arctan x$，求 $f^{(n)}(0)$。
\end{example}
\paragraph{思路分析}先由已知一阶导数写 $(1+x^2)f'=1$。对这个等式求高阶导数，二次多项式使 Leibniz 求和很短；无需预用无穷级数。
\paragraph{解答}$f'=1/(1+x^2)$ 光滑，因此所需高阶导数存在。对等式求 $m\ge1$ 次导数，得
\[
 (1+x^2)f^{(m+1)}+2mx f^{(m)}+m(m-1)f^{(m-1)}=0.
\]
在 $0$ 得 $f^{(m+1)}(0)=-m(m-1)f^{(m-1)}(0)$。$m=1$ 给出 $f''(0)=0$；结合 $f(0)=0,f'(0)=1$，逐次递推得到
\[
 f^{(2k)}(0)=0\quad(k\ge0),\qquad
 f^{(2k+1)}(0)=(-1)^k(2k)!\quad(k\ge0).
\]
奇偶两条递推链应分别看待。

\begin{example}[归纳得到指数与三角函数的高阶导数]
求 $f(x)=\ee^x\cos x$ 的 $n$ 阶导数。
\end{example}
\paragraph{思路分析}求一两阶观察变化：指数因子保留，三角部分相位增加、振幅放大。观察只是猜测，再用归纳法证明。
\paragraph{解答}由乘积法则
\[
 f'(x)=\ee^x(\cos x-\sin x)=\sqrt2\,\ee^x\cos(x+\pi/4).
\]
猜测 $f^{(n)}(x)=2^{n/2}\ee^x\cos(x+n\pi/4)$。$n=0$ 成立；若第 $n$ 阶成立，则再求导为
\[
 2^{n/2}\ee^x[\cos(x+n\pi/4)-\sin(x+n\pi/4)]
 =2^{(n+1)/2}\ee^x\cos(x+(n+1)\pi/4).
\]
所以公式对所有 $n\ge0$ 成立。也可用 Leibniz 公式，但归纳式更紧凑，说明两种方法应按最终表达式的结构选择。

\section{微分、线性近似与误差}
\label{lec:05:differential}

\begin{definition}[一元函数的可微性]
若函数增量可表示为
\[
 f(x_0+h)-f(x_0)=Ah+r(h),\qquad\frac{r(h)}h\to0,
\]
其中 $A$ 不依赖 $h$，则称 $f$ 在 $x_0$ 可微。线性项 $Ah$ 称为微分。取自变量微分 $\dd x=h$，记作 $\dd y=A\dd x$。
\end{definition}

\begin{theorem}
一元函数在一点可微当且仅当在该点可导，且 $A=f'(x_0)$。
\end{theorem}
\begin{proof}
若可微，把增量式除以 $h$，极限为 $A$，故可导。反之，若差商趋于 $f'(x_0)$，令 $\varepsilon(h)$ 为差商与此数之差，则 $\varepsilon(h)\to0$；增量恰为 $f'(x_0)h+\varepsilon(h)h$，余项除以 $h$ 趋于零，故可微。
\end{proof}

于是 $\Delta y=\dd y+o(h)$，切线给出函数增量的一阶线性近似。
当 $f'(x_0)\ne0$ 时，$\Delta y/\dd y\to1$，可以说两者等价；当 $f'(x_0)=0$ 时，微分恒为零，不能写此比值。例如 $y=x^2$ 在零点的 $\dd y=0$，实际增量却是 $h^2$。

链式法则给出一阶微分的形式不变性：若 $y=f(u),u=g(t)$ 可导，则
\[
 \dd y=f'(u)g'(t)\dd t=f'(u)\dd u.
\]
无论 $u$ 是自变量还是中间变量，一阶表达式都相同。但 $u=g(t)$ 时 $\dd u=g'(t)\dd t$，一般并不等于有限增量 $\Delta u$。同理
\[
 \dd(uv)=u\dd v+v\dd u,\qquad
 \dd(u/v)=\frac{v\dd u-u\dd v}{v^2},\quad v\ne0.
\]
二阶微分不能直接照搬一阶形式不变性；换自变量后二阶导数包含内函数的二阶导数。

\begin{example}[近似值与可验证误差]
用微分近似 $\sqrt{4.04}$，并直接估计本次近似的误差。
\end{example}
\paragraph{思路分析}选择已知平方根的基点 $4$，增量为 $h=0.04$。仅有 $o(h)$ 说明相对自变量增量的误差趋零，不能自动给出固定 $h$ 的数值误差；本例用代数恒等式补充估计。
\paragraph{解答}对 $f(x)=\sqrt{x}$，$f'(4)=1/4$，因此
\[
 \sqrt{4.04}\approx2+\tfrac14\cdot0.04=2.01.
\]
令 $s=\sqrt{4+h}$，对 $h>0$ 有 $s>2$，并且
\[
 s-2-\frac h4
 =\frac h{s+2}-\frac h4
 =-\frac{h^2}{4(s+2)^2}.
\]
所以近似偏大，且在 $h=0.04$ 时
$0<2.01-\sqrt{4.04}<h^2/64=0.000025$。这个界由精确恒等式得到，不依赖尚未介绍的泰勒余项。

\paragraph{误差传播的基本含义}
若测量量 $x$ 有小误差 $h$，输出 $f(x)$ 的误差一阶近似为 $f'(x)h$。若 $x_0,f(x_0)$ 均非零，则相对误差近似为
\[
 \frac{\Delta y}{f(x_0)}\approx
 \frac{x_0f'(x_0)}{f(x_0)}\frac h{x_0}.
\]
例如球体积 $V=\tfrac43\pi r^3$，半径相对误差约 $1\%$ 时，体积相对误差一阶约 $3\%$。精确比值为 $(1+h/r)^3-1$，展开这个有限多项式即可看出还存在 $3(h/r)^2+(h/r)^3$，并评估线性近似是否满足精度要求。

\section{分层练习与解析}
\label{lec:05:exercises}

\begin{exercise}[基础：定义与左右导数]
设 $f(0)=0$，$f(x)=x\sin(1/x)$（$x\ne0$）。判断零点连续性与可导性。
\end{exercise}
\paragraph{解析}$|f(x)|\le|x|\to0$，所以连续。差商为 $\sin(1/h)$；沿 $h_n=1/(\pi/2+2n\pi)$ 值为 $1$，沿 $k_n=1/(3\pi/2+2n\pi)$ 值为 $-1$，故导数不存在。比较前面的 $x^2\sin(1/x)$，差商中多一个趋零因子是关键。

\begin{exercise}[基础：链式与对数求导]
求 $y=(\sin x)^x$ 在 $0<x<\pi$ 的导数。
\end{exercise}
\paragraph{解析}该区间 $\sin x>0$，所以 $\ln y=x\ln(\sin x)$ 合法。求导得
\[
 \frac{y'}y=\ln(\sin x)+x\frac{\cos x}{\sin x},\qquad
 y'=(\sin x)^x[\ln(\sin x)+x\cot x].
\]
首项来自指数变化，次项来自底数变化。

\begin{exercise}[迁移：隐函数与参数方程]
（1）由 $x^2+xy+y^2=3$ 确定经过 $(1,1)$ 的分支，求该点的 $y',y''$。
（2）设 $x=t-\sin t,y=1-\cos t$，在 $0<t<2\pi$ 求一、二阶导数。
\end{exercise}
\paragraph{解析}（1）$F_y=x+2y=3\ne0$。一阶关系为
$2x+y+(x+2y)y'=0$，代点得 $y'=-1$。继续求导：
\[
 2+2y'+2(y')^2+(x+2y)y''=0,
\]
故 $2-2+2+3y''=0$，得到 $y''=-2/3$。
（2）$x'=1-\cos t>0$，故
\[
 \frac{\dd y}{\dd x}=\frac{\sin t}{1-\cos t}=\cot(t/2),\qquad
 \frac{\dd^2y}{\dd x^2}
 =\frac{-\tfrac12\csc^2(t/2)}{1-\cos t}
 =-\frac1{4\sin^4(t/2)}.
\]
区间条件保证分母不为零。

\begin{exercise}[迁移：降幂与高阶导数]
求 $f(x)=\cos^4x$ 的 $n$ 阶导数（$n\ge1$）。
\end{exercise}
\paragraph{解析}先用 $\cos^2x=(1+\cos2x)/2$ 两次降幂，得
\[
 \cos^4x=\frac38+\frac12\cos2x+\frac18\cos4x.
\]
常数项在 $n\ge1$ 时消失，故
\[
 f^{(n)}(x)=2^{n-1}\cos(2x+n\pi/2)
            +2^{2n-3}\cos(4x+n\pi/2).
\]
若把 $n=0$ 代入末式会漏掉 $3/8$，范围必须保留。

\begin{exercise}[综合：微分近似的边界]
设 $f$ 在 $a$ 可导且 $f'(a)=0$。是否能由 $f(a+h)-f(a)=o(h)$ 推出该增量恒为零？是否能由 $f(a+h)=f(a)+ch^2+o(h^2)$ 推出 $f''(a)$ 存在？
\end{exercise}
\paragraph{解析}两个结论均不能推出。第一问取 $f(x)=(x-a)^2$。第二问令 $a=0$，
\[
 f(x)=\begin{cases}cx^2+x^3\sin(1/x^2),&x\ne0,\\0,&x=0.\end{cases}
\]
差商 $f(x)/x=cx+x^2\sin(1/x^2)\to0$，故 $f'(0)=0$，并有所给二阶渐近式。但是对 $x\ne0$，
\[
 \frac{f'(x)-f'(0)}x
 =2c+3x\sin(1/x^2)-\frac2x\cos(1/x^2).
\]
沿 $x_n=1/\sqrt{2n\pi}$ 最后一项趋于 $-\infty$；沿 $z_n=1/\sqrt{(2n+1)\pi}$ 趋于 $+\infty$。二阶导数不存在。函数值的余项很小，不能保证余项的导数也很小。

\paragraph{章末自查}能否从题面判断应使用定义、链式、反函数、隐函数、参数方程或 Leibniz 公式？每次除以某个量前，是否确认其非零？取对数是否缩小了原定义域？近似答案是否附有满足本次精度的误差说明？这些检查与求导运算本身同样必要。


\chapter{微分中值定理}
\label{lec:06:chapter}

导数描述一个点附近的变化率，实际问题却常给出两个点的函数值，或者要求控制整个区间的变化。微分中值定理把这两类信息联系起来。本章的前置知识是导数定义与运算法则、闭区间连续函数的最值定理和介值定理。我们先从极值处的导数出发，建立 Rolle、Lagrange、Cauchy 定理，再学习零点、不等式和辅助函数的证明方法。所有论证均不需要泰勒公式。

\section{从极值到 Rolle 定理}
\subsection{Fermat 定理：内部极值的必要条件}
\begin{definition}
设 $x_0$ 是函数定义区间的内点。若存在 $\delta>0$，使该区间内满足 $|x-x_0|<\delta$ 的点均有 $f(x)\le f(x_0)$，称 $x_0$ 为局部极大值点；把不等号反向便得到局部极小值点。二者统称局部极值点。
\end{definition}
局部极值只比较附近的函数值，最值则比较指定集合上的全部函数值。内部最值点一定是局部极值点，端点则只有一侧可比较。
\begin{theorem}[Fermat 定理]
若 $x_0$ 是内部局部极值点，且 $f'(x_0)$ 存在，则 $f'(x_0)=0$。
\end{theorem}
\begin{proof}
先设 $x_0$ 是局部极大值点。对充分小的 $h>0$，
\[
\frac{f(x_0+h)-f(x_0)}h\le0;
\]
对充分小的 $h<0$，分子仍非正，分母变为负，故同一差商非负。取单侧极限，得到 $f'_+(x_0)\le0$、$f'_-(x_0)\ge0$。可导性保证二者相等，公共值只能是零。局部极小值情形的不等号相反，结论相同。
\end{proof}
\begin{example}[检查必要条件的边界]
分别判断 $x^3$ 在零点、$|x|$ 在零点、$x$ 在 $[0,1]$ 的左端点是否适用 Fermat 定理。
\end{example}
\paragraph{解答}
$x^3$ 在零点导数为零，但零点两侧函数值分别为负、正，因此不是极值点；导数为零不能反推极值。$|x|$ 在零点取极小值，但左右导数为 $-1,1$，不可导。$x$ 在 $[0,1]$ 的零点取最小值，但零点是端点，其右导数为 $1$。后两例分别说明可导与内点条件不能省略。

\subsection{Rolle 定理及条件的用途}
\begin{theorem}[Rolle 定理]\label{lec:06:rolle}
设 $a<b$。若 $f$ 在 $[a,b]$ 连续，在 $(a,b)$ 可导，且 $f(a)=f(b)$，则存在 $\xi\in(a,b)$，使 $f'(\xi)=0$。
\end{theorem}
\begin{proof}
闭区间连续性保证存在最大值 $M$ 和最小值 $m$。若 $M=m$，则 $f$ 为常数，任取内部点均满足结论。若 $M>m$，记两端共同函数值为 $c$。由于两个不同最值不可能都等于 $c$，至少一个最值与 $c$ 不同，其取得点便不能是端点。该点是内部极值点，可导性和 Fermat 定理给出导数为零。
\end{proof}
这里没有要求 $f'$ 连续，也没有要求端点可导。中值点只保证存在，不保证唯一或等于中点。
\begin{example}[逐一检查 Rolle 的条件]
说明以下三种情况为什么不能直接应用 Rolle 定理：
\[
f(x)=\begin{cases}x,&0\le x<1,\\0,&x=1;\end{cases}
\qquad g(x)=|x|\quad(-1\le x\le1),\qquad h(x)=x\quad(0\le x\le1).
\]
\end{example}
\paragraph{解答}
$f(0)=f(1)=0$，内部导数恒为 $1$，但在右端点不连续；$g(-1)=g(1)$ 且连续，但在内部的零点不可导，其他点导数为 $\pm1$；$h$ 连续可导，却不满足两端等值。这些函数都没有内部导数零点，分别显示三个条件的作用。

\begin{example}[零点之间的导数零点]
若 $f$ 在 $[a,b]$ 连续、在 $(a,b)$ 二阶可导，且有三个互异零点 $x_1<x_2<x_3$，证明 $f''$ 在 $(x_1,x_3)$ 有零点。
\end{example}
\paragraph{思路分析}
每两个相邻零点之间用一次 Rolle，可以先得到两个互异的 $f'$ 零点；再在这两个点之间用 Rolle。第二次需要的是 $f'$ 的连续性，由二阶可导保证。
\paragraph{证明}
在 $[x_1,x_2]$、$[x_2,x_3]$ 上分别应用 Rolle，得到
\[
u\in(x_1,x_2),\quad v\in(x_2,x_3),\quad f'(u)=f'(v)=0.
\]
因为 $[u,v]\subset(a,b)$，$f'$ 在该闭区间可导，因而连续；在其内部亦可导。再用 Rolle，得到 $w\in(u,v)$，使 $f''(w)=0$。注意三个原零点可以含原区间端点，而新得到的 $u,v$ 均在原区间内部。
\paragraph{方法回顾}
更一般地，$n+1$ 个不同零点逐层产生 $n,n-1,\ldots,1$ 个导数零点；每一层必须有相应阶数的可导性。反向使用时，若 $f^{(n)}$ 无零点，则 $f$ 至多有 $n$ 个不同零点。

\section{Lagrange 与 Cauchy 中值定理}
\subsection{平均变化率与瞬时变化率}
\begin{theorem}[Lagrange 中值定理]\label{lec:06:lagrange}
若 $f$ 在 $[a,b]$ 连续，在 $(a,b)$ 可导，$a<b$，则存在 $\xi\in(a,b)$，使
\[
f(b)-f(a)=f'(\xi)(b-a).
\]
\end{theorem}
\paragraph{构造的来源}
目标是让 $f'(\xi)-k=0$，其中 $k=[f(b)-f(a)]/(b-a)$ 为割线斜率。导数中的常数 $k$ 可由减去直线 $k(x-a)$ 产生。减去整条端点割线还会同时制造两个零点。
\begin{proof}
令
\[
F(x)=f(x)-f(a)-\frac{f(b)-f(a)}{b-a}(x-a).
\]
$F$ 在闭区间连续、内部可导，并且 $F(a)=F(b)=0$。Rolle 定理给出 $F'(\xi)=0$，即所求公式。
\end{proof}
\begin{figure}[htbp]
\centering
\begin{tikzpicture}[scale=1.05]
\draw[->] (-.3,0)--(4.5,0) node[right] {$x$};
\draw[->] (0,-.2)--(0,3.4) node[above] {$y$};
\draw[thick,domain=.5:4,samples=60] plot (\x,{.12*\x*\x+.45});
\draw (.5,.48)--(4,2.37);
\draw[dashed] (1.1,.437)--(3.4,1.679);
\fill (.5,.48) circle (1.4pt) node[above left] {$A$};
\fill (4,2.37) circle (1.4pt) node[above] {$B$};
\fill (2.25,1.0575) circle (1.4pt) node[below right,fill=white,inner sep=1pt] {$P$};
\draw[dotted] (.5,0)--(.5,.48) (2.25,0)--(2.25,1.0575) (4,0)--(4,2.37);
\node[below] at (.5,0) {$a$}; \node[below] at (2.25,0) {$\xi$}; \node[below] at (4,0) {$b$};
\end{tikzpicture}
\caption{割线 $AB$ 的斜率等于内部某点 $P$ 的切线斜率；图中虚线是该切线。}
\end{figure}
若把 $\xi$ 写成 $a+\theta(b-a)$，则 $0<\theta<1$。这只改变中值点的记法，不能把 $\theta$ 当作固定常数。

\begin{example}[从导数范围估计增量]
证明对 $x>0$ 有
\[
\frac{x}{1+x}<\ln(1+x)<x.
\]
\end{example}
\paragraph{思路分析}
对数差 $\ln(1+x)-\ln1$ 的两端自变量之差是 $x$；导数 $1/t$ 在两端之间有简单上下界，因此先用中值定理，再对导数估计。
\paragraph{证明}
函数 $\ln t$ 在 $[1,1+x]$ 连续、内部可导。存在 $\xi\in(1,1+x)$，使 $\ln(1+x)=x/\xi$。由
\[
\frac1{1+x}<\frac1\xi<1
\]
并乘以正数 $x$，即得结论。若 $-1<x<0$，在 $[1+x,1]$ 上应用同一定理，仍得 $\ln(1+x)=x/\xi$，但乘以负数时必须反向，最终同样得到 $x/(1+x)<\ln(1+x)<x$。

\subsection{Cauchy 定理：两个增量的比值}
\begin{theorem}[Cauchy 中值定理]\label{lec:06:cauchy}
设 $a<b$。若 $f,g$ 在 $[a,b]$ 连续、在 $(a,b)$ 可导，则存在同一个 $\xi\in(a,b)$，使
\[
[f(b)-f(a)]g'(\xi)=[g(b)-g(a)]f'(\xi).
\]
若另有 $g'(x)\ne0$ 对所有 $x\in(a,b)$ 成立，则 $g(b)\ne g(a)$，并可写成
\[
\frac{f(b)-f(a)}{g(b)-g(a)}=\frac{f'(\xi)}{g'(\xi)}.
\]
\end{theorem}
\begin{proof}
记 $A=f(b)-f(a)$、$B=g(b)-g(a)$，令
\[
F(x)=A[g(x)-g(a)]-B[f(x)-f(a)].
\]
此组合由目标等式的两项反向求导而来，且系数恰使 $F(a)=F(b)=0$。连续与可导条件由 $f,g$ 继承。Rolle 定理给出 $Ag'(\xi)-Bf'(\xi)=0$。若 $g'$ 内部处处非零，假设 $g(a)=g(b)$ 会由 Rolle 推出某点 $g'=0$，矛盾；所以两处分母均非零，可除法得商式。
\end{proof}
分别对 $f,g$ 使用 Lagrange 通常得到不同的中值点，不能把两个导数直接相除。Cauchy 的价值正在于得到同一个点。
\begin{example}[确定同一个中值点]
对 $f(t)=t^2$、$g(t)=t^3$ 在 $[1,2]$ 使用 Cauchy 定理，求满足商式的中值点。
\end{example}
\paragraph{解答}
两函数连续可导，且 $g'(t)=3t^2>0$。因此
\[
\frac{4-1}{8-1}=\frac{2\xi}{3\xi^2}=\frac2{3\xi},
\qquad \xi=\frac{14}{9}\in(1,2).
\]
若分别用 Lagrange，$f$ 的中值点是 $3/2$，$g$ 的中值点是 $\sqrt{7/3}$，确实不同。

\section{导数怎样控制整个区间}
\subsection{常数、单调性和导数界}
\begin{proposition}\label{lec:06:derivative-control}
设 $f$ 在区间 $I$ 内可导；若 $I$ 含端点，则另要求 $f$ 在这些端点连续。
若内部处处 $f'=0$，则 $f$ 为常数；若内部处处 $f'\ge0$，则 $f$ 单调不减；若内部处处 $f'>0$，则 $f$ 严格递增。若内部处处 $|f'|\le M$，则任意 $x,y\in I$ 都有
\[
|f(x)-f(y)|\le M|x-y|.
\]
\end{proposition}
\begin{proof}
任取 $x<y$，在 $[x,y]$ 上使用 Lagrange，得到 $f(y)-f(x)=f'(\xi)(y-x)$。因 $y-x>0$，导数的零值、非负或正号分别给出三个结论；取绝对值则给出最后的估计。$x=y$ 时估计直接成立。
\end{proof}
严格递增不要求导数处处正，例如 $x^3$ 在零点导数为零。若 $f'\ge0$，且任何非空开子区间上 $f'$ 都不恒为零，也能推出严格递增：否则存在 $x<y$ 使 $f(x)=f(y)$，单调不减会迫使 $f$ 在 $[x,y]$ 恒为常数，与假设矛盾。若 $f$ 严格递增且某点可导，由该点左右差商均非负可知导数非负。

\begin{example}[恒等式与端点延伸]
证明 $\arcsin x+\arccos x=\pi/2$ 对 $-1\le x\le1$ 成立，并证明任意实数 $x,y$ 有 $|\sin x-\sin y|\le|x-y|$。
\end{example}
\paragraph{解答}
令 $H(x)=\arcsin x+\arccos x$。在 $(-1,1)$ 上
\[
H'(x)=\frac1{\sqrt{1-x^2}}-\frac1{\sqrt{1-x^2}}=0.
\]
所以 $H$ 在该开区间为常数，取 $x=0$ 得常数为 $\pi/2$。两函数在 $\pm1$ 连续，取单侧极限将等式延伸到端点；不需要端点导数。第二式中，$x=y$ 时成立；否则中值定理给出 $\sin x-\sin y=\cos\xi(x-y)$，利用 $|\cos\xi|\le1$ 得证。

\subsection{零点存在和唯一必须分开证明}
\begin{example}[严格单调函数的唯一根]
证明 $\ee^x+x=0$ 在实数范围有且仅有一根。
\end{example}
\paragraph{思路分析}
连续性加异号端点负责存在；导数正号负责唯一。单调性本身不保证有根，例如 $\ee^x$ 处处正。
\paragraph{证明}
令 $F(x)=\ee^x+x$。$F(-1)=\ee^{-1}-1<0$，$F(0)=1>0$，连续性保证 $(-1,0)$ 内有根。又 $F'(x)=\ee^x+1>0$，故严格递增，不能有两个不同零点。
\paragraph{另一种唯一性论证}
若有根 $u<v$，在 $[u,v]$ 用 Rolle 会得到内部点 $F'=0$，与导数正号矛盾。这一论证适合只需要排除多个根的题目。

\begin{example}[由导数下界找到异号端点]
设 $f\in C[0,+\infty)$，在 $(0,+\infty)$ 可导，$f(0)<0$，且 $f'(x)\ge m>0$。证明 $f$ 在正半轴有唯一零点。
\end{example}
\paragraph{思路分析}
无穷区间不能直接用闭区间介值定理；先用导数下界把函数值抬到正值，再选一个有限区间。
\paragraph{证明}
对任意 $x>0$，在 $[0,x]$ 用 Lagrange，得 $f(x)-f(0)=xf'(\xi)\ge mx$。取 $B=(1-f(0))/m>0$，则 $f(B)\ge1>0$。由介值定理在 $(0,B)$ 有根；由 $f'>0$ 得严格递增，故唯一。统一正下界 $m$ 很关键：仅有 $f'>0$ 不保证函数最终为正，例如 $f(x)=-\ee^{-x}$。

\subsection{不等式与数列收敛}
\begin{example}[把差转化为单调函数]
证明当 $b>a\ge\ee$ 时 $a^b>b^a$。
\end{example}
\paragraph{思路分析}
底数与指数交叉出现，取对数可把乘幂变为乘积，再除以正数 $ab$，目标变为同一个函数在两点的比较。
\paragraph{证明}
目标等价于 $b\ln a>a\ln b$，也就是 $(\ln a)/a>(\ln b)/b$。令 $H(t)=(\ln t)/t$，则 $H'(t)=(1-\ln t)/t^2$。在 $[a,b]$ 用 Lagrange，得
\[
H(b)-H(a)=\frac{1-\ln\xi}{\xi^2}(b-a)<0,
\]
因为 $\xi>a\ge\ee$。这也覆盖 $a=\ee$ 时端点导数为零的情况。

\begin{example}[调和和与对数之差]
证明数列 $u_n=\sum_{k=1}^n1/k-\ln n$ 收敛。
\end{example}
\paragraph{思路分析}
无需猜测极限的具体值。相邻项之差包含对数增量，用前面的对数不等式检查单调性；同一不等式还能制造可望远镜相消的下界。
\paragraph{证明}
对 $n\ge2$，取 $x=1/(n-1)$，得到 $\ln[n/(n-1)]>1/n$，故
\[
u_n-u_{n-1}=\frac1n-\ln\frac n{n-1}<0.
\]
另一方面，$\ln(1+1/k)<1/k$，所以
\[
u_n>\sum_{k=1}^n\ln\frac{k+1}{k}-\ln n
=\ln(n+1)-\ln n>0.
\]
$u_n$ 单调递减且以下界零有界，由单调有界定理收敛。此证明不使用尚未介绍的积分。

\section{怎样构造辅助函数}
\subsection{先识别导数形式，再制造端点条件}
存在性目标常为某个表达式在一点等于零。先把它看成一个导数的组成部分，再检查新函数的两个端值是否相等；如果原端点不能使用，可先找两个内部零点。
\begin{center}
\begin{tabular}{lll}
\toprule
目标表达式 & 候选辅助函数 & 需要检查的条件\\
\midrule
$f'-c$ & $f-cx$ & 两端值相等\\
$f'+pf$，$p$ 为常数 & $\ee^{px}f$ & 两端加权值相等\\
$f'+g'f$ & $\ee^gf$ & $g$ 连续、内部可导\\
$f'g-fg'$ & $f/g$ & 全区间 $g\ne0$\\
$xf'+nf$ & $x^nf$ & 权函数定义与端点值\\
\bottomrule
\end{tabular}
\end{center}
这张表来自乘积和商的求导法则。选择辅助函数后仍须完整检查 Rolle 条件，不能只写出函数名称。

\begin{example}[带可变系数的导数关系]
设 $f,g$ 在 $[a,b]$ 连续、内部可导，且 $f(a)=f(b)=0$。证明存在 $\xi\in(a,b)$，使 $f'(\xi)+g'(\xi)f(\xi)=0$。
\end{example}
\paragraph{思路分析}
乘积导数含 $f'$ 和 $f$ 两项，希望第二项的系数是 $g'$。权函数 $\ee^{g(x)}$ 的导数恰为 $g'(x)\ee^{g(x)}$，从而两项可提出共同非零因子。
\paragraph{证明}
令 $F=\ee^gf$。它在闭区间连续、内部可导，两端值均为零。因此存在 $\xi$ 使
\[
0=F'(\xi)=\ee^{g(\xi)}[f'(\xi)+g'(\xi)f(\xi)].
\]
指数因子严格正，括号必须为零。若题目只说 $g$ 在内部可导，未保证端点连续，则不能无检查地使用原端点；当前假设明确补足了辅助函数的闭区间连续性。

\begin{example}[先找内部零点，再加权]
设 $f$ 在 $[0,3]$ 连续、内部可导，$f(0)f(2)<0$、$f(2)f(3)<0$。给定实数 $p$，证明存在 $\xi\in(0,3)$，使 $f'(\xi)=pf(\xi)$。
\end{example}
\paragraph{思路分析}
目标对应 $\ee^{-px}f(x)$ 的导数，但原端点未必等值。两个异号条件先给出分离的零点，它们乘以非零权函数后仍为零。
\paragraph{证明}
介值定理分别给出 $c_1\in(0,2)$、$c_2\in(2,3)$，使 $f(c_1)=f(c_2)=0$。令 $F(x)=\ee^{-px}f(x)$，在 $[c_1,c_2]$ 连续、内部可导且两端为零。Rolle 给出 $\xi\in(c_1,c_2)$，使
\[
0=F'(\xi)=\ee^{-p\xi}[f'(\xi)-pf(\xi)].
\]
除以正权函数即得结论。

\begin{example}[寻找已知表达式的原函数形式]
给定实数 $a_1,\ldots,a_n$，证明 $\sum_{k=1}^na_k\cos(kx)$ 在 $(0,\pi)$ 内有零点。
\end{example}
\paragraph{思路分析}
余弦和不一定容易找到异号端点，但它是正弦和的导数；正弦在 $0,\pi$ 的整数倍角处均为零。这使 Rolle 比直接的介值法更自然。
\paragraph{证明}
令 $F(x)=\sum_{k=1}^n(a_k/k)\sin(kx)$，则 $F$ 连续可导，$F(0)=F(\pi)=0$。Rolle 给出 $\xi\in(0,\pi)$，使 $F'(\xi)=\sum_{k=1}^na_k\cos(k\xi)=0$。这里仅用已知三角函数的求导，不需要一般积分理论。

\begin{example}[商函数的适用边界]
设 $f,g$ 在 $[a,b]$ 连续、内部可导，$g$ 在整个闭区间无零点，且 $f(a)g(b)=f(b)g(a)$。证明存在 $\xi\in(a,b)$，使 $f'(\xi)g(\xi)-f(\xi)g'(\xi)=0$。
\end{example}
\paragraph{证明}
无零点条件使 $H=f/g$ 在闭区间连续、内部可导。端点关系除以 $g(a)g(b)\ne0$ 后得到 $H(a)=H(b)$。Rolle 得到
\[
0=H'(\xi)=\frac{f'(\xi)g(\xi)-f(\xi)g'(\xi)}{g(\xi)^2},
\]
即得结论。不能仅检查两个端点分母非零：若 $f=1$、$g=x^2$ 在 $[-1,1]$ 上，则端点比值相等，但商函数在零点无定义，不能对它使用 Rolle。

\subsection{二阶差分也可以用两次 Rolle}
\begin{example}[二次插值辅助函数]\label{lec:06:second-difference}
设 $f$ 在 $[a,b]$ 连续、内部二阶可导，$c=(a+b)/2$。证明存在 $\xi\in(a,b)$，使
\[
f(a)-2f(c)+f(b)=\frac{(b-a)^2}{4}f''(\xi).
\]
\end{example}
\paragraph{思路分析}
希望两次求导后只留下 $f''$ 和一个常数，因此减去二次多项式。让该多项式同时匹配 $a,c,b$ 三点的函数值，就制造三个零点；二次项系数由二阶差分决定。
\paragraph{证明}
记 $h=(b-a)/2>0$，设
\[
A=\frac{f(a)-2f(c)+f(b)}{2h^2},\qquad
B=\frac{f(b)-f(a)}{2h},
\quad P(x)=f(c)+B(x-c)+A(x-c)^2.
\]
代入 $x=c-h,c,c+h$，分别得到 $P(a)=f(a)$、$P(c)=f(c)$、$P(b)=f(b)$。令 $F=f-P$，则 $F(a)=F(c)=F(b)=0$。两次 Rolle 产生 $\xi\in(a,b)$，使 $F''(\xi)=0$。所以
\[
f''(\xi)=2A=\frac{f(a)-2f(c)+f(b)}{h^2}.
\]
由于 $h^2=(b-a)^2/4$，结论成立。此方法直接得到同一个 $f''(\xi)$，不把不同点的二阶导数擅自合并，也无需先用泰勒公式。

\section{导数的介值性：Darboux 定理}
\begin{theorem}[Darboux 定理]\label{lec:06:darboux}
设 $f$ 在 $[a,b]$ 连续、内部可导，端点具有有限单侧导数。若
\[
f'_+(a)<\lambda<f'_-(b),
\]
则存在 $x_0\in(a,b)$，使 $f'(x_0)=\lambda$。两端导数大小反向时，同样达到二者之间的每一个数。
\end{theorem}
\paragraph{证明思路}
不能直接对 $f'$ 使用连续函数的介值定理，因为题目没有保证 $f'$ 连续。把目标导数值 $\lambda$ 减去后，转为 $g=f-\lambda x$ 的导数零点。两端导数一负一正，将迫使 $g$ 的最小值在内部取得。
\begin{proof}
令 $g(x)=f(x)-\lambda x$。它在闭区间连续，且 $g'_+(a)<0$、$g'_-(b)>0$。由单侧导数定义，存在靠近 $a$ 的 $u>a$，满足
\[
\frac{g(u)-g(a)}{u-a}<0,
\quad\text{故 }g(u)<g(a).
\]
也存在靠近 $b$ 的 $v<b$，满足
\[
\frac{g(v)-g(b)}{v-b}>0.
\]
由于 $v-b<0$，第二式给出 $g(v)<g(b)$。因此 $a,b$ 都不能取得全区间最小值。闭区间最值定理保证存在内部最小值点 $x_0$；Fermat 定理得 $g'(x_0)=0$，即 $f'(x_0)=\lambda$。导数大小反向时对 $-f$ 使用上述证明，或将最小值改为最大值。
\end{proof}
若 $f$ 在一个开区间可导，任取两个内部点，把它们作为闭区间端点，即知 $f'$ 具有介值性。介值性弱于连续性：导数可以不连续，却不能从一个值直接跳到另一个值而漏掉中间值。
\begin{example}[不连续但具有介值性的导数]
设 $f(x)=x^2\sin(1/x)$（$x\ne0$），$f(0)=0$。说明 $f$ 处处可导而 $f'$ 在零点不连续。
\end{example}
\paragraph{解答}
零点差商为 $x\sin(1/x)$，其绝对值不超过 $|x|$，所以 $f'(0)=0$。非零点求导得
\[
f'(x)=2x\sin(1/x)-\cos(1/x).
\]
取 $x_n=1/(2n\pi)$、$y_n=1/[(2n+1)\pi]$，则两列趋零，但 $f'(x_n)=-1$、$f'(y_n)=1$，故不连续。Darboux 定理仍保证在每个 $[y_n,x_n]$ 内，导数达到 $-1$ 与 $1$ 之间的全部值；不连续来自振荡，而非跳跃。

\section{分层练习与完整解析}
以下均为自编练习，依次训练条件检查、方法迁移和综合构造。
\begin{exercise}
求 $f(x)=x^2$ 在 $[1,3]$ 上 Lagrange 定理的中值点，并解释它是否必为所有函数的区间中点。
\end{exercise}
\paragraph{解析}
割线斜率为 $(9-1)/(3-1)=4$，$f'(\xi)=2\xi=4$，故 $\xi=2$。这只是本函数的计算结果。例如 $x^3$ 在 $[0,1]$ 的中值点满足 $3\xi^2=1$，为 $1/\sqrt3$，不是 $1/2$。

\begin{exercise}
设 $f$ 在 $[a,b]$ 连续、内部可导，且 $m\le f'(x)\le M$。证明 $m(b-a)\le f(b)-f(a)\le M(b-a)$。
\end{exercise}
\paragraph{解析}
Lagrange 得 $f(b)-f(a)=f'(\xi)(b-a)$，因为 $b-a>0$，将导数上下界乘以区间长度即可。$m,M$ 可为负数，论证仍成立。

\begin{exercise}
证明 $x^3+x-2=0$ 有唯一实根，并求出该根。
\end{exercise}
\paragraph{解析}
代入 $x=1$ 得零，已证明存在。导数为 $3x^2+1>0$，函数严格递增，所以不能有第二个根。唯一根为 $1$。若已找到根，存在部分不必再重复使用介值定理。

\begin{exercise}
设 $f$ 在 $[1,2]$ 连续、内部可导，$f(1)=f(2)=0$。证明存在 $\xi\in(1,2)$，使 $\xi f'(\xi)+2f(\xi)=0$。
\end{exercise}
\paragraph{解析}
$(x^2f)'=x[xf'+2f]$，故选 $F=x^2f$。它连续、内部可导且两端为零。Rolle 得 $0=F'(\xi)=\xi[\xi f'(\xi)+2f(\xi)]$。因 $\xi>1$ 非零，可除去该因子。

\begin{exercise}
设 $f$ 在 $[-1,1]$ 连续、内部二阶可导，$f(-1)=f(0)=f(1)=0$。证明 $f''$ 内部有零点；若只给 $f(-1)=f(1)=0$，结论是否成立？
\end{exercise}
\paragraph{解析}
第一问在 $[-1,0]$、$[0,1]$ 分别用 Rolle 得 $u<0<v$ 且 $f'(u)=f'(v)=0$，再对 $f'$ 在 $[u,v]$ 用 Rolle 即得结论。第二问不成立：$f(x)=1-x^2$ 满足两端为零，但 $f''=-2$ 处处非零。

\begin{exercise}
设 $f$ 在 $[0,1]$ 连续、内部可导，具有单侧端点导数，且 $f'_+(0)>0>f'_-(1)$。证明存在内部点 $c$，使 $f'(c)=0$；说明能否直接断定 $f'$ 连续。
\end{exercise}
\paragraph{解析}
由 $f'_+(0)>0$，靠近零点的内部函数值大于 $f(0)$；由 $f'_-(1)<0$，靠近一的内部函数值也大于 $f(1)$（差商负、分母负）。因此最大值不在两个端点取得，必在内部 $c$ 取得，Fermat 得 $f'(c)=0$。可导性不推出导数连续，上述振荡导数便是反例。

\paragraph{章末自查}
遇到增量先考虑 Lagrange，遇到两种增量的比值考虑 Cauchy；目标是导数组合为零时，逆向检查乘积、商与线性组合的求导。每次应用定理都写出实际的闭区间，确认连续、内部可导和必要的端点关系。证明唯一根要同时完成存在与唯一；求到某个中值点后，不将它与另一次定理产生的点视为同一个点。


\chapter{洛必达法则}
\label{lec:07:chapter}

本章需要函数极限、导数运算法则以及柯西中值定理。我们要解决的问题是：当商的分子、分母都趋于零，或者分母的绝对值趋于无穷时，怎样借助导数比较它们变化的快慢？洛必达法则把一个商的极限问题转化为导数之商的极限问题。它有明确的条件，适合用作计算工具，也可以用来建立对数、幂函数与指数函数的增长比较。

\section{从未定式到导数之商}
\label{lec:07:indeterminate}

\subsection{未定式表达的是信息不足}

设 $x\to0$。仅知道 $f(x)\to0$、$g(x)\to0$，无法决定 $f(x)/g(x)$ 的极限。例如
\[
\frac{x}{x}=1,\qquad
\frac{x^2}{x}=x\to0,\qquad
\frac{x}{x^2}=\frac1x.
\]
第三个商甚至没有双侧极限。这里的 $0/0$ 是一种\emph{未定式类型}，不表示可以对两个零做除法；它提醒我们必须比较趋于零的速度。同样，当 $x\to+\infty$ 时，$x/x$、$x/x^2$ 与 $x^2/x$ 都是 $\infty/\infty$ 型，极限却分别为 $1$、$0$ 与 $+\infty$。

导数给出局部变化率。对于同时趋于零的两个函数，柯西中值定理可以把原商表示成某个中间点处的导数之商；这正是洛必达法则的依据。注意，法则中的新表达式为
\[
\frac{f'(x)}{g'(x)},
\]
而非原商的导数
\[
\left(\frac fg\right)'(x)
=\frac{f'(x)g(x)-f(x)g'(x)}{g(x)^2}.
\]
两者解决的问题不同，不能混用。

\subsection{先判断，后求导}

计算前先求分子、分母各自的极限。若分母趋于非零有限数，直接应用商的极限法则；若分子趋于非零数而分母趋于零，应检查符号和单侧极限。只有形式和条件匹配后，才考虑洛必达法则。

\begin{example}[不属于未定式的商]
求 $\displaystyle\lim_{x\to0}\frac{1+x}{2+x}$，并说明为什么不能分别求导。
\end{example}
\paragraph{思路分析}
分子趋于 $1$，分母趋于 $2$，商的极限法则已经给出答案。分别求导会丢掉这两个常数，改变极限。
\paragraph{解答}
由连续性，原极限为 $1/2$。导数之商恒为 $1/1=1$，与原答案不同。原式既非 $0/0$ 型，也不满足分母绝对值趋于无穷的条件，故不能使用洛必达法则。

\section{零比零型的定理与证明}
\label{lec:07:zero}

\begin{theorem}[有限端点的零比零型洛必达法则]
\label{lec:07:zero-theorem}
设 $a\in\R$，$f,g$ 在某个右侧区间 $(a,a+\delta)$ 内可导，并且对该区间内的每个 $x$ 都有 $g'(x)\ne0$。若
\[
\lim_{x\to a^+}f(x)=\lim_{x\to a^+}g(x)=0,
\qquad
\lim_{x\to a^+}\frac{f'(x)}{g'(x)}=L,
\]
其中 $L$ 为实数、$+\infty$ 或 $-\infty$，则
\[
\lim_{x\to a^+}\frac{f(x)}{g(x)}=L.
\]
将右侧区间改成 $(a-\delta,a)$，相同结论对 $x\to a^-$ 成立。若左右两侧满足条件且得到同一个 $L$，便得到双侧极限。
\end{theorem}

\paragraph{条件怎样理解}
可导性要求在\emph{整个足够小的去心单侧邻域}内成立；只知道 $f'(a),g'(a)$ 存在不足以直接应用本定理。$g'(x)\ne0$ 同样是邻域条件，而非只检查 $g'(a)$。反过来，定理不要求 $f,g$ 在 $a$ 处有定义或可导；端点处补成连续函数就够了。最后，导数之商要有确定的有限或无穷极限；如果它振荡，定理没有给出原商的结论。

\paragraph{证明思路}
若把端点函数值补为零，$f(x)/g(x)$ 就成为以 $a$ 为一个端点的增量之商。柯西中值定理可以在 $a$ 与 $x$ 之间找到 $\xi$，使增量之商等于 $f'(\xi)/g'(\xi)$。关键是先保证分母非零，再证明 $x\to a^+$ 时 $\xi\to a^+$。

\begin{proof}
定义 $F(a)=G(a)=0$，并在 $x>a$ 时令 $F(x)=f(x)$、$G(x)=g(x)$。由两函数趋于零，$F,G$ 在 $a$ 处右连续；它们在 $[a,x]$ 上连续、在 $(a,x)$ 内可导。

先证明 $g(x)\ne0$。若某个 $x>a$ 满足 $g(x)=0$，则 $G(a)=G(x)=0$。罗尔定理会给出某个 $\eta\in(a,x)$ 满足 $g'(\eta)=0$，与假设矛盾。因此原商在该右侧区间有定义。

对 $F,G$ 在 $[a,x]$ 上应用柯西中值定理，存在 $\xi\in(a,x)$，满足
\[
[F(x)-F(a)]G'(\xi)
=[G(x)-G(a)]F'(\xi).
\]
由于 $G(x)-G(a)=g(x)\ne0$ 且 $G'(\xi)=g'(\xi)\ne0$，可除得
\begin{equation}
\frac{f(x)}{g(x)}=\frac{f'(\xi)}{g'(\xi)},
\qquad a<\xi<x.
\label{lec:07:cauchy-ratio}
\end{equation}
由 $0<\xi-a<x-a$，当 $x\to a^+$ 时 $\xi\to a^+$。若 $L$ 有限，对任意 $\varepsilon>0$，在足够小的右侧区间内都有
\[
\left|\frac{f'(t)}{g'(t)}-L\right|<\varepsilon.
\]
当 $x$ 位于此区间时，$\xi$ 也在此区间，故由式 \eqref{lec:07:cauchy-ratio} 得 $|f(x)/g(x)-L|<\varepsilon$。

若 $L=+\infty$，对任意 $M>0$，在足够小的右侧区间内有 $f'(t)/g'(t)>M$，再由同一等式得 $f(x)/g(x)>M$。$L=-\infty$ 的证明把不等号改为小于 $-M$。左侧证明只需交换区间端点；双侧结论由左右极限相同得到。
\end{proof}

\begin{example}[一次求导后停止]
求 $\displaystyle\lim_{x\to0}\frac{\ln(1+x)}x$。
\end{example}
\paragraph{思路分析}
分子、分母均趋于零；对数求导后变成有理式，可直接代入。
\paragraph{解答}
取 $0<|x|<1/2$，则 $1+x>0$，两函数可导，分母导数为 $1\ne0$。又有
\[
\lim_{x\to0}\frac{[\ln(1+x)]'}{x'}
=\lim_{x\to0}\frac1{1+x}=1.
\]
由定理，原极限为 $1$。导数之商已经可直接求极限，不需要继续求导。

\begin{example}[无穷结论与单侧方向]
求 $\displaystyle\lim_{x\to0^+}\frac{\sin x}{x^2}$ 和 $\displaystyle\lim_{x\to0^-}\frac{\sin x}{x^2}$。
\end{example}
\paragraph{思路分析}
虽然两边都是 $0/0$ 型，导数之商含 $1/x$，必须保留趋近方向。
\paragraph{解答}
在 $0<|x|<\pi/4$ 内，分子、分母可导且 $(x^2)'=2x\ne0$。导数之商为
\[
\frac{\cos x}{2x}.
\]
右侧 $\cos x\to1$、$2x\to0^+$，所以该商趋于 $+\infty$；左侧则趋于 $-\infty$。分别应用单侧定理，原商的右极限为 $+\infty$，左极限为 $-\infty$，因此原商没有双侧极限，也不能把答案写成无方向的“无穷”。

\section{分母趋于无穷与无穷远形式}
\label{lec:07:infinity}

\begin{theorem}[分母绝对值趋于无穷的洛必达法则]
\label{lec:07:infinity-theorem}
设 $f,g$ 在 $(a,a+\delta)$ 内可导，$g'(x)\ne0$，且
\[
\lim_{x\to a^+}|g(x)|=+\infty,
\qquad
\lim_{x\to a^+}\frac{f'(x)}{g'(x)}=L,
\]
其中 $L\in\R\cup\{+\infty,-\infty\}$。则 $f(x)/g(x)\to L$。
左侧趋近有相应结论。这里不要求分子本身趋于无穷，甚至不要求分子单独有极限；通常所说的 $\infty/\infty$ 型是本定理的一个常用情形。
\end{theorem}

\paragraph{为何不能照搬零比零型的证明}
此时无法在端点补一个有限函数值。需要在邻域内固定一点 $c$，比较 $f(x)-f(c)$ 与 $g(x)-g(c)$。分母趋于无穷保证固定的 $g(c)$ 相对于 $g(x)$ 可以忽略，这样就能从增量之商回到原商。

\begin{proof}
先设 $L$ 有限。给定 $\varepsilon>0$，缩小右侧区间，使其中处处满足
\[
\left|\frac{f'(t)}{g'(t)}-L\right|<\varepsilon.
\]
在此区间固定 $c$，再取 $a<x<c$。由罗尔定理和 $g'\ne0$ 可知 $g(x)\ne g(c)$。柯西中值定理给出某个 $\xi\in(x,c)$，满足
\[
R_x:=\frac{f(x)-f(c)}{g(x)-g(c)}
=\frac{f'(\xi)}{g'(\xi)},\qquad |R_x-L|<\varepsilon.
\]
当 $x$ 充分接近 $a$ 时，$g(x)\ne0$，且
\begin{align*}
\frac{f(x)}{g(x)}
&=R_x\left(1-\frac{g(c)}{g(x)}\right)+\frac{f(c)}{g(x)},\\
\left|\frac{f(x)}{g(x)}-L\right|
&\le\varepsilon\left|1-\frac{g(c)}{g(x)}\right|
+\left|\frac{f(c)-Lg(c)}{g(x)}\right|.
\end{align*}
因为 $|g(x)|\to\infty$，右端第一项的上极限不超过 $\varepsilon$，第二项趋于零。更具体地，在足够靠近 $a$ 时两项之和小于 $3\varepsilon$。$\varepsilon$ 任意，故原商趋于 $L$。

再设 $L=+\infty$。给定 $M>0$，选取邻域使 $f'(t)/g'(t)>2(M+1)$，并在其中固定 $c$。上述中值定理给出 $R_x>2(M+1)$。又因 $g(c)/g(x)\to0$、$f(c)/g(x)\to0$，当 $x$ 足够接近 $a$ 时，
\[
1-\frac{g(c)}{g(x)}>\frac12,
\qquad \left|\frac{f(c)}{g(x)}\right|<1.
\]
因此 $f(x)/g(x)>M$。$L=-\infty$ 时对 $-f$ 应用已经证明的正无穷情形。左侧的证明相同。
\end{proof}

\subsection{把无穷远化为有限端点}

以上两个定理也适用于 $x\to+\infty$ 或 $x\to-\infty$，此时可导性与 $g'(x)\ne0$ 要在某个尾区间内成立。以 $x\to+\infty$ 为例，令 $t=1/x$，则 $t\to0^+$。设
\[
F(t)=f(1/t),\qquad G(t)=g(1/t).
\]
在 $t>0$ 的足够小区间内，
\[
\frac{F'(t)}{G'(t)}
=\frac{f'(1/t)(-1/t^2)}{g'(1/t)(-1/t^2)}
=\frac{f'(1/t)}{g'(1/t)}.
\]
原定理的条件和导数比极限都得到保留，故可以使用有限端点的结果。$x\to-\infty$ 对应 $t\to0^-$。这一换元也说明：无穷远形式不是另一套求导规则。

\begin{example}[对数与任意正幂的比较]
证明对任意 $p>0$，$\displaystyle\lim_{x\to+\infty}\frac{\ln x}{x^p}=0$。
\end{example}
\paragraph{思路分析}
原式是 $\infty/\infty$ 型。一次求导后，$1/x$ 与 $x^{p-1}$ 的商可以合成一个负幂，无须对非整数幂反复求导。
\paragraph{解答}
在 $x>1$ 内，两函数可导，$(x^p)'=px^{p-1}>0$。导数之商满足
\[
\frac{(\ln x)'}{(x^p)'}
=\frac{1/x}{px^{p-1}}=\frac1{px^p}\to0.
\]
由无穷远形式的定理，结论成立。若把 $x$ 换成正整数 $n$，由函数极限的归结原则得到 $\ln n/n^p\to0$；不能对离散指标 $n$ 直接求导。

\begin{example}[幂函数与指数函数的比较]
设 $\alpha>0$、$b>1$，证明 $\displaystyle\lim_{x\to+\infty}\frac{x^\alpha}{b^x}=0$。
\end{example}
\paragraph{思路分析}
非整数幂经过多次求导后可能变成负幂，逐步判断会较繁。先用更大的整数幂控制它，再对整数幂进行有限次求导，计算更清楚。
\paragraph{解答}
取整数 $m>\alpha$。当 $x\ge1$ 时，
\[
0\le\frac{x^\alpha}{b^x}\le\frac{x^m}{b^x}.
\]
对右端，分母的各阶导数为 $(\ln b)^j b^x>0$。在求导 $j$ 次之后（$0\le j<m$），商为
\[
\frac{m!}{(m-j)!}\frac{x^{m-j}}{(\ln b)^j b^x},
\]
分子与分母均趋于 $+\infty$，可导性和下一次的分母导数非零都成立。求导 $m$ 次后的商为
\[
\frac{m!}{(\ln b)^m b^x}\to0.
\]
从最后一步向前逐次应用洛必达法则，得到 $x^m/b^x\to0$，再由夹逼定理得到所求结论。

\paragraph{方法回顾}
比较结论“对数慢于正幂，正幂慢于底数大于 $1$ 的指数”有明确参数条件。本章的证明只涉及这些函数，不把它无条件推广到任意增长函数，也不由此直接处理阶乘。

\section{连续使用：每一步都是一次新的判断}
\label{lec:07:repeated}

一次求导后仍为 $0/0$ 或满足分母趋于无穷的条件，可以考虑再次使用；但必须重新检查新的分母导数不为零、函数可导以及导数比的极限。不能仅凭第一次合法就把后面所有求导都视为合法。

书写上最好先算出最后一个容易判断的极限，再由后往前援引定理。这样每次使用所需的“导数比极限存在”都有依据。下面的等式链用于简洁展示，随后的条件检查说明每个等号成立的理由。

\begin{example}[两次消去]
求 $\displaystyle\lim_{x\to0}\frac{\ee^x-1-x}{1-\cos x}$。
\end{example}
\paragraph{思路分析}
指数函数减去常数项和线性项后，一次求导仍有消去；分母一次求导是 $\sin x$，再求导后趋于 $1$。
\paragraph{解答}
选择 $0<|x|<\pi/4$。第一次原式为 $0/0$ 型，分母导数为 $\sin x\ne0$；导数之商为 $(\ee^x-1)/\sin x$，它仍是 $0/0$ 型。第二次使用时，新的分母导数为 $\cos x>0$。所有涉及的函数均可导，而
\[
\lim_{x\to0}\frac{\ee^x}{\cos x}=1.
\]
先对 $(\ee^x-1)/\sin x$ 应用定理，得到其极限为 $1$；再对原商应用定理，得到
\[
\lim_{x\to0}\frac{\ee^x-1-x}{1-\cos x}
=\lim_{x\to0}\frac{\ee^x-1}{\sin x}
=\lim_{x\to0}\frac{\ee^x}{\cos x}=1.
\]

\begin{example}[三阶消去与及时停止]
求 $\displaystyle\lim_{x\to0}\frac{x-\sin x}{x^3}$。
\end{example}
\paragraph{思路分析}
求导后分子依次为 $1-\cos x$、$\sin x$；第二次求导后已能使用 $\sin x/x\to1$。若选第三次求导，也要重新核对条件。
\paragraph{解答}
原式与第一次导数之商都为 $0/0$ 型。在 $x\ne0$ 的足够小邻域内，前两次所需的分母导数分别为 $3x^2\ne0$、$6x\ne0$。于是
\[
\lim_{x\to0}\frac{x-\sin x}{x^3}
=\lim_{x\to0}\frac{1-\cos x}{3x^2}
=\lim_{x\to0}\frac{\sin x}{6x}
=\frac16.
\]
最后一步用已经建立的重要极限。也可第三次使用洛必达法则：$\sin x/(6x)$ 为 $0/0$ 型，分母导数 $6\ne0$，而 $\cos x/6\to1/6$。两种停止方式给出同一结果。

\section{把其他未定式化成商}
\label{lec:07:transforms}

\subsection{乘积与差：先做代数整理}

若 $u\to0$、$|v|\to\infty$，乘积 $uv$ 为 $0\cdot\infty$ 型。在相应因子非零的邻域内，可以写成
\[
uv=\frac{v}{1/u}\qquad\text{或}\qquad uv=\frac{u}{1/v}.
\]
第一种通常成为分母趋于无穷的商，第二种通常成为 $0/0$ 型。选择哪一种，应看求导后哪种更简单，而不是固定把某个因子移到分母。

\begin{example}[幂与对数的乘积]
设 $p>0$，求 $\displaystyle\lim_{x\to0^+}x^p\ln x$。
\end{example}
\paragraph{思路分析}
$x^p\to0$、$\ln x\to-\infty$。将 $x^p$ 放入分母的倒数，求导后可约成 $x^p$。
\paragraph{解答}
在 $x>0$ 时有
\[
x^p\ln x=\frac{\ln x}{x^{-p}}.
\]
分母趋于 $+\infty$，两函数可导，$(x^{-p})'=-px^{-p-1}\ne0$；导数之商为
\[
\frac{1/x}{-px^{-p-1}}=-\frac{x^p}{p}\to0.
\]
所以原极限为 $0$。原乘积在 $0<x<1$ 时为负，趋于 $0$ 与这一符号相符。

差式中的 $\infty-\infty$ 通常先通分、提取公因子或有理化。不能对两个无穷分别求极限后相减，也不能直接“对差使用洛必达”。

\begin{example}[通分保留相消后的信息]
求 $\displaystyle\lim_{x\to0}\left(\frac1x-\frac1{\ee^x-1}\right)$。
\end{example}
\paragraph{思路分析}
两个倒数都有无穷量，相减后的主导部分会消去。通分使消去显现在分子 $\ee^x-1-x$ 中；分母则可拆出一个已知极限因子，减少求导次数。
\paragraph{解答}
在 $x\ne0$ 时，$\ee^x-1\ne0$，故
\begin{align*}
\frac1x-\frac1{\ee^x-1}
&=\frac{\ee^x-1-x}{x(\ee^x-1)}\\
&=\frac{\ee^x-1-x}{x^2}\cdot\frac{x}{\ee^x-1}.
\end{align*}
第一个因子为 $0/0$ 型，分母导数 $2x\ne0$；其导数之商 $(\ee^x-1)/(2x)$ 仍为 $0/0$ 型，下一分母导数为 $2\ne0$。所以
\[
\lim_{x\to0}\frac{\ee^x-1-x}{x^2}
=\lim_{x\to0}\frac{\ee^x-1}{2x}
=\lim_{x\to0}\frac{\ee^x}{2}=\frac12.
\]
第二个因子由一次洛必达（分母导数 $\ee^x>0$）得到极限 $1$。因此原极限为 $1/2$。

\paragraph{方法回顾}
也可在通分后对整个商求导，但分母会产生多个项。先把 $x(\ee^x-1)$ 写成 $x^2$ 与 $(\ee^x-1)/x$ 的乘积，只在乘除因子中使用已知极限，保留了分子中的相消信息。

\begin{example}[有理化比求导更短]
求 $\displaystyle\lim_{x\to+\infty}\bigl(\sqrt{x^2+x}-x\bigr)$。
\end{example}
\paragraph{思路分析}
原式是 $\infty-\infty$ 型。根式与多项式相减时，共轭乘积可以直接消去最高次项。
\paragraph{解答}
对 $x>0$，
\[
\sqrt{x^2+x}-x
=\frac{x}{\sqrt{x^2+x}+x}
=\frac1{\sqrt{1+1/x}+1}\to\frac12.
\]
这里没有必要使用洛必达法则。有理化后的商已经能由连续性求极限。

\subsection{幂指式：取对数后恢复指数}
\label{lec:07:powers}

对于 $u(x)^{v(x)}$，先检查底数在所讨论邻域内为正，从而有
\[
u(x)^{v(x)}=\exp\bigl(v(x)\ln u(x)\bigr).
\]
求出 $H(x)=v(x)\ln u(x)$ 的极限后再恢复指数：若 $H\to h\in\R$，原式趋于 $\ee^h$；若 $H\to+\infty$ 或 $-\infty$，原式分别趋于 $+\infty$ 或 $0$。常见的 $1^\infty$、$0^0$、$\infty^0$ 都可转成乘积，再选合适的商。

\begin{example}[三种幂指型的基本模型]
求以下极限，其中 $a,b\in\R$：
\[
\lim_{x\to0}(1+ax)^{b/x},\qquad
\lim_{x\to0^+}x^x,\qquad
\lim_{x\to+\infty}x^{1/x}.
\]
\end{example}
\paragraph{思路分析}
三题的底数最终都为正，取对数后分别得到对数差商、幂与对数的乘积、对数与幂的商。后两种已经有结论。
\paragraph{解答}
第一题中，取足够小的邻域使 $1+ax>0$。若 $a=0$ 或 $b=0$，原式恒为 $1=\ee^{ab}$。其余情形下，
\[
\ln\bigl((1+ax)^{b/x}\bigr)
=b\frac{\ln(1+ax)}x.
\]
这是 $0/0$ 型，分母导数为 $1$，而导数之商 $a/(1+ax)\to a$。所以对数趋于 $ab$，原极限为 $\ee^{ab}$。

第二题中，$\ln(x^x)=x\ln x\to0$，因此 $x^x\to1$。这说明底数趋于零、指数也趋于零时，不能仅由底数判断幂的极限。

第三题中，$\ln(x^{1/x})=(\ln x)/x\to0$，因此 $x^{1/x}\to1$。三种结果都来自\emph{整个指数}的极限。

\begin{example}[底数趋于一时的精度]
求 $\displaystyle\lim_{x\to0}\left(\frac{\sin x}{x}\right)^{1/x^2}$。
\end{example}
\paragraph{思路分析}
底数趋于 $1$，但指数趋于无穷；直接把底数换成 $1$ 会丢失被指数放大的误差。取对数后，要把 $\ln(\sin x/x)$ 与 $\sin x/x-1$ 联系起来，再使用三阶消去例题。
\paragraph{解答}
在 $0<|x|<\pi/2$ 内，$\sin x/x>0$。设 $u(x)=\sin x/x-1$，则 $u\to0$。并且此区间内 $u\ne0$：对 $0<x<\pi/2$，
$x-\sin x$ 的导数为 $1-\cos x>0$，从而 $\sin x<x$；负半轴由偶性得到相同的比值不等式。于是
\begin{align*}
\frac{\ln(\sin x/x)}{x^2}
&=\frac{\ln(1+u(x))}{u(x)}\cdot\frac{\sin x-x}{x^3}\\
&\longrightarrow 1\cdot\left(-\frac16\right)=-\frac16.
\end{align*}
第一因子用对数差商的极限及复合极限，第二因子用本章三阶消去的结果。因此所求极限为 $\ee^{-1/6}$。全部计算只用重要极限和洛必达法则，没有使用泰勒公式。

\begin{example}[整体取对数，避免不同步取极限]
求 $\displaystyle\lim_{x\to+\infty}\frac{(1+1/x)^{x^2}}{\ee^x}$。
\end{example}
\paragraph{思路分析}
把 $(1+1/x)^x\to\ee$ 先代入，再保留外层指数 $x$，会放大被省略的误差。应对整个商取对数，把两个大的指数放到同一个差式中。
\paragraph{解答}
原式为正，其对数为
\[
H(x)=x^2\ln(1+1/x)-x.
\]
令 $t=1/x\to0^+$，得到
\[
H(x)=\frac{\ln(1+t)-t}{t^2}.
\]
该商为 $0/0$ 型，函数在 $t>0$ 可导，分母导数 $2t\ne0$。导数之商可以先化简：
\[
\frac{1/(1+t)-1}{2t}
=-\frac1{2(1+t)}\to-\frac12.
\]
故 $H(x)\to-1/2$，原极限为 $\ee^{-1/2}$。

\paragraph{易错说明}
虽然 $\ln(1+t)\sim t$，但在 $\ln(1+t)-t$ 中直接替换会把分子误判为零。相消后的量还要除以 $t^2$，必须保留足以确定这个差的精度。化为 $0/0$ 型后，一次求导并化简已经提供所需精度。

\section{使用边界与方法选择}
\label{lec:07:boundaries}

\subsection{导数比没有极限，原商仍可能有极限}

洛必达法则是一个充分条件定理，不能反向使用。考虑
\[
f(x)=x+\sin x,\qquad g(x)=x,\qquad x\to+\infty.
\]
原商可写成 $1+\sin x/x$，由 $|\sin x/x|\le1/x$ 得到极限 $1$。但导数之商 $1+\cos x$ 沿 $x=2k\pi$ 为 $2$，沿 $x=(2k+1)\pi$ 为 $0$，没有极限。这不是原商发散的证据，而是此条求导路径不能给出结论。

有限端点也有同样现象。令 $f(x)=x+x^2\sin(1/x)$（$x\ne0$）、$g(x)=x$。则
\[
\frac{f(x)}{g(x)}=1+x\sin(1/x)\to1,
\]
而
\[
\frac{f'(x)}{g'(x)}
=1+2x\sin(1/x)-\cos(1/x)
\]
没有极限。分子、分母趋于零和邻域可导，并不自动保证导数比极限存在。

\subsection{分母导数非零也要单独核对}

对
\[
\frac{x-\sin x}{x+\cos x}\qquad(x\to+\infty),
\]
若分别求导，会得到 $(1-\cos x)/(1-\sin x)$。除了振荡问题，新的分母 $1-\sin x$ 在 $x=\pi/2+2k\pi$ 处为零，因此原分母的导数在任何尾区间都不能处处非零，定理条件已不成立。正确做法是同除以 $x$：
\[
\frac{x-\sin x}{x+\cos x}
=\frac{1-\sin x/x}{1+\cos x/x}\to1.
\]
分母自身在足够大的 $x$ 上非零，只能保证原商有定义，不能代替分母\emph{导数}非零的要求。

\subsection{求导不一定使问题简化}

对于 $\ee^x$ 与有界振荡项相加的商，例如
\[
\frac{\ee^x-\sin x}{\ee^x+\cos x},
\]
同除以 $\ee^x$ 就得到极限 $1$。求导后仍有指数和三角函数，多次求导会继续出现类似结构，计算没有缩短。这里不能笼统地说“含振荡项就不能用洛必达”：指数占主导时，导数比也可能有极限。应比较具体表达式，选能直接闭合论证的方法。

\begin{center}
\begin{tabular}{p{0.29\textwidth}p{0.59\textwidth}}
\toprule
题目特征 & 优先动作及检查 \\
\midrule
分母趋于非零有限数 & 直接用极限四则运算；无需求导。 \\
$0/0$ 型或分母趋于无穷 & 核对邻域可导、分母导数非零；看导数比是否更简单。 \\
根式之差 & 先有理化，再判断所得商。 \\
若干大项或小项相消 & 先整理成一个差或一个商，保留相消后的信息。 \\
幂指型且底数为正 & 整体取对数，计算指数极限后恢复指数。 \\
主导项加有界振荡 & 同除以主导项，使用有界量与无穷小的乘积。 \\
\bottomrule
\end{tabular}
\end{center}

\paragraph{一套可执行的检查顺序}
先确定定义域和趋近方向，再判断原式的类型。若选择洛必达，逐项写出可导区间和分母导数不为零的理由；计算导数之商后，先做约分、因式整理等代数化简。若新商仍满足条件且求导能降低复杂度，再重复；若振荡、循环或出现更繁的项，改用已有的极限方法。最后检查答案的符号、左右侧是否一致，以及参数的临界值是否被单独处理。

\section{分层练习与解析}
\label{lec:07:exercises}

以下题目均为本章自编练习，不标为真实考试题。

\begin{exercise}[基础巩固]
设 $a\in\R$，求 $\displaystyle\lim_{x\to0}\frac{\ee^{ax}-1}{x}$，并写出适用条件。
\end{exercise}
\paragraph{解析}
分子、分母均趋于零，两函数在实轴可导，分母导数为 $1\ne0$，导数之商 $a\ee^{ax}\to a$。故极限为 $a$，包括 $a=0$ 时分子恒为零的情形。

\begin{exercise}[基础巩固]
求 $\displaystyle\lim_{x\to0}\frac{\ln(1+x)-x}{x^2}$。
\end{exercise}
\paragraph{解析}
取 $0<|x|<1/2$，原式是 $0/0$ 型，两函数可导，分母导数 $2x\ne0$。导数之商为
\[
\frac{1/(1+x)-1}{2x}=-\frac1{2(1+x)}\to-\frac12.
\]
所以极限为 $-1/2$。先化简就能结束，不必再次求导。

\begin{exercise}[方法迁移]
设 $p>0$、$q>0$，求 $\displaystyle\lim_{x\to0^+}x^p|\ln x|^q$。
\end{exercise}
\paragraph{解析}
令 $t=-\ln x\to+\infty$，则 $x=\ee^{-t}$，原式等于 $t^q/\ee^{pt}$。幂与指数比较例题适用：$q>0$，指数底数 $\ee^p>1$。因此极限为 $0$。该换元把非整数次的对数幂变成了普通幂，避免直接对复杂复合式反复求导。

\begin{exercise}[方法迁移]
求 $\displaystyle\lim_{x\to0}\left(\frac1{\sin x}-\frac1x\right)$。
\end{exercise}
\paragraph{解析}
在 $0<|x|<\pi/2$ 内，
\[
\frac1{\sin x}-\frac1x
=\frac{x-\sin x}{x\sin x}
=\frac{x-\sin x}{x^3}\cdot\frac{x}{\sin x}\cdot x.
\]
三个因子的极限依次为 $1/6$、$1$、$0$，故所求极限为 $0$。通分后利用已知极限因子，避免扩大求导表达式。

\begin{exercise}[综合提升]
求 $\displaystyle\lim_{x\to0}\frac{(1+x)^{1/x}-\ee}{x}$。
\end{exercise}
\paragraph{解析}
取 $0<|x|<1/2$，底数为正。令
\[
h(x)=\frac{\ln(1+x)}x-1
=\frac{\ln(1+x)-x}{x}.
\]
由对数差商，$h(x)\to0$，而前面的基础练习给出
\[
\frac{h(x)}x=\frac{\ln(1+x)-x}{x^2}\to-\frac12.
\]
还需从指数误差恢复原误差。$\ee^t-1$ 与 $t$ 都趋于零，分母导数为 $1$，导数之商 $\ee^t\to1$，故 $(\ee^t-1)/t\to1$。在所取去心邻域中 $h(x)\ne0$：$\ln(1+x)<x$ 对 $x\ne0$ 成立，可由 $x-\ln(1+x)$ 的导数 $x/(1+x)$ 及零点处的最小值证明。因此
\[
\frac{(1+x)^{1/x}-\ee}{x}
=\ee\frac{\ee^{h(x)}-1}{h(x)}\frac{h(x)}x
\longrightarrow-\frac{\ee}{2}.
\]
只知道 $(1+x)^{1/x}\to\ee$ 不足以求此差商；需要同时控制趋近 $\ee$ 的一阶误差。

\begin{exercise}[边界辨析]
设 $x\to0$，考察 $\displaystyle\frac{x^2\sin(1/x)}x$。能否因导数之商没有极限而判定原商没有极限？
\end{exercise}
\paragraph{解析}
原商为 $x\sin(1/x)$，其绝对值不超过 $|x|$，因此极限为 $0$。导数之商为
\[
2x\sin(1/x)-\cos(1/x).
\]
沿 $x_k=1/(2k\pi)$，它等于 $-1$；沿 $y_k=1/((2k+1)\pi)$，它等于 $1$，所以导数比没有极限。只能得出本次洛必达法则不能给出答案，不能得出原商发散。

\paragraph{章末自查}
应能回答：为什么定理检查的是邻域内的可导性？为什么 $0/0$ 型证明中的中间点会趋于端点？为什么分母趋于无穷时固定常数可以忽略？重复求导时要重新检查哪些条件？取对数后为什么必须先求整个指数的极限？如果导数比没有极限，下一步应怎样重新选择方法？


\chapter{泰勒公式}
\label{lec:08:chapter}
本章以导数、中值定理和无穷小记号为前置知识。导数用直线描述局部变化，泰勒公式则用多项式记录更多阶的变化。我们要回答三个问题：多项式系数如何确定，舍去的部分有多小，怎样据此计算极限或证明估计。

\section{从线性近似到多项式近似}
可导性给出 $f(a+h)=f(a)+f'(a)h+o(h)$。如果两个表达式的一次项相同，它们相减后的一次信息消失，必须研究二次乃至更高次的变化。例如 $\ee^h-1-h$ 的大小不能由 $\ee^h-1\sim h$ 决定。
设 $P_n(a+h)=c_0+c_1h+\cdots+c_nh^n$，要求它在 $a$ 处与 $f$ 的各阶导数一致。因为 $P_n^{(k)}(a)=k!c_k$，系数只能取为 $c_k=f^{(k)}(a)/k!$。
\begin{definition}[泰勒多项式与余项]
具有相应导数的函数 $f$ 在 $a$ 处的 $n$ 阶泰勒多项式是
\[
T_n(x)=\sum_{k=0}^n\frac{f^{(k)}(a)}{k!}(x-a)^k,
\qquad R_n(x)=f(x)-T_n(x).
\]
取 $a=0$ 时称为麦克劳林多项式。这里 $f^{(0)}=f$。
\end{definition}
余项是公式的一部分。写出多项式以后，仍需证明余项达到所需精度。$R_n=o((x-a)^n)$ 表示其与 $(x-a)^n$ 的商趋于零；它没有直接给出固定区间上的数值界。
\begin{figure}[htbp]
\centering
\begin{tikzpicture}[x=2.1cm,y=1.7cm]
\draw[->] (-1.3,0)--(1.35,0) node[right] {$x$};
\draw[->] (0,-0.2)--(0,3) node[above] {$y$};
\draw[thick,domain=-1.1:1.05,samples=70] plot (\x,{exp(\x)});
\draw[dashed,domain=-1.1:1.05,samples=2] plot (\x,{1+\x});
\draw[densely dotted,thick,domain=-1.1:1.05,samples=50] plot (\x,{1+\x+\x*\x/2});
\fill (0,1) circle (1pt) node[left] {$1$};
\node[right] at (1.05,2.86) {$\ee^x$};
\node[right] at (1.05,2.60) {$T_2$};
\node[right] at (1.05,2.05) {$T_1$};
\end{tikzpicture}
\caption{指数函数与零点的一次、二次泰勒多项式。三者在零点有相同函数值与斜率，二次多项式还匹配二阶变化。}
\label{lec:08:fig-approx}
\end{figure}

\section{两种余项及其证明}
\subsection{局部比较使用皮亚诺余项}
\begin{theorem}[带皮亚诺余项的泰勒公式]
设 $n\ge1$，$f$ 在 $a$ 的某邻域内具有直到 $n-1$ 阶导数，且 $f^{(n-1)}$ 在 $a$ 可导。则
\[
f(a+h)=\sum_{k=0}^n\frac{f^{(k)}(a)}{k!}h^k+o(h^n),\qquad h\to0.
\]
这是本章使用的一组充分条件；最高阶导数无须在邻域内连续。
\end{theorem}
\begin{proof}
令 $F(h)=f(a+h)-T_n(a+h)$。有 $F^{(k)}(0)=0$，$0\le k\le n$。最高阶导数在零点存在的含义给出
\[
F^{(n-1)}(h)=F^{(n-1)}(0)+F^{(n)}(0)h+o(h)=o(h).
\]
给定 $\varepsilon>0$，取 $\delta>0$，使 $0<|t|<\delta$ 时 $|F^{(n-1)}(t)|\le\varepsilon|t|$。对于 $n\ge2$，在连接 $0,h$ 的区间上对 $F^{(n-2)}$ 应用拉格朗日中值定理：某 $t$ 位于其间，且
\[
|F^{(n-2)}(h)|=|hF^{(n-1)}(t)|\le\varepsilon |h|^2.
\]
依次向下应用同一定理，得到 $|F(h)|\le\varepsilon |h|^n$。这些导函数在零点连续，在区间内部可导，所以每次中值定理的条件均满足。$n=1$ 时第一步已是结论。因此 $F(h)/h^n\to0$。
\end{proof}
\subsection{误差与区间不等式使用拉格朗日余项}
\begin{theorem}[带拉格朗日余项的泰勒公式]
设 $a\ne x$，$f$ 在连接 $a,x$ 的闭区间上属于 $C^n$，且在内部存在 $n+1$ 阶导数。则存在严格位于 $a,x$ 之间的 $\xi$，使
\[
f(x)=T_n(x)+\frac{f^{(n+1)}(\xi)}{(n+1)!}(x-a)^{n+1}.
\]
这里 $n=0$ 也是允许的，公式成为拉格朗日中值定理。
\end{theorem}
\begin{proof}
固定 $x$，令 $K=(f(x)-T_n(x))/(x-a)^{n+1}$，构造
\[
G(t)=f(t)-T_n(t)-K(t-a)^{n+1}.
\]
选取 $K$ 的目的是令 $G(x)=0$；同时 $G(a)=G'(a)=\cdots=G^{(n)}(a)=0$。先对两个零点 $a,x$ 用罗尔定理，得到内部的 $G'$ 零点，再把它与 $a$ 处的 $G'$ 零点配对。如此反复，得到 $G^{(n+1)}(\xi)=0$。各次使用所需的闭区间连续性来自 $f\in C^n$，最后一次的内部可导性来自 $f^{(n+1)}$ 的存在。于是
\[
0=f^{(n+1)}(\xi)-(n+1)!K,
\]
代回 $K$ 即得公式。$x<a$ 时把区间端点按大小排列，论证相同。
\end{proof}
若 $|f^{(n+1)}(t)|\le M$ 对整段成立，则
\[
|R_n(x)|\le\frac{M}{(n+1)!}|x-a|^{n+1}.
\]
同样，余项的符号由导数符号和 $(x-a)^{n+1}$ 的符号共同决定。不要漏掉增量为负的情形。

\begin{example}[指数近似及误差]用 $1+x+x^2/2$ 近似 $\ee^x$，给出 $|x|\le0.1$ 时的统一误差界。\end{example}
\paragraph{思路分析}数值近似需要可计算的上界。指数函数三阶导数仍是自身，在连接零点与 $x$ 的区间内不超过 $\ee^{0.1}$。
\paragraph{解答}由拉格朗日余项，
\[
\ee^x=1+x+\frac{x^2}{2}+\frac{\ee^\xi x^3}{6},\qquad
|R_2(x)|\le\frac{\ee^{0.1}}6|x|^3\le\frac{\ee^{0.1}}{6000}.
\]
当 $x\ge0$ 时该近似偏小；当 $x<0$ 时偏大。绝对误差界同时覆盖两侧。

\section{常用展开怎样得到}
以下均是 $x\to0$ 时的有限阶公式，不是省略了收敛论证的无穷级数。
指数函数的每阶导数都是 $\ee^x$，正弦、余弦导数每四阶循环。因此
\begin{align*}
\ee^x&=1+x+\frac{x^2}{2!}+\frac{x^3}{3!}+\frac{x^4}{4!}+o(x^4),\\
\sin x&=x-\frac{x^3}{3!}+\frac{x^5}{5!}+o(x^5),\\
\cos x&=1-\frac{x^2}{2!}+\frac{x^4}{4!}+o(x^4).
\end{align*}
正弦偶次系数、余弦奇次系数为零，可直接由各阶导数在零点的值验证。
对 $x>-1$，逐次求导得
\[
\frac{\dd^k}{\dd x^k}\ln(1+x)=\frac{(-1)^{k-1}(k-1)!}{(1+x)^k},\quad k\ge1,
\]
所以
\[
\ln(1+x)=x-\frac{x^2}{2}+\frac{x^3}{3}-\frac{x^4}{4}+o(x^4).
\]
对固定实数 $\alpha$，在 $1+x>0$ 的邻域内，幂函数逐次求导得到下降乘积，故
\[
(1+x)^\alpha=1+\alpha x+\frac{\alpha(\alpha-1)}2x^2
+\frac{\alpha(\alpha-1)(\alpha-2)}6x^3+o(x^3).
\]
令 $\alpha=-1$ 可得 $1/(1+x)=1-x+x^2-x^3+o(x^3)$；令 $\alpha=1/2$ 可得 $\sqrt{1+x}=1+x/2-x^2/8+x^3/16+o(x^3)$。

\begin{example}[正切展开的系数比较]求 $\tan x$ 到五次项的展开。\end{example}
\paragraph{思路分析}直接计算五阶导数繁琐。零点附近 $\cos x\ne0$，正切光滑且为奇函数，可以设奇次多项式，再用 $\tan x\cos x=\sin x$ 确定系数。
\paragraph{解答}设 $\tan x=ax+bx^3+cx^5+o(x^5)$。乘以余弦展开，得到
\[
ax+\left(b-\frac a2\right)x^3
+\left(c-\frac b2+\frac a{24}\right)x^5+o(x^5)
=x-\frac{x^3}6+\frac{x^5}{120}+o(x^5).
\]
依次比较系数，$a=1$，$b=1/3$，$c=1/120+1/6-1/24=2/15$。于是
\[
\tan x=x+\frac{x^3}3+\frac{2x^5}{15}+o(x^5).
\]
系数比较的依据将在下一节证明。

反三角函数可从其导数在零点处的值求出。例如 $g=\arctan x$ 有 $g'=1/(1+x^2)$，该有理函数的各阶导数给出 $g'(0)=1$、$g'''(0)=-2$、$g^{(5)}(0)=24$；奇偶性使偶数阶导数为零。因此
\[
\arctan x=x-\frac{x^3}3+\frac{x^5}5+o(x^5).
\]
类似地，$(\arcsin x)'=(1-x^2)^{-1/2}$ 的二项式系数给出
\[
\arcsin x=x+\frac{x^3}6+\frac{3x^5}{40}+o(x^5).
\]
这里使用导函数的泰勒系数确定原函数的导数值，没有对任意 $o$ 余项直接求导。

\section{展开阶数与余项运算}
\label{lec:08:order}
\subsection{先判断目标精度}
若分母与 $cx^m$ 等价（$c\ne0$），且要求有限非零极限，通常须把整个分子算到 $x^m$ 并留下 $o(x^m)$。若低次项未消去，商可能发散；若到 $m$ 次仍全部消去，商可能为零。先通分、提因子或取对数，往往能减少所需展开项。
\begin{proposition}[有限展开系数的唯一性]
若 $\sum_{k=0}^na_kx^k-\sum_{k=0}^nb_kx^k=o(x^n)$，则 $a_k=b_k$ 对 $0\le k\le n$ 全部成立。
\end{proposition}
\begin{proof}若有不同，取最低次不同系数的编号 $j$。两边除以 $x^j$，左边趋于非零常数 $a_j-b_j$，右边为 $o(x^{n-j})$，趋于零，矛盾。\end{proof}
\subsection{余项如何传播}
同阶 $o$ 相加仍是同阶 $o$。若 $A=O(x^p)$、$B=o(x^q)$，则 $AB=o(x^{p+q})$；这是因为 $A/x^p$ 有界而 $B/x^q\to0$。除以 $x^r$ 后，$o(x^m)/x^r=o(x^{m-r})$。作除法时须先检查除数的首项不为零。
复合时若 $u(x)=O(x^p)\to0$，则外层 $o(u^q)$ 可写为 $o(x^{pq})$：写成 $u^q\varepsilon(u)$，其中 $\varepsilon(u)\to0$ 即可。若 $u$ 可能为零，把余项在零点定义为零。若只知 $u\to0$ 而不知道其阶，不能擅自换成 $o(x^q)$。

\begin{example}[抵消后的四阶极限]计算 $\displaystyle\lim_{x\to0}\frac{\cos x-\ee^{-x^2/2}}{x^4}$。\end{example}
\paragraph{思路分析}常数项和二次项会抵消。余弦须到四次，指数的内部变量是二阶，外层只须到二次。
\paragraph{解答}
\[
\cos x=1-\frac{x^2}2+\frac{x^4}{24}+o(x^4),\qquad
\ee^{-x^2/2}=1-\frac{x^2}2+\frac{x^4}8+o(x^4).
\]
两式相减得 $-x^4/12+o(x^4)$，除以 $x^4$，极限为 $-1/12$。单独把两项都替换成 $1$ 会丢失全部需要的信息。

\begin{example}[乘积各因子的精度]计算 $\displaystyle\lim_{x\to0}\frac{\ee^x\sin x-x-x^2}{x^3}$。\end{example}
\paragraph{思路分析}正弦从一次项起，故指数到二次项即可；指数有常数项，故正弦必须到三次。
\paragraph{解答}
\begin{align*}
\ee^x\sin x
&=\left(1+x+\frac{x^2}2+o(x^2)\right)
\left(x-\frac{x^3}6+o(x^3)\right)\\
&=x+x^2+\left(\frac12-\frac16\right)x^3+o(x^3).
\end{align*}
指数余项乘以 $O(x)$ 成为 $o(x^3)$。减去 $x+x^2$ 后极限为 $1/3$。

\begin{example}[复合展开的三次项]把 $\sqrt{1+\sin x}$ 展开到三次项。\end{example}
\paragraph{思路分析}$u=\sin x\sim x$，目标为 $x^3$，外层必须保留 $u^3$。把外层仅写到 $u^2$ 再丢弃 $O(u^3)$，不能确定三次系数。
\paragraph{解答}由 $u=x-x^3/6+o(x^3)$，有 $u^2=x^2+o(x^3)$、$u^3=x^3+o(x^3)$。因此
\begin{align*}
\sqrt{1+\sin x}
&=1+\frac u2-\frac{u^2}8+\frac{u^3}{16}+o(u^3)\\
&=1+\frac x2-\frac{x^2}8+\left(-\frac1{12}+\frac1{16}\right)x^3+o(x^3)\\
&=1+\frac x2-\frac{x^2}8-\frac{x^3}{48}+o(x^3).
\end{align*}

\begin{example}[商的递次系数]求 $x/(\ee^x-1)$ 的二阶展开，并在零点连续补值。\end{example}
\paragraph{思路分析}先消去分子分母共有的 $x$，使分母常数项为 $1$，再用乘法反求商的系数。
\paragraph{解答}零点外，
\[
\frac{x}{\ee^x-1}=\frac1{1+x/2+x^2/6+o(x^2)}.
\]
设商为 $A+Bx+Cx^2+o(x^2)$，比较乘积与 $1$，得到 $A=1$、$B+A/2=0$、$C+B/2+A/6=0$。于是
\[
\frac{x}{\ee^x-1}=1-\frac x2+\frac{x^2}{12}+o(x^2).
\]
取零点函数值 $1$ 可连续补值。系数计算也可用 $1/(1+v)=1-v+v^2+o(v^2)$，其适用性由 $v\to0$ 保证。

\begin{example}[幂指表达式先取对数]求 $\displaystyle\lim_{n\to\infty}n\bigl[(1+1/n)^n-\ee\bigr]$。\end{example}
\paragraph{思路分析}先让指数变成可展开的量。不能先假定原极限存在后写它的对数。
\paragraph{解答}令 $t=1/n$。由对数三阶公式，
\[
\frac{\ln(1+t)}t=1-\frac t2+O(t^2).
\]
所以 $(1+t)^{1/t}=\ee\exp(-t/2+O(t^2))=\ee(1-t/2+O(t^2))$。减去 $\ee$ 后除以 $t$ 得极限 $-\ee/2$。正底数条件在 $t>0$ 时成立。

\section{系数、高阶导数与极值}
\label{lec:08:derivative}
满足皮亚诺定理条件时，由系数唯一性，若 $f(a+h)=\sum_{k=0}^nc_kh^k+o(h^n)$，则 $f^{(k)}(a)=k!c_k$。先检查导数存在，再读取系数。任意渐近式本身不能保证高阶可导。
例如令 $f(0)=0$、$f(x)=x^3\sin(1/x^2)$（$x\ne0$）。虽有 $f(x)=o(x^2)$ 且 $f'(0)=0$，但
\[
f'(x)=3x^2\sin(1/x^2)-2\cos(1/x^2)
\]
使 $[f'(x)-f'(0)]/x$ 无极限，故 $f''(0)$ 不存在。不能从零二次系数断言二阶导数为零，也不能直接对 $o(x^2)$ 求导。

\begin{example}[由隐式关系确定展开]设零点附近的函数 $y=y(x)$ 在 $0$ 连续，满足 $y=x\ee^y$ 且 $y(0)=0$，求其三阶展开。\end{example}
\paragraph{思路分析}先验证这个关系确实确定一个光滑函数，再设有限展开。把 $x$ 作为一次因子，右侧指数只须保留二次项。
\paragraph{解答}将关系写成 $x=\phi(y)=y\ee^{-y}$。有 $\phi'(y)=\ee^{-y}(1-y)$，在 $|y|<1/2$ 上严格为正。因此 $\phi$ 严格递增，零点附近有连续反函数。由于给定 $y(x)$ 在 $0$ 连续且 $y(0)=0$，充分小的 $|x|$ 保证 $|y(x)|<1/2$，故它恰为这个局部反函数。反函数求导公式给出 $y'=1/\phi'(y)$；因分母不为零，反复求导可知 $y$ 在该邻域任意阶可导。这验证了设三阶展开的前提，也排除了正侧可能存在的另一实解分支。
设 $y=ax+bx^2+cx^3+o(x^3)$，代入 $y=x(1+y+y^2/2+o(y^2))$。由 $y=O(x)$，右侧为
\[
x+ax^2+\left(b+\frac{a^2}2\right)x^3+o(x^3).
\]
依次比较得 $a=1$、$b=a=1$、$c=b+a^2/2=3/2$，故 $y=x+x^2+3x^3/2+o(x^3)$。仅从形式方程设多项式而不验证反函数的存在与光滑性，会遗漏泰勒公式的前提。

\begin{example}[复合系数求导]设 $f(x)=\ee^{x^2}\cos x$，求 $f^{(4)}(0)$；再求 $g(x)=\ee^x$ 的 $g'''(1)$。\end{example}
\paragraph{思路分析}第一问函数光滑且为偶函数，只需要四次系数。第二问把展开点平移到零点。
\paragraph{解答}
\[
f(x)=\left(1+x^2+\frac{x^4}2\right)
\left(1-\frac{x^2}2+\frac{x^4}{24}\right)+o(x^4)
=1+\frac{x^2}2+\frac{x^4}{24}+o(x^4).
\]
故 $f^{(4)}(0)=4!/24=1$。令 $h=x-1$，$g(1+h)=\ee(1+h+h^2/2+h^3/6)+o(h^3)$，故 $g'''(1)=3!\ee/6=\ee$。

\begin{theorem}[首个非零高阶导数判定极值]
设 $f$ 满足 $a$ 处 $m$ 阶皮亚诺公式的条件，$f'(a)=\cdots=f^{(m-1)}(a)=0$，且 $f^{(m)}(a)\ne0$。若 $m$ 为偶数，则该导数为正时 $a$ 是严格极小值点，为负时是严格极大值点；若 $m$ 为奇数，则 $a$ 不是极值点。
\end{theorem}
\begin{proof}记 $c=f^{(m)}(a)/m!\ne0$。有 $f(a+h)-f(a)=h^m(c+o(1))$。取足够小邻域使 $c+o(1)$ 与 $c$ 同号。偶次 $h^m$ 在两侧均为正，奇次在两侧异号，分别得到所述结论。\end{proof}
\begin{example}[二阶检验退化后继续展开]判定 $f(x)=\cos x-1+x^2/2$ 和 $g(x)=\sin x-x$ 在零点是否有极值。\end{example}
\paragraph{解答}$f(x)=x^4/24+o(x^4)$，零点附近非零 $x$ 时 $f(x)>f(0)=0$，为严格极小值点。$g(x)=-x^3/6+o(x^3)$，两侧差值异号，故无极值。二阶导数为零只表示二阶检验失效。

\section{有限公式、无穷级数与区间证明}
\label{lec:08:global}
\subsection{阶数趋于无穷需要额外论证}
固定 $n$、令 $x\to a$ 的皮亚诺公式，与固定 $x$、令 $n\to\infty$ 的级数表示是两种极限过程。即使每阶导数存在，也不能自动断言函数等于其泰勒级数。
对指数函数，固定 $x$ 后，拉格朗日余项满足
\[
|R_n(x)|\le\ee^{|x|}\frac{|x|^{n+1}}{(n+1)!}\longrightarrow0.
\]
最后的收敛可用相邻项比值 $|x|/(n+2)\to0$ 证明。这样才得到 $\ee^x=\sum_{k=0}^\infty x^k/k!$。几何级数则从有限恒等式
\[
\frac1{1-x}=\sum_{k=0}^nx^k+\frac{x^{n+1}}{1-x}
\]
出发，仅当 $|x|<1$ 余项趋零。级数的适用区间须随证明给出。

\begin{example}[余项证明指数函数的不等式与无理性]证明 $\ee^x>1+x$ 对 $x\ne0$ 成立，并证明 $\ee$ 为无理数。\end{example}
\paragraph{思路分析}第一问使用一阶多项式，二阶导数恒正。第二问希望得到一个严格位于 $0$ 与 $1$ 之间的整数，需乘以阶乘消去分母，再控制余项。
\paragraph{解答}由拉格朗日公式 $\ee^x=1+x+\ee^\xi x^2/2>1+x$。
先建立粗界 $\ee<4$：在 $x=1/2$ 处用零阶公式，$\ee^{1/2}=1+\ee^\xi/2<1+\ee^{1/2}/2$，故 $\ee^{1/2}<2$，平方得结论。
假设 $\ee=p/q$，其中 $p,q$ 为正整数。取 $n\ge q$ 且 $n+1>4$，指数函数在 $[0,1]$ 光滑，故
\[
\ee=\sum_{k=0}^n\frac1{k!}+\frac{\ee^\xi}{(n+1)!},\qquad0<\xi<1.
\]
乘以 $n!$，左侧 $n!p/q-\sum_{k=0}^nn!/k!$ 为整数，右侧为 $\ee^\xi/(n+1)$，满足 $0<\ee^\xi/(n+1)<4/(n+1)<1$，矛盾。因此 $\ee$ 无理。

\begin{example}[有界函数控制导数]设 $f\in C^2(\R)$，且 $|f|\le A$、$|f''|\le B$，其中 $A,B\ge0$。证明 $|f'(x)|\le\sqrt{2AB}$。\end{example}
\paragraph{思路分析}在同一点向左右展开，相减消去常数项；增量 $h$ 可自由选取，再平衡函数值误差与二阶余项。
\paragraph{解答}对任意 $h>0$，两个拉格朗日公式给出
\[
f(x+h)-f(x-h)=2hf'(x)+\frac{h^2}2\bigl(f''(\xi_+)-f''(\xi_-)\bigr).
\]
取绝对值并移项，$2h|f'(x)|\le2A+Bh^2$，从而 $|f'(x)|\le A/h+Bh/2$。若 $A,B>0$，取 $h=\sqrt{2A/B}$ 得所需界。若 $A=0$，函数恒零；若 $B=0$，上式对所有 $h$ 成立，令 $h\to\infty$ 得 $f'=0$。退化情形也满足结论。
\paragraph{方法回顾}这里两个余项的中间点通常不同，不能把它们随意合成同一个点。取绝对值估计不要求这两个点相同。

\section{分层练习与解析}
\begin{exercise}[基础巩固]求 $\ln(1+2x)$ 的三阶展开，以及 $\displaystyle\lim_{x\to0}\frac{\tan x-\sin x}{x^3}$。\end{exercise}
\noindent\textbf{解析}\quad 将 $2x$ 代入对数公式得 $2x-2x^2+8x^3/3+o(x^3)$。第二问分子为 $(1/3+1/6)x^3+o(x^3)$，极限为 $1/2$。
\begin{exercise}[方法迁移]求常数 $a,b$，使 $\ee^x-1-ax-bx^2=o(x^2)$；再求 $\ln(\cos x)$ 的四阶展开。\end{exercise}
\noindent\textbf{解析}\quad 展开后为 $(1-a)x+(1/2-b)x^2+o(x^2)$，由系数唯一性得 $a=1,b=1/2$。零点附近余弦为正，令 $u=\cos x-1=-x^2/2+x^4/24+o(x^4)$，
\[
\ln(\cos x)=u-\frac{u^2}2+o(u^2)
=-\frac{x^2}2-\frac{x^4}{12}+o(x^4).
\]
\begin{exercise}[综合提升]设 $f\in C^2[a,b]$，$f'(a)=f'(b)=0$，证明存在 $c\in(a,b)$，使 $|f''(c)|\ge4|f(b)-f(a)|/(b-a)^2$。\end{exercise}
\noindent\textbf{解析}\quad 取中点 $m=(a+b)/2$。分别从 $a,b$ 展开到 $m$，得到
\[
f(m)-f(a)=\frac{(b-a)^2}8f''(\xi_1),\qquad
f(m)-f(b)=\frac{(b-a)^2}8f''(\xi_2).
\]
相减并取绝对值，$|f(b)-f(a)|\le(b-a)^2\max(|f''(\xi_1)|,|f''(\xi_2)|)/4$。在两点中取绝对值较大者即可；无需假设最大值在整个开区间取得。
\begin{exercise}[误差检验]在 $0\le x\le1/2$ 上，用 $x-x^2/2$ 近似 $\ln(1+x)$，证明误差在 $0$ 与 $x^3/3$ 之间。\end{exercise}
\noindent\textbf{解析}\quad 三阶导数为 $2/(1+t)^3$。拉格朗日公式给出 $R_2=x^3/[3(1+\xi)^3]$，其中 $0<\xi<x$（$x=0$ 时余项为零），故 $0\le R_2\le x^3/3\le1/24$。
\paragraph{自查}能否解释为何选展开点与阶数，逐个跟踪乘除或复合后的余项，并在使用系数求导、全区间估计和无穷级数时写出各自条件？这些条件分别解决不同问题，不能互相代替。


\chapter{导数的应用}
\label{lec:09:chapter}
导数把函数在一点附近的变化转化为可计算的量。本章进一步用一阶导数组织增减，用二阶导数描述弯曲，再把这些信息用于不等式、作图和优化。需要的前置知识是连续函数的介值与最值定理、求导法则、中值定理和泰勒公式。曲率中的弧长微分将在第十二章得到积分形式；此处从局部长度变化推导。

\section{切线与法线：先确定点，再确定方向}
若 $f$ 在 $x_0$ 可导，则图像在 $P=(x_0,f(x_0))$ 的切线为
\[y-f(x_0)=f'(x_0)(x-x_0).\]
法线是过 $P$ 且垂直于切线的直线。若 $f'(x_0)\ne0$，法线斜率为 $-1/f'(x_0)$；若切线水平，则法线为 $x=x_0$。对于正则参数曲线，切向量是 $(x'(t_0),y'(t_0))$，即使 $x'(t_0)=0$ 也可以使用向量表示切线；此时 $y'(t_0)\ne0$，切线竖直、法线水平。
\begin{example}[水平切线与竖直法线]
求 $y=x^2$ 在 $x=1$ 和 $x=0$ 处的切线、法线。
\end{example}
\paragraph{解答}$y'=2x$。在 $(1,1)$ 的切线为 $y-1=2(x-1)$，法线为 $y-1=-(x-1)/2$。在 $(0,0)$，切线为 $y=0$，法线为 $x=0$；零斜率时不能用 $-1/0$ 表示法线斜率。

\section{单调性：从局部斜率到区间变化}
\label{lec:09:monotonicity}
\begin{theorem}[导数与单调性]
设 $f$ 在区间 $I$ 上连续，在内部可导。若内部各点 $f'\ge0$，则 $f$ 单调不减；若 $f'>0$，则严格递增。相应的负号给出单调不增和严格递减。反过来，可导的单调不减函数满足 $f'\ge0$。
\end{theorem}
\begin{proof}
任取 $u<v$ 属于 $I$，在 $[u,v]$ 上应用拉格朗日中值定理，存在 $\xi\in(u,v)$，使 $f(v)-f(u)=f'(\xi)(v-u)$，其符号由导数决定。反向结论由差商得到：对正负两种增量 $h$，单调不减性均给出 $(f(x+h)-f(x))/h\ge0$，取极限得 $f'(x)\ge0$。
\end{proof}
严格递增并不要求处处 $f'>0$，例如 $x^3$ 在零点导数为零。若 $f'\ge0$ 且导数只在有限个点为零，可以把任意 $[u,v]$ 分成若干段；每个非退化段内有导数严格为正的小区间，故总增量严格为正。求单调区间时先写定义域，再用驻点、不可导点与间断点分段；不能跨过定义域的缺口使用中值定理。

\begin{example}[导数零点并不破坏严格单调]
求 $f(x)=x^3-3x$ 的单调区间。
\end{example}
\paragraph{思路分析}导数可因式分解，符号的变化只可能发生在因子的零点。
\paragraph{解答}定义域为 $\R$，$f'(x)=3(x-1)(x+1)$。当 $x<-1$ 或 $x>1$ 时两因子同号，导数为正；当 $-1<x<1$ 时导数为负。因此 $f$ 在 $(-\infty,-1]$、$[1,+\infty)$ 严格递增，在 $[-1,1]$ 严格递减。端点可以并入相应区间，因为函数连续。

\begin{example}[构造差函数证明不等式]
证明对 $x>0$，$\dfrac{x}{1+x}<\ln(1+x)<x$。
\end{example}
\paragraph{思路分析}两边都在 $x=0$ 取零。把待证差的函数构造成从零开始严格增加的函数，可以证明整个半轴上的不等式。
\paragraph{证明}令 $A(x)=x-\ln(1+x)$，则 $A(0)=0$，$A'(x)=x/(1+x)>0$。令 $B(x)=\ln(1+x)-x/(1+x)$，则 $B(0)=0$，
\[ B'(x)=\frac1{1+x}-\frac1{(1+x)^2}=\frac{x}{(1+x)^2}>0. \]
因此 $A(x),B(x)>0$。局部泰勒展开只能直接控制零点附近；这里的导数符号覆盖全部 $x>0$。

\section{极值与最值：候选点和判别}
\label{lec:09:extrema}
\begin{definition}
若内点 $x_0$ 的某邻域内 $f(x)\le f(x_0)$，称 $x_0$ 为局部极大值点，$f(x_0)$ 为局部极大值；反向不等式定义局部极小值。若不等式对整个定义域成立，则称全局最大值或最小值。除特别说明，本章的极值指内点的局部极值。
\end{definition}
费马引理指出：可导的内点极值必有 $f'(x_0)=0$。证明时，在局部极大值点的右差商非正、左差商非负，二者若有同一极限，只能等于零。由此得到候选点清单：内部驻点、内部不可导点，以及求闭区间最值时的端点。驻点只是必要候选，$x^3$ 的零点就没有极值；$|x|$ 的极小值点又没有导数。

\begin{theorem}[第一导数判别]
设 $f$ 在 $x_0$ 连续，在其两侧某邻域可导。若左侧 $f'>0$、右侧 $f'<0$，则 $x_0$ 为严格极大值点；若左负右正，则为严格极小值点。若两侧同为正或同为负，则无极值。
\end{theorem}
\begin{proof}
左正意味着从左侧向 $x_0$ 递增，故左侧函数值小于 $f(x_0)$；右负意味着从 $x_0$ 向右递减，故右侧也小于 $f(x_0)$。其他情形以相同的单调性比较得到。
\end{proof}
\begin{theorem}[第二导数判别]
若 $f'(x_0)=0$ 且 $f''(x_0)>0$，则 $x_0$ 为严格极小值点；若 $f''(x_0)<0$，则为严格极大值点。
\end{theorem}
\begin{proof}
由二阶皮亚诺公式，$f(x_0+h)-f(x_0)=\tfrac12f''(x_0)h^2+o(h^2)$。主项在 $h\ne0$ 充分小时的符号固定，给出结论。
\end{proof}
若 $f''(x_0)=0$，判别没有结论。更一般地，若最低非零阶为 $m$，即 $f^{(1)}(x_0)=\cdots=f^{(m-1)}(x_0)=0$、$f^{(m)}(x_0)\ne0$，且可使用 $m$ 阶皮亚诺公式，则最低项为 $f^{(m)}(x_0)h^m/m!$：偶数阶给出极值，奇数阶不给出极值。这解释了 $x^4$ 与 $x^3$ 在零点的不同表现。

\begin{example}[不可导点与闭区间最值]
求 $f(x)=|x|-x^2$ 在 $[-2,2]$ 上的极值和最值。
\end{example}
\paragraph{思路分析}绝对值分界点必须保留；闭区间端点也要参加最值比较。
\paragraph{解答}当 $x<0$，$f'=-1-2x$；当 $x>0$，$f'=1-2x$。导数在 $(-2,-1/2)$、$(0,1/2)$ 为正，在 $(-1/2,0)$、$(1/2,2)$ 为负。因此 $x=\pm1/2$ 是极大值点，值为 $1/4$；零点左右由减到增，是不可导的极小值点，值为零。候选点的值为
\[f(-2)=f(2)=-2,\qquad f(\pm1/2)=1/4,\qquad f(0)=0.\]
最大值为 $1/4$，最小值为 $-2$。局部极小值零并不是全局最小值。

\begin{example}[开区间上的最大值]
求 $f(x)=x^{1/x}$（$x>0$）的最大值。
\end{example}
\paragraph{思路分析}底数与指数同时变化，直接比较困难。对数严格递增，因此最大化 $f$ 等价于最大化 $\ln f$。
\paragraph{解答}令 $H(x)=\ln x/x$，则 $H'=(1-\ln x)/x^2$。在 $(0,\ee)$ 上为正，在 $(\ee,+\infty)$ 上为负。每个 $x<\ee$ 都因递增而有 $H(x)<H(\ee)$，每个 $x>\ee$ 都因递减而有同样结论，故最大值是 $\ee^{1/\ee}$。这个结论来自两侧单调性，并非只由唯一驻点推出。开区间上连续函数未必取到最值，例如 $f(x)=x$ 在 $(0,1)$ 上只有上确界 $1$。

\section{凸、凹与拐点}
\label{lec:09:convexity}
\begin{definition}[凸性约定]
若对区间 $I$ 内任意 $u,v$ 和 $0\le\lambda\le1$，有
\[ f((1-\lambda)u+\lambda v)\le(1-\lambda)f(u)+\lambda f(v), \]
则称 $f$ 为凸函数（下凸）；图像位于弦的下方。反向不等式定义凹函数（上凸）。若 $u\ne v$、$0<\lambda<1$ 时均为严格不等式，则称严格凸或严格凹。不同教材对中文凹凸名称的用法不同，本书以此不等式为准。
\end{definition}
\begin{theorem}[可导与二阶可导函数的凸性]
可导函数 $f$ 在区间上凸，当且仅当 $f'$ 单调不减，也等价于每一点的切线不等式
\[ f(v)\ge f(u)+f'(u)(v-u). \]
若 $f$ 二阶可导，则凸性等价于 $f''\ge0$。$f''>0$ 是严格凸的充分条件。
\end{theorem}
\begin{proof}
凸性定义经移项给出三点割线斜率的次序：若 $u<w<v$，则
\[\frac{f(w)-f(u)}{w-u}\le\frac{f(v)-f(u)}{v-u}\le\frac{f(v)-f(w)}{v-w}.\]
令 $w\to u^+$、$w\to v^-$，可得 $f'(u)\le[f(v)-f(u)]/(v-u)\le f'(v)$，故导数单调，并得到切线不等式。反之，若导数单调，分别在 $[u,w]$、$[w,v]$ 用中值定理，两段割线斜率有序，移项即为凸性定义。若各点切线不等式成立，令 $w=(1-\lambda)u+\lambda v$，把 $f(u)\ge f(w)+f'(w)(u-w)$ 与 $f(v)\ge f(w)+f'(w)(v-w)$ 加权相加，导数项消去，得到凸性。最后把导数与单调性定理用于 $f'$，得到二阶判别；若 $f''>0$，两段斜率严格有序，故严格凸。
\end{proof}
严格凸不能反推 $f''>0$：$x^4$ 严格凸，但零点二阶导数为零。凸也不能反推严格凸，直线就满足凸性的等号。

\begin{figure}[htbp]
\centering
\begin{tikzpicture}[scale=1.05]
\draw[->] (-.4,0)--(3.5,0) node[right] {$x$};
\draw[->] (0,-.3)--(0,3.3) node[above] {$y$};
\draw[thick,domain=.2:3,samples=60] plot (\x,{.27*\x*\x+.3});
\draw[dashed] (.5,.3675)--(2.8,2.4168);
\draw (0,-.3075)--(2.8,1.9605);
\fill (1.5,.9075) circle(1.3pt);
\node[below right] at (1.5,.9075) {$P$};
\node[above left] at (2,1.704) {弦};
\node[below right] at (2.5,1.758) {切线};
\end{tikzpicture}
\caption{凸（下凸）图像位于弦下方、切线上方。实线切线取于点 $P$。}
\end{figure}

\begin{example}[切线与加权不等式]
证明 $\ee^x\ge1+x$（$x\in\R$），以及对正数 $a,b$ 和 $0\le\lambda\le1$，$a^{1-\lambda}b^\lambda\le(1-\lambda)a+\lambda b$。
\end{example}
\paragraph{思路分析}第一个右端是指数函数在零点的切线。第二个把乘积取对数，正好出现凹函数的加权定义。
\paragraph{证明}$f''(x)=\ee^x>0$，零点切线为 $1+x$，切线不等式给出第一式，严格凸保证等号仅在零点。$g(x)=\ln x$ 的二阶导数为 $-1/x^2<0$，故
\[\ln((1-\lambda)a+\lambda b)\ge(1-\lambda)\ln a+\lambda\ln b.\]
两边取指数即得第二式；当 $0<\lambda<1$ 时等号当且仅当 $a=b$，端点权重则总有等号。
凸性的定义可用归纳法推广为有限加权 Jensen 不等式：若 $\lambda_i\ge0$、$\sum\lambda_i=1$，则 $f(\sum\lambda_i x_i)\le\sum\lambda_i f(x_i)$。归纳时把前 $n-1$ 项先按总权重归一化，再与最后一项使用二点不等式；总权重为零时直接去掉零权项。

\begin{definition}[拐点]
若 $f$ 在内点 $x_0$ 连续，且图像在该点左右由凸变凹或由凹变凸，则 $(x_0,f(x_0))$ 是拐点。
\end{definition}
在两侧二阶导数存在且符号固定的常见情形，查 $f''$ 的符号改变即可。$f''=0$ 和二阶导数不存在的内点都是候选，必须检查两侧；没有定义的点不能成为图像的拐点。
若 $f''(x_0)=0$、$f'''(x_0)\ne0$，则由 $f''(x_0+h)/h\to f'''(x_0)$ 的保号性，$f''$ 在左右异号，故是拐点。这是充分条件。

\begin{example}[零二阶导数与拐点]
判断 $x^3$、$x^4$、$\sqrt[3]{x}$ 在零点是否有拐点。
\end{example}
\paragraph{解答}$x^3$ 的二阶导数 $6x$ 左负右正，零点是拐点。$x^4$ 的二阶导数 $12x^2$ 两侧均正，零点不是拐点。$\sqrt[3]{x}$ 连续，而 $x\ne0$ 时二阶导数为 $-2/(9x^{5/3})$，左正右负，零点也是拐点，虽然该点的一阶、二阶导数都不存在。

\section{渐近线与完整作图}
\label{lec:09:asymptotes}
若在 $x\to a$ 的至少一侧 $f(x)\to+\infty$ 或 $-\infty$，则 $x=a$ 为垂直渐近线。仅分母为零不足以保证这一点，例如 $\sin x/x$ 在零点有有限极限。
在某个无穷方向，直线 $y=kx+b$ 是渐近线的准确含义是 $f(x)-kx-b\to0$，因为曲线点到该直线的距离为 $|f(x)-kx-b|/\sqrt{1+k^2}$。由此除以 $x$，可得必要的 $k=\lim f(x)/x$；再相减得 $b=\lim(f(x)-kx)$。两个有限极限同时存在也充分。$k=0$ 给水平渐近线，$k\ne0$ 给斜渐近线。正负无穷必须分开。

\begin{example}[两端不同的斜渐近线]
求 $f(x)=\sqrt{x^2+1}$ 的渐近线。
\end{example}
\paragraph{思路分析}根号的主项是 $|x|$，负无穷时不能替换成 $x$。
\paragraph{解答}正无穷时 $k=1$，并且
\[f(x)-x=\frac1{\sqrt{x^2+1}+x}\longrightarrow0.\]
负无穷时 $k=-1$，且
\[f(x)+x=\frac1{\sqrt{x^2+1}-x}\longrightarrow0.\]
故两端分别为 $y=x$ 和 $y=-x$。定义域为全实轴且处处连续有限，无垂直渐近线。

\begin{example}[由信息表组织图像]
分析并描绘 $f(x)=x+1/x$。
\end{example}
\paragraph{思路分析}先定两支的位置，再查增减和弯曲，最后用渐近线约束远端走向。
\paragraph{解答}定义域为 $\R\setminus\{0\}$，函数为奇函数，无实零点。零点右侧趋正无穷、左侧趋负无穷，故 $x=0$ 为垂直渐近线；两端 $f(x)-x=1/x\to0$，故斜渐近线为 $y=x$。
\[f'(x)=1-\frac1{x^2},\qquad f''(x)=\frac2{x^3}.\]
导数在 $(-\infty,-1)$、$(1,+\infty)$ 为正，在 $(-1,0)$、$(0,1)$ 为负。因此 $(-1,-2)$ 是局部极大点，$(1,2)$ 是局部极小点。负半轴凹（上凸），正半轴凸（下凸）；零点不在图像上，不能叫拐点。正支在直线 $y=x$ 上方，负支在下方。
\begin{figure}[htbp]
\centering
\begin{tikzpicture}[scale=.8]
\draw[->] (-4,0)--(4,0) node[right] {$x$};
\draw[->] (0,-4.8)--(0,4.8) node[above] {$y$};
\draw[dashed] (-4,-4)--(4,4) node[right] {$y=x$};
\draw[thick,domain=.23:3.6,samples=90] plot (\x,{\x+1/\x});
\draw[thick,domain=-3.6:-.23,samples=90] plot (\x,{\x+1/\x});
\fill (1,2) circle(1.5pt) node[above right] {$(1,2)$};
\fill (-1,-2) circle(1.5pt) node[below left,fill=white,inner sep=2pt] {$(-1,-2)$};
\end{tikzpicture}
\caption{两支分别沿垂直渐近线与斜渐近线伸展。}
\end{figure}
一般作图顺序是：定义域与连续性，奇偶或周期，截距及端点极限，渐近线，一阶导数符号与极值，二阶导数符号与拐点，最后连接各段。图像应满足全部已知信息，而不是仅穿过几个算出的点。

\section{曲率：单位长度中的切线转动}
\label{lec:09:curvature}
坡度描述方向，曲率描述方向改变的快慢。直线即使斜率很大，曲率仍为零；圆沿各点的切线不断转动，曲率恒定。
\begin{definition}
对二阶光滑的正则曲线，选局部连续的切线方向角 $\theta$ 与有向弧长参数 $s$，曲率定义为 $\kappa=|\dd\theta/\dd s|$。正则指速度不为零。
\end{definition}
对于 $y=f(x)$，有 $\tan\theta=f'(x)$，所以 $\dd\theta/\dd x=f''/(1+(f')^2)$。局部弧长变化满足 $\dd s/\dd x=\sqrt{1+(f')^2}$，由链式法则得到
\[\kappa=\frac{|f''(x)|}{(1+(f'(x))^2)^{3/2}}.\]
对于参数曲线 $(x(t),y(t))$，设速度 $v=\sqrt{(x')^2+(y')^2}>0$，局部写 $x'=v\cos\theta$、$y'=v\sin\theta$。求导再作差消去 $v'$，得到
\[x'y''-y'x''=v^2\theta',\qquad \frac{\dd s}{\dd t}=v,\qquad
\kappa=\frac{|x'y''-y'x''|}{((x')^2+(y')^2)^{3/2}}.\]
这种推导也适用于 $x'=0$ 的竖直切线处；只要 $x',y'$ 不同时为零，公式就有意义。若 $\kappa>0$，曲率半径 $\rho=1/\kappa$；曲率为零时没有有限的曲率圆。

\begin{example}[椭圆曲率及退化检查]
求椭圆 $x=a\cos t$、$y=b\sin t$（$a\ge b>0$）的曲率及最大、最小曲率。
\end{example}
\paragraph{解答}四个导数为 $x'=-a\sin t$、$y'=b\cos t$、$x''=-a\cos t$、$y''=-b\sin t$。分子行列式为 $ab(\sin^2t+\cos^2t)=ab$，故
\[\kappa(t)=\frac{ab}{(a^2\sin^2t+b^2\cos^2t)^{3/2}}.\]
当 $a>b$，分母底数为 $b^2+(a^2-b^2)\sin^2t$，介于 $b^2$ 与 $a^2$ 之间。最大曲率 $a/b^2$ 在长轴两端取得，最小曲率 $b/a^2$ 在短轴两端取得。若 $a=b=R$，则每一点都是最大与最小曲率点，曲率均为 $1/R$。

\subsection{极坐标形式与曲率圆}
令 $x=r(\theta)\cos\theta$、$y=r(\theta)\sin\theta$。先求出
\begin{align*}
x'&=r'\cos\theta-r\sin\theta,&y'&=r'\sin\theta+r\cos\theta,\\
x''&=(r''-r)\cos\theta-2r'\sin\theta,&y''&=(r''-r)\sin\theta+2r'\cos\theta.
\end{align*}
平方相加使交叉项抵消，得到 $(x')^2+(y')^2=(r')^2+r^2$；行列式展开后得到 $x'y''-y'x''=r^2+2(r')^2-rr''$。因此在 $r^2+(r')^2>0$ 处
\[\kappa=\frac{|r^2+2(r')^2-rr''|}{(r^2+(r')^2)^{3/2}}.\]
对于图像 $y=f(x)$ 且 $f''\ne0$，曲率圆圆心坐标为
\[\left(x-\frac{f'(1+(f')^2)}{f''},\quad f(x)+\frac{1+(f')^2}{f''}\right).\]
这是沿单位法向量 $(-f',1)/\sqrt{1+(f')^2}$ 移动有向长度 $(1+(f')^2)^{3/2}/f''$ 的结果；符号决定圆心在法线哪一侧。曲率圆也称密切圆，它与原曲线在该点具有相同的切线和曲率。曲率圆圆心的轨迹称为渐屈线，可能出现退化情形，例如圆的渐屈线退化为圆心。一般渐开线关系还需讨论轨迹的正规性；本章只使用上述曲率及曲率圆公式。
\begin{example}[抛物线的曲率圆]
求 $y=x^2/2$ 在原点的曲率圆，并求极坐标圆 $r=R>0$ 的曲率。
\end{example}
\paragraph{解答}抛物线有 $f'(0)=0$、$f''(0)=1$，故 $\kappa=1$、$\rho=1$，圆心为 $(0,1)$，圆方程为 $x^2+(y-1)^2=1$。极坐标圆的 $r'=r''=0$，代入公式得 $\kappa=R^2/(R^2)^{3/2}=1/R$，与圆的几何直观一致。

\section{几何优化：约束、边界与结论}
\label{lec:09:optimization}
几何优化先把图形条件译成变量的可行范围，再建立目标函数。消元后寻找内部驻点与不可导点，最后检查边界或端点极限。只有连续性和闭有界可行区间才自动保证最值存在；可行区间为开区间时，要用全区间单调性或其他全局估计证明。唯一驻点不说明它一定是最大值，例如 $x^2$ 在 $\R$ 上只有一个驻点，但没有最大值。

\begin{example}[矩形面积与可行域]
周长为 $P>0$ 的矩形，求最大面积。
\end{example}
\paragraph{解答}设一边 $x$，另一边为 $P/2-x$，非退化条件是 $0<x<P/2$。面积 $A=x(P/2-x)$，$A'=P/2-2x$ 在 $x=P/4$ 左正右负，故全区间最大值为 $P^2/16$，由正方形取得。也可配方 $A=P^2/16-(x-P/4)^2$，直接给出全局上界和等号条件。两端极限为零，不能把不存在的退化矩形误称为最小面积的取到点。

\begin{example}[固定母线的圆锥]
母线长为 $L>0$ 的圆锥，求最大容积以及展开侧面扇形的角度。
\end{example}
\paragraph{思路分析}以高为变量可以避免体积函数中出现根号；底半径用勾股约束消去。
\paragraph{解答}设高为 $h$、底半径为 $r$，$r^2+h^2=L^2$、$0<h<L$。于是
\[V(h)=\frac\pi3(L^2h-h^3),\qquad V'(h)=\frac\pi3(L^2-3h^2).\]
导数在 $h=L/\sqrt3$ 前正后负，两端容积极限为零，故最大容积 $2\pi L^3/(9\sqrt3)$ 在该点取得。此时 $r=L\sqrt{2/3}$。展开扇形的弧长等于底周长，若其角为 $\beta$，则 $L\beta=2\pi r$，故 $\beta=2\pi\sqrt{2/3}$。若从整圆剪裁，剪去角为 $2\pi-\beta$，两种角必须区分。

\section{分层练习与解析}
以下均为自编练习。
\begin{exercise}[基础巩固]
求 $f(x)=x^4-2x^2$ 在 $[-2,2]$ 上的极值、最大值和最小值。
\end{exercise}
\paragraph{解析}$f'=4x(x-1)(x+1)$，依次在 $(-2,-1)$、$(-1,0)$、$(0,1)$、$(1,2)$ 为负、正、负、正。因此 $\pm1$ 是极小值点，值 $-1$；零点为极大值点，值零。端点值均为 $8$，故最大值 $8$ 在 $\pm2$ 取得，最小值 $-1$ 在 $\pm1$ 取得。
\begin{exercise}[方法迁移]
分析 $f(x)=\ln x/x$ 的单调性、凹凸性和渐近线，并求极值与拐点。
\end{exercise}
\paragraph{解析}定义域为 $(0,+\infty)$。$f'=(1-\ln x)/x^2$，在 $(0,\ee)$ 正，在 $(\ee,+\infty)$ 负，最大值为 $1/\ee$。$f''=(2\ln x-3)/x^3$，在 $(0,\ee^{3/2})$ 负，在 $(\ee^{3/2},+\infty)$ 正，故先凹（上凸）后凸（下凸），拐点为 $(\ee^{3/2},\tfrac32\ee^{-3/2})$。在零点右侧趋负无穷，故 $x=0$ 为垂直渐近线；无穷远处趋零，故 $y=0$ 为水平渐近线。没有最小值，因为零点右侧无下界。
\begin{exercise}[综合提升]
求 $f_a(x)=x^3-3ax$ 在 $[-1,1]$ 上的最大值，其中 $a\ge0$。
\end{exercise}
\paragraph{解析}导数为 $3(x^2-a)$。若 $a=0$，函数严格递增，最大值为 $1$。若 $0<a<1$，内部候选点为 $\pm\sqrt a$，候选值为
\[f_a(-1)=3a-1,\quad f_a(1)=1-3a,\quad f_a(-\sqrt a)=2a^{3/2},\quad f_a(\sqrt a)=-2a^{3/2}.\]
令 $s=\sqrt a$，有 $2s^3-(3s^2-1)=(s-1)^2(2s+1)\ge0$，故左端点不会胜过内部极大值；而 $2s^3-(1-3s^2)=2(s-1/2)(s+1)^2$，分界是 $a=1/4$。若 $a\ge1$，导数在内部非正，最大值在 $-1$。因此
\[\max f_a=\begin{cases}1-3a,&0\le a\le1/4,\\2a^{3/2},&1/4\le a\le1,\\3a-1,&a\ge1.\end{cases}\]
在 $a=1/4$，$x=1$ 与 $x=-1/2$ 同时取到最大值；在 $a=1$ 两个描述对应同一个点 $x=-1$。参数改变时，驻点进出区间和候选值交叉都可能形成分界。
\begin{exercise}[曲率与优化]
半径为 $R>0$ 的圆内接矩形，求最大面积；求圆上任意点的曲率。
\end{exercise}
\paragraph{解析}设矩形半边长为 $u,v>0$，则 $u^2+v^2=R^2$，面积 $4uv$。由 $(u-v)^2\ge0$ 得 $2uv\le R^2$，故最大面积为 $2R^2$，在 $u=v=R/\sqrt2$ 时取得。参数化圆为 $(R\cos t,R\sin t)$，曲率分子为 $R^2$，分母为 $R^3$，故处处 $\kappa=1/R$。面积最大与曲率常量是两种不同问题，各自的条件与计算应分别检查。

章末自查：能够区分驻点、极值点与最值点；能够解释端点和不可导点为何必须检查；能够用弦不等式固定凹凸术语；能够说明 $f''=0$ 为什么不足以判拐点；能够分别计算正负无穷的渐近线；能够把优化模型的每个候选点放回可行范围并检验全局结论。


\chapter{不定积分}
\label{lec:10:chapter}

本章所需的前置知识是导数、微分、链式法则、乘积求导和中值定理。求导回答函数怎样变化，本章反过来问：已知变化率，怎样找回函数？计算目标是一族函数，因而最后既要检查导数，也要检查定义区间和积分常数。

\section{原函数与积分常数}
\begin{definition}
设函数 $f$ 在开区间 $I$ 上有定义。若 $F$ 在 $I$ 上可导且对每个 $x\in I$ 都有 $F'(x)=f(x)$，则称 $F$ 是 $f$ 在 $I$ 上的一个原函数。$f$ 在该区间上的所有原函数组成不定积分，记作 $\int f(x)\,\dd x$。
\end{definition}
例如 $(x^2)'=2x$，而 $(x^2+7)'=2x$。导数丢掉了常数信息，这正是答案必须带常数的原因。
\begin{theorem}
若 $F$ 是 $f$ 在区间 $I$ 上的一个原函数，则全部原函数恰为 $F+C$，其中 $C$ 为任意实常数。
\end{theorem}
\begin{proof}
常数的导数为零，所以 $F+C$ 均为原函数。反过来，若 $G'=f$，则 $(G-F)'=0$。对任意 $a<b$ 且 $a,b\in I$，$G-F$ 在 $[a,b]$ 连续、在 $(a,b)$ 可导，中值定理给出
\[
(G-F)(b)-(G-F)(a)=(G-F)'(\xi)(b-a)=0.
\]
任意两点的值相同，所以 $G-F$ 是常数。
\end{proof}
这段证明用到了两点之间整段属于定义域。若定义域断开，每个连通区间可以有自己的常数。例如 $1/x$ 在 $\R\setminus\{0\}$ 上的全部原函数为
\[
F(x)=\begin{cases}\ln(-x)+C_-,&x<0,\\ \ln x+C_+,&x>0,\end{cases}
\]
两常数没有衔接条件。写 $\ln|x|+C$ 时通常是在某一个固定区间上讨论。

\begin{proposition}[线性性质]
若 $F'=f$、$G'=g$，则 $aF+bG$ 是 $af+bg$ 的一个原函数，其中 $a,b$ 为常数。因此可逐项积分并在最后统一补上任意常数 $C$。
\end{proposition}
\begin{proof}
直接求导：$(aF+bG)'=af+bg$。
\end{proof}
这里没有乘法或除法的对应公式。例如 $\int x^2\,\dd x=x^3/3+C$，而把它写成 $(\int x\,\dd x)^2$ 会得到四次函数，其导数并非 $x^2$。

连续函数在区间上有原函数，这将在下一章由变上限积分证明。本章计算中涉及的连续函数都可据此确认存在性；但存在原函数不等于能用有限个初等函数写出，例如 $\ee^{-x^2}$ 的原函数通常用积分定义。

\section{基本公式与恒等变形}
以下列的是一个原函数；使用时补上 $C$，并限制在被积函数和结果都有意义的区间内。对任意实指数 $\alpha$，幂公式首先在 $x>0$ 使用；若负半轴上的幂也有实意义且满足相同求导公式，则可在相应区间使用。
\begin{align*}
\int x^\alpha\,\dd x&=\frac{x^{\alpha+1}}{\alpha+1}+C\quad(\alpha\ne-1),&
\int\frac{\dd x}{x}&=\ln|x|+C\quad(x\ne0),\\
\int\ee^x\,\dd x&=\ee^x+C,&
\int b^x\,\dd x&=\frac{b^x}{\ln b}+C\quad(b>0,
\ b\ne1),\\
\int\sin x\,\dd x&=-\cos x+C,&
\int\cos x\,\dd x&=\sin x+C,\\
\int\sec^2x\,\dd x&=\tan x+C,&
\int\csc^2x\,\dd x&=-\cot x+C,\\
\int\sec x\tan x\,\dd x&=\sec x+C,&
\int\csc x\cot x\,\dd x&=-\csc x+C.
\end{align*}
三角公式分别避开分母为零的点；当 $b=1$ 时，被积函数为 $1$，原函数为 $x$，不能代入含 $\ln b$ 的公式。
对 $a>0$，另有
\begin{align*}
\int\frac{\dd x}{x^2+a^2}&=\frac1a\arctan\frac xa+C,\\
\int\frac{\dd x}{\sqrt{a^2-x^2}}&=\arcsin\frac xa+C\quad(|x|<a),\\
\int\frac{\dd x}{x^2-a^2}&=\frac1{2a}\ln\left|\frac{x-a}{x+a}\right|+C\quad(x\ne\pm a).
\end{align*}
第一式的 $1/a$ 来自链式法则。第三式来自
\[
\frac1{x^2-a^2}=\frac1{2a}\left(\frac1{x-a}-\frac1{x+a}\right).
\]
对双曲函数，按定义 $\sinh x=(\ee^x-\ee^{-x})/2$、$\cosh x=(\ee^x+\ee^{-x})/2$，逐项积分即得 $\int\sinh x\,\dd x=\cosh x+C$ 和 $\int\cosh x\,\dd x=\sinh x+C$。

\begin{example}[先化简再积分]
计算 $\displaystyle\int\frac{x^4}{1+x^2}\,\dd x$ 和 $\displaystyle\int\tan^2x\,\dd x$。
\end{example}
\paragraph{思路分析}第一个分子次数较高，除法可把它降为整式与基本分式之和；第二个用恒等式把幂转为已知导数。
\paragraph{解答}
\[
x^4=(x^2-1)(1+x^2)+1,
\qquad \frac{x^4}{1+x^2}=x^2-1+\frac1{1+x^2}.
\]
因此第一积分为 $x^3/3-x+\arctan x+C$。在 $\cos x\ne0$ 的区间上，$\tan^2x=\sec^2x-1$，第二积分为 $\tan x-x+C$。分别求导得到原来的两个被积函数。

\section{换元：逆用链式法则}
\label{lec:10:substitution}
\subsection{凑微分与第一类换元}
若 $F'=f$，$g$ 可导，且 $g(x)$ 落在 $F$ 的定义区间内，则
\[
\int f(g(x))g'(x)\,\dd x=F(g(x))+C.
\]
依据就是 $(F(g(x)))'=f(g(x))g'(x)$。实际操作先找内层 $g$，算出 $\dd g=g'\dd x$，检查剩余因子能否完全写成 $g$ 的函数，再统一用一个新变量积分，最后回代。只差常数可以补偿；差一个变量因子时，不能凭空补上。

\begin{example}[常数补偿与对数内层]
计算 $\displaystyle\int\frac{x}{1+x^2}\,\dd x$ 和 $\displaystyle\int\frac{\sqrt{1+\ln x}}x\,\dd x$。
\end{example}
\paragraph{思路分析}第一题分子的 $x\dd x$ 是分母微分的一半。第二题 $\dd x/x$ 正好是对数内层的微分，含对数并不意味着必须分部积分。
\paragraph{解答}令 $u=1+x^2$，则 $\dd u=2x\dd x$，第一题为 $\frac12\int\dd u/u=\frac12\ln(1+x^2)+C$。第二题在 $x>\ee^{-1}$ 上令 $v=1+\ln x$，有 $\dd v=\dd x/x$，所以
\[
\int\frac{\sqrt{1+\ln x}}x\,\dd x
=\int v^{1/2}\,\dd v=\frac23(1+\ln x)^{3/2}+C.
\]
求导时系数 $\frac23\cdot\frac32$ 相消，再乘 $1/x$，得到原式。端点 $x=\ee^{-1}$ 可另讨论右导数；通常原函数公式写在开区间上。

\begin{example}[正割和余割的原函数]
推导 $\int\sec x\,\dd x$ 和 $\int\csc x\,\dd x$。
\end{example}
\paragraph{思路分析}目标是构造 $h'/h$。求导 $\sec x+\tan x$，会产生它自身乘 $\sec x$。
\paragraph{解答}在 $\cos x\ne0$ 的区间，$(\sec x+\tan x)(\sec x-\tan x)=1$，故前者不为零；并且
\[
(\sec x+\tan x)'=\sec x(\sec x+\tan x).
\]
因而 $\int\sec x\,\dd x=\ln|\sec x+\tan x|+C$。同理，$\csc x-\cot x$ 不为零且其导数等于 $\csc x(\csc x-\cot x)$，故
\[
\int\csc x\,\dd x=\ln|\csc x-\cot x|+C.
\]
另由 $\tan x=\sin x/\cos x$、$\cot x=\cos x/\sin x$ 得
\[
\int\tan x\,\dd x=-\ln|\cos x|+C,
\qquad\int\cot x\,\dd x=\ln|\sin x|+C.
\]

\paragraph{错误辨析}把 $\int(1-x)^{-1/2}\,\dd x$ 写成 $\arcsin\sqrt{x}$，既引入了额外的 $x\ge0$ 限制，又遗漏 $\dd\sqrt{x}=\dd x/(2\sqrt{x})$。正确换元 $u=1-x$ 给出 $-2\sqrt{1-x}+C$，适用于 $x<1$。

\subsection{第二类换元与根式有理化}
令 $x=\varphi(t)$。在选定的区间上要求该变换可逆，且 $\varphi$ 和逆函数可导；则求出 $\int f(\varphi(t))\varphi'(t)\,\dd t=H(t)+C$ 后回代 $t=\varphi^{-1}(x)$。回代结果的导数为 $f(x)$，这证明了公式。选取单调支的目的，是让回代唯一并确定根式符号。

\begin{example}[不同根指数的统一代换]
计算 $\displaystyle\int\frac{\dd x}{\sqrt{x}+\sqrt[3]{x}}$。
\end{example}
\paragraph{思路分析}实定义域是 $x>0$。根指数 $2,3$ 的最小公倍数是 $6$，取 $x=t^6$ 能同时消去两种根号。
\paragraph{解答}取 $t=x^{1/6}>0$，则 $\sqrt{x}=t^3$、$\sqrt[3]{x}=t^2$、$\dd x=6t^5\dd t$，所以
\begin{align*}
\int\frac{\dd x}{\sqrt{x}+\sqrt[3]{x}}
&=6\int\frac{t^3}{t+1}\,\dd t
=6\int\left(t^2-t+1-\frac1{t+1}\right)\dd t\\
&=2\sqrt{x}-3x^{1/3}+6x^{1/6}
-6\ln(1+x^{1/6})+C.
\end{align*}
除法的依据是 $t^3=(t+1)(t^2-t+1)-1$。验算可先对 $t$ 求导得到 $6t^3/(t+1)$，再乘 $\dd t/\dd x=1/(6t^5)$，即为 $1/(t^3+t^2)$。

对 $R\bigl(x,\sqrt[n]{(ax+b)/(cx+d)}\bigr)$，若 $ad-bc\ne0$，取根式为 $t$ 后可解出
\[
x=\frac{dt^n-b}{a-ct^n},\qquad
\dd x=\frac{n(ad-bc)t^{n-1}}{(a-ct^n)^2}\,\dd t.
\]
先检查根式实定义域及分母，再选择可逆支；此时所有项都成为 $t$ 的有理函数。含 $\sqrt{(x-a)(x-b)}$ 时应写成 $|x-a|\sqrt{(x-b)/(x-a)}$，不能在未检查区间时去掉绝对值。

\subsection{三类二次根式}
先配方，再辨认 $a^2-u^2$、$a^2+u^2$、$u^2-a^2$，其中 $a>0$。对应可取 $u=a\sin t$、$u=a\tan t$、$u=a\sec t$。关键不是记住代换名，而是用恒等式使根号内成为平方，并选择符号已知的支。

\begin{example}[半圆根式的三角代换]
求 $\displaystyle\int\sqrt{a^2-x^2}\,\dd x$，其中 $a>0$。
\end{example}
\paragraph{思路分析}取 $x=a\sin t$，$t\in(-\pi/2,\pi/2)$，此时 $\cos t>0$，所以根式就是 $a\cos t$，无需带未定符号的绝对值。
\paragraph{解答}
\begin{align*}
\int\sqrt{a^2-x^2}\,\dd x
&=a^2\int\cos^2t\,\dd t
=\frac{a^2}2t+\frac{a^2}4\sin2t+C\\
&=\frac12x\sqrt{a^2-x^2}+\frac{a^2}2\arcsin\frac xa+C.
\end{align*}
回代用 $\sin t=x/a$、$\cos t=\sqrt{a^2-x^2}/a$。对结果求导，除 $\frac12\sqrt{a^2-x^2}$ 外的两项合为
\[
-\frac{x^2}{2\sqrt{a^2-x^2}}+
\frac{a^2}{2\sqrt{a^2-x^2}}=\frac12\sqrt{a^2-x^2},
\]
因此验算通过，公式在 $(-a,a)$ 上成立。例如 $4x-x^2=4-(x-2)^2$，先取 $u=x-2$ 再用本式。

\begin{example}[平方和与平方差根式]
求 $\displaystyle\int\frac{\dd x}{\sqrt{x^2+a^2}}$ 与 $\displaystyle\int\frac{\dd x}{\sqrt{x^2-a^2}}$，$a>0$。
\end{example}
\paragraph{思路分析}第一题用正切代换，第二题先在 $x>a$ 用正割代换；负半轴另处理，防止开方时丢掉符号。
\paragraph{解答}第一题取 $x=a\tan t$，$t\in(-\pi/2,\pi/2)$，根式为 $a\sec t>0$，得
\[
\int\sec t\,\dd t=\ln|\sec t+\tan t|+C
=\ln\bigl(x+\sqrt{x^2+a^2}\bigr)+C.
\]
常数 $-\ln a$ 已吸收进 $C$，且根号大于 $|x|$，对数内始终正。
第二题在 $x>a$ 取 $x=a\sec t$，$t\in(0,\pi/2)$，则根式为 $a\tan t$，也化为 $\int\sec t\,\dd t$，结果是 $\ln(x+\sqrt{x^2-a^2})+C_+$。在 $x<-a$ 取 $u=-x>a$，积分变成 $-\int\dd u/\sqrt{u^2-a^2}$。利用
\[
(-x+\sqrt{x^2-a^2})(x+\sqrt{x^2-a^2})=-a^2,
\]
可将负支结果写成 $\ln|x+\sqrt{x^2-a^2}|+C_-$。两支常数独立。直接求导统一表达式：记 $s=\sqrt{x^2\pm a^2}$，则
\[
\frac{\dd}{\dd x}\ln|x+s|=\frac{1+x/s}{x+s}=\frac1s.
\]

\section{分部积分：逆用乘积法则}
\label{lec:10:parts}
由 $(uv)'=u'v+uv'$ 得到
\[
\int u\,\dd v=uv-\int v\,\dd u.
\]
选 $u$ 要使求导后的结构简化，选 $\dd v$ 要能求出 $v$，还要检查新积分是否更容易。对数、反三角通常适合求导，多项式反复求导会归零；但 $\ln x/x$ 应优先凑微分。这些判断比固定排序更可靠。

\begin{example}[对数与反三角的分部积分]
计算 $\int x^{1/2}\ln x\,\dd x$ 和 $\int\arctan x\,\dd x$。
\end{example}
\paragraph{思路分析}对数求导后成为 $1/x$，反三角求导后成为有理函数；这两种变化使剩余积分变简单。
\paragraph{解答}第一题 $x>0$，取 $u=\ln x$、$\dd v=x^{1/2}\dd x$，得 $\dd u=\dd x/x$、$v=2x^{3/2}/3$，所以
\[
\int x^{1/2}\ln x\,\dd x
=\frac23x^{3/2}\ln x-\frac23\int x^{1/2}\,\dd x
=\frac23x^{3/2}\ln x-\frac49x^{3/2}+C.
\]
求导出现的 $\frac23\sqrt{x}$ 与第二项导数相消。第二题取 $u=\arctan x$、$v=x$，于是
\[
\int\arctan x\,\dd x=x\arctan x-\int\frac{x}{1+x^2}\,\dd x
=x\arctan x-\frac12\ln(1+x^2)+C.
\]
求导后 $x/(1+x^2)$ 相消，余下 $\arctan x$。
一般地，在 $x>0$ 上，当 $\mu\ne-1$ 时同样得到
\[
\int x^\mu\ln x\,\dd x=
\frac{x^{\mu+1}\ln x}{\mu+1}-\frac{x^{\mu+1}}{(\mu+1)^2}+C;
\]
当 $\mu=-1$ 时改用 $\dd\ln x$，结果为 $(\ln x)^2/2+C$。

\subsection{反复分部与表格记录}
\begin{example}[多项式乘三角函数]
计算 $\displaystyle\int(x^2+1)\cos2x\,\dd x$。
\end{example}
\paragraph{思路分析}多项式求导两次变成常数，第三次归零；三角函数可以不断积分。表格只是记录连续分部的符号与系数。
\paragraph{解答}第一次分部得
\[
I=\frac12(x^2+1)\sin2x-\int x\sin2x\,\dd x.
\]
对余项再次分部：
\[
\int x\sin2x\,\dd x=-\frac12x\cos2x+\frac14\sin2x.
\]
因此
\[
I=\left(\frac{x^2}2+\frac14\right)\sin2x+\frac x2\cos2x+C.
\]
也可用下表逐行记录：每一行求导栏与下一行积分栏相乘，符号交替。
\begin{center}
\begin{tabular}{ccc}
\toprule
符号 & 连续求导 & 连续积分\\
\midrule
$+$ & $x^2+1$ & $\cos2x$\\
$-$ & $2x$ & $\frac12\sin2x$\\
$+$ & $2$ & $-\frac14\cos2x$\\
$-$ & $0$ & $-\frac18\sin2x$\\
\bottomrule
\end{tabular}
\end{center}
三个乘积为 $(x^2+1)\sin2x/2$、$x\cos2x/2$、$-\sin2x/4$。验算中余弦项合为 $(x^2+1)\cos2x$，正弦项 $x\sin2x-x\sin2x$ 相消。

\subsection{循环出现的原积分}
\begin{example}[指数乘正弦]
计算 $I=\int\ee^{ax}\sin bx\,\dd x$，其中 $a\ne0$。
\end{example}
\paragraph{思路分析}求导正弦、余弦会相互转换，指数积分保持自身。两次分部后会出现 $I$ 的常数倍，可把计算变成代数方程。
\paragraph{解答}选定一个原函数代表并暂时省略常数，有
\begin{align*}
I&=\frac{\ee^{ax}}a\sin bx-\frac baJ,\\
J=\int\ee^{ax}\cos bx\,\dd x
&=\frac{\ee^{ax}}a\cos bx+\frac baI.
\end{align*}
代回第一式，移项得
\[
\left(1+\frac{b^2}{a^2}\right)I=
\frac{\ee^{ax}}{a^2}(a\sin bx-b\cos bx).
\]
所以全部原函数为
\[
\frac{\ee^{ax}}{a^2+b^2}(a\sin bx-b\cos bx)+C.
\]
求导时余弦项 $-ab\cos bx+ab\cos bx$ 相消，正弦系数为 $a^2+b^2$。该答案也适用于 $a=0,b\ne0$，可直接求导确认；$a=b=0$ 时原被积函数为零，原函数为常数。

\section{有理函数与重因子的部分分式}
\label{lec:10:rational}
对 $P(x)/Q(x)$，先检查实定义域 $Q(x)\ne0$，再做多项式除法使余式次数小于 $Q$。实系数多项式可分解为一次因子和不可约二次因子的乘积，这是本节引用的代数结论。
真分式分解中，一次因子 $(x-a)^m$ 对应
\[
\frac{A_1}{x-a}+\frac{A_2}{(x-a)^2}+\cdots+\frac{A_m}{(x-a)^m};
\]
不可约二次因子 $q(x)=x^2+px+r$ 的 $m$ 次幂对应
\[
\frac{B_1x+D_1}{q(x)}+\frac{B_2x+D_2}{q(x)^2}+\cdots+
\frac{B_mx+D_m}{q(x)^m}.
\]
必须列出每一阶，二次分母的分子必须允许一次项。通分后的多项式恒等式可用代值和比较系数解出常数；代入分母零点应在去分母之后进行。

\begin{example}[一次重因子]
计算 $\displaystyle\int\frac{\dd x}{x(x-1)^2}$。
\end{example}
\paragraph{思路分析}定义域避开 $0,1$；$(x-1)^2$ 需要两项。只设最高阶一项无法表示所有真分式。
\paragraph{解答}设
\[
\frac1{x(x-1)^2}=\frac A x+\frac B{x-1}+\frac D{(x-1)^2}.
\]
去分母：$1=A(x-1)^2+Bx(x-1)+Dx$。取 $x=0$ 得 $A=1$，取 $x=1$ 得 $D=1$，比较 $x^2$ 系数得 $B=-1$。所以
\[
\int\frac{\dd x}{x(x-1)^2}=\ln|x|-\ln|x-1|-\frac1{x-1}+C.
\]
求导恢复三项，通分后分子为 $(x-1)^2-x(x-1)+x=1$。三个连通区间的常数可独立。

\subsection{二次分母和二次重因子}
分子先写成
\[
Ax+B=\frac A2(2x+p)+\left(B-\frac{Ap}2\right),
\]
第一部分凑 $\dd q$，第二部分配方 $q=(x+p/2)^2+a^2$，$a^2=r-p^2/4>0$。此写法对 $A=0$ 也成立。
高次二次分母可归结为 $I_n=\int\dd u/(u^2+a^2)^n$。对 $n\ge2$，由
\[
\left(\frac{u}{(u^2+a^2)^{n-1}}\right)'
=\frac{2(n-1)a^2}{(u^2+a^2)^n}
-\frac{2n-3}{(u^2+a^2)^{n-1}}
\]
积分并整理得到递推式
\[
I_n=\frac{u}{2(n-1)a^2(u^2+a^2)^{n-1}}
+\frac{2n-3}{2(n-1)a^2}I_{n-1},
\qquad I_1=\frac1a\arctan\frac ua.
\]
等式按原函数代表理解，最后补上常数。

\begin{example}[不可约二次重因子]
计算 $\displaystyle\int\frac{\dd x}{x(x^2+1)^2}$。
\end{example}
\paragraph{思路分析}设式中必须同时有 $(x^2+1)^{-1}$ 与 $(x^2+1)^{-2}$，每一项先允许一次分子，再由恒等式确定哪些系数为零。
\paragraph{解答}设
\[
\frac1{x(x^2+1)^2}=\frac A x+\frac{Bx+D}{x^2+1}
+\frac{Ex+G}{(x^2+1)^2}.
\]
去分母得 $1=A(x^2+1)^2+x(Bx+D)(x^2+1)+x(Ex+G)$。常数项给 $A=1$；$x^3$ 与 $x$ 系数给 $D=G=0$；$x^4$ 系数给 $B=-1$；$x^2$ 系数给 $E=-1$。故
\begin{align*}
\int\frac{\dd x}{x(x^2+1)^2}
&=\int\left(\frac1x-\frac{x}{x^2+1}-\frac{x}{(x^2+1)^2}\right)\dd x\\
&=\ln|x|-\frac12\ln(x^2+1)+\frac1{2(x^2+1)}+C.
\end{align*}
最后一项求导为 $-x/(x^2+1)^2$，结合分解即可验算。
若分解出现常数分子，则用上面的递推式；例如
\[
\int\frac{\dd x}{(x^2+1)^2}=\frac{x}{2(x^2+1)}+\frac12\arctan x+C.
\]

\section{三角函数的有理式}
\label{lec:10:trig}
先检查能否利用奇次幂凑微分或偶次幂降幂；只含二次齐次组合时，$t=\tan x$ 常比半角代换简单。复杂的 $R(\sin x,\cos x)$ 可用万能代换，但所得有理函数有时次数较高。

\begin{example}[奇次幂与偶次幂]
计算 $\int\sin^4x\cos^3x\,\dd x$ 与 $\int\sin^2x\cos^2x\,\dd x$。
\end{example}
\paragraph{思路分析}第一题保留一个 $\cos x\dd x=\dd\sin x$，其余余弦平方改写；第二题没有奇次因子，先用二倍角降幂。
\paragraph{解答}
\begin{align*}
\int\sin^4x\cos^3x\,\dd x
&=\int\sin^4x(1-\sin^2x)\,\dd\sin x
=\frac15\sin^5x-\frac17\sin^7x+C,\\
\int\sin^2x\cos^2x\,\dd x
&=\frac14\int\sin^22x\,\dd x
=\frac18\int(1-\cos4x)\,\dd x
=\frac x8-\frac{\sin4x}{32}+C.
\end{align*}
第一答案求导为 $\sin^4x(1-\sin^2x)\cos x$；第二答案求导为 $(1-\cos4x)/8$，两者恢复原式。

\begin{example}[正切代换]
求 $\displaystyle\int\frac{\dd x}{\sin x\cos^3x}$。
\end{example}
\paragraph{思路分析}提出 $\dd x/\cos^2x=\dd\tan x$，其余 $1/(\sin x\cos x)$ 能由正切表示，不需要半角代换。
\paragraph{解答}在 $\sin x\cos x\ne0$ 的区间上令 $t=\tan x$，则
\[
\frac1{\sin x\cos x}=\frac{1+t^2}{t},\qquad
\int\frac{\dd x}{\sin x\cos^3x}
=\int\left(\frac1t+t\right)\dd t
=\ln|\tan x|+\frac12\tan^2x+C.
\]
求导为 $(\sec^2x/\tan x)+\tan x\sec^2x=\sec^2x(1+\tan^2x)/\tan x$，即原式。

\begin{example}[万能代换]
计算 $\displaystyle\int\frac{\dd x}{1+\sin x+\cos x}$。
\end{example}
\paragraph{思路分析}分母同时线性含正弦、余弦，半角代换能将其变为一次分式。先选避开 $\tan(x/2)$ 极点的区间。
\paragraph{解答}令 $t=\tan(x/2)$。由半角恒等式和求导，
\[
\sin x=\frac{2t}{1+t^2},\qquad
\cos x=\frac{1-t^2}{1+t^2},\qquad
\dd x=\frac2{1+t^2}\,\dd t.
\]
故
\[
\int\frac{\dd x}{1+\sin x+\cos x}
=\int\frac{\dd t}{1+t}
=\ln|1+\tan(x/2)|+C.
\]
验算为 $\frac{1+t^2}{2(1+t)}$，正是原分母的倒数。分母为零的点不在定义域；本题半角代换的极点 $x=(2k+1)\pi$ 也恰是原分母零点，所以表达式覆盖全部连通定义区间。一般题目若代换产生原题没有的极点，须在其他局部区间求解并通过连续、可导条件衔接，不能把该点当成原题奇点。

\section{绝对值与分段衔接}
\begin{example}[全区间原函数]
求 $\displaystyle\int\ee^{|x|}\,\dd x$ 在整个 $\R$ 上的原函数。
\end{example}
\paragraph{思路分析}两侧先分别积分。原函数需要在零点连续且可导；这里零点属于定义域，两侧常数不能独立。
\paragraph{解答}两侧表达式为 $-\ee^{-x}+C_-$（$x<0$）和 $\ee^x+C_+$（$x>0$）。连续性要求 $-1+C_-=1+C_+$，所以 $C_+=C_--2$。设 $C=C_-$，得到
\[
F(x)=\begin{cases}-\ee^{-x}+C,&x<0,\\ \ee^x+C-2,&x\ge0.\end{cases}
\]
零点两侧导数均趋于 $1$；直接用差商也得 $F'(0)=1$，与 $\ee^{|0|}=1$ 相符。仅拼接连续还不足以保证任意分段答案可导，每次都应检查分界点差商。

\section{分层练习与完整解析}
以下均为自编练习，用于检验方法迁移。
\begin{exercise}[基础巩固]
计算 $\int x\ee^{x^2}\,\dd x$、$\int\ln x\,\dd x$（$x>0$）和 $\int\dd x/(4+x^2)$。
\end{exercise}
\paragraph{解析}第一题取 $u=x^2$，$\dd u=2x\dd x$，得 $\ee^{x^2}/2+C$。第二题取 $u=\ln x$、$v=x$，得 $x\ln x-\int1\,\dd x=x\ln x-x+C$。第三题写成 $\frac12\int\dd(x/2)/(1+(x/2)^2)$，得 $\frac12\arctan(x/2)+C$。三个答案分别求导即恢复被积函数。

\begin{exercise}[方法迁移]
计算 $\int x^2(1-x)^5\,\dd x$ 与 $\int(2x+3)/(x^2+2x+2)\,\dd x$。
\end{exercise}
\paragraph{解析}第一题取 $t=1-x$，$x=1-t$、$\dd x=-\dd t$，只展开低次因子：
\[
-\int(1-2t+t^2)t^5\,\dd t
=-\frac{t^6}6+\frac{2t^7}7-\frac{t^8}8+C.
\]
回代 $t=1-x$。求导为 $t^5(1-2t+t^2)=x^2(1-x)^5$。第二题分子为 $(2x+2)+1$，分母为 $(x+1)^2+1$，得 $\ln(x^2+2x+2)+\arctan(x+1)+C$；求导通分后的分子为 $2x+2+1$。

\begin{exercise}[综合提升]
计算 $\int\sqrt{1-x^2}\,\dd x$，并用分部积分得到与三角代换相同的结果。
\end{exercise}
\paragraph{解析}在 $|x|<1$ 记一个原函数为 $S$，取 $u=\sqrt{1-x^2}$、$v=x$，得
\[
S=x\sqrt{1-x^2}+\int\frac{x^2}{\sqrt{1-x^2}}\,\dd x.
\]
将 $x^2=1-(1-x^2)$ 代入，余项为 $\arcsin x-S$，所以 $2S=x\sqrt{1-x^2}+\arcsin x$，最后补 $C$。此法的构造来自目标：把余项重新拆出原根式，形成可解的循环等式。

\begin{exercise}[边界检查]
函数 $F(x)=\ln|x|+C$ 是否表示 $1/x$ 在 $\R\setminus\{0\}$ 上的全部原函数？阶梯函数 $f(x)=0$（$x<0$）、$f(x)=1$（$x\ge0$）是否在整个 $\R$ 上有原函数？
\end{exercise}
\paragraph{解析}第一问若 $C$ 被理解为一个全域固定常数，就漏掉两支常数不同的情况。第二问若有原函数，则在 $x<0$ 必为 $C_-$，在 $x>0$ 必为 $x+C_+$。连续性迫使两常数及零点值相同，但零点左右差商分别为 $0,1$，不可导，矛盾。因此没有全域原函数。

\paragraph{方法自查}先确定区间；再做恒等化简并寻找复合函数的微分；含乘积时比较分部前后的复杂度；有理式先除法、后分解；根式先配方和选支；三角式先检查奇偶幂，再考虑正切或半角。每次回代后求导，检查系数、符号、绝对值以及是否遗漏定义域的一支。


\chapter{定积分}
\label{lec:11:chapter}
本章把“变化率”与“累积量”连接起来。前置知识是连续、导数、微分中值定理与不定积分。先由分割求和定义累积量，再证明它与原函数的联系；最后学习计算、离散和式极限与积分证明。有限区间上的定积分首先指黎曼积分。无穷区间或无界被积函数的积分需要另取极限，其收敛理论见第十三章。

\section{从矩形近似到黎曼积分}
\label{lec:11:riemann}
若速度为常数，路程等于速度乘时间；若速度连续变化，可以把时间划成短段，用段内一个时刻的速度近似该段速度。曲线下方的有向面积也采用相同结构：每段宽度乘代表高度，再把全部小段相加。关键问题是：分割愈细时，结果能否摆脱取样点的选择？
\begin{definition}[黎曼积分]
设 $a<b$，$f$ 为 $[a,b]$ 上的有界实函数。分割 $P$ 为
\[a=x_0<x_1<\cdots<x_n=b,\qquad \Delta x_i=x_i-x_{i-1},\qquad |P|=\max_i\Delta x_i.\]
每段任取 $\xi_i\in[x_{i-1},x_i]$，黎曼和为 $S(P,\xi)=\sum_{i=1}^n f(\xi_i)\Delta x_i$。若存在 $I\in\R$，使对每个 $\varepsilon>0$ 都存在 $\delta>0$，只要 $|P|<\delta$，对所有分割和所有取样点都有 $|S(P,\xi)-I|<\varepsilon$，则称 $f$ 黎曼可积，并记 $I=\int_a^b f(x)\dd x$。
\end{definition}
这里控制的是最大段长，并非仅要求段数趋于无穷；也不是只检查某一组右端点和。定义保证在已经证明可积后，可以自由选择方便的分割与取样点计算积分。
\begin{figure}[htbp]
\centering
\begin{tikzpicture}[x=1.5cm,y=1.2cm]
\draw[->] (0,0)--(3.4,0) node[right] {$x$};
\draw[->] (0,0)--(0,2.3) node[above] {$y$};
\foreach \a/\b/\h in {0/0.6/0.5288,0.6/1.2/0.7592,1.2/1.8/1.22,1.8/2.4/1.9112}{\draw[fill=gray!15] (\a,0) rectangle (\b,\h);}
\draw[thick,domain=0:2.5,samples=60] plot (\x,{0.5+0.32*\x*\x});
\node at (2.8,1.8) {$y=f(x)$};
\draw[<->] (0.6,-0.7)--(1.2,-0.7) node[midway,below] {$\Delta x_i$};
\draw[dashed] (0.9,0)--(0.9,0.7592);
\node[below] at (0.9,0) {$\xi_i$};
\end{tikzpicture}
\caption{每个矩形贡献 $f(\xi_i)\Delta x_i$；最大宽度趋于零时考察总和。}
\end{figure}

\subsection{怎样判断可积}
\begin{theorem}[达布判据]
设 $f$ 有界，记 $M_i=\sup_{[x_{i-1},x_i]}f$，$m_i=\inf_{[x_{i-1},x_i]}f$，上、下和为 $U(P)=\sum M_i\Delta x_i$，$L(P)=\sum m_i\Delta x_i$。$f$ 可积当且仅当对每个 $\varepsilon>0$，存在分割 $P$ 使
\[U(P)-L(P)=\sum_i(M_i-m_i)\Delta x_i<\varepsilon.\]
\end{theorem}
\begin{proof}
细分使下和增大、上和减小；任意两个分割取共同细分，可得任一下和不超过任一上和。因此 $\sup_P L(P)\le\inf_P U(P)$，判据中的条件正是二者相等。说明这一等式为什么对应定义：固定一个上、下和之差很小的分割 $P_0$。对任意细网格 $Q$，只有跨越 $P_0$ 内部分点的小段不能直接落入某个固定段；若 $P_0$ 有 $r$ 个内部分点且 $|f|\le B$，这些小段的总长至多 $2r|Q|$，对上下界的额外影响至多 $4Br|Q|$。其余小段的任意黎曼和夹在 $P_0$ 的上、下和之间。先选 $P_0$，再令 $|Q|$ 足够小，便使所有取样点和靠近共同值。反过来，若定义成立，在细分割中分别取接近各段上确界、下确界的点，两个和都靠近 $I$，故上下和之差可任意小。
\end{proof}
\begin{proposition}[常用充分条件]
闭区间上的连续函数、单调函数，以及只有有限个间断点的有界函数均黎曼可积。
\end{proposition}
\begin{proof}
连续函数在闭区间一致连续，故网格足够细时，各段振幅 $M_i-m_i<\varepsilon/(b-a)$，上下和差小于 $\varepsilon$。若 $f$ 递增，等分的上下和差为 $(b-a)[f(b)-f(a)]/n$，趋于零；递减情形对 $-f$ 使用此结论。若有有限个间断点，先用总长很小的有限邻域包住这些点，其贡献由有界性控制；余下是有限个闭区间，分别使用连续函数的一致连续性。两部分误差均可任意小，达布判据即得可积。
\end{proof}
\begin{example}[可积与连续的边界]
比较阶梯函数与狄利克雷函数在 $[-1,1]$ 上的可积性。
\end{example}
\paragraph{解答} 阶梯函数 $H(x)=0$（$x<0$）、$H(x)=1$（$x\ge0$）只有一个间断点，故可积，积分为 $1$。狄利克雷函数 $D(x)$ 在有理点取 $1$、无理点取 $0$。每一小段同时含有两类点；全部取有理点得和 $2$，全部取无理点得和 $0$，所以不可积。可积并不推出连续，有界也不推出可积。若仅在单点 $c$ 取 $1$、其余点取 $0$，其积分为零：包含 $c$ 的小段总贡献至多 $2|P|$。因此改变有限个点不改变可积函数的积分。

\begin{example}[直接用定义计算]
由黎曼和求 $\int_0^1x^2\dd x$。
\end{example}
\paragraph{思路分析} 已由连续性知道可积，所以可选等分右端点。取样点和变成已知平方和。
\paragraph{解答}
\[S_n=\frac1n\sum_{k=1}^n\left(\frac kn\right)^2
=\frac{n(n+1)(2n+1)}{6n^3}\longrightarrow\frac13.\]
所以积分为 $1/3$。这里使用一组和计算的依据是预先证明了可积，不能用这一组和收敛反推任意有界函数可积。

\section{性质、估值与积分中值}
\label{lec:11:mean}
约定 $\int_a^a f=0$，$\int_b^a f=-\int_a^b f$。以下保序与估值公式按 $a<b$ 使用；反向积分需要先调整方向。若 $f,g$ 可积，则
\begin{align*}
\int_a^b(\alpha f+\beta g)&=\alpha\int_a^b f+\beta\int_a^b g,\\
\int_a^b f&=\int_a^c f+\int_c^b f,\\
f\le g&\ \Longrightarrow\ \int_a^b f\le\int_a^b g,\\
\left|\int_a^b f\right|&\le\int_a^b|f|.
\end{align*}
线性与保序来自黎曼和及极限，区间可加性来自在 $c$ 处分割。$|f|$ 可积可由 $||u|-|v||\le|u-v|$ 与达布判据证明；最后一个不等式由 $-|f|\le f\le|f|$ 得到。若 $m\le f\le M$，则
\[m(b-a)\le\int_a^b f\le M(b-a).\]
若连续非负的 $f$ 在某点为正，连续性使它在一段正长度区间上保持正下界，故积分严格为正。仅有“可积、非负、不恒为零”不足以得严格正值，上一节的单点函数即为反例。有向积分可能正负抵消，几何面积应计算 $\int|f|$。

\begin{theorem}[加权积分中值定理]
若 $f\in C[a,b]$，$g$ 在 $[a,b]$ 可积且 $g\ge0$，则存在 $\xi\in[a,b]$，使
\[\int_a^b f(x)g(x)\dd x=f(\xi)\int_a^b g(x)\dd x.\]
取 $g=1$，得到平均值公式 $\int_a^b f=f(\xi)(b-a)$。$g\le0$ 可对 $-g$ 应用同一结论。
\end{theorem}
\begin{proof}
设 $m=\min f$、$M=\max f$。由 $mg\le fg\le Mg$ 得 $mG\le\int fg\le MG$，其中 $G=\int g\ge0$。若 $G>0$，商 $\int fg/G$ 在 $[m,M]$ 内，由介值定理得到 $f(\xi)$。若 $G=0$，同一夹逼直接得 $\int fg=0$，任意 $\xi$ 皆可；此时若 $g$ 仅可积，不能声称 $g\equiv0$。
\end{proof}
连续性与权不变号各有用途。取 $f(x)=g(x)=x$，区间 $[-1,1]$ 上 $\int g=0$ 而 $\int fg=2/3$，说明变号权会使结论失败。

\begin{example}[不求原函数的估计]
证明 $1<\int_0^1\sqrt{1+x^4}\dd x<11/10$。
\end{example}
\paragraph{思路分析} 根式原函数复杂，而比较函数 $1$ 与 $1+x^4/2$ 都易积。需要的是整个区间上的不等式。
\paragraph{证明} 当 $x>0$，$\sqrt{1+x^4}>1$；又 $(1+x^4/2)^2=1+x^4+x^8/4>1+x^4$，两边为正，所以 $\sqrt{1+x^4}<1+x^4/2$。连续差函数在内部严格为正，积分得
\[1<\int_0^1\sqrt{1+x^4}\dd x<\int_0^1(1+x^4/2)\dd x=11/10.\]
不能把仅在 $x\to0$ 成立的渐近关系作为整个 $[0,1]$ 上的大小比较。

\begin{example}[先找到相同函数值再用罗尔定理]
设 $f\in C[0,1]$，在 $(0,1)$ 可导，且 $3\int_{2/3}^1 f(x)\dd x=f(0)$，证明存在 $\xi\in(0,1)$ 使 $f'(\xi)=0$。
\end{example}
\paragraph{思路分析} 罗尔定理需要两个相同函数值。积分给的是平均值，先将其化成点值。
\paragraph{证明} 积分中值定理给出 $c\in[2/3,1]$，使 $3\int_{2/3}^1 f=f(c)$。故 $f(c)=f(0)$ 且 $c>0$。在 $[0,c]$ 上，连续、内部可导与端点等值三个条件均成立，由罗尔定理得 $\xi\in(0,c)\subset(0,1)$。

\section{累积函数与微积分基本定理}
\label{lec:11:ftc}
\begin{theorem}[积分上限函数]
若 $f$ 在 $[a,b]$ 可积，$F(x)=\int_a^x f(t)\dd t$，则 $F$ 连续。若 $f$ 在内部点 $x$ 连续，则 $F'(x)=f(x)$。在端点连续时分别得到相应单侧导数。
\end{theorem}
\begin{proof}
由 $|f|\le B$，$|F(y)-F(x)|=|\int_x^y f|\le B|y-x|$，故连续。对非零且足够小的 $h$，
\[\frac{F(x+h)-F(x)}h-f(x)=\frac1h\int_x^{x+h}[f(t)-f(x)]\dd t.\]
不论 $h$ 正负，右侧绝对值均不超过连接 $x$ 与 $x+h$ 的区间上 $|f(t)-f(x)|$ 的上确界。连续性使这个量趋于零，故导数等于 $f(x)$。
\end{proof}
证明揭示了“短段上的平均变化率趋于当地高度”，也避免先假设未知原函数存在。连续函数因而一定有原函数；但原函数未必能用有限个初等函数写出。
\begin{theorem}[牛顿--莱布尼茨公式]
若 $f\in C[a,b]$，且 $G'=f$，则 $\int_a^b f(x)\dd x=G(b)-G(a)$。
\end{theorem}
\begin{proof}
积分上限函数 $F$ 满足 $F'=f=G'$，中值定理说明 $F-G$ 为常数。由 $F(a)=0$ 得 $F(x)=G(x)-G(a)$，令 $x=b$ 即得。
\end{proof}
分段连续函数应在分界点拆开使用公式。例如 $f(x)=x^2+4$（$x<0$）、$\ee^x$（$x\ge0$），则 $\int_{-1}^1f=13/3+\ee-1$。跳跃点不影响积分，但不能把不连续 $f$ 的积分上限函数误判为处处满足 $F'=f$。

\subsection{变限求导与含参数表达式}
若 $f$ 连续，$u,v$ 可导且值落在定义区间内，由链式法则得
\[\frac{\dd}{\dd x}\int_{u(x)}^{v(x)}f(t)\dd t=f(v(x))v'(x)-f(u(x))u'(x).\]
积分变量 $t$ 是内部的哑变量。先代入上、下限，再乘对应导数；不要把外层自变量当成积分变量。
\begin{example}[两端都变化]
求 $H(x)=\int_{x^2}^{\sin x}\ee^{t^2}\dd t$ 的导数。
\end{example}
\paragraph{解答} 被积函数连续，因此
\[H'(x)=\ee^{\sin^2x}\cos x-2x\ee^{x^4}.\]
即使 $x^2>\sin x$，有向积分与同一公式仍成立。
\begin{example}[被积式也含有外层变量]
设 $f$ 连续，求 $K(x)=\int_a^{x^2}(x^2-t)f(t)\dd t$ 的导数。
\end{example}
\paragraph{思路分析} 单独对上限求导会遗漏 $x^2$ 的变化。先用线性把外层变量提出来。
\paragraph{解答}
\begin{align*}
K(x)&=x^2\int_a^{x^2}f(t)\dd t-\int_a^{x^2}tf(t)\dd t,\\
K'(x)&=2x\int_a^{x^2}f(t)\dd t+x^2f(x^2)2x-x^2f(x^2)2x
=2x\int_a^{x^2}f(t)\dd t.
\end{align*}
若表达式无法拆开，可引用含参积分求导定理：在相关紧矩形上 $h(t,x)$ 与 $\partial h/\partial x$ 均连续时，$\frac{\dd}{\dd x}\int_a^xh(t,x)\dd t=h(x,x)+\int_a^x\partial_xh(t,x)\dd t$。这是比普通变限公式更强的结论，使用时须检查条件。

\begin{example}[缩小区间上的积分极限]
求 $\lim_{x\to0}\frac{\int_0^{x^4}\sin(t^2)\dd t}{x^{12}}$。
\end{example}
\paragraph{解答} 分子、分母均趋于零，分母导数在 $x\ne0$ 时非零。变限求导及洛必达法则给出
\[\lim_{x\to0}\frac{4x^3\sin(x^8)}{12x^{11}}=\frac13\lim_{x\to0}\frac{\sin(x^8)}{x^8}=\frac13.\]
也可用一致估计：当 $0\le t\le x^4$ 时，$\sin(t^2)/t^2$ 一致接近 $1$，故积分夹在 $(1\pm\varepsilon)\int_0^{x^4}t^2\dd t$ 之间。这个估计才使积分内替换合法。

\section{换元、分部与累积量再积分}
\label{lec:11:calculation}
\begin{theorem}[定积分换元与分部公式]
若 $f$ 连续、$\varphi\in C^1[\alpha,\beta]$ 且其值在 $f$ 的定义区间内，则
\[\int_{\varphi(\alpha)}^{\varphi(\beta)}f(x)\dd x=\int_\alpha^\beta f(\varphi(t))\varphi'(t)\dd t.\]
若 $u,v\in C^1[a,b]$，则
\[\int_a^b u(x)v'(x)\dd x=[u(x)v(x)]_a^b-\int_a^b u'(x)v(x)\dd x.\]
\end{theorem}
\begin{proof}
取 $f$ 的原函数 $G$，复合求导给出 $(G\circ\varphi)'=f(\varphi)\varphi'$，应用基本公式得换元式。由 $(uv)'=u'v+uv'$ 两边积分并移项得分部式。
\end{proof}
上述换元充分条件不要求 $\varphi$ 单调；但实际用 $x=\varphi(t)$ 解出新上下限时，选单调分支能避免漏区间。根式换元必须保留绝对值，例如 $\sqrt{\cos^2t}=|\cos t|$，只有所选区间保证 $\cos t\ge0$ 时才可去掉绝对值。
\begin{example}[根式与三角代换]
计算 $\int_0^a\sqrt{a^2-x^2}\dd x$，其中 $a>0$。
\end{example}
\paragraph{解答} 令 $x=a\sin t$，$0\le t\le\pi/2$，则根式为 $a\cos t$、$\dd x=a\cos t\dd t$。所以
\[I=a^2\int_0^{\pi/2}\cos^2t\dd t=\frac{a^2}2\int_0^{\pi/2}(1+\cos2t)\dd t=\frac{\pi a^2}4.\]
结果与半径 $a$ 的四分之一圆面积相符，提供独立的数量级检验。
\begin{example}[使求导后更简单]
计算 $\int_0^1x\ln(1+x)\dd x$。
\end{example}
\paragraph{思路分析} 对数求导成为有理函数，$x$ 易于积分，故取 $u=\ln(1+x)$、$v=x^2/2$。
\paragraph{解答}
\begin{align*}
I&=\frac12\ln2-\frac12\int_0^1\frac{x^2}{1+x}\dd x,\\
\frac{x^2}{1+x}&=x-1+\frac1{1+x},\\
I&=\frac12\ln2-\frac12\left(\frac12-1+\ln2\right)=\frac14.
\end{align*}

\begin{proposition}[消去一层累积积分]
若 $f$ 连续，则
\[\int_a^b\left(\int_a^xf(t)\dd t\right)\dd x=\int_a^b(b-t)f(t)\dd t.\]
更一般地，若 $w$ 连续、$W'=w$，则
\[\int_a^bw(x)\left(\int_a^xf(t)\dd t\right)\dd x=\int_a^b[W(b)-W(t)]f(t)\dd t.\]
\end{proposition}
\begin{proof}
令 $F(x)=\int_a^xf$，对 $\int wF$ 分部积分，得 $W(b)F(b)-\int_a^bW(x)f(x)\dd x$，整理便得。取 $w=1$ 即得第一式。
\end{proof}
三角区域 $a\le t\le x\le b$ 中，固定 $t$ 后 $x$ 的长度为 $b-t$，说明公式的几何来源。本章用一元分部积分给出证明，不以未学的二重积分换序代替论证。
\begin{example}[内层原函数无法初等表示]
求 $\int_0^1x\left(\int_0^x\ee^{-t^2}\dd t\right)\dd x$ 的单重积分表达式。
\end{example}
\paragraph{解答} 取 $W(x)=x^2/2$，上式给出 $\frac12\int_0^1(1-t^2)\ee^{-t^2}\dd t$。再由 $(t\ee^{-t^2})'=\ee^{-t^2}-2t^2\ee^{-t^2}$ 可得结果为 $\frac14\int_0^1\ee^{-t^2}\dd t+\frac1{4\ee}$。原函数非初等并不妨碍化简；不应捏造一个初等闭式。

\section{对称、周期与华里士递推}
\label{lec:11:symmetry}
对可积函数，换元 $x=a+b-t$ 给出
\[\int_a^bf(x)\dd x=\int_a^bf(a+b-x)\dd x.\]
先写出变换后的完整被积式，再判断相加、相减是否简化。奇偶性是其在 $[-a,a]$ 的特例：奇函数积分为零，偶函数积分为正半区间积分的两倍。
\begin{example}[配对消去复杂分母]
求 $I=\int_0^{\pi/2}\frac{\sin x}{\sin x+\cos x}\dd x$ 及 $J=\int_{-1}^1\frac{x^{10}}{1+\ee^x}\dd x$。
\end{example}
\paragraph{解答} 分母在区间上恒正。第一式换元 $x=\pi/2-t$ 得另一表达式 $I=\int_0^{\pi/2}\frac{\cos x}{\sin x+\cos x}\dd x$。相加得 $2I=\pi/2$，故 $I=\pi/4$。第二式换元 $x=-t$ 得 $J=\int_{-1}^1x^{10}\ee^x/(1+\ee^x)\dd x$，相加得 $2J=\int_{-1}^1x^{10}\dd x=2/11$，故 $J=1/11$。
\begin{proposition}[对称权中的一次因子]
若 $q$ 可积且 $q(a+b-x)=q(x)$，则 $\int_a^b xq(x)\dd x=\frac{a+b}2\int_a^bq(x)\dd x$。
\end{proposition}
\begin{proof}
变换后积分为 $\int_a^b(a+b-x)q(x)\dd x$，与原式相加即可。
\end{proof}
例如 $q(x)=f(\sin x)$ 在 $[0,\pi]$ 满足条件。但 $q(x)=\sin x/(2+\cos x)$ 变换后分母为 $2-\cos x$，不能使用这个公式。检查的是整体对称关系，而不是表达式中是否出现正弦。

\begin{proposition}[周期积分与平均值]
设 $T>0$，$f(x+T)=f(x)$，且 $f$ 在一个周期上黎曼可积，则任何整周期积分均为 $A=\int_0^Tf$，并且
\[\lim_{x\to+\infty}\frac1x\int_0^xf(t)\dd t=\frac AT.\]
\end{proposition}
\begin{proof}
周期延拓使各有界区间可积且 $|f|\le B$。拆分并平移积分得 $\int_c^{c+T}f=\int_0^Tf$。写 $x=nT+r$、$0\le r<T$，则
\[\frac1x\int_0^xf=\frac nxA+\frac1x\int_0^rf.\]
余项绝对值不超过 $BT/x$，趋于零，而 $n/x=(x-r)/(Tx)\to1/T$。
\end{proof}
\begin{example}[长时间平均]
求 $\lim_{x\to\infty}x^{-1}\int_0^x|\sin t|\dd t$。
\end{example}
\paragraph{解答} $|\sin t|$ 连续、周期为 $\pi$，一个周期积分为 $2$，所以平均值极限为 $2/\pi$。对这个 $\infty/\infty$ 型式子直接使用洛必达无法得到极限，因为分子导数 $|\sin x|$ 不趋于常数。

\begin{theorem}[华里士积分]
令 $I_n=\int_0^{\pi/2}\sin^nx\dd x$，$n$ 为非负整数，则
\[I_n=\int_0^{\pi/2}\cos^nx\dd x,\quad I_0=\frac\pi2,\quad I_1=1,\quad I_n=\frac{n-1}n I_{n-2}\ (n\ge2).\]
因而
\[I_{2m}=\frac{(2m-1)!!}{(2m)!!}\frac\pi2,\qquad I_{2m+1}=\frac{(2m)!!}{(2m+1)!!}.\]
其中双阶乘每次减 $2$，并约定 $0!!=(-1)!!=1$。
\end{theorem}
\begin{proof}
正余弦互换来自 $x=\pi/2-t$。对 $n\ge2$，将一因子 $\sin x\dd x=-\dd(\cos x)$，分部得
\[I_n=[-\sin^{n-1}x\cos x]_0^{\pi/2}+(n-1)\int_0^{\pi/2}\sin^{n-2}x\cos^2x\dd x
=(n-1)(I_{n-2}-I_n).\]
边界项为零，移项得递推。偶数阶最终归至 $I_0$，奇数阶归至 $I_1$，分别乘上各递推系数即得通项。
\end{proof}
\begin{example}[根式转成递推差]
求 $\int_0^1x^8\sqrt{1-x^2}\dd x$。
\end{example}
\paragraph{解答} 令 $x=\sin t$、$t\in[0,\pi/2]$，根式和微分各贡献一个 $\cos t$。所以积分为
\[I_8-I_{10}=\left(1-\frac9{10}\right)I_8
=\frac1{10}\frac{7\cdot5\cdot3\cdot1}{8\cdot6\cdot4\cdot2}\frac\pi2=\frac{7\pi}{512}.\]
该计算也提示不能直接照抄参考资料的数值：课程资料第115页把 $7!!/10!!$ 算错，其所列 $7\pi/768$ 不正确。

\subsection{正弦与余弦混合幂}
令 $J_{m,n}=\int_0^{\pi/2}\sin^m x\cos^n x\dd x$，其中 $m,n$ 为非负整数。若有奇数次幂，留出一个正弦或余弦作微分，用 $\sin^2x+\cos^2x=1$ 将其余偶次幂改为多项式；若全为偶数，则可逐步降阶。具体地，对 $m\ge2$，取 $u=\sin^{m-1}x$、$\dd v=\sin x\cos^nx\dd x$，有 $v=-\cos^{n+1}x/(n+1)$，两端乘积为零，因此
\[J_{m,n}=\frac{m-1}{n+1}\int_0^{\pi/2}\sin^{m-2}x\cos^{n+2}x\dd x
=\frac{m-1}{n+1}(J_{m-2,n}-J_{m,n}).\]
整理得 $J_{m,n}=(m-1)J_{m-2,n}/(m+n)$；也可互换 $m,n$ 降低余弦次数。
\begin{example}[奇次幂留一个微分]
计算 $\int_0^{\pi/2}\sin^3x\cos^2x\dd x$。
\end{example}
\paragraph{解答} 写成 $(1-\cos^2x)\cos^2x\sin x\dd x$。令 $u=\cos x$，上下限从 $1$ 到 $0$，$\dd u=-\sin x\dd x$，积分为 $\int_0^1(u^2-u^4)\dd u=1/3-1/5=2/15$。递推也给 $J_{3,2}=(2/5)J_{1,2}=(2/5)(1/3)$。奇次幂适合直接换元，偶次幂适合递推，方法依据是能否留出单个微分因子。

\section{和式极限与积分极限}
\label{lec:11:limits}
识别和式时依次确定步长、取样点、固定函数与区间。例如 $\frac1n\sum f(2k/n)$ 对应 $\frac12\int_0^2f$，因为真正步长为 $2/n$。求和函数随 $n$ 改变时，还须控制总误差，不能以“每项近似”跳过这一条件。
\begin{example}[标准和与扰动和]
求 $\lim_{n\to\infty}\sum_{k=1}^n n/[n^2+k(k-1)]$。
\end{example}
\paragraph{解答} 用 $(k-1)^2\le k(k-1)\le k^2$ 直接夹逼：
\[\frac1n\sum_{k=1}^n\frac1{1+(k/n)^2}\le S_n\le\frac1n\sum_{k=1}^n\frac1{1+((k-1)/n)^2}.\]
两边为连续函数 $1/(1+t^2)$ 的右端、左端黎曼和，都趋于 $\pi/4$。夹逼便得同一极限。另一通用原则是若 $\sup_k|a_{n,k}-b_{n,k}|\le\eta_n\to0$，则 $|n^{-1}\sum(a_{n,k}-b_{n,k})|\le\eta_n$。每个固定 $k$ 的差趋于零不足以保证求和后误差消失。
\begin{proposition}[积分的高次根]
若 $h\ge0$ 在 $[a,b]$ 连续且 $a<b$，则
\[\lim_{n\to\infty}\left(\int_a^bh(t)^n\dd t\right)^{1/n}=\max_{[a,b]}h=M.\]
\end{proposition}
\begin{proof}
$M=0$ 时直接成立。若 $M>0$，上界为 $M(b-a)^{1/n}$。任给 $0<\varepsilon<M$，最大值点附近有正长度的子区间 $J$，在其中 $h>M-\varepsilon$。记其长为 $\ell>0$，下界为 $(M-\varepsilon)\ell^{1/n}$。先令 $n\to\infty$，再令 $\varepsilon\downarrow0$，得到结论；邻域长度必须固定，不能随 $n$ 任意缩小后忽略。
\end{proof}
\begin{example}[高次根模型]
求 $\lim_n(\int_1^2\ee^{-nt^2}\dd t)^{1/n}$。
\end{example}
\paragraph{解答} 取 $h(t)=\ee^{-t^2}$，最大值为 $\ee^{-1}$，即为极限。具体地，上界为 $\ee^{-n}$；在 $[1,1+\eta]$ 上积分得下界 $\eta\ee^{-n(1+\eta)^2}$，开根、先取 $n$ 极限再取 $\eta\downarrow0$ 得同一结果。
\begin{remark}[交换极限与积分的边界]
若 $f_n\to f$ 在 $[a,b]$ 一致收敛且各 $f_n$ 可积，则 $f$ 可积且 $|\int f_n-\int f|\le(b-a)\sup|f_n-f|\to0$。可积性可由达布判据和 $f_n$ 的一致逼近证明。仅逐点收敛不够：连续三角尖峰 $f_n(x)=n\max\{1-n|x-1/n|,0\}$（$n\ge2$）在 $[0,1]$ 逐点趋零，积分却均为 $1$。尖峰高度增大而底宽减小，面积保持不变。
\end{remark}

\section{积分不等式与辅助函数}
\label{lec:11:inequality}
常用路线是先逐点比较，再积分；若已知累积积分的界，先构造积分上限函数，再分部积分提取矩；若目标含平方，考虑非负平方积分。存在唯一性则把等式两端相减成连续函数，用端点符号和严格单调性。
\begin{theorem}[积分型柯西--施瓦茨不等式]
对连续实函数 $f,g$，
\[\left(\int_a^bfg\right)^2\le\left(\int_a^bf^2\right)\left(\int_a^bg^2\right).\]
等号当且仅当存在不全为零的常数 $\alpha,\beta$，使 $\alpha f+\beta g\equiv0$。
\end{theorem}
\begin{proof}
记 $A=\int f^2$、$B=\int fg$、$C=\int g^2$。若 $C=0$，连续非负性使 $g\equiv0$，故 $B=0$。若 $C>0$，$\int(f-\lambda g)^2=A-2\lambda B+\lambda^2C\ge0$，取 $\lambda=B/C$ 得 $A-B^2/C\ge0$。等号恰使该平方积分为零，连续性使 $f=(B/C)g$；退化情形也包含在所给线性相关条件内。
\end{proof}
\begin{example}[由累积量取得平方下界]
设 $f\in C[0,1]$ 且 $\int_x^1f(t)\dd t\ge(1-x^3)/2$。证明 $\int_0^1f^2\ge5/12$。
\end{example}
\paragraph{思路分析} 不能对不等式两边求导来比较 $f$。将给定的累积函数再积分，通过分部得到 $f$ 的一次矩，随后与 $x$ 配对使用柯西不等式。
\paragraph{证明} 令 $F(x)=\int_x^1f$，$F'=-f$、$F(1)=0$。所以
\[\int_0^1xf(x)\dd x=[-xF(x)]_0^1+\int_0^1F(x)\dd x\ge\frac38.\]
再由 $(\int xf)^2\le(\int x^2)(\int f^2)=\frac13\int f^2$ 得 $\int f^2\ge27/64>5/12$。

\begin{example}[存在性与唯一性分开证明]
设连续 $f>0$ 于 $[0,1]$，证明存在唯一 $c\in(0,1)$ 使 $\int_0^cf=\int_c^1 1/f$。
\end{example}
\paragraph{证明} 连续性及恒正保证 $1/f$ 连续。令 $H(x)=\int_0^xf-\int_x^1 1/f$。$H$ 连续，$H(0)<0<H(1)$，介值定理给出内部零点。又 $H'=f+1/f>0$，$H$ 严格递增，所以零点至多一个。

\begin{proposition}[凸函数积分的两侧界]
若 $f$ 在 $[a,b]$ 连续且凸，则
\[f((a+b)/2)\le\frac1{b-a}\int_a^bf(t)\dd t\le\frac{f(a)+f(b)}2.\]
\end{proposition}
\begin{proof}
凸性给出 $2f((a+b)/2)\le f(t)+f(a+b-t)$，积分并用对称换元得左界。右界由 $f(t)\le[(b-t)f(a)+(t-a)f(b)]/(b-a)$ 积分得到。严格凸且 $a<b$ 时两界均严格，因为内部弦不等式严格，连续差在正长度区间为正。
\end{proof}
\begin{proposition}[单调权的第二中值公式]
若 $f$ 连续、$g\in C^1[a,b]$ 且单调递增，则存在 $\xi\in[a,b]$ 使
\[\int_a^bfg=g(a)\int_a^\xi f+g(b)\int_\xi^bf.\]
\end{proposition}
\begin{proof}
令 $F(x)=\int_a^xf$，则 $F(a)=0$、$F'=f$。分部得 $\int_a^bfg=g(b)F(b)-\int_a^bFg'$。单调递增保证 $g'\ge0$，对连续 $F$ 与权 $g'$ 使用加权中值定理，得 $\int Fg'=F(\xi)[g(b)-g(a)]$，代回整理便得结论；若 $g'\equiv0$，同一证明仍成立。这里只陈述并证明连续可导单调权的版本。
\end{proof}
\begin{example}[反函数面积与杨格不等式]
设 $f\in C^1[0,\infty)$ 严格递增、$f(0)=0$。若 $a\ge0$ 且 $b\ge0$ 属于 $f$ 的值域，证明 $\int_0^af+\int_0^bf^{-1}\ge ab$，并确定等号。
\end{example}
\paragraph{证明} 先令 $c=f^{-1}(b)$；这个定义需要值域条件。连续严格单调使反函数连续。换元 $y=f(x)$，再分部，得 $\int_0^bf^{-1}(y)\dd y=\int_0^cxf'(x)\dd x=cb-\int_0^cf$。于是目标差为 $\int_c^a[f(x)-b]\dd x$。$a>c$ 时被积式在 $(c,a]$ 为正；$a<c$ 时它在 $[a,c)$ 为负，反向积分后为正。故差非负，等号恰为 $a=c$，即 $b=f(a)$。例如 $f(x)=1-\ee^{-x}$ 的值域是 $[0,1)$，不能对 $b=2$ 写反函数积分。

\section{分层练习与完整解析}
\begin{exercise}[基础巩固]
求 $\int_{-1}^1|x|\dd x$，并求 $\frac{\dd}{\dd x}\int_{\cos x}^{x^3}\sqrt{1+t^2}\dd t$。
\end{exercise}
\paragraph{解析} 面积需用绝对值，偶性给出 $2\int_0^1x\dd x=1$。第二问上下限求导后为 $3x^2\sqrt{1+x^6}+\sin x\sqrt{1+\cos^2x}$；下限导数为 $-\sin x$，与公式中的减号抵消。
\begin{exercise}[方法迁移]
求 $\lim_n n^{-1}\sum_{k=1}^n\sqrt{1+2k/n}$，及 $\int_0^\pi x\sin^3x\dd x$。
\end{exercise}
\paragraph{解析} 第一问是 $f(t)=\sqrt{1+2t}$ 的黎曼和，积分为 $[(1+2t)^{3/2}/3]_0^1=\sqrt3-1/3$。第二问 $\sin^3x$ 关于 $\pi/2$ 对称，因此积分为 $(\pi/2)\int_0^\pi\sin^3x\dd x=\pi I_3=2\pi/3$，其中 $I_3=(2/3)I_1=2/3$。
\begin{exercise}[综合提升]
设 $f\in C^2[0,1]$，$f(0)=f(1)=f'(0)=0$、$f'(1)=1$，证明 $\int_0^1(f'')^2\dd x\ge4$，并给出等号函数。
\end{exercise}
\paragraph{思路分析} 端点给出 $f''$ 的两个矩。寻找一次函数权重，使配方或柯西不等式给出尽可能清楚的界。
\paragraph{证明} 基本公式与分部得 $\int_0^1f''=1$、$\int_0^1xf''=[xf']_0^1-\int_0^1f'=1$，故 $\int_0^1(x-1/3)f''=2/3$。而 $\int_0^1(x-1/3)^2=1/9$，柯西不等式给出 $4/9\le(1/9)\int(f'')^2$。等号要求 $f''=6x-2$，由 $f'(0)=f(0)=0$ 得 $f=x^3-x^2$，它满足另两端点条件且平方积分为 $4$。
\begin{exercise}[证明边界]
设 $f\in C^1[a,b]$，$f(a)=0$。证明 $\int_a^bf^2\le\frac12(b-a)^2\int_a^b(f')^2$。
\end{exercise}
\paragraph{证明} $f(x)=\int_a^xf'$，柯西不等式给出 $f(x)^2\le(x-a)\int_a^x(f')^2\le(x-a)\int_a^b(f')^2$。再对 $x$ 积分，系数为 $\int_a^b(x-a)\dd x=(b-a)^2/2$。这里给的是有效界，并未宣称常数最优。
\paragraph{章末自查} 能否解释定义中的“所有分割与所有取样点”？能否在求导前区分积分变量、积分限与外层参数？能否说明扰动和式的总误差为什么消失？证明严格不等式时，能否指出提供正长度区间的条件？这些问题比单独背出公式更能检验方法是否可迁移。


\chapter{定积分的应用}
\label{lec:12:chapter}

本章以前面的定积分、换元、分部积分和导数最值为基础。问题的核心是把一个几何量或物理量分成小部分，找到每一部分的线性主部，再累加。积分技巧通常在建模完成后才进入计算；若微元、边界或重叠关系错误，计算再准确也得不到所求的量。

\section{微元法：先解释被积函数}
\label{lec:12:element}
设总量在区间上可相加，$q$ 在 $[a,b]$ 上黎曼可积。在分割 $a=x_0<\cdots<x_n=b$ 上，若小部分满足
\[
 \Delta Q_i=q(\xi_i)\Delta x_i+\varepsilon_i\Delta x_i,
 \qquad \max_i|\varepsilon_i|\longrightarrow0
\]
（当分割的最大长度趋于零），则误差总和的绝对值不超过
$(b-a)\max_i|\varepsilon_i|$，从而
\[
 Q=\lim\sum_i q(\xi_i)\Delta x_i=\int_a^b q(x)\,\dd x.
\]
微元记号 $\dd Q=q(x)\dd x$ 表示这种线性主部，并不声称一个有宽度的小部分恰好等于其近似值。

建模时依次回答：所求量能否相加；选哪个变量使薄片内的关键量近似不变；薄片的形状是什么；微元中的每一因子表示什么；积分限如何覆盖整个对象且不重复。最后用量纲和非负性检查。例如 $\rho gA(y)h(y)\dd y$ 的单位为力乘长度，适合做功；$\rho gh(y)w(y)\dd y$ 的单位为力，适合静水压力。

\section{平面图形的面积}
\label{lec:12:area}
\subsection{直角坐标：上下边界与左右边界}
若连续边界在 $[a,b]$ 上满足 $f(x)\ge g(x)$，竖直薄条的高度为 $f(x)-g(x)$，宽度为 $\dd x$，所以
\[
 A=\int_a^b[f(x)-g(x)]\,\dd x.
\]
连续性保证在足够窄的条内，高度的最大值与最小值之差一致趋于零，故上下矩形面积夹逼到同一积分。这也给出了面积公式的依据。若上下位置交换，应先求交点并分段，或者写 $\int_a^b|f-g|\dd x$ 后再去绝对值。对水平薄条，公式变为
\[
 A=\int_c^d[x_{\text{右}}(y)-x_{\text{左}}(y)]\,\dd y.
\]
这里左右边界不一定都是同一个函数的反函数。选择变量的标准是边界易写、分段较少，而非必须对 $x$ 积分。

\begin{example}[竖条与横条互相检验]
求 $y=x$ 与 $y=x^2$ 围成的面积。
\end{example}
\paragraph{思路分析}联立得到 $x=0,1$。在 $0<x<1$ 上 $x>x^2$，所以竖条为直线减抛物线；横条则要解出 $x=y$ 与 $x=\sqrt y$。
\paragraph{解答}
\[
 A=\int_0^1(x-x^2)\dd x=\frac12-\frac13=\frac16.
\]
也有 $A=\int_0^1(\sqrt y-y)\dd y=2/3-1/2=1/6$。两种选择描述同一区域。
\begin{figure}[htbp]
\centering
\begin{tikzpicture}[scale=3]
\draw[->](-.1,0)--(1.2,0) node[right]{$x$};
\draw[->](0,-.1)--(0,1.2) node[above]{$y$};
\fill[gray!15] (0,0)--(1,1)--plot[domain=1:0,samples=40](\x,{\x*\x})--cycle;
\draw (0,0)--(1,1) node[above left]{$y=x$};
\draw[domain=0:1,samples=40]plot(\x,{\x*\x}) node[right]{$y=x^2$};
\draw[thick](.55,.3025)--(.55,.55);
\draw[dashed](.55,0)--(.55,.3025);
\node[below] at(.55,0){$x$}; \node[below]at(1,0){$1$};
\node[right]at(.55,.42){$x-x^2$};
\end{tikzpicture}
\caption{竖条微元的高度由同一个横坐标处的两条边界决定。}
\end{figure}

\begin{example}[边界换位必须分段]
求 $y=x^3$、$y=x$ 所围全部有限区域的总面积。
\end{example}
\paragraph{思路分析}交点为 $-1,0,1$，包含两个区域。负半轴上 $x^3>x$，正半轴上 $x>x^3$。
\paragraph{解答}
\[
 A=\int_{-1}^0(x^3-x)\dd x+\int_0^1(x-x^3)\dd x
 =2\left(\frac12-\frac14\right)=\frac12.
\]
直接算 $\int_{-1}^1(x-x^3)\dd x=0$ 得到的是有向积分，正负抵消后不能表示几何面积。

\subsection{参数方程与行进方向}
若 $x=x(t)$ 在一段上严格单调，且曲线在 $x$ 轴上方，则其下方面积可由换元算出。参数从左到右时为 $\int y(t)x'(t)\dd t$；从右到左时该积分为负，要反转上下限。若 $x(t)$ 改变单调方向，先分成能表示单值边界的弧段，分别按区域计算，不能对一个任意闭合参数轨迹直接套“曲线下方面积”。

\begin{example}[椭圆的参数面积]
求 $x^2/a^2+y^2/b^2=1$ 所围面积，$a,b>0$。
\end{example}
\paragraph{思路分析}利用四象限对称，只算第一象限。参数 $x=a\cos t,y=b\sin t$ 在 $0\le t\le\pi/2$ 从右向左行进，因此面积需要负号。
\paragraph{解答}
\[
 A=4\left[-\int_0^{\pi/2}b\sin t(-a\sin t)\dd t\right]
 =4ab\int_0^{\pi/2}\frac{1-\cos2t}{2}\dd t=\pi ab.
\]
当 $a=b$ 时回到圆面积 $\pi a^2$，给出了一个几何检验。

\subsection{极坐标的扇形微元}
在极角区间 $[\alpha,\beta]$ 内，设 $r(\theta)\ge0$ 连续，射线扫过的区域单次覆盖。角宽 $\Delta\theta$ 的小区域由最小、最大半径的扇形夹住，扇形面积为半径平方乘角宽的一半。于是
\[
 A=\frac12\int_\alpha^\beta r(\theta)^2\dd\theta.
\]
若同一射线上 $0\le r_{\text{内}}\le r_{\text{外}}$，则
\[
 A=\frac12\int_\alpha^\beta(r_{\text{外}}^2-r_{\text{内}}^2)\dd\theta.
\]
这些公式要求描述实际径向区间并单次扫描。负极径表示反向射线上的点；把任意角区间代入平方公式可能重复覆盖。玫瑰线先查 $r=0$，再判断相邻零点间是否恰好一瓣。

\begin{example}[一瓣与两条极坐标边界]
求 $r=\sin3\theta$ 在 $0\le\theta\le\pi/3$ 的一瓣面积；再求 $1\le r\le2\cos\theta$ 所描述区域的面积。
\end{example}
\paragraph{思路分析}第一问半径非负且两端为零。第二问要先使外半径不小于内半径，即 $\cos\theta\ge1/2$，故角区间为 $[-\pi/3,\pi/3]$。
\paragraph{解答}
\[
 A_1=\frac12\int_0^{\pi/3}\sin^23\theta\dd\theta=\frac\pi{12}.
\]
第二问由对称性，
\[
 A_2=\int_0^{\pi/3}(4\cos^2\theta-1)\dd\theta
 =\left[\theta+\sin2\theta\right]_0^{\pi/3}
 =\frac\pi3+\frac{\sqrt3}{2}.
\]
积分限来自两边界的交点，不能沿用整个圆的角区间。

\section{体积：截面法与柱壳法}
\label{lec:12:volume}
\subsection{截面、圆盘与垫片}
设垂直于 $x$ 轴的截面面积为连续函数 $A(x)$，薄片厚度为 $\dd x$，则
\[
 V=\int_a^b A(x)\dd x.
\]
旋转轴为水平直线 $y=c$ 时，取竖条旋转，得到垂直于轴的截面，沿轴方向用 $\dd x$ 累加。若竖条 $g(x)\le y\le f(x)$ 完全在轴同侧，截面为垫片，外、内半径分别为两端到轴距离中的最大、最小值：
\[
 V=\pi\int_a^b(R^2-r^2)\dd x.
\]
若 $g(x)\le c\le f(x)$，条跨轴，截面是圆盘，半径为
$R=\max\{c-g(x),f(x)-c\}$，内半径为零。不能把两个半径的平方相减。绕竖直轴 $x=c$ 时，用水平条得到垫片，以 $\dd y$ 累加。

\begin{example}[非旋转体的已知截面]
某立体在 $0\le x\le1$ 的垂直截面是正方形，其边长为 $x-x^2$，求体积。
\end{example}
\paragraph{思路分析}截面面积是边长平方，不含 $\pi$。薄片体积由“截面积乘厚度”决定。
\paragraph{解答}
\[
 V=\int_0^1(x-x^2)^2\dd x=\frac13-\frac12+\frac15=\frac1{30}.
\]

\begin{example}[垫片与偏移轴]
区域 $D$ 由 $y=x,y=x^2$ 围成，求绕 $x$ 轴与绕 $y=-1$ 旋转的体积。
\end{example}
\paragraph{思路分析}两轴都是水平轴，用竖条。对 $x$ 轴半径为 $x,x^2$；对 $y=-1$ 半径为 $x+1,x^2+1$。
\paragraph{解答}
\begin{align*}
 V_x&=\pi\int_0^1(x^2-x^4)\dd x=\frac{2\pi}{15},\\
 V_{-1}&=\pi\int_0^1[(1+x)^2-(1+x^2)^2]\dd x\\
 &=\pi\int_0^1(2x-x^2-x^4)\dd x=\frac{7\pi}{15}.
\end{align*}
平移旋转轴改变的是每个点到轴的距离，不能只给最后体积加一个常数。

\subsection{柱壳：平行轴的薄条}
绕 $y$ 轴时，位于右半平面的竖条旋转成薄柱壳。半径为 $x$，高度为 $h(x)$，厚度为 $\dd x$。圆柱体积差
\[
 \pi[(x+\Delta x)^2-x^2]h(x)
 =2\pi xh(x)\Delta x+\pi h(x)(\Delta x)^2
\]
的线性主部为 $2\pi xh(x)\dd x$。若区域在轴同侧、薄条构成的壳互不重复，则
\[
 V=2\pi\int \rho h\,\dd\rho.
\]
更具体地，绕 $x=c$ 且区域在其右侧时，$\rho=x-c$，竖条高度为上减下；绕 $y=c$ 且区域在其上侧时，用水平条，半径 $y-c$，壳高为右减左。

\begin{center}
\begin{tabular}{lll}
\toprule
方法 & 薄条与旋转轴 & 积分推进方向\\
\midrule
圆盘、垫片 & 垂直 & 平行于轴\\
柱壳 & 平行 & 垂直于轴\\
\bottomrule
\end{tabular}
\end{center}
例如绕 $x$ 轴对 $x$ 积分用垫片，绕 $y$ 轴对 $x$ 积分用柱壳。说清“薄条方向”与“变量推进方向”可避免把两者混淆。

\begin{example}[同一旋转体的两种建模]
区域 $D$ 由 $y=x,y=x^2$ 围成，求绕 $y$ 轴的体积。
\end{example}
\paragraph{思路分析}竖条平行于轴，直接用壳高 $x-x^2$；横条垂直于轴，半径为 $\sqrt y$ 与 $y$。
\paragraph{解答}
\[
 V=2\pi\int_0^1x(x-x^2)\dd x=2\pi\left(\frac13-\frac14\right)=\frac\pi6.
\]
垫片法给 $V=\pi\int_0^1(y-y^2)\dd y=\pi/6$。选择方法时不仅看能否反解，还应比较最终积分难度。

\begin{example}[跨轴重叠]
求 $y=1-x^2$ 与 $x$ 轴围成区域绕 $y$ 轴的体积。
\end{example}
\paragraph{思路分析}区域跨过 $y$ 轴。左右两侧旋转后产生同一个实体；若对 $[-1,1]$ 积分 $2\pi|x|(1-x^2)$，会把每个壳算两次。横条旋转成圆盘最直接。
\paragraph{解答}在 $0\le y\le1$ 上，横条从 $-\sqrt{1-y}$ 到 $\sqrt{1-y}$，圆盘半径为 $\sqrt{1-y}$，故
\[
 V=\pi\int_0^1(1-y)\dd y=\frac\pi2.
\]
只用右半区域做柱壳也有 $2\pi\int_0^1x(1-x^2)\dd x=\pi/2$。一般跨轴且不对称时，应对同一半径处的高度区间取并集，先合并重叠再积分；不能盲目只取某一侧。

\section{曲线的弧长}
\label{lec:12:length}
\subsection{从折线逼近到弧长积分}
设 $\gamma(t)=(x(t),y(t))$ 在 $[\alpha,\beta]$ 上连续可微。曲线长度定义为沿参数顺序内接折线长度的上确界。记速度 $v(t)=\sqrt{x'(t)^2+y'(t)^2}$，则
\begin{theorem}[参数弧长]
连续可微曲线的行进长度为
\[
 L=\int_\alpha^\beta v(t)\dd t.
\]
\end{theorem}
\begin{proof}
对一小段 $[u,u+\Delta t]$，微积分基本定理给出
$\gamma(u+\Delta t)-\gamma(u)=\int_u^{u+\Delta t}\gamma'(t)\dd t$。
由向量三角不等式，弦长不超过 $\int_u^{u+\Delta t}v(t)\dd t$，故任意折线长不超过全积分。
另一方面，$\gamma'$ 在闭区间上一致连续。若分割网格足够小，则这段增量与 $\gamma'(u)\Delta t$ 的距离不超过 $\varepsilon\Delta t$。由长度差不超过向量差，弦长与 $v(u)\Delta t$ 的差也不超过 $\varepsilon\Delta t$。求和后总误差不超过 $\varepsilon(\beta-\alpha)$，折线长趋于速度积分。上、下两方面合并即得结论。
\end{proof}
速度可在孤立点为零，例如摆线尖点，不影响上述证明。若参数重复遍历同一曲线，积分计算的是全部行进长度；求几何曲线长度须选单次遍历区间。

取参数 $t=x$，得到
\[
 L=\int_a^b\sqrt{1+f'(x)^2}\dd x.
\]
若曲线用 $x=g(y)$ 表示，则改为 $\int_c^d\sqrt{1+g'(y)^2}\dd y$。极坐标参数满足
\[
 x'=r'\cos\theta-r\sin\theta,\qquad
 y'=r'\sin\theta+r\cos\theta,
\]
交叉项相消后 $x'^2+y'^2=r'^2+r^2$，因而
\[
 L=\int_\alpha^\beta\sqrt{r^2+r'^2}\dd\theta.
\]
弧长微元 $\dd s$ 总为非负；对增大的参数写 $\dd s=v(t)\dd t$。根号开方必须保留绝对值。

\begin{example}[直角坐标弧长]
求 $y=\frac23x^{3/2}$ 在 $0\le x\le1$ 上的弧长。
\end{example}
\paragraph{思路分析}导数为 $\sqrt x$，所以根号内部只含一次项，直接换元即可。
\paragraph{解答}
\[
 L=\int_0^1\sqrt{1+x}\dd x
 =\left[\frac23(1+x)^{3/2}\right]_0^1
 =\frac23(2\sqrt2-1).
\]
端点间直线距离为 $\sqrt{13}/3$，结果不小于该距离，符合几何下界。

\begin{example}[摆线与半角绝对值]
求摆线 $x=a(t-\sin t),y=a(1-\cos t)$ 的一拱长度，$a>0,0\le t\le2\pi$。
\end{example}
\paragraph{思路分析}平方和先用半角公式化简，再查半角在参数区间的符号；不要分开积分两个导数。
\paragraph{解答}
\[
 v=a\sqrt{(1-\cos t)^2+\sin^2t}
 =a\sqrt{2-2\cos t}=2a|\sin(t/2)|.
\]
在本区间 $\sin(t/2)\ge0$，故
$L=2a\int_0^{2\pi}\sin(t/2)\dd t=8a$。

\begin{example}[极坐标指数曲线]
求 $r=\ee^{k\theta}$ 在 $0\le\theta\le b$ 上的弧长，$b>0$，$k$ 为实数。
\end{example}
\paragraph{思路分析}极径及其导数只差常数倍，所以速度根号可直接提出极径。单独处理 $k=0$ 可防止除零。
\paragraph{解答}若 $k\ne0$，
\[
 L=\sqrt{1+k^2}\int_0^b\ee^{k\theta}\dd\theta
 =\frac{\sqrt{1+k^2}}{k}(\ee^{kb}-1).
\]
该式即使 $k<0$ 也为正。若 $k=0$，曲线为单位圆上的行进轨迹，长度为 $b$；当 $b>2\pi$ 时有重复遍历，整个几何圆的长度仍为 $2\pi$。

\section{旋转曲面的侧面积}
\label{lec:12:surface}
曲线一小段用线段近似，绕轴形成圆台。圆台的侧面积为 $\pi(R+r)l$，其中 $l$ 是斜母线长度，而非轴向高度。令小段趋于零，得到
\[
 \dd S=2\pi\rho\,\dd s,\qquad S=2\pi\int\rho\,\dd s,
\]
其中 $\rho$ 是曲线点到旋转轴的距离。对 $y=f(x)$ 绕 $y=c$ 旋转，
\[
 S=2\pi\int_a^b|f(x)-c|\sqrt{1+f'(x)^2}\dd x.
\]
绕 $x=c$ 则将半径换成 $|x-c|$。参数曲线可直接代入 $\dd s=v(t)\dd t$。公式计算扫过的曲面面积，重复生成同一面也会重复计数，因此要选不重复的母线；侧面积也不包含封闭端面的面积。

\begin{example}[圆台与曲面计算]
求 $y=x^3$ 在 $0\le x\le1$ 上绕 $x$ 轴旋转的侧面积。
\end{example}
\paragraph{思路分析}半径是 $x^3$，母线微元含 $\sqrt{1+9x^4}$，恰好能用 $u=1+9x^4$ 换元。
\paragraph{解答}
\[
 S=2\pi\int_0^1x^3\sqrt{1+9x^4}\dd x
 =\frac\pi{18}\int_1^{10}\sqrt u\dd u
 =\frac\pi{27}(10\sqrt{10}-1).
\]
若要求封闭实体的总表面积，还须加 $x=1$ 处圆盘的面积 $\pi$。误用 $2\pi\int x^3\dd x$ 会漏掉斜母线因素。

\begin{example}[球面与端点的参数选择]
半径为 $R>0$ 的球，其表面积是多少？
\end{example}
\paragraph{思路分析}以 $y=\sqrt{R^2-x^2}$ 为母线，端点导数无界。改用上半圆参数 $x=R\cos t,y=R\sin t$，$0\le t\le\pi$，可以在本章连续可微条件下直接计算。
\paragraph{解答}速度为 $R$，旋转半径为 $R\sin t$，母线绕 $x$ 轴生成整个球面且不重复，因此
\[
 S=2\pi R^2\int_0^\pi\sin t\dd t=4\pi R^2.
\]
若将整个圆周旋转，同一球面被上下半圆分别生成，会重复计数。

\section{物理应用与质量中心}
\label{lec:12:physics}
\subsection{变力做功}
一维运动中 $F(x)$ 表示沿正向的力分量。微小位移的功为 $F(x)\dd x$，从 $a$ 到 $b$ 的功为 $W=\int_a^bF(x)\dd x$。有向功可能为负；缓慢拉伸弹簧所需外力与弹簧恢复力的方向相反，必须说明计算哪一个力的功。

\begin{example}[弹簧常数由观测确定]
将弹簧从自然长度拉长 $0.1\,\mathrm m$ 需力 $20\,\mathrm N$。求继续从拉长 $0.1\,\mathrm m$ 到 $0.3\,\mathrm m$ 时外力的功。
\end{example}
\paragraph{思路分析}胡克定律 $F=kx$ 中 $x$ 是相对自然长度的伸长量。先由观测求 $k$，再在给定伸长区间积分。
\paragraph{解答}$k=20/0.1=200\,\mathrm{N/m}$，故
\[
 W=\int_{0.1}^{0.3}200x\dd x=100(0.3^2-0.1^2)=8\,\mathrm J.
\]
不能以最终力 $60\,\mathrm N$ 乘总位移 $0.2\,\mathrm m$，因为力不是常量。

\subsection{抽水做功：重量乘提升距离}
忽略摩擦及出口动能，设液体密度 $\rho$ 恒定，重力加速度为 $g$。高度 $y$ 处水平水层的截面积为 $A(y)$，体积 $A(y)\dd y$，重力 $\rho gA(y)\dd y$。若出口高度为 $h_0$，这层水须提升 $h_0-y$，故
\[
 W=\rho g\int_{y_1}^{y_2}A(y)(h_0-y)\dd y.
\]
积分限覆盖原来有水的高度，提升距离按每层原始位置计算。这里无需把水位下降当作新的积分变量，因为每一层水的势能增加可以直接累加。

\begin{example}[圆锥容器抽水]
顶点向下的圆锥容器高 $H$、口半径 $R$，盛满水。将水抽到高出口 $H+h$（$h\ge0$）处，求功。
\end{example}
\paragraph{思路分析}原点设在锥顶，向上为 $y$ 轴。相似三角形给 $r(y)/y=R/H$；每层不是同样重，也不是提升同样距离。
\paragraph{解答}
\begin{align*}
 W&=\rho g\frac{\pi R^2}{H^2}\int_0^Hy^2(H+h-y)\dd y\\
 &=\rho g\frac{\pi R^2}{H^2}
 \left(\frac{(H+h)H^3}{3}-\frac{H^4}{4}\right)\\
 &=\pi\rho gR^2\left(\frac{H^2}{12}+\frac{hH}{3}\right).
\end{align*}
当 $h=0$，得到 $\pi\rho gR^2H^2/12$。增加出口高度 $h$ 后多做的功等于全部水的重力乘 $h$，与结果的第二项一致。
\begin{figure}[htbp]
\centering
\begin{tikzpicture}[scale=1.1]
\draw(-2,3)--(0,0)--(2,3);\draw(-2,3)--(2,3);
\draw[->](0,-.2)--(0,4)node[above]{$y$};
\draw[dashed](-2.4,3.7)--(2.7,3.7) node[right]{出口 $H+h$};
\draw[thick](-1,1.5)--(1,1.5);
\draw[<->](0,1.65)--(1,1.65)node[midway,above]{$r(y)$};
\draw[<->](2.3,1.5)--(2.3,3.7)node[midway,right]{$H+h-y$};
\node[left]at(0,1.5){$y$};\node[above left]at(0,3){$H$};
\end{tikzpicture}
\caption{轴截面中的水平水层：相似关系确定半径，出口高度确定提升距离。}
\end{figure}

\subsection{静水压力：深度乘面积}
液面下深度 $h$ 处的表压为 $p=\rho gh$。对竖直平板，取水平窄条，使同一条上的压强相同；若宽度为 $w(y)$、深度为 $h(y)$，则一侧水产生的合力大小为
\[
 F=\rho g\int h(y)w(y)\dd y.
\]
此式使用相对于大气的表压；两侧都有液体时要积分两侧压强之差。板只有部分浸水时，积分限只覆盖浸水部分。

\begin{example}[圆形观测窗]
竖直圆窗半径 $R$，圆心在液面下 $H>R$，求一侧水压力。
\end{example}
\paragraph{思路分析}以圆心为原点，向下为 $y$ 正向，深度为 $H+y$。由圆方程，水平条的宽度为 $2\sqrt{R^2-y^2}$。奇偶性可消去偏离圆心的项。
\paragraph{解答}
\begin{align*}
 F&=\rho g\int_{-R}^R(H+y)2\sqrt{R^2-y^2}\dd y\\
 &=\rho gH\int_{-R}^R2\sqrt{R^2-y^2}\dd y
 =\rho gH\pi R^2.
\end{align*}
最后的积分是圆面积；含 $y$ 的项为奇函数，积分为零。结果等于圆心深度处的压强乘面积，这里由对称性证明，而不是假设整个窗的压强相同。

\subsection{密度、质量与质心}
线密度 $\lambda(x)\ge0$ 的细杆质量为 $m=\int_a^b\lambda(x)\dd x$。当 $m>0$，由各微段的力矩相加，质心坐标为
\[
 \bar x=\frac{\int_a^bx\lambda(x)\dd x}{m}.
\]
均匀平面区域 $g(x)\le y\le f(x)$ 的面积为 $A>0$，竖条中心纵坐标为 $(f+g)/2$，因此
\[
 \bar x=\frac1A\int_a^bx(f-g)\dd x,
 \qquad
 \bar y=\frac1{2A}\int_a^b(f^2-g^2)\dd x.
\]
这些公式的依据是关于坐标轴的面积力矩，不能把 $\bar y$ 误写成边界函数平均值的普通平均。

\begin{example}[非均匀杆与曲边区域的质心]
求 $0\le x\le1$、$\lambda(x)=1+x$ 的细杆质心；求 $0\le y\le x,0\le x\le1$ 的均匀三角形区域质心。
\end{example}
\paragraph{思路分析}先算总质量或面积，再用相应力矩除以总量。杆右侧更重，质心应在 $1/2$ 的右边。
\paragraph{解答}杆的质量为 $m=1+1/2=3/2$，力矩为 $1/2+1/3=5/6$，故 $\bar x=5/9$。
三角形面积为 $A=1/2$，所以
\[
 \bar x=2\int_0^1x^2\dd x=\frac23,
 \qquad \bar y=\int_0^1x^2\dd x=\frac13.
\]

\section{含参几何量的最值}
\label{lec:12:optimization}
几何优化的顺序是：写出闭合区域，确定参数可行范围，建立面积或体积函数，化为单变量，寻找驻点并比较端点或边界极限。单有 $F'(t)=0$ 不足以证明全局最值。对连续函数 $f$ 和可导边界 $u,v$，由原函数链式法则，
\[
 \frac{\dd}{\dd t}\int_{u(t)}^{v(t)}f(x)\dd x
 =f(v(t))v'(t)-f(u(t))u'(t).
\]

\begin{example}[移动窗口的最大面积]
对 $t>0$，求曲线 $y=\ee^{-x}$、$x$ 轴、$x=t$ 和 $x=2t$ 围成的面积最大值。
\end{example}
\paragraph{思路分析}窗口随参数变宽，同时移向函数更小的位置。先将两种效应包含在同一个积分中，再用导数比较。
\paragraph{解答}
\[
 A(t)=\int_t^{2t}\ee^{-x}\dd x=\ee^{-t}-\ee^{-2t},
 \quad A'(t)=\ee^{-t}(2\ee^{-t}-1).
\]
在 $0<t<\ln2$ 导数为正，$t>\ln2$ 导数为负，故全局最大值在 $t=\ln2$ 取得，$A_{\max}=1/2-1/4=1/4$。边界极限 $t\to0^+$ 及 $t\to\infty$ 都给零，与判定一致。

\section{分层练习与解析}
以下均为自编练习，用于检验建模步骤。
\begin{exercise}[基础：水平边界]
求 $x=y^2$ 与 $x=2-y^2$ 围成的面积。
\end{exercise}
\paragraph{解析}交点满足 $y^2=1$，故 $-1\le y\le1$。右减左得
$A=\int_{-1}^1(2-2y^2)\dd y=8/3$。若使用竖条，需要分别描述上下边界，横条更直接。

\begin{exercise}[基础：心形线面积与弧长]
求 $r=a(1+\cos\theta)$（$a>0$）的面积与整条曲线弧长。
\end{exercise}
\paragraph{解析}区间 $[0,2\pi]$ 单次描出心形线，面积为
\[
 A=\frac{a^2}{2}\int_0^{2\pi}(1+2\cos\theta+\cos^2\theta)\dd\theta
 =\frac32\pi a^2.
\]
速度为 $a\sqrt{2+2\cos\theta}=2a|\cos(\theta/2)|$，在 $\pi$ 处分段，
\[
 L=2a\int_0^\pi\cos(\theta/2)\dd\theta
 -2a\int_\pi^{2\pi}\cos(\theta/2)\dd\theta=8a.
\]

\begin{exercise}[迁移：偏移的竖直轴]
求 $y=x,y=x^2$ 围成区域绕 $x=2$ 的体积。
\end{exercise}
\paragraph{解析}区域在轴左侧，竖条半径为 $2-x$、高度为 $x-x^2$。故
\[
 V=2\pi\int_0^1(2-x)(x-x^2)\dd x
 =2\pi\left(1-1+\frac14\right)=\frac\pi2.
\]
半径不能写成 $x-2$，否则体积为负。

\begin{exercise}[迁移：侧面积与总面积]
求线段 $y=x,0\le x\le h$（$h>0$）绕 $x$ 轴形成圆锥的侧面积与封闭表面积。
\end{exercise}
\paragraph{解析}半径为 $x$，$\dd s=\sqrt2\dd x$，故侧面积
$S=2\pi\sqrt2\int_0^hx\dd x=\pi\sqrt2h^2$。底圆面积为 $\pi h^2$，总面积为 $\pi h^2(1+\sqrt2)$。这与圆锥公式 $\pi r l$（$r=h,l=\sqrt2h$）一致。

\begin{exercise}[迁移：部分装水的圆柱]
半径 $R$、高 $H$ 的竖直圆柱容器内水深为 $d$（$0\le d\le H$），将水抽至容器顶口，求功。
\end{exercise}
\paragraph{解析}底面为原点，水层截面积恒为 $\pi R^2$，提升距离为 $H-y$，但积分上限是 $d$：
\[
 W=\pi\rho gR^2\int_0^d(H-y)\dd y
 =\pi\rho gR^2(Hd-d^2/2).
\]
当 $d=0$ 时功为零，$d=H$ 时为 $\pi\rho gR^2H^2/2$。

\begin{exercise}[迁移：三角形闸门]
竖直三角形闸门顶边位于水面，顶边宽 $B$，下顶点深度为 $H$。求一侧水压力。
\end{exercise}
\paragraph{解析}以深度 $y$ 为变量，由相似三角形 $w(y)=B(1-y/H)$。因此
\[
 F=\rho gB\int_0^Hy(1-y/H)\dd y
 =\frac16\rho gBH^2.
\]
全板面积为 $BH/2$，面积中心深度为 $H/3$，相乘也给同一结果。最深处压强乘全板面积会高估。

\begin{exercise}[综合：面积约束与可行区间]
设 $f(x)=ax+b\ge0$（$0\le x\le1$），且曲线下方面积为 $1$。求绕 $x$ 轴的最小体积。
\end{exercise}
\paragraph{解析}面积约束为 $a/2+b=1$，所以 $b=1-a/2$。线性函数非负等价于端点非负，即 $-2\le a\le2$。体积为
\[
 \frac V\pi=\int_0^1[1+a(x-1/2)]^2\dd x
 =1+a^2\int_0^1(x-1/2)^2\dd x=1+\frac{a^2}{12}.
\]
故 $a=0,b=1$ 时最小，$V_{\min}=\pi$。零交叉项来自 $\int_0^1(x-1/2)\dd x=0$，必须先保留并检查该项。

\paragraph{章末自查}能够解释面积条的高度、垫片的两个半径、柱壳的半径与高度、侧面积的母线长度、抽水的重量与提升距离、静水压力的深度与条宽，就能从同一个微元原则重建这些公式。每题还须检查扫描是否重复、参数范围是否合理、根号是否漏绝对值、结果单位是否与所求量一致。


\chapter{反常积分}
\label{lec:13:chapter}

本章使用定积分、原函数、换元、分部积分及函数极限。定积分原来处理有限区间上有界函数的累积量；这里要回答两个新问题：区间无限延伸时，累积量能否有限？函数在某点无界时，该点附近的累积量能否有限？共同做法是先避开反常位置，在普通定积分上计算，再考察截断极限。

\section{截断定义与反常位置}
\label{lec:13:definition}
\begin{definition}[无穷区间积分]
设实值函数 $f$ 在 $[a,+\infty)$ 的每个有限闭子区间上黎曼可积。若有限极限
\[
I=\lim_{T\to+\infty}\int_a^T f(x)\,\dd x
\]
存在，称反常积分 $\int_a^{+\infty}f(x)\,\dd x$ 收敛，其值为 $I$；否则称发散。类似地，$\int_{-\infty}^a f$ 用 $\lim_{A\to-\infty}\int_A^a f$ 定义。
\end{definition}
定义要求极限为实数。截断积分趋于 $+\infty$、$-\infty$ 或振荡没有极限，都属于发散。符号 $+\infty$ 表示极限方向，不能作为普通数代入原函数。

\begin{definition}[有限端点的瑕积分]
设 $f$ 在 $(a,b]$ 的每个闭子区间上可积。若 $a$ 附近有无界等奇异现象，以
\[
\int_a^b f(x)\,\dd x
=\lim_{\varepsilon\to0^+}\int_{a+\varepsilon}^b f(x)\,\dd x
\]
定义瑕积分，有限极限存在时称收敛。右端点 $b$ 的定义为
$\lim_{\varepsilon\to0^+}\int_a^{b-\varepsilon}f(x)\,\dd x$。
\end{definition}
端点无界不要求函数一定趋于无穷：例如 $\sin(1/x)/x$ 在零点附近无界振荡，也要用截断极限。若端点极限有限，可以补定义后按普通定积分处理。

\subsection{多个反常位置分别判断}
若 $c\in(a,b)$ 是内部奇点，定义
\[
\int_a^b f=\int_a^c f+\int_c^b f
\]
的前提是两项分别收敛。若两个端点都反常，先取普通内点 $d$，分别定义 $\int_a^d f$ 和 $\int_d^b f$；若积分区间为整条实轴，也分别检查 $\int_{-\infty}^d f$ 与 $\int_d^{+\infty}f$。有限个反常位置都用此法拆分。分点的选择不影响结果，因为改变分点只增减一段普通定积分。

\begin{example}[独立截断与主值]
判断 $\int_{-1}^1\dd x/x$ 是否收敛，并计算其对称截断极限。
\end{example}
\paragraph{思路分析} 零点是内部奇点。普通反常积分要求两侧各自有限；把两侧的无穷量抵消，会改变定义。
\paragraph{解答} 对 $\varepsilon,\delta>0$，
\[
\int_{-1}^{-\varepsilon}\frac{\dd x}{x}=\ln\varepsilon\to-\infty,
\qquad
\int_\delta^1\frac{\dd x}{x}=-\ln\delta\to+\infty.
\]
所以普通反常积分发散。取相同截断尺度则两项之和恒为零，因此 Cauchy 主值为零。一般截断和为 $\ln(\varepsilon/\delta)$；取 $\delta=\varepsilon^2$ 时它趋于 $+\infty$。这说明对称截断隐藏了路径依赖。
\begin{remark}
主值是另一个定义，本章除明确写出“主值”外均讨论普通反常积分。奇函数在对称区间上的积分为零，只有在积分本身已存在时才能直接使用。
\end{remark}

\begin{example}[无界端点的有限面积]
计算 $\int_0^1\dd x/\sqrt{1-x^2}$。
\end{example}
\paragraph{思路分析} 只有右端点反常。原函数为 $\arcsin x$，但先在 $[0,1-\varepsilon]$ 上使用牛顿--莱布尼茨公式。
\paragraph{解答}
\[
\int_0^{1-\varepsilon}\frac{\dd x}{\sqrt{1-x^2}}
=\arcsin(1-\varepsilon),\qquad
\lim_{\varepsilon\to0^+}\arcsin(1-\varepsilon)=\frac\pi2.
\]
故积分收敛，值为 $\pi/2$。函数高度趋于无穷，仍可能有有限累积量。

\section{基准幂函数积分}
\label{lec:13:p-integrals}
\begin{proposition}[两种幂积分的临界指数]
对实数 $p$ 和 $a>0$，
\[
\int_a^{+\infty}x^{-p}\,\dd x\text{ 收敛}\iff p>1;
\qquad
\int_0^a x^{-p}\,\dd x\text{ 收敛}\iff p<1.
\]
收敛时积分值依次为 $a^{1-p}/(p-1)$、$a^{1-p}/(1-p)$。
\end{proposition}
\begin{proof}
当 $p\ne1$ 时，截断计算给出
\[
\int_a^T x^{-p}\,\dd x=\frac{T^{1-p}-a^{1-p}}{1-p},
\qquad
\int_\varepsilon^a x^{-p}\,\dd x
=\frac{a^{1-p}-\varepsilon^{1-p}}{1-p}.
\]
在无穷远处，$T^{1-p}$ 只有 $p>1$ 时趋零；$p<1$ 时趋于正无穷，截断积分发散。在零点附近，$\varepsilon^{1-p}$ 只有 $p<1$ 时趋零；$p>1$ 时趋于正无穷，截断积分同样发散。$p=1$ 时，两式分别为 $\ln(T/a)$、$\ln(a/\varepsilon)$，均趋于正无穷。
\end{proof}
有限端点的 $\int_a^b(x-a)^{-p}\,\dd x$ 和 $\int_a^b(b-x)^{-p}\,\dd x$ 都收敛当且仅当 $p<1$。前者令 $t=x-a$，后者令 $t=b-x$；后一代换有 $\dd x=-\dd t$，反向积分限与负号相抵消。$p\le0$ 时端点可补成连续函数，仍符合此结论。

\paragraph{如何记住方向} 无穷远需要衰减得快，故 $p$ 要大于 $1$；有限端点要求奇性增长得慢，故 $p$ 要小于 $1$。应先写清反常端点，再使用这两个条件。

\section{Cauchy 收敛准则：控制尾部积分}
\label{lec:13:cauchy}
\begin{theorem}[反常积分的 Cauchy 准则]
设 $f$ 在 $[a,+\infty)$ 的每个有限闭子区间上可积。积分 $\int_a^{+\infty}f$ 收敛，当且仅当对每个 $\varepsilon>0$，存在 $A>a$，使任意 $T>S>A$ 都满足
\[
 \left|\int_S^T f(x)\,\dd x\right|<\varepsilon.
\]
\end{theorem}
\paragraph{直观理解} 收敛要求截断位置足够远以后，再增加任意一段尾部都只能改变很小的积分值。这里要控制的是所有尾部区间，不能只检查某一列选定的区间。
\begin{proof}
记 $F(T)=\int_a^T f$。若 $F(T)\to I\in\R$，取 $A$ 使 $T>A$ 时 $|F(T)-I|<\varepsilon/2$。于是 $T>S>A$ 时，
\[
 |F(T)-F(S)|\le |F(T)-I|+|F(S)-I|<\varepsilon.
\]
反过来，假设尾部条件成立。取正整数 $n$ 并令 $u_n=F(a+n)$，给定任意正误差，当 $m,n$ 足够大时，尾部条件给出 $|u_m-u_n|$ 小于该误差。因此 $\{u_n\}$ 是 Cauchy 数列，由第2章的准则，存在 $I\in\R$ 使 $u_n\to I$。
再给定 $\varepsilon>0$，用尾部条件选取 $A$，使任意两个大于 $A$ 的截断点所对应的 $F$ 值之差小于 $\varepsilon/2$。固定一个足够大的 $n$，使 $a+n>A$ 且 $|F(a+n)-I|<\varepsilon/2$。对任意 $T>A$，无论 $T$ 与 $a+n$ 谁更大，都有
\[
 |F(T)-I|\le |F(T)-F(a+n)|+|F(a+n)-I|<\varepsilon.
\]
所以 $F(T)\to I$，反常积分收敛。
\end{proof}
对于左端点 $a$ 为瑕点的积分 $\int_a^b f$，对应条件是：对每个 $\varepsilon>0$，存在 $\delta>0$，使 $a<u<v<a+\delta$ 时 $|\int_u^v f|<\varepsilon$。证明仍考察两个截断积分的差，充分性可先取截断点 $a+(b-a)/n$ 再控制任意截断点。右端点情形同理。
\paragraph{如何用于证明发散} 找到一列越来越靠近反常位置的区间，使各区间上积分的绝对值始终不小于同一个正常数，就违反必要条件。本章后面的振荡积分练习会使用这种方法。

\section{非负函数的比较判别}
\label{lec:13:comparison}
以下“局部可积”表示在避开反常位置的每个有限闭子区间上可积。本节函数均为实值。

\begin{theorem}[非负积分的单调有界准则]
若 $f\ge0$ 在 $[a,+\infty)$ 局部可积，则 $\int_a^{+\infty}f$ 收敛当且仅当截断积分 $F(T)=\int_a^T f$ 有上界。
\end{theorem}
\begin{proof}
对 $T_2>T_1$，$F(T_2)-F(T_1)=\int_{T_1}^{T_2}f\ge0$，所以 $F$ 单调不减。有上界时，由单调有界原理有有限极限；若有限极限存在，单调性给出 $F(T)\le\lim_{S\to\infty}F(S)$。有限反常端点也可用越来越接近该端点的截断区间证明。
\end{proof}

\begin{theorem}[比较判别法]
设 $f,g$ 非负、局部可积，并在所研究的反常端点附近满足 $f\le g$。则 $\int g$ 收敛推出 $\int f$ 收敛；$\int f$ 发散推出 $\int g$ 发散。此处各积分只针对同一个反常端点。
\end{theorem}
\begin{proof}
先把端点附近以外的普通积分分离。对剩余截断积分，$0\le\int f\le\int g$。若后者有上界，前者也有上界，由上一准则收敛。第二个结论是第一个结论的逆否命题。
\end{proof}
只需端点附近比较，因为有限普通区间不会影响敛散性。证明收敛应找可积的上界，证明非负积分发散应找不可积的下界；找到不可积上界或可积下界，都不能完成判断。

\begin{theorem}[极限比较判别法]
设 $f\ge0$、$g>0$ 在反常端点附近局部可积，且该端点处
\[
\lim\frac{f(x)}{g(x)}=c,\qquad 0<c<+\infty.
\]
则两积分同敛散。若极限为 $0$，$\int g$ 收敛可推出 $\int f$ 收敛；若极限为 $+\infty$，$\int g$ 发散可推出 $\int f$ 发散。
\end{theorem}
\begin{proof}
当 $0<c<+\infty$ 时，在端点的足够小邻域内有
$(c/2)g\le f\le(3c/2)g$，应用两次比较判别法。极限为零时在该邻域内 $f\le g$；极限为正无穷时 $f\ge g$，分别得到所述单向结论。
\end{proof}
因此，非负函数的渐近等价可以用于敛散判断，不能直接用来计算积分值。比值趋零也不自动代表收敛：$1/(x\ln x)=o(1/x)$，但它在无穷远处仍不可积。

\begin{example}[从分母提取主项]
判断 $\int_1^{+\infty}\dd x/[x\sqrt[3]{x^2+1}]$ 的敛散性。
\end{example}
\paragraph{思路分析} 反常位置只有无穷远，分母的主增长为 $x\cdot x^{2/3}=x^{5/3}$。
\paragraph{解答} 被积函数非负连续，而且
\[
\frac{1/[x\sqrt[3]{x^2+1}]}{x^{-5/3}}
=\frac1{\sqrt[3]{1+x^{-2}}}\to1.
\]
由于 $5/3>1$，比较的幂积分收敛，原积分收敛。

\begin{example}[两端约束取交集]
求 $\int_0^1 x^\beta/\sqrt{1-x^2}\,\dd x$ 收敛的实参数范围。
\end{example}
\paragraph{思路分析} 零点的幂次随参数变化，右端点的根式固定产生奇性，应取 $1/2$ 分段。
\paragraph{解答} 在零点处，
\[
\frac{x^\beta/\sqrt{1-x^2}}{x^\beta}\to1,
\]
故左侧收敛当且仅当 $\beta>-1$。在右端点处，利用 $1-x^2=(1-x)(1+x)$，
\[
\frac{x^\beta/\sqrt{1-x^2}}{(1-x)^{-1/2}}
=\frac{x^\beta}{\sqrt{1+x}}\to\frac1{\sqrt2}.
\]
右侧对每个实数 $\beta$ 都收敛。两侧条件取交集，得到 $\beta>-1$。

\begin{example}[分子抵消降低奇性]
求 $\int_0^1(1-\cos x)x^{-m}\,\dd x$ 收敛的实参数范围。
\end{example}
\paragraph{思路分析} 分母单独看似有 $x^{-m}$ 的奇性，但分子在零点有二阶零点。判断实际阶数必须先考虑这项抵消。
\paragraph{解答} 被积函数非负，利用 $1-\cos x=2\sin^2(x/2)$ 及 $\sin t/t\to1$，
\[
\lim_{x\to0^+}\frac{(1-\cos x)x^{-m}}{x^{2-m}}
=\lim_{x\to0^+}\frac{1-\cos x}{x^2}=\frac12.
\]
幂指数为 $m-2$，有限端点条件是 $m-2<1$，故恰在 $m<3$ 时收敛。$m=3$ 与 $1/(2x)$ 同阶，仍发散。

\subsection{临界幂次中的对数}
\label{lec:13:logarithm}
\begin{example}[对数代换处理临界项]
判断 $\int_\ee^{+\infty}\dd x/[x(\ln x)^q]$ 与
$\int_0^{\ee^{-1}}\dd x/[x(\ln(1/x))^q]$ 的敛散性。
\end{example}
\paragraph{思路分析} 两题的幂次都是临界的 $x^{-1}$。$\dd x/x$ 正好是对数微分，用代换使对数成为新的幂函数。
\paragraph{解答} 第一题在 $[\ee,T]$ 上令 $t=\ln x$，得到
\[
\int_\ee^T\frac{\dd x}{x(\ln x)^q}
=\int_1^{\ln T}t^{-q}\,\dd t.
\]
$T\to\infty$ 时 $\ln T\to\infty$，故收敛当且仅当 $q>1$。
第二题在 $[\varepsilon,\ee^{-1}]$ 上令 $t=\ln(1/x)=-\ln x$，有 $\dd x/x=-\dd t$，于是
\[
\int_\varepsilon^{\ee^{-1}}\frac{\dd x}{x(\ln(1/x))^q}
=\int_1^{\ln(1/\varepsilon)}t^{-q}\,\dd t.
\]
$\varepsilon\to0^+$ 时新上限趋于正无穷，条件同为 $q>1$。若 $q=1$，原函数转为 $\ln\ln x$ 或 $\ln\ln(1/x)$，发散非常缓慢，仍是发散。

\section{换元与分部积分必须保留极限}
\label{lec:13:calculation}
反常积分的换元规则来自普通定积分：先固定所有截断端点，在有限区间上做一一对应的光滑代换，记录新积分限，再令原截断端点趋向反常位置。分部积分也先写有限区间公式；把边界项与剩余积分分别取极限前，应证明它们有相应极限。不得把两个发散项形式相减后当作有限答案。

\begin{example}[平方代换与两端极限]
计算 $\int_0^{+\infty}\dd x/[\sqrt{x}(1+x)]$。
\end{example}
\paragraph{思路分析} 零点与无穷远都要检查。平方代换同时消去根式，产生 $1/(1+t^2)$ 的标准积分。
\paragraph{解答} 零点附近与 $x^{-1/2}$ 同阶，无穷远与 $x^{-3/2}$ 同阶，两端分别收敛。在 $[\varepsilon,T]$ 令 $x=t^2$，其中 $t>0$，
\[
\int_\varepsilon^T\frac{\dd x}{\sqrt{x}(1+x)}
=2\int_{\sqrt\varepsilon}^{\sqrt T}\frac{\dd t}{1+t^2}
=2\bigl(\arctan\sqrt T-\arctan\sqrt\varepsilon\bigr).
\]
分别令 $\varepsilon\to0^+$、$T\to\infty$，得到积分值 $\pi$。这里两个端点分别存在极限，所以可以合写截断式；合写本身不能代替两端审敛。

\begin{example}[分部积分中的消失边界项]
计算 $\int_0^1\ln x\,\dd x$，并判断 $\int_0^1\ln x/(1-x^2)\,\dd x$ 的敛散性。
\end{example}
\paragraph{思路分析} 对数无界但比任意负幂增长慢。第一题可以直接截断分部积分，第二题先检查右端是否真有奇性，再对绝对值比较。
\paragraph{解答} 在 $[\varepsilon,1]$ 分部积分，
\[
\int_\varepsilon^1\ln x\,\dd x
=[x\ln x-x]_\varepsilon^1
=-1-\varepsilon\ln\varepsilon+\varepsilon.
\]
令 $t=1/\varepsilon$，$\varepsilon\ln\varepsilon=-\ln t/t\to0$，所以积分值为 $-1$；又因 $\ln x\le0$，其绝对值积分为 $1$。
第二题在右端点有
\[
\lim_{x\to1^-}\frac{\ln x}{1-x^2}
=\lim_{x\to1^-}\frac{1/x}{-2x}=-\frac12,
\]
可补成连续函数。在零点处
\[
\lim_{x\to0^+}x^{1/2}\left|\frac{\ln x}{1-x^2}\right|=0,
\]
因为令 $t=1/x$ 后 $-x^{1/2}\ln x=(\ln t)/\sqrt t\to0$。故绝对值最终小于 $x^{-1/2}$，零点处绝对收敛，整个积分绝对收敛。

\section{变号函数、绝对收敛与振荡}
\label{lec:13:oscillation}
\begin{definition}
若 $\int|f|$ 收敛，称 $\int f$ 绝对收敛。若 $\int f$ 收敛而 $\int|f|$ 发散，称 $\int f$ 条件收敛。多个反常位置仍分别使用这些定义。
\end{definition}
\begin{theorem}[绝对收敛准则]
局部可积实函数的绝对收敛推出收敛。
\end{theorem}
\begin{proof}
在每个普通截断区间上，$0\le f+|f|\le2|f|$。若 $\int|f|$ 收敛，比较判别法推出 $\int(f+|f|)$ 收敛。两者的截断积分都有有限极限，相减得 $\int f$ 收敛。
\end{proof}
若绝对值比较失败，应进一步检查振荡抵消；$\int|f|$ 发散不能推出 $\int f$ 发散。以下给出一个足以覆盖常用振荡模型的判别法。

\begin{theorem}[Dirichlet 判别法的可微版本]
\label{lec:13:dirichlet}
设 $h$ 在 $[a,+\infty)$ 连续，$H(x)=\int_a^x h(t)\,\dd t$ 有界；设 $g\in C^1[a,+\infty)$ 非负、单调不增，且 $g(x)\to0$。则 $\int_a^{+\infty}g(x)h(x)\,\dd x$ 收敛。
\end{theorem}
\begin{proof}
设 $|H(x)|\le M$。在有限区间上分部积分，
\[
\int_a^T gh=g(T)H(T)-g(a)H(a)-\int_a^T g'H
=g(T)H(T)-\int_a^T g'H.
\]
第一项趋零。因 $g'\le0$，
\[
\int_a^T|g'(x)H(x)|\,\dd x
\le M\int_a^T(-g'(x))\,\dd x
=M\bigl(g(a)-g(T)\bigr)\le Mg(a).
\]
非负积分单调有界，所以剩余积分绝对收敛。两个极限均已验证，原积分收敛。此证明说明“振荡原函数有界”与“振幅单调趋零”各自承担什么作用。
\end{proof}
此处给的是有充分条件的可微版本；去掉单调性，仅有 $g\to0$ 时不能直接引用该结论。

\begin{example}[同一振荡函数的两种收敛]
判断 $\int_1^{+\infty}\sin x/x^2\,\dd x$ 和 $\int_1^{+\infty}\sin x/x\,\dd x$ 是否绝对收敛。
\end{example}
\paragraph{思路分析} 第一题的绝对值有可积幂函数上界。第二题振幅衰减较慢，要把普通收敛与绝对收敛分开证明。
\paragraph{解答} 第一题满足 $|\sin x/x^2|\le x^{-2}$，故绝对收敛。第二题中 $h(x)=\sin x$ 的累积函数为 $\cos1-\cos x$，有界；$g(x)=1/x$ 非负单调趋零，故由上述定理收敛。也可直接检查
\[
\int_1^T\frac{\sin x}{x}\,\dd x
=\cos1-\frac{\cos T}{T}-\int_1^T\frac{\cos x}{x^2}\,\dd x,
\]
其中边界项有极限，末项绝对收敛。另一方面，对 $k\ge1$，
\[
\int_{k\pi}^{(k+1)\pi}\frac{|\sin x|}{x}\,\dd x
\ge\frac1{(k+1)\pi}\int_{k\pi}^{(k+1)\pi}|\sin x|\,\dd x
=\frac2{(k+1)\pi}.
\]
右侧累加为发散的调和级数，故绝对值积分发散。第二题条件收敛。

\begin{example}[倒数代换把瑕点变为无穷远]
判断 $\int_0^1\sin(1/x)/x\,\dd x$ 的收敛类型。
\end{example}
\paragraph{思路分析} 零点处振荡加快，直接做幂函数比较丢失抵消。取 $t=1/x$ 把高频振荡变为 $\sin t$。
\paragraph{解答} 在 $[\varepsilon,1]$ 上有 $\dd x=-t^{-2}\dd t$，所以
\[
\int_\varepsilon^1\frac{\sin(1/x)}x\,\dd x
=\int_1^{1/\varepsilon}\frac{\sin t}{t}\,\dd t.
\]
当 $\varepsilon\to0^+$ 时上限趋于正无穷，上一例证明右侧收敛。对绝对值作同一代换，变成发散的 $\int_1^\infty|\sin t|/t\,\dd t$，故原瑕积分条件收敛。

\subsection{不要把级数必要条件照搬给积分}
级数收敛要求通项趋零；反常积分收敛一般不要求被积函数趋零。可构造连续非负函数：在每个整数 $n\ge2$ 处放一个高度为 $1$、底边长为 $2^{-n}$ 的三角形，其余处为零。各支撑区间互不相交，每个三角形面积为 $2^{-n-1}$，总面积有限，但 $f(n)=1$，函数不趋零。即使连续、非负，仍不能以“被积函数不趋零”作为一般发散判据。要证明发散，需展示不能消失的尾部积分、不可积下界或截断积分无界。

\section{补充：Gamma 函数}
\label{lec:13:gamma}
本节保留课程资料中的特殊函数内容，作为反常积分、指数衰减及分部积分的综合应用。高斯积分的精确求值涉及后续多元积分工具，此处明确作为引用，不把它当作已证明的本章基础结论。

\begin{definition}[Gamma 函数的积分定义]
对 $s>0$，定义
\[
\Gamma(s)=\int_0^{+\infty}x^{s-1}\ee^{-x}\,\dd x.
\]
\end{definition}
\paragraph{为何恰取正参数} 在零点处，$x^{s-1}\ee^{-x}\sim x^{s-1}$，所以左侧积分收敛当且仅当 $s>0$。无穷远处，对任意固定实数 $s$，
$x^{s+1}\ee^{-x}\to0$，故充分大时 $x^{s-1}\ee^{-x}\le x^{-2}$。这个极限可用整数幂上界证明：选整数 $m>\max(s+1,0)$，$x\ge1$ 时 $x^{s+1}/\ee^x\le x^m/\ee^x\to0$，后者用有限次洛必达法则。于是两端分别收敛。$s\le0$ 时原积分在零点发散。

\begin{proposition}[递推关系]
对 $s>0$，$\Gamma(s+1)=s\Gamma(s)$，并且 $\Gamma(1)=1$、$\Gamma(n+1)=n!$（$n$ 为非负整数）。
\end{proposition}
\begin{proof}
在 $[\varepsilon,T]$ 上分部积分得到
\[
\int_\varepsilon^T x^s\ee^{-x}\,\dd x
=[-x^s\ee^{-x}]_\varepsilon^T
+s\int_\varepsilon^T x^{s-1}\ee^{-x}\,\dd x.
\]
$s>0$ 保证 $\varepsilon^s\ee^{-\varepsilon}\to0$；指数衰减保证 $T^s\ee^{-T}\to0$；末项两端已经证明收敛。分别取极限得到递推式。直接积分得 $\Gamma(1)=\lim_{T\to\infty}(1-\ee^{-T})=1$，反复递推得到整数值。
\end{proof}
引用后续课程的高斯积分 $\int_0^\infty\ee^{-t^2}\,\dd t=\sqrt\pi/2$，在截断区间令 $x=t^2$ 可得 $\Gamma(1/2)=\sqrt\pi$。由递推关系得到 $\Gamma(3/2)=\sqrt\pi/2$。实负参数可通过递推关系另行延拓至非整数点，但延拓值不等于上面已经发散的积分值；本章不讨论复分析中的解析延拓与极点。

\begin{example}[缩放代换与 Gamma 值]
设 $a>0$，$s>0$，计算 $\int_0^{+\infty}x^{s-1}\ee^{-ax}\,\dd x$；求 $\int_0^\infty x^2\ee^{-2x}\,\dd x$。
\end{example}
\paragraph{思路分析} 指数中的系数 $a$ 可由线性缩放消去。需同时变换微分与幂函数，并检查 $a$ 为正保证方向不变。
\paragraph{解答} 在 $[\varepsilon,T]$ 上令 $t=ax$，有
\[
\int_\varepsilon^T x^{s-1}\ee^{-ax}\,\dd x
=a^{-s}\int_{a\varepsilon}^{aT}t^{s-1}\ee^{-t}\,\dd t.
\]
$a\varepsilon\to0^+$、$aT\to\infty$，所以所求为 $a^{-s}\Gamma(s)$。取 $s=3,a=2$，得到 $2^{-3}\Gamma(3)=2!/8=1/4$。

\section{分层练习与完整解析}
以下均为自编练习，用于检验定义、方法选择和参数边界。

\begin{exercise}[基础：多个无穷端点]
计算 $\int_{-\infty}^{+\infty}\dd x/(1+x^2)$；判断 $\int_{-\infty}^{+\infty}x/(1+x^2)\,\dd x$ 的普通收敛性和对称主值。
\end{exercise}
\noindent\textbf{解析}\quad 第一题在零点分段。两侧的截断积分分别为 $\arctan T$ 和 $-\arctan A$；当 $T\to+\infty$、$A\to-\infty$ 时，两者均趋于 $\pi/2$，故总积分为 $\pi$。第二题右側截断为 $\tfrac12\ln(1+T^2)\to\infty$，左侧为 $-\tfrac12\ln(1+A^2)\to-\infty$，故普通积分发散。对称区间 $[-T,T]$ 上被积函数为可积奇函数，积分为零，主值为零。

\begin{exercise}[基础：直接比较]
判断 $\int_1^\infty\ee^{-x^2}\,\dd x$ 的敛散性，并给出截断尾部的上界。
\end{exercise}
\noindent\textbf{解析}\quad  $x\ge1$ 时 $x^2\ge x$，故 $0\le\ee^{-x^2}\le\ee^{-x}$，后者积分收敛。若 $T\ge1$，则 $\int_T^\infty\ee^{-x^2}\,\dd x\le\int_T^\infty\ee^{-x}\,\dd x=\ee^{-T}$。尾部上界趋零，同时量化了截断误差。

\begin{exercise}[迁移：参数与两端]
求 $\int_0^\infty x^\alpha/(1+x)^\beta\,\dd x$ 收敛的实参数条件。
\end{exercise}
\noindent\textbf{解析}\quad  在零点处与 $x^\alpha$ 同阶，收敛要求 $\alpha>-1$；在无穷远处与 $x^{\alpha-\beta}$ 同阶，收敛要求 $\beta-\alpha>1$。被积函数正，极限比较双向成立，所以条件是 $\alpha>-1$ 且 $\beta>\alpha+1$。两条等号边界各与 $1/x$ 同阶而发散。

\begin{exercise}[迁移：更慢的对数]
求 $\int_{\ee^\ee}^\infty\dd x/[x\ln x(\ln\ln x)^r]$ 收敛的条件。
\end{exercise}
\noindent\textbf{解析}\quad  在截断区间令 $t=\ln\ln x$，有 $\dd t=\dd x/(x\ln x)$，下限变为 $1$，上限变为 $\ln\ln T\to\infty$。所得 $\int_1^\infty t^{-r}\,\dd t$ 恰在 $r>1$ 时收敛。此题说明幂次和第一个对数均临界时，还要继续考察下一层对数。

\begin{exercise}[综合：振荡的参数分支]
对 $\alpha>0$，判断 $\int_1^\infty\sin x/x^\alpha\,\dd x$ 的收敛类型。
\end{exercise}
\noindent\textbf{解析}\quad  对每个 $\alpha>0$，$x^{-\alpha}$ 属于 Dirichlet 判别法中的单调振幅，故原积分收敛。$\alpha>1$ 时 $|\sin x|/x^\alpha\le x^{-\alpha}$，所以绝对收敛。$0<\alpha\le1$ 时，在半周期上
\[
\int_{k\pi}^{(k+1)\pi}\frac{|\sin x|}{x^\alpha}\,\dd x
\ge\frac{2}{((k+1)\pi)^\alpha}.
\]
右侧累加发散：$\alpha\le1$ 时 $(k+1)^{-\alpha}\ge(k+1)^{-1}$。故此时条件收敛。若扩展到 $\alpha=0$，截断积分 $\cos1-\cos T$ 没有极限，不能纳入上述收敛范围。

\begin{exercise}[综合：瑕点处振荡]
求 $\int_0^1 x^{-r}\sin(1/x)\,\dd x$ 在实参数 $r$ 下的收敛类型。
\end{exercise}
\noindent\textbf{解析}\quad  令 $t=1/x$，截断积分变为 $\int_1^{1/\varepsilon}t^{r-2}\sin t\,\dd t$。$r<2$ 时振幅 $t^{r-2}$ 单调趋零，故收敛。$r<1$ 时绝对值有可积上界 $t^{r-2}$，故绝对收敛；$1\le r<2$ 时令 $\alpha=2-r\in(0,1]$，由上一练习的半周期下界可知条件收敛。若 $r\ge2$，在区间 $[2k\pi+\pi/6,2k\pi+5\pi/6]$ 上 $\sin t\ge1/2$ 且 $t^{r-2}\ge1$，因此每段积分至少为 $\pi/3$。若截断积分有有限极限，则两个都趋向无穷的端点对应的积分差必须趋零，与该下界矛盾，故发散。完整分类为：$r<1$ 绝对收敛，$1\le r<2$ 条件收敛，$r\ge2$ 发散。

\paragraph{章末自查} 写下所有反常位置，分别截断；确认比较对象在同一端点满足正确的幂次条件；对变号函数区分绝对值与原函数；换元时记录微分、方向和极限；分部积分时逐项核验边界。参数题的临界值必须单独检查，渐近式的使用必须先确认自变量处于其适用极限。

\chapter*{资料与编写说明}
\markboth{资料与编写说明}{}
\addcontentsline{toc}{chapter}{资料与编写说明}
课程覆盖范围参考用户提供的《微积分（甲）I 讲义》；讲解组织及视觉风格参考
《微积分（甲）I 期末考试高频考点与备考策略》（页眉署名 Ephemeral）。
正文结合初学者的前置知识重新组织，基础和经典例题独立推导，说明定理的条件与使用边界。
原始资料仅在本地保存，未随知识库公开上传。

全文提供自包含的 LaTeX 源文件及对应 PDF。构建方法见知识库的课程入口。

\end{document}
